\documentclass[12pt]{article}

\usepackage{amsmath,amssymb,amsthm,mathtools}
\usepackage{enumitem}
\usepackage{booktabs}
\usepackage{array}
\usepackage[margin=1in]{geometry}
\usepackage{parskip}
\usepackage{xcolor}
\usepackage[colorlinks=true,linkcolor=blue!55!black,citecolor=blue!55!black]{hyperref}

%----------------------------------------------------------------------
% Theorem environments
%----------------------------------------------------------------------
\newtheorem{theorem}{Theorem}[section]
\newtheorem{lemma}[theorem]{Lemma}
\newtheorem{proposition}[theorem]{Proposition}
\newtheorem{corollary}[theorem]{Corollary}
\newtheorem{metatheorem}[theorem]{Meta-Theorem}

\theoremstyle{definition}
\newtheorem{definition}[theorem]{Definition}
\newtheorem{remark}[theorem]{Remark}
\newtheorem{example}[theorem]{Example}
\newtheorem{openproblem}[theorem]{Open Problem}
\newtheorem{heuristic}[theorem]{Heuristic}
\newtheorem{conjecture}[theorem]{Conjecture}
\newtheorem{assumption}[theorem]{Hypothesis}

%----------------------------------------------------------------------
% Operators
%----------------------------------------------------------------------
\DeclareMathOperator{\rank}{rank}
\DeclareMathOperator{\tr}{tr}
\DeclareMathOperator{\Tr}{Tr}
\DeclareMathOperator{\supp}{supp}
\DeclareMathOperator{\PoS}{PoS}
\DeclareMathOperator{\Ldim}{Ldim}
\DeclareMathOperator*{\argmin}{arg\,min}
\DeclareMathOperator*{\argmax}{arg\,max}

%----------------------------------------------------------------------
% Macros
%----------------------------------------------------------------------
\newcommand{\eqdef}{\vcentcolon=}
\newcommand{\norm}[1]{\left\lVert#1\right\rVert}
\newcommand{\vertiii}[1]{%
  \left\lvert\left\lvert\left\lvert#1\right\rvert\right\rvert\right\rvert%
}


\newcommand{\Hcal}{\mathcal H}

\newcommand{\RetKL}{\Delta_{\mathrm{ret}}^{\mathrm{KL}}}
\newcommand{\RetQuad}{\Delta_{\mathrm{ret}}^{\mathrm{quad}}}
\newcommand{\GrdLin}{\Delta_{\mathrm{grd}}^{\mathrm{lin}}}
\newcommand{\Dunres}{\Delta_{\mathrm{grd}}^{\mathrm{unres}}}
\newcommand{\DHankstr}{\Delta_{\mathrm{grd}}^{\mathrm{Hank,str}}}
\newcommand{\Com}{\Delta_{\mathrm{com}}}
\newcommand{\Comex}{\Com}
\newcommand{\ComRD}{\Delta_{\mathrm{com}}^{\mathrm{RD}}}
\newcommand{\ComGame}{\Delta_{\mathrm{com}}^{\mathrm{game}}}

\newcommand{\Sp}[1]{\mathcal S_{#1}}

\newcommand{\rankBdet}{\operatorname{rank}_{\mathbb B}^{\mathrm{det}}}
\newcommand{\kappaMN}{\kappa_{\mathrm{MN}}}

\newcommand{\Reg}{\mathrm{Reg}}
\newcommand{\Est}{\mathsf E}
\newcommand{\Aapp}{A_T}
\newcommand{\Fbudget}{\mathsf F}
\newcommand{\PsiAgg}{\Psi}

\newcommand{\EA}{\operatorname{Pinch}_{\mathcal A}}
\newcommand{\KyF}[2]{\left\|#1\right\|_{(#2)}}
\newcommand{\PoSlin}{\PoS_{\mathrm{lin}}^{\mathrm{loc}}}
\newcommand{\PoSquad}{\PoS_{\mathrm{quad}}}

\newcommand{\Ren}[1]{D_{#1}}
\newcommand{\sRen}[1]{\widetilde D_{#1}}
\newcommand{\KL}{D_{\mathrm{KL}}}
\newcommand{\Dmax}{D_{\max}}
\newcommand{\Dmaxsym}{D_{\max}^{\mathrm{sym}}}

\newcommand{\Splus}{\mathcal S^{+}}
\newcommand{\Lsync}{L_{\mathrm{sync}}^{\mathrm{adapt}}}
\newcommand{\Lsyncu}{L_{\mathrm{sync}}^{\mathrm{univ}}}
\newcommand{\Esync}{E_{\mathrm{sync}}}
\newcommand{\EsyncSI}{E_{\mathrm{sync}}^{\mathrm{SI}}}
\newcommand{\MistRI}{\operatorname{Mistakes}_{\mathrm{active,RI}}}
\newcommand{\Valt}{V_{\mathrm{alt}}}
\newcommand{\Vsim}{V_{\mathrm{sim}}}
\newcommand{\Gact}{\mathcal G^{\mathrm{act}}}
\newcommand{\free}{\mathsf f}
\newcommand{\rd}{\mathsf r}
\newcommand{\rot}{\mathrm{rot}}
\newcommand{\flipl}{\mathrm{flip}}
\newcommand{\mistk}{\mathrm{mis}}
\newcommand{\ind}[1]{\mathbf 1\{#1\}}
\newcommand{\RetKLr}[1]{\Delta_{\mathrm{ret}}^{\mathrm{KL},(#1)}}
\newcommand{\RetKLc}{\Delta_{\mathrm{ret}}^{\mathrm{KL,ctrl}}}
\newcommand{\RetQuadc}{\Delta_{\mathrm{ret}}^{\mathrm{quad,ctrl}}}
\newcommand{\RetKLg}{\Delta_{\mathrm{ret}}^{\mathrm{KL,gen}}}

\title{
\bfseries
The Rate--Distortion Theory of Bounded Sequential Transduction:\\
A Comparative Syntax for Finite-State Approximation
}
\author{}
\date{}

\begin{document}
\maketitle

%======================================================================
\begin{abstract}
We develop a comparative framework for finite-state approximation of
sequential objects.  The framework uses histories, finite-index right
congruences, state-count resources, enriched cost profunctors, stationary
aggregation, and variational gaps.  Commitment uses deterministic Mealy
machines and Myhill--Nerode residuals.  Retention uses stationary controlled
unifilar causal machines --- controlled $\epsilon$-machines, whose state update
$S_{t+1}=\tau(S_t,X_t,Y_t)$ may read the realized output; the input-driven
synchronized subclass, in which causal states factor through input histories,
contains all complexity constructions.
Grounding uses linear finite-rank Hankel realizations, whose resource is
operator rank rather than quotient index; a closed-form finite-budget
characterization of exact symbolic deterministic grounding remains open, but
an observable deterministic floor, its finite-state limit, and certified
finite-horizon and finite-section intervals are established below.  Task theories use stationary
Ces\`aro aggregation or discounted prefix measures.  Finite prefixes do not
define a probability measure over all lengths.

The retention full-KL gap satisfies an exact finite-state
information-bottleneck decomposition over stationary lumpable quotients, whose
controlled analogue, conditional on the input, holds for general unifilar
machines (Theorem~\ref{thm:controlled-ib}):
\[
\RetKL(M)
=
I(S;Z)
-
\max_{\substack{
\phi\ \text{lumpable}\\
|\mathcal K|\le M
}}
I(K_\phi;Z).
\]
Zero retention occurs exactly when $M$ is at least the stationary support size
of distinct predictive states.  The Gaussian quadratic restricted gap
satisfies a covariance-tail lower bound with effective rank $M-1$:
\[
\RetQuad(M)
\ge
\frac12
\sum_{i\ge M}
\lambda_i(\Sigma_\pi).
\]
Full KL admits a precise local Fisher asymptotic lower bound under regularity
assumptions.

The unrestricted linear finite-rank Hankel relaxation gap satisfies
\[
\Dunres(M)
=
\sigma_{M+1}(H_\nu)
\]
by Eckart--Young--Mirsky \cite{eckart1936,mirsky1960}.  The Hankel-structured finite-rank realization gap
satisfies the lower bound
\[
\DHankstr(M)
\ge
\sigma_{M+1}(H_\nu),
\]
with equality only under additional structural hypotheses.  For a single-letter
alphabet, $|\Sigma|=1$, the shift-intertwining Hardy-space unitary of
Theorem~\ref{thm:aak-equality} is the canonical basis-to-monomial
correspondence and is therefore automatic, not a separately supplied
hypothesis; equality then follows from Adamjan--Arov--Krein provided the one
remaining hypothesis holds, namely that the transported symbol is compact.
For $|\Sigma|>1$ the free monoid carries no shift of multiplicity one, the
scalar theorem does not apply, and equality remains conditional on a
Nehari/AAK-type finite-rank approximation theorem for the multi-shift
setting, which is not established here.

For finite output alphabets a global converse holds in probability
coordinates: $\RetKL(M)\ge\sum_{i\ge M}\lambda_i(\Sigma_p)$, where
$\Sigma_p$ is the covariance of the predictive probability vectors, with no
regularity or small-radius hypothesis.  The constant $1$ is optimal, being
approached by antipodal binary perturbations of the uniform law.  Under uniform
interiority a corresponding bound holds in a fixed minimal natural-parameter
chart, with a constant governed by the least eigenvalue of the log-partition
Hessian on that chart.  That interiority hypothesis cannot be removed: in the
natural-parameter chart no universal constant exists, since two Bernoulli laws
approaching opposite faces of the simplex keep the retention cost below
$\log2$ while their parameter covariance diverges.  There is likewise no
canonical global ``Fisher covariance'' independent of the choice of chart.

The computational-complexity theory gives NP-completeness for Gaussian
quadratic restricted retention already for minimal definite reset machines of
synchronization depth one.  A tagged depth-two construction also gives
NP-hardness.  For unrestricted full-KL retention over a finite output alphabet,
embedding $k$-means into a small interior neighborhood of the uniform law makes
the Jensen--Shannon block cost a controlled perturbation of the Euclidean one
and yields NP-hardness under a promise, again for depth-one reset machines in
which every partition is lumpable.  Because the embedding controls the
objective relatively and not merely across a promise gap, the reduction is
approximation preserving, and the problem is APX-hard: no polynomial-time
approximation scheme exists unless $\mathsf P=\mathsf{NP}$.  Both statements
concern an output alphabet whose size grows with the dimension of the embedded
instance; hardness for a fixed alphabet is not established.  The result there is
hardness rather than completeness, since deciding inequalities between
rationals and sums of logarithms of rationals is not known to be polynomial.

The temporal theory distinguishes three protocols: reset-word prediction,
persistent-stream prediction, and active query learning.  In the reset-word
realizable protocol, standard automata VC/Littlestone estimates give
$\Theta(M\log M)$ mistake complexity over fixed nontrivial alphabets.  For a
single persistent transducer stream, finite-class counting gives the
unconditional upper bound $O(M\log M)$, and an explicit stream-realizable
mistake tree gives the matching $\Omega(M\log M)$ lower bound, so the
persistent-stream complexity is also $\Theta(M\log M)$ for every fixed
alphabet pair with $|\mathcal I|\ge2$, $|\mathcal O|\ge2$.  In the active protocol the learner chooses its own inputs, so mistakes are
counted against an explicit identification objective: $\MistRI(M)$ denotes the
mistakes incurred before the learner determines the Myhill--Nerode residual
class of the current history.  Some such objective is unavoidable, since an
unconstrained active learner may repeat an input of known output forever and
make no mistakes at all.  Against that objective,
\[
\MistRI(M)
\le
O(\Esync(M)),
\]
and on the full class $\mathcal H_M$ this is $\Theta(M\log M)$
\emph{unconditionally}: an active halving learner attains the objective within
$\log_2|\mathcal H_M\times Q|=O(M\log M)$ mistakes, and an explicit gated
family forces $\Omega(M\log M)$, so
\[
\MistRI(M)
=
\Theta(M\log M),
\]
with no synchronizability or operational-equivalence hypothesis.  This is a
statement about active residual identification, not about unrestricted active
prediction.
State identification in a known minimal skeleton costs
$\Theta(\log M)$ mistakes --- the halving upper bound is matched
by a cyclic-shift family on which each round exposes one fresh bit of the
initial state --- so a decomposition whose second term is a
state-identification mistake complexity is degenerate for every subclass; a
direct-sum instance does force its two components on disjoint rounds, but both
components realized here are model-identification components, and no matching
upper bound in the two components separately is claimed.
Finite-class aggregation gives agnostic regret $O(\sqrt{T M\log M})$; matching
lower bounds are inherited only in protocols admitting an $\Omega(M\log M)$
sequential shattering construction.  In the active protocol neither bound
follows merely from letting the learner choose inputs: the $\Omega(M\log M)$
lower bound is supplied by the explicit gated construction rather than by
transfer from the passive setting, and separating the two terms of the additive
form still requires a direct-sum instance.

A conditional Schatten template states that a rank-controlled response
operator together with a Schatten-$p$ domination inequality yields a
Schatten-$p$ tail lower bound.  An oracle inequality composes offline
approximation and online estimation within type-correct hypothesis classes.
A conditional continuous bias--variance sandwich is proved for regular
approximation and estimation envelopes, and, under the explicit approximation
and estimation floors of Assumption~\ref{ass:oracle-floors}, a matching
minimax lower bound is proved, so the envelope is optimal up to the
aggregation penalty.  Whether that penalty is itself necessary remains open.

A linear second-order safety surrogate based on pinching satisfies Ky Fan
majorization and is distinct from the discrete right-congruence Price of
Safety.  A vertex correspondence labels retention by $\alpha=1$, commitment by
a Booleanized symmetrized $\alpha=0$ support vertex, and grounding by an
$\alpha=\infty$ operator-norm face; this is an organizing analogy, recorded as
a formal consistency rather than a derivation, and none of the analytic
theorems depend on it.  Every approximation claim carries an explicit type
signature fixing regime, feasible set, aggregation, support, protocol, and
analytical bridge, so that statements from different regimes are not
conflated.  The framework separates shared finite-state syntax from regime-specific
analysis: a typed variational schema, structural monotonicity and zero
thresholds, a conditional response-operator converse, and regime-specific
analytical engines.  It identifies a small set of sharply posed open research
programs in place of a single undifferentiated open problem of exact
symbolic deterministic grounding.
\end{abstract}
\tableofcontents

%======================================================================
\section{Introduction: The Bounded Sequential Generator and Its Regimes}
\label{sec:intro}
%======================================================================

This manuscript develops a comparative theory of finite-state approximation
for sequential objects.  The resource is memory, measured by the number of
internal states or, equivalently, by the index of a finite right congruence on
histories.  The cost is a regime-specific distortion functional: Boolean
mismatch for commitment, KL divergence for retention, operator-norm distance
for linear grounding relaxations, and related quantitative costs.

The framework may be viewed as a rate--distortion theory of bounded sequential
transduction.  The rate is the finite-state budget $M$.  The distortion is the
task loss induced by the regime.

%----------------------------------------------------------------------
\subsection{Bounded sequential generators}
\label{subsec:bsg-intro}
%----------------------------------------------------------------------

A \textbf{Bounded Sequential Generator} (BSG) is a deterministic finite-state
transducer, i.e.\ a Mealy machine,
\[
\mathcal A=(\mathcal S,s_0,\delta,\lambda),
\qquad
|\mathcal S|\le M,
\]
with finite input alphabet $\mathcal I$, finite output alphabet $\mathcal O$,
transition map
\[
\delta:\mathcal S\times\mathcal I\to\mathcal S,
\]
and output map
\[
\lambda:\mathcal S\times\mathcal I\to\mathcal O.
\]
It defines a synchronous causal function
\[
F_{\mathcal A}:\mathcal I^{\mathbb N}\to\mathcal O^{\mathbb N}.
\]

The deterministic Mealy BSG is the machine model for the commitment regime.
In that regime, the target is a deterministic specification, a safety
property, or a deterministic strategic rule, and the question is whether a
finite-state machine implements or enforces it exactly or approximately.  The
exact zero-error theory is governed by Myhill--Nerode residual equivalence.

%----------------------------------------------------------------------
\subsection{Three approximation regimes}
\label{subsec:three-regimes}
%----------------------------------------------------------------------

The manuscript studies three regimes.  They share a finite-state syntax but
differ in their semantic objects.  The framework does not reduce the three
regimes to a single machine category.

\begin{enumerate}[label=(\arabic*),leftmargin=2.2em]
\item \textbf{Commitment.}
The target is a deterministic residual function.  A history $u$ determines a
residual continuation map
\[
F_u(w)=F(uw)_{|u|+1:|u|+|w|}.
\]
Two histories are equivalent if they induce the same residual function.  The
minimal exact state complexity is the Myhill--Nerode index
\[
\kappa_{\det}(F)
=
\operatorname{index}(\sim_F).
\]
The exact commitment gap is a threshold phenomenon:
\[
\Com(M)=0
\quad\Longleftrightarrow\quad
M\ge\kappa_{\det}(F).
\]

\item \textbf{Retention.}
The target is a stochastic controlled process whose predictive structure is
encoded by causal states.  Each positive-mass causal state $s\in\Splus$
carries a predictive distribution $P_s$ over future outputs.  An $M$-state
approximation merges causal states into at most $M$ lumped machine states
through a stationary lumpable quotient.

The full-KL retention gap is
\[
\RetKL(M)
=
\min_{\substack{
\phi:\Splus\to\mathcal K\ \text{lumpable}\\
|\mathcal K|\le M
}}
\sum_{k\in\mathcal K}
\sum_{s:\phi(s)=k}
\pi_s
D_{\mathrm{KL}}(P_s\|\bar P_k),
\]
where
\[
\bar P_k
=
\frac{
\sum_{s:\phi(s)=k}\pi_s P_s
}{
\sum_{s:\phi(s)=k}\pi_s
}.
\]
In synchronized machines, lumpable quotients of $\Splus$ correspond to
finite-index right congruences on sufficiently long histories.  In the general
unifilar model the causal-state update reads the realized output, the feasible
quotients are the unifilar-lumpable ones, and the input-driven synchronized
subclass recovers right congruences on input histories.

When the predictive family is restricted to common-covariance Gaussians, one
obtains a quadratic restricted gap
\[
\RetQuad(M).
\]
The full-KL gap and the quadratic restricted gap are distinct objects.

\item \textbf{Grounding.}
The target is a stochastic channel or residual kernel.  The spectral converse
uses linear finite-rank realizations rather than deterministic Mealy
machines.  The relevant object is a Hankel operator $H_\nu$ built from an
impulse-response function.

The unrestricted linear finite-rank Hankel relaxation gap is
\[
\Dunres(M)
=
\inf_{\operatorname{rank}B\le M}
\norm{H_\nu-B}_{\mathrm{op}},
\]
and satisfies
\[
\Dunres(M)
=
\sigma_{M+1}(H_\nu)
\]
by Eckart--Young--Mirsky \cite{eckart1936,mirsky1960}.

The Hankel-structured finite-rank realization gap is
\[
\DHankstr(M)
=
\inf_{\substack{
\operatorname{rank}B\le M\\
B\ \text{Hankel}
}}
\norm{H_\nu-B}_{\mathrm{op}},
\]
and satisfies
\[
\DHankstr(M)
\ge
\sigma_{M+1}(H_\nu).
\]
Equality in the Hankel-structured problem requires additional structure
(Theorems~\ref{thm:aak-equality} and~\ref{thm:aak-multiletter}).  The
resource $M$ here is a finite-rank budget on a bounded operator, not the
index of a right congruence; Section~\ref{subsec:grounding-instantiation} and
the master scope of Section~\ref{subsec:meta-scope} record this distinction
explicitly, and Definition~\ref{def:symbolic-grounding-gap} below supplies a
quotient-typed \emph{symbolic} grounding gap, over deterministic Mealy
machines, whose resource is again a right-congruence index and which is
therefore the object genuinely comparable to commitment and retention.  A
closed-form finite-budget value for that symbolic gap remains open; certified
bounds are given in
Propositions~\ref{prop:grounding-finite-section},~\ref{prop:symbolic-finite-horizon}
and Theorem~\ref{thm:observable-floor}.
\end{enumerate}

The three regimes share a \textbf{typed finite-memory budget} rather than a
single literal resource.  For commitment and retention (and for the symbolic
grounding gap of Definition~\ref{def:symbolic-grounding-gap}), the budget $M$
is the index of a finite right congruence on histories.  For the
\emph{unrestricted linear} grounding relaxation, the budget $M$ is instead the
rank of a finite-rank operator approximating a Hankel operator.  These two
notions of budget are related --- both cap the ``resolution'' of the
approximant --- but they are not asserted equivalent: a rank-$M$ operator is
not in general reachable through an $M$-state quotient, and an $M$-state
quotient does not in general act on the joint alphabet the way a rank-$M$
Hankel truncation does.  Section~\ref{subsec:meta-scope} makes this typed
picture precise as four pillars common to all three regimes.

%----------------------------------------------------------------------
\subsection{The shared finite-state schema}
\label{subsec:shared-schema-intro}
%----------------------------------------------------------------------

The shared combinatorial resource is the index of a finite quotient of a
regime-typed history system (Definition~\ref{def:history-system}, developed in
Section~\ref{subsec:history-actions}).
Commitment quotients live on input histories $\mathcal I^*$ with a total right
action; unifilar retention quotients live on feasible joint input--output
histories with a pruned action, in the support-relative sense of
Definition~\ref{def:support-right-cong}; grounding quotients live on the joint
alphabet $\Sigma=\mathcal I\times\mathcal O$.  For the commitment regime, and
for the input-driven retention subclass, the classical form is the following;
Remark~\ref{rem:regime-history-instances} tabulates the three instantiations.

Let $\mathcal I^*$ denote the set of finite input histories.  A right
congruence $\sim$ on $\mathcal I^*$, in the standard sense of automata theory
\cite{sakarovitch2009}, is an equivalence relation satisfying
\[
u\sim v
\implies
ux\sim vx
\qquad
\forall x\in\mathcal I.
\]
Equivalently, the quotient $\mathcal I^*/\!\sim$ carries a well-defined
deterministic transition
\[
[u]\cdot x=[ux].
\]
The index of $\sim$ is the number of equivalence classes.  This index is the
finite-state budget.

A target sequential object may be represented as a trajectory
\[
\delta:\mathcal I^*\to\operatorname{ob}(\mathbf R),
\]
where $\mathbf R$ is a regime-specific residual category.  An $M$-state
approximation factors $\delta$ through a finite quotient:
\[
\mathcal I^*
\longrightarrow
\mathcal I^*/\!\sim
\longrightarrow
\mathbf R,
\qquad
\operatorname{index}(\sim)\le M.
\]
The associated variational gap is
\[
\Delta_{\mathbb T}(M)
=
\inf_{\substack{\sim,\Psi\\ \operatorname{index}(\sim)\le M}}
\mathcal L(\delta,\Psi\circ Q).
\]

\begin{remark}[Scope of the Quotient Display]
\label{rem:quotient-display-scope}
The display above is the quotient-typed instance of the schema, and it is the
literal resource for commitment, for retention, and for the symbolic
grounding gap of Definition~\ref{def:symbolic-grounding-gap}.  The
\emph{unrestricted linear} grounding relaxation of
Theorem~\ref{thm:spectral-grounding} instead instantiates the schema with a
finite-rank feasible set,
\[
\Delta_{\mathbb T}(M)
=
\inf_{\operatorname{rank}(B)\le M}
\mathcal L(\delta,B),
\]
a genuinely different budget type.  Section~\ref{subsec:meta-scope} records
both instances as coordinates of one typed variational schema without
asserting that the two budgets coincide.
\end{remark}

Each regime attaches its own per-state semantic object:
\begin{itemize}[nosep]
\item \textbf{commitment:} a residual function;
\item \textbf{retention:} a predictive distribution or exponential-family
parameter;
\item \textbf{grounding:} an operator, kernel, or linear realization
coordinate.
\end{itemize}

The shared theorems concern the quotient structure and the index of right
congruences.  The analytical engines---Myhill--Nerode theory, information
bottleneck, Hankel approximation, game-theoretic commitment, online learning,
and operator algebra---are regime-specific.

%----------------------------------------------------------------------
\subsection{Stationary aggregation and type-correct costs}
\label{subsec:stationary-intro}
%----------------------------------------------------------------------

The average-case theory uses stationary aggregation.  For each fixed length
$n$, a stationary process assigns total mass one to the prefixes of length $n$:
\[
\sum_{|u|=n}\mu(u)=1.
\]
Consequently,
\[
\sum_{u\in\mathcal I^*}\mu(u)=\infty,
\]
so the prefix weights do not define a probability measure on the whole free
monoid $\mathcal I^*$.  The theory uses stationary aggregation schemes.

Two canonical average-case aggregations are admitted:
\begin{enumerate}[nosep,label=(\alph*)]
\item a stationary Ces\`aro average;
\item a discounted prefix measure.
\end{enumerate}

The worst-case aggregation remains the supremum over all finite histories:
\[
\mathcal L_\infty(\delta,\widehat\delta)
=
\sup_{u\in\mathcal I^*}
\mathcal E(\delta(u),\widehat\delta(u)).
\]

The cost profunctor $\mathcal E$ is regime-specific.  Commitment uses a
Boolean or $0/1$ equality cost.  Retention uses KL divergence or a quadratic
Fisher surrogate.  Grounding uses operator-norm distance or a related linear
relaxation cost.  The structural Boolean cost $\{0,\infty\}$ is distinguished
from the normalized operational cost $\{0,1\}$.

%----------------------------------------------------------------------
\subsection{Analytical boundaries}
\label{subsec:boundaries-intro}
%----------------------------------------------------------------------

The framework separates shared syntax from regime-specific analysis.

\begin{itemize}[leftmargin=2.0em]
\item \textbf{Shared syntax versus regime-specific semantics.}
The schema provides histories, right congruences, finite quotients, and
variational gaps.  It does not derive the Myhill--Nerode theorem, the
Eckart--Young--Mirsky theorem, or the information-bottleneck identity from
categorical axioms alone.

\item \textbf{Exact theorems versus local theorems.}
The commitment zero-error threshold, the full-KL finite-state
information-bottleneck identity, and the unrestricted linear Hankel spectral
converse are exact within their domains.  The full-KL spectral lower bound is
local and requires explicit regularity assumptions.  The Hankel-structured
finite-rank realization gap satisfies the unrestricted singular-value lower
bound, but equality requires additional structure.

\item \textbf{Converses versus achievability.}
Spectral converses are lower bounds on approximation gaps.  They become
online regret upper rates only when paired with explicit approximation
achievability bounds.

\item \textbf{Linear relaxations versus discrete problems.}
The linear Hankel gap and the linear Price-of-Safety surrogate are exact
linear objects.  They do not automatically solve the corresponding discrete
right-congruence optimization problems.

\item \textbf{Heuristics versus theorems.}
The spectral-tail picture is a heuristic.  The tail shape depends on
aggregation, operator type, affine versus linear approximation, rank versus
state count, compactness, domination, and quadratic-regime assumptions.

\item \textbf{No unconditional cross-regime ordering.}
The framework does not assert a universal inequality among grounding,
retention, and commitment gaps.  The Independence Theorem shows that the
gaps can be varied independently.
\end{itemize}

%----------------------------------------------------------------------
\subsection{Components of the framework}
\label{subsec:integrated-components-intro}
%----------------------------------------------------------------------

The framework comprises the following components.

\paragraph{Schatten conditional converse template.}
A task theory with a rank-controlled response operator $A_\delta$ and a
Schatten-$p$ domination inequality satisfies a unified converse:
\[
\Delta_{\mathbb T}(M)
\ge
\frac1{c_p}
\Phi^{(p)}_{r_{\mathbb T}(M)}(A_\delta),
\]
where $r_{\mathbb T}(M)$ is the effective rank budget and
$\Phi^{(p)}_r$ is the Schatten-$p$ tail.

Grounding is the Schatten-$\infty$ instance for the unrestricted linear
finite-rank relaxation; retention is the Schatten-$1$ instance with effective
rank $M-1$; commitment is the Boolean-rank vertex.  Contextwise $\ell^p$
aggregation does not by itself imply Schatten-$p$ domination.

\paragraph{Adaptive two-axis oracle inequality.}
Within each type-correct regime, offline approximation and online estimation
compose into a bias--variance decomposition:
\[
\Reg_T(\mathcal A_M)
\le
\Aapp^{\mathsf r}(M)
+
\Est_M^{\mathsf r}(T).
\]
Adaptive aggregation over the state budget $M$ yields an oracle inequality
without knowledge of the optimal budget.  For bounded agnostic losses, the
adaptive penalty is
\[
\PsiAgg_M(T)
=
O\!\left(
\sqrt{T\ln(eM)}
+
\ln(eM)
\right).
\]
Realizable $0/1$ and mixable log-loss model selection have logarithmic
overhead.  Spectral tails enter online rates only under explicit spectral-rate
assumptions.

Estimation rates may be supplied by a Littlestone bound, a multiclass
Littlestone bound, a mixability bound, or a linear-regret bound.  The oracle
inequality does not assume a single estimation rate across regimes.

\paragraph{Linear Price-of-Safety surrogate.}
For a quadratic/Fisher retention embedding with state-indexed centered Gram
matrix
\[
G_{st}
=
\sqrt{\pi_s\pi_t}\,
\langle y_s-\bar y,\;y_t-\bar y\rangle,
\]
and a safety partition inducing a pinching conditional expectation
$\EA$, the linear block-local Price-of-Safety surrogate is
\[
\PoSlin(M)
=
\frac12\KyF{G}{M-1}
-
\frac12\KyF{\EA G}{M-1}.
\]
Pinching majorization gives nonnegativity:
\[
\PoSlin(M)\ge0.
\]
The coupling bound
\[
\PoSlin(M)
\le
\frac12\KyF{G-\EA G}{M-1}
\]
is an upper bound, not an equality.  This surrogate is exact for the linear
block-local problem but is not the discrete right-congruence Price of Safety
without relaxation-error control.

\paragraph{Exponent vertex correspondence.}
The manuscript separates three exponent-like layers:
\begin{enumerate}[nosep,label=(\arabic*)]
\item R\'enyi order $\alpha$;
\item Schatten exponent $p$;
\item aggregation or vertex label.
\end{enumerate}

Retention is the $\alpha=1$ KL vertex.  Commitment is recovered from the
$\alpha=0$ support vertex after symmetrization and Boolean valuation.
Grounding is organized by the $\alpha=\infty$ supremum vertex, but the exact
unrestricted linear finite-rank grounding cost is operator-norm distance, not
max-divergence.

The equality of vertex labels
\[
\alpha=p\in\{0,1,\infty\}
\]
is a formal consistency, not a derivation of the Mirsky converses from
R\'enyi identities.

%----------------------------------------------------------------------
\subsection{Roadmap}
\label{subsec:roadmap}
%----------------------------------------------------------------------

Section~\ref{sec:schema} develops the shared finite-state schema.
Section~\ref{sec:specialized} defines specialized task theories.
Section~\ref{sec:meta} states the structural meta-theorems.
Section~\ref{sec:retention} develops retention.
Section~\ref{sec:commitment} develops commitment.
Section~\ref{sec:grounding} develops grounding.
Section~\ref{sec:schatten} develops the conditional Schatten-converse
template.
Section~\ref{sec:oracle} develops the adaptive two-axis oracle inequality and
the associated budget laws.
Section~\ref{sec:pos-linear} develops the linear Price-of-Safety surrogate and
the pinching law.
Section~\ref{sec:exponent} develops the exponent vertex correspondence.
Section~\ref{sec:joint} develops joint, comparative, interaction, and
independence results.
Section~\ref{sec:temporal} develops the temporal axis.
Section~\ref{sec:spectral-heuristic} states the spectral-tail heuristic.
Section~\ref{sec:conditional-rep} states the Conditional Representation
Theorem.
Section~\ref{sec:type-discipline} fixes the type signature that every
approximation statement carries.
Section~\ref{sec:epistemic} records epistemic boundaries.
Section~\ref{sec:conclusion} concludes.

%======================================================================
\section{The Schema of Finite-State Approximation}
\label{sec:schema}
%======================================================================

This section develops the shared finite-state schema.  The schema supplies a
common syntax for all regimes: histories, enriched cost categories, finite
right congruences, quotient machines, stationary aggregation, and profinite
completion.  The analytical meaning of the residual objects is supplied later
by each regime.

%----------------------------------------------------------------------
\subsection{The Enriching Poset}
\label{subsec:enriching-poset}
%----------------------------------------------------------------------

All costs take values in the monoidal poset
\[
\mathbf V=([0,\infty],\ge,+,0).
\]
The unit is $0$, and the order is reversed from the usual real order: a
morphism $v\to v'$ exists iff $v\ge v'$ in the usual order.  Thus $0$ is
terminal and $\infty$ is initial.  A morphism $v\to v'$ records that the cost
has decreased from $v$ to $v'$.

The Boolean sub-poset is
\[
\mathbf V_{\mathbb B}=\{0,\infty\}\subset[0,\infty].
\]
It corresponds to exact equality versus definite mismatch.  A normalized
Boolean mismatch cost with values $\{0,1\}$ is treated separately as an
operational corollary.  The structural zero-error theory uses
$\{0,\infty\}$; the normalized decision-theoretic version uses $\{0,1\}$.

The reversed-order convention is the standard Lawvere metric convention
\cite{lawvere1973}.

%----------------------------------------------------------------------
\subsection{$\mathbf V$-Categories and $\mathbf V$-Profunctors}
\label{subsec:V-cats-profunctors}
%----------------------------------------------------------------------

A \textbf{$\mathbf V$-category} $\mathbf R$ consists of a set of objects
$\operatorname{ob}(\mathbf R)$ and hom-objects
\[
\mathbf R(X,Y)\in\mathbf V
\]
satisfying:
\begin{enumerate}[label=(\roman*),nosep]
\item \textbf{Identity.}
\[
0\ge \mathbf R(X,X),
\]
hence $\mathbf R(X,X)=0$.

\item \textbf{Composition.}
\[
\mathbf R(X,Y)+\mathbf R(Y,Z)\ge \mathbf R(X,Z).
\]
\end{enumerate}

A \textbf{$\mathbf V$-profunctor}
\[
\mathcal E:\mathbf R^{\mathrm{op}}\times\mathbf R\to\mathbf V
\]
assigns costs $\mathcal E(X,Y)$ with left and right actions satisfying
\[
\mathbf R(X',X)+\mathcal E(X,Y)\ge \mathcal E(X',Y),
\]
and
\[
\mathcal E(X,Y)+\mathbf R(Y,Y')\ge \mathcal E(X,Y').
\]

The profunctor $\mathcal E$ is the cost bimodule.  It measures the mismatch
between residual objects $X$ and $Y$, while the left and right actions express
compatibility with the internal geometry of $\mathbf R$.

\begin{remark}[Discrete versus enriched residuals]
\label{rem:discrete-vs-enriched}
For retention and commitment, $\mathbf R$ is usually a discrete
$\mathbf V$-category:
\[
\mathbf R(X,Y)=
\begin{cases}
0,&X=Y,\\
\infty,&X\neq Y.
\end{cases}
\]
Then the action inequalities impose no triangle inequality on $\mathcal E$.
KL divergence is not a metric, and Boolean mismatch is not a metric.

For grounding, $\mathbf R$ may be a metric $\mathbf V$-category of operators,
with
\[
\mathbf R(X,Y)=
\norm{X-Y}_{\mathrm{op}}.
\]
If $\mathcal E$ is the identity profunctor,
\[
\mathcal E(X,Y)=\mathbf R(X,Y),
\]
then the action inequalities reduce to the triangle inequality for the
operator norm.
\end{remark}

%----------------------------------------------------------------------
\subsection{The History $\mathbf V$-Category}
\label{subsec:history-category}
%----------------------------------------------------------------------

Let $\mathcal I$ be a finite input alphabet.  The \textbf{history
$\mathbf V$-category} $\mathbf H$ is the discrete $\mathbf V$-category on the
free monoid $\mathcal I^*$:
\begin{itemize}[nosep]
\item objects are finite words $u\in\mathcal I^*$;
\item hom-objects are
\[
\mathbf H(u,v)=
\begin{cases}
0,&u=v,\\
\infty,&u\neq v.
\end{cases}
\]
\end{itemize}

Because $\mathbf H$ is discrete, any object-level map
\[
\delta:\mathcal I^*\to\operatorname{ob}(\mathbf R)
\]
is automatically a $\mathbf V$-functor.  No nontrivial morphism inequalities
need to be checked.

The prefix structure of histories is not encoded in the hom-objects of
$\mathbf H$.  It is encoded externally in the right-congruence condition,
which ensures that deterministic finite-state approximations evolve
consistently under input extensions.

%----------------------------------------------------------------------
\subsection{Task Theory with Stationary Aggregation}
\label{subsec:task-theory}
%----------------------------------------------------------------------

\begin{definition}[Task Theory]
\label{def:task-theory}
A \textbf{task theory} is a tuple
\[
\mathbb T=(\mathbf H,\mathbf R,\mathcal E,\delta,\mathrm{Agg},\mathcal L)
\]
where:
\begin{enumerate}[nosep,label=(\arabic*)]
\item $\mathbf H$ is the history $\mathbf V$-category.
\item $\mathbf R$ is a $\mathbf V$-category of residual objects.
\item $\mathcal E:\mathbf R^{\mathrm{op}}\times\mathbf R\to\mathbf V$ is a
cost profunctor which is \textbf{reflexive}: $\mathcal E(X,X)=0$ for every
$X\in\mathbf R$.  Reflexivity is not automatic for a $\mathbf V$-profunctor.
In the Lawvere quantale $([0,\infty],\ge,+)$ the constant profunctor
$\mathcal E\equiv1$ on a one-object $\mathbf R$ satisfies every profunctor
axiom, and for it the trajectory cannot be realized at any budget even though
its Nerode index is $1$.  Every cost used in this manuscript --- KL divergence,
operator norm, and the $0/1$ commitment loss --- is reflexive, so the axiom
constrains nothing in the instances but is required for the generic statements
of Meta-Theorem~\ref{meta:monotone}.
\item $\delta:\mathbf H\to\mathbf R$ is the true trajectory.
\item $\mathrm{Agg}$ is an aggregation scheme, one of:
\begin{itemize}[nosep]
\item a stationary Ces\`aro average;
\item a discounted prefix measure; or
\item the worst-case supremum
$\mathcal L_\infty(\delta,\widehat\delta)=\sup_{u\in\mathcal I^*}
\mathcal E(\delta(u),\widehat\delta(u))$.
\end{itemize}
The first two are \emph{average-case} schemes and require the stationary or
discounted datum; the third requires none.  A task theory carrying the
worst-case scheme is called \emph{uniform}.  Results are stated against a
declared scheme, since the Boolean dichotomy in particular takes different
forms under worst-case and average-case aggregation
(Meta-Theorem~\ref{meta:boolean}).
\item $\mathcal L$ is the resulting cost functional.
\end{enumerate}
\end{definition}

\begin{definition}[Stationary Ces\`aro aggregation]
\label{def:cesaro-agg}
Let $U_t$ denote the stationary past at time $t$.  For a cost process
\[
c_t=\mathcal E(\delta(U_t),\widehat\delta(U_t)),
\]
define
\[
\overline{\mathcal L}_1(\delta,\widehat\delta)
=
\lim_{T\to\infty}
\frac1T
\sum_{t=0}^{T-1}
\mathbb E[c_t],
\]
whenever the limit exists.
\end{definition}

For finite-state stationary ergodic processes and finite-state
approximations, the limit exists under standard ergodicity assumptions.  In
non-ergodic settings, one may replace the limit by a limsup or by ergodic
decomposition.

\begin{definition}[Discounted prefix aggregation]
\label{def:discounted-agg}
For $0<\gamma<1$, define a probability measure on $\mathcal I^*$ by
\[
\mu_\gamma(u)
=
(1-\gamma)\gamma^{|u|}
\mathbb P(X_{1:|u|}=u).
\]
Then
\[
\mathcal L_1^\gamma(\delta,\widehat\delta)
=
\mathbb E_{u\sim\mu_\gamma}
\mathcal E(\delta(u),\widehat\delta(u)).
\]
\end{definition}

The discounted prefix measure is a probability measure because
\[
\sum_{u\in\mathcal I^*}\mu_\gamma(u)
=
(1-\gamma)
\sum_{n=0}^\infty
\gamma^n
\sum_{|u|=n}
\mathbb P(X_{1:n}=u)
=
(1-\gamma)
\sum_{n=0}^\infty
\gamma^n
=
1.
\]

The worst-case aggregation remains
\[
\mathcal L_\infty(\delta,\widehat\delta)
=
\sup_{u\in\mathcal I^*}
\mathcal E(\delta(u),\widehat\delta(u)).
\]
It is not an average-case aggregation and should not be conflated with the
stationary Ces\`aro or discounted prefix schemes.

\begin{remark}[Finite Prefixes versus Stationary States]
\label{rem:prefixes-versus-states}
Finite prefixes do not define a probability measure on all of
$\mathcal I^*$.  For stationary retention, the canonical object is the
stationary state process $(S_t)$ and its stationary support $\Splus$.  Finite
histories enter directly in synchronized machines, where $S_t$ is a function
of a finite suffix of the input past.
\end{remark}

%----------------------------------------------------------------------
\subsection{Regime-Typed History Actions}
\label{subsec:history-actions}
%----------------------------------------------------------------------

The finite-state resource shared by the three regimes is the index of a finite
quotient of a history system.  Which history system is regime-dependent: the
regimes share the \emph{form} of a finite-index right-action quotient, not a
single literal free monoid.  Forcing every regime onto $\mathcal I^*$ would
misdescribe predictive machines whose state update reads the realized output.

\begin{definition}[History System]
\label{def:history-system}
A \textbf{history system} is a tuple
\[
\mathsf H=(H,h_\varnothing,\mathcal A,\cdot),
\]
where $H$ is a set of finite histories with root $h_\varnothing$, $\mathcal A$
is an event alphabet, and $\cdot$ is a partial right action
\[
H\times\mathcal A\rightharpoonup H,
\qquad
(h,a)\mapsto h\cdot a,
\]
defined exactly when $a$ is \emph{feasible} from $h$; write
$\mathrm{Feas}(h)\subseteq\mathcal A$ for the set of feasible events.  The
action is \emph{pruned}: whenever $h\cdot a$ is defined it lies in $H$.  If
$\mathrm{Feas}(h)=\mathcal A$ for every $h$, the action is \emph{total}.
\end{definition}

\begin{definition}[Support-Relative Right Congruence]
\label{def:support-right-cong}
Let $\mathsf H=(H,h_\varnothing,\mathcal A,\cdot)$ be a history system.  A
\textbf{support-relative right congruence} on $\mathsf H$ is an equivalence
relation $\sim$ on $H$ such that $u\sim v$ implies
\begin{enumerate}[label=(\roman*),nosep]
\item $\mathrm{Feas}(u)=\mathrm{Feas}(v)$; and
\item $u\cdot a\sim v\cdot a$ for every $a\in\mathrm{Feas}(u)$.
\end{enumerate}
Its \textbf{index} is the number of equivalence classes.  The quotient
$H/\!\sim$ carries a well-defined partial transition $[u]\cdot a=[u\cdot a]$,
with $\mathrm{Feas}([u])=\mathrm{Feas}(u)$; clause~(i) is exactly what makes
the feasible-event set of a class well defined.  For a total action clause~(i)
is vacuous and the definition reduces to Definition~\ref{def:right-cong}.
\end{definition}

\begin{remark}[Clause~(i) Is Not Redundant]
\label{rem:support-clause-i-needed}
Clause~(ii) alone, read as a condition on commonly feasible events, does not
give a quotient history system.  Take $\mathcal A=\{a,b\}$ and two histories
$u,v$ with $\mathrm{Feas}(u)=\{a\}$ and $\mathrm{Feas}(v)=\{b\}$.  The
relation identifying $u$ with $v$ satisfies clause~(ii) vacuously, there being
no commonly feasible event, yet the class $[u]$ has no well-defined feasible
set and the quotient carries no transition structure.  This is the mechanism
behind Remark~\ref{rem:unifilar-support-not-automatic} below.
\end{remark}

\begin{remark}[The Regimes as History Systems]
\label{rem:regime-history-instances}
The principal instantiations are the following.

\begin{center}
\begin{tabular}{@{}llp{3.1cm}l@{}}
\toprule
regime & $H$ & $\mathcal A$ & action \\
\midrule
commitment & $\mathcal I^*$ & $\mathcal I$ & total \\
retention, unifilar & feasible joint histories
& $\mathcal I\times\mathcal O$ & pruned \\
retention, input-driven & $\mathcal I^*$ & $\mathcal I$ & total \\
grounding & $\Sigma^*$, $\Sigma=\mathcal I\times\mathcal O$ & $\Sigma$ & total \\
\bottomrule
\end{tabular}
\end{center}

For commitment and grounding the action is total and
Definition~\ref{def:right-cong} applies verbatim.  For unifilar retention the
feasible joint histories form a pruned tree under the stationary law, and
compatibility with the action is understood in the support-relative sense.  The
input-driven model of Definition~\ref{def:controlled-markov} is the case in
which the state process factors through the input projection, and the
input-history syntax is recovered.  The history $\mathbf V$-category of
Section~\ref{subsec:history-category} is built on whichever $H$ the regime
supplies; being discrete, it makes every object-level trajectory a
$\mathbf V$-functor regardless of which system is used.
\end{remark}

\begin{lemma}[Cofilteredness for Support-Relative Congruences]
\label{lem:cofiltered-support}
Let $\sim_1,\sim_2$ be finite-index support-relative right congruences on a
history system $\mathsf H$.  Then $\sim_1\cap\sim_2$ is one as well, refines
both, and has index at most
$\operatorname{index}(\sim_1)\operatorname{index}(\sim_2)$.  The finite
quotients of $\mathsf H$ therefore form a cofiltered system.
\end{lemma}

\begin{proof}
The intersection of equivalence relations is an equivalence relation.  If
$u\,(\sim_1\cap\sim_2)\,v$ then $\mathrm{Feas}(u)=\mathrm{Feas}(v)$ already
from $\sim_1$, giving clause~(i); and for $a$ in that common feasible set,
$u\cdot a\sim_iv\cdot a$ for $i=1,2$, giving clause~(ii).  The map
$[u]_{\sim_1\cap\sim_2}\mapsto([u]_{\sim_1},[u]_{\sim_2})$ is well defined and
injective, which is the index bound.  The argument is that of
Lemma~\ref{lem:cofiltered}, with the feasible-set clause carried along.
\end{proof}

%----------------------------------------------------------------------
\subsection{Right Congruences and Quotients}
\label{subsec:right-cong}
%----------------------------------------------------------------------

\begin{definition}[Right Congruence]
\label{def:right-cong}
A \textbf{right congruence} $\sim$ on $\mathbf H$ is an equivalence relation
on $\mathcal I^*$ such that
\[
u\sim v
\implies
uw\sim vw
\quad
\forall w\in\mathcal I^*.
\]
Equivalently, it suffices to require
\[
u\sim v
\implies
ux\sim vx
\quad
\forall x\in\mathcal I.
\]
The \textbf{index} of $\sim$ is the number of equivalence classes.
\end{definition}

This is the total-action instance of Definition~\ref{def:support-right-cong};
pruned retention history systems use the support-relative form.

The quotient $\mathbf Q=\mathbf H/\!\sim$ has objects $[u]$ and hom-objects
\[
\mathbf Q([u],[v])=
\begin{cases}
0,&[u]=[v],\\
\infty,&\text{otherwise}.
\end{cases}
\]
The quotient map
\[
Q:\mathbf H\to\mathbf Q,
\qquad
u\mapsto[u],
\]
is a $\mathbf V$-functor.  The right-congruence condition ensures that the
quotient carries a well-defined deterministic transition
\[
[u]\cdot x=[ux].
\]

\begin{definition}[$M$-State Approximation]
\label{def:M-state-approx}
An \textbf{$M$-state approximation} of a task theory $\mathbb T$ is a pair
$(\sim,\Psi)$ where $\sim$ is a right congruence of index at most $M$ and
\[
\Psi:\mathbf Q\to\mathbf R
\]
is a $\mathbf V$-functor.  The approximate trajectory is
\[
\Psi\circ Q:\mathbf H\to\mathbf R.
\]
\end{definition}

\begin{definition}[$M$-State Gap]
\label{def:M-state-gap}
The \textbf{$M$-state gap} is
\[
\Delta_{\mathbb T}(M)
=
\inf_{\substack{\sim,\Psi\\ \operatorname{index}(\sim)\le M}}
\mathcal L(\delta,\Psi\circ Q).
\]
\end{definition}

\begin{remark}[Machine type depends on $\mathbf R$]
\label{rem:realization-types}
The composite $\Psi\circ Q$ has deterministic state-transition structure, but
its per-state content depends on $\mathbf R$:
\begin{itemize}[nosep]
\item \textbf{Commitment:} each state stores a residual function; the
approximant is a deterministic Mealy BSG.
\item \textbf{Retention:} each state stores a predictive distribution or
kernel; the approximant is a unifilar lumped controlled predictor
(Definition~\ref{def:unifilar-lumpable}), and in the input-driven synchronized
subclass a deterministic-transition, stochastic-emission machine, in the sense
of Definition~\ref{def:unifilar-machine} and
Remark~\ref{rem:unifilar-proper-subclass}.
\item \textbf{Grounding spectral converse:} each state stores an operator or
linear coordinate; the approximant is a linear finite-rank realization.
\end{itemize}
The shared theorems concern $\operatorname{index}(\sim)$.  The
regime-specific theorems concern the semantic objects attached to the states.
\end{remark}

\begin{remark}[Stationary Retention Quotients]
\label{rem:schema-stationary-retention}
For stationary controlled unifilar machines, the primary feasible objects are
unifilar-lumpable quotients
\[
\phi:\Splus\to\mathcal K
\]
(Definition~\ref{def:unifilar-lumpable}).  In a machine with joint-history
synchronization depth $L$, each such quotient with block-uniform support
induces a support-relative right congruence on feasible joint histories of
length at least $L$ by
\[
u\sim_\phi v
\iff
\phi(\sigma(u))=\phi(\sigma(v))
\]
(Proposition~\ref{prop:unifilar-lumpability}); in the input-driven
synchronized subclass this projects to the input-history right congruence of
Proposition~\ref{prop:lumpability}.  Thus the finite-state budget remains an
index resource, now on regime-typed histories.  For non-synchronized
stationary machines, finite-index right congruences on finite histories should
be replaced by finite-state factors of the stationary state process.
\end{remark}

%----------------------------------------------------------------------
\subsection{Profinite Completion}
\label{subsec:profinite-completion}
%----------------------------------------------------------------------

The finite quotients of $\mathbf H$ form a cofiltered system; the
support-relative analogue, for pruned history systems, is
Lemma~\ref{lem:cofiltered-support}.  Let $\mathcal C$ be the cofiltered poset
of finite-index right congruences on $\mathcal I^*$, ordered by refinement:
\[
\sim_1\le\sim_2
\quad\Longleftrightarrow\quad
\sim_1\text{ refines }\sim_2.
\]

\begin{lemma}[$\mathcal C$ Is Cofiltered]
\label{lem:cofiltered}
If $\sim_1$ and $\sim_2$ are finite-index right congruences on $\mathcal I^*$,
then so is $\sim_1\cap\sim_2$, and it refines both.
\end{lemma}

\begin{proof}
The intersection is an equivalence relation.  It is a right congruence: if
$u\sim_1v$ and $u\sim_2v$, then for every $w\in\mathcal I^*$ we have
$uw\sim_1vw$ and $uw\sim_2vw$, hence
$uw\,(\sim_1\cap\sim_2)\,vw$.  The map
\[
[u]_{\sim_1\cap\sim_2}
\mapsto
\bigl([u]_{\sim_1},\,[u]_{\sim_2}\bigr)
\]
is well defined and injective, so
\[
\operatorname{index}(\sim_1\cap\sim_2)
\le
\operatorname{index}(\sim_1)\,
\operatorname{index}(\sim_2)
<
\infty.
\]
Thus $\sim_1\cap\sim_2\in\mathcal C$ refines both, so $\mathcal C$ is
cofiltered.
\end{proof}

\begin{definition}[Profinite completion]
\label{def:profinite-completion}
The \textbf{profinite completion} of $\mathbf H$ is
\[
\widehat{\mathbf H}
=
\varprojlim_{\sim\in\mathcal C}
\mathbf H/\!\sim.
\]
\end{definition}

An object of $\widehat{\mathbf H}$ is a compatible family
\[
(q_\sim)_\sim,
\qquad
q_\sim\in\mathbf H/\!\sim,
\]
such that whenever $\sim_1$ refines $\sim_2$, the projection of
$q_{\sim_1}$ to $\mathbf H/\!\sim_2$ equals $q_{\sim_2}$.

There is a canonical dense embedding
\[
\iota:\mathbf H\hookrightarrow\widehat{\mathbf H},
\qquad
u\mapsto([u]_\sim)_\sim.
\]
Injectivity is not formal: it holds because the free monoid $\mathcal I^*$ is
\emph{residually finite} with respect to finite-index right congruences.
Given distinct $u\ne v$, set $N>\max\{|u|,|v|\}$ and let $\sim_N$ identify two
words when they are equal, or when both have length at least $N$.  This is a
right congruence: if $w\sim_Nw'$ and $x\in\mathcal I$ then either $w=w'$, whence
$wx=w'x$, or both $|w|,|w'|\ge N$, whence both $|wx|,|w'x|\ge N$.  It has
finite index $|\mathcal I^{<N}|+1$, and separates $u$ from $v$.  Hence the
components of $\iota(u)$ and $\iota(v)$ differ at $\sim_N$.

The same separation argument applies to pruned history systems, once
feasibility is carried along.

\begin{proposition}[Residual Finiteness of Synchronized Unifilar History Trees]
\label{prop:unifilar-residually-finite}
Let a unifilar machine (Definition~\ref{def:unifilar-machine}) have
joint-history synchronization, so that its feasible histories of length at
least $L$ form a pruned history system $F$ carrying a state readout
$\sigma:F\to\Splus$ with $\sigma(u\cdot a)=\tau(\sigma(u),a)$ on feasible
events.  Then $F$ is residually finite with respect to finite-index
support-relative right congruences: the canonical map from $F$ into the inverse
limit of its finite quotients is injective.
\end{proposition}

\begin{proof}
Let $u\neq v$ in $F$.  Choose $N>\max\{|u|,|v|\}$ and define $\sim_N$ by
declaring $w\sim_Nw'$ when either $w=w'$, or $|w|,|w'|\ge N$ and
$\sigma(w)=\sigma(w')$.

\emph{Finite index.}  There are at most
$\sum_{n<N}|\mathcal I\times\mathcal O|^{\,n}$ singleton classes of length
below $N$ and at most $|\Splus|$ state-labelled classes at length $N$ or
above.

\emph{Clause~(i).}  Feasibility of $(x,y)$ from a history $w$ means
$P_{\sigma(w)}^x(\{y\})>0$, so it is a function of $\sigma(w)$ alone.  Two
histories identified by $\sim_N$ are either equal or have equal $\sigma$, and
in both cases have the same feasible events.

\emph{Clause~(ii).}  In the identity case $w\cdot a=w'\cdot a$.  Otherwise
$|w|,|w'|\ge N$ and $\sigma(w)=\sigma(w')$, so for a feasible $(x,y)$,
\[
\sigma\bigl(w\cdot(x,y)\bigr)
=
\tau(\sigma(w),x,y)
=
\tau(\sigma(w'),x,y)
=
\sigma\bigl(w'\cdot(x,y)\bigr),
\]
and both extensions have length at least $N+1$, hence
$w\cdot(x,y)\sim_Nw'\cdot(x,y)$.

Since $N$ exceeds $|u|$ and $|v|$, both are singleton classes of $\sim_N$ and
are therefore separated, so the corresponding components of their images in
the inverse limit differ.
\end{proof}

For total-action regimes the completion is that of the free monoid; for
synchronized unifilar retention it is the completion of the feasible
joint-history tree, residually finite by
Proposition~\ref{prop:unifilar-residually-finite}.

The \textbf{profinite uniformity} on $\mathbf H$ is the initial uniformity
with respect to the projections
\[
\mathbf H\to\mathbf H/\!\sim.
\]
Equivalently, its basic entourages are
\[
U_\sim
=
\{(u,v)\in\mathcal I^*\times\mathcal I^*:u\sim v\},
\]
where $\sim$ ranges over finite-index right congruences.

\begin{remark}[Cofiltered limits and Lawvere's framework]
\label{rem:lawvere-limit}
The enriched poset
\[
\mathbf V=([0,\infty],\ge,+,0)
\]
is complete symmetric monoidal closed in the sense of Lawvere
\cite{lawvere1973}.  The relevant enriched cofiltered limits therefore exist.
Because $\mathbf V$ is posetal, the enriched inverse limit coincides with the
ordinary compatible-family limit on objects, equipped with the induced
profinite uniformity.
\end{remark}

\begin{remark}[Profinite uniformity for discrete residual categories]
\label{rem:profinite-discrete}
Let $\mathbf R$ be a discrete $\mathbf V$-category.  Uniform continuity of a
trajectory
\[
\delta:\mathbf H\to\mathbf R
\]
in the profinite uniformity is equivalent to $\delta$ factoring through some
finite quotient:
\[
\mathbf H
\longrightarrow
\mathbf H/\!\sim
\longrightarrow
\mathbf R
\]
for some finite-index right congruence $\sim$.  Indeed, in a discrete
$\mathbf V$-category the only entourages are those containing the diagonal, so
uniform continuity requires a finite-index congruence $\sim$ with
$u\sim v\Rightarrow\delta(u)=\delta(v)$, which is exactly factorization through
$\mathbf H/\!\sim$.
\end{remark}

\begin{remark}[Compactness and convergence]
\label{rem:profinite-compact}
Each finite quotient $\mathbf H/\!\sim$ is a finite discrete space.  The
inverse limit $\widehat{\mathbf H}$ is compact in the profinite topology.
This compactness underlies the profinite convergence meta-theorem: under
compact-uniform hypotheses on $\mathbf R$ and $\mathcal E$, uniform continuity
of $\delta$ is equivalent to continuous extension to $\widehat{\mathbf H}$,
and hence to approximability by finite-index right congruences with vanishing
gap.
\end{remark}

The profinite completion is not an analytic Hilbert-space completion such as
an $L^2$ completion.  It is a logical or Stone-type completion generated by
finite deterministic quotients.

%======================================================================
\section{Specialized Task Theories}
\label{sec:specialized}
%======================================================================

The generic schema becomes operative after the residual category, cost
profunctor, and aggregation scheme are specialized.  This section defines the
principal specialized task theories used later: separated task theories,
compact uniform task theories, Hilbert-module task theories, stationary
controlled causal machines, full-KL retention, and Gaussian quadratic
restricted retention.

%----------------------------------------------------------------------
\subsection{Separated Task Theories}
\label{subsec:separated}
%----------------------------------------------------------------------

\begin{definition}[Separated Task Theory]
\label{def:separated-task}
A task theory $\mathbb T$ is \textbf{separated} if:
\[
\mathbf R(X,Y)=\mathbf R(Y,X)=0
\implies
X=Y,
\]
and
\[
\mathcal E(X,Y)=0
\implies
X=Y.
\]
\end{definition}

\begin{definition}[Nerode Equivalence]
\label{def:nerode}
For any task theory $\mathbb T$, the \textbf{Nerode equivalence}
$\sim_\delta$ is
\[
u\sim_\delta v
\iff
\forall w\in\mathcal I^*,
\quad
\mathbf R(\delta(uw),\delta(vw))
=
\mathbf R(\delta(vw),\delta(uw))
=
0.
\]
\end{definition}

The two-sided condition is used because $\mathbf R$ need not be symmetric.  If
$\mathbf R$ is symmetric, the condition reduces to
\[
\mathbf R(\delta(uw),\delta(vw))=0
\quad
\forall w.
\]

The Nerode equivalence is always a right congruence.  Indeed, if
$u\sim_\delta v$, then for every $x\in\mathcal I$ and every
$w\in\mathcal I^*$,
\[
\delta(uxw)=\delta(u(xw)),
\qquad
\delta(vxw)=\delta(v(xw)),
\]
so the defining zero-distance condition for $u$ and $v$ applied to the
continuation $xw$ gives $ux\sim_\delta vx$.

In the commitment regime, $\sim_\delta$ is the usual Myhill--Nerode residual
equivalence.  In the retention regime, when predictive distributions are
distinct on the stationary support, it coincides with causal-state
equivalence.  In grounding regimes, it may correspond to equality of residual
kernels or impulse-response functionals, depending on the chosen residual
category.

%----------------------------------------------------------------------
\subsection{Compact Uniform Task Theories}
\label{subsec:compact-uniform}
%----------------------------------------------------------------------

\begin{definition}[Compact Metrizable $\mathbf V$-Category]
\label{def:compact-metrizable}
A $\mathbf V$-category $\mathbf R$ is \textbf{compact metrizable} if
$\operatorname{ob}(\mathbf R)$ carries a compact metric topology and the
hom-object map
\[
\mathbf R(-,-):
\operatorname{ob}(\mathbf R)\times\operatorname{ob}(\mathbf R)
\to
[0,\infty]
\]
is continuous in the metric topology, with the usual extended-real topology on
$[0,\infty]$.
\end{definition}

More generally, one may use a compact uniform space in place of a compact
metric space.  The profinite convergence meta-theorem is stated under the
hypothesis that $\mathbf R$ is compact metrizable or compact uniform and that
the cost profunctor $\mathcal E$ is uniformly continuous and separated.

\begin{remark}[Separate Continuity Is Insufficient]
\label{rem:separate-continuity-insufficient}
Separate continuity alone is insufficient for the profinite equivalence.  The
meta-theorem uses joint uniform continuity, or compact-uniform continuity, of
the cost profunctor.
\end{remark}

%----------------------------------------------------------------------
\subsection{Hilbert-Module Task Theories and Linear Hankel Gaps}
\label{subsec:hilbert-module}
%----------------------------------------------------------------------

\begin{definition}[Hilbert-Module Task Theory]
\label{def:hilbert-module}
A \textbf{Hilbert-module task theory} is a separated task theory on a finite
joint alphabet
\[
\Sigma=\mathcal I\times\mathcal O
\]
equipped with the following additional data:
\begin{enumerate}[label=(\arabic*)]
\item a Banach space $B$ of residual observables;
\item a bounded linear functional
\[
\phi:B\to\mathbb C;
\]
\item an impulse-response function
\[
h:\Sigma^*\to\mathbb C
\]
obtained from the true trajectory by
\[
h(u)=\phi(\delta(u));
\]
\item a Hankel operator
\[
H_\delta:\ell^2(\Sigma^*)\to\ell^2(\Sigma^*)
\]
specified on the standard basis by
\[
(H_\delta e_u)(v)=h(uv),
\]
\emph{provided} this densely defined formula extends to a bounded operator on
$\ell^2(\Sigma^*)$.  Boundedness is not automatic: for general $h$ the vector
$v\mapsto h(uv)$ need not be square summable, and even when it is, the induced
matrix need not be bounded.  A sufficient condition is the geometric decay
hypothesis of Proposition~\ref{prop:stable-decay}, which yields boundedness by
the Schur test;
\item compactness of the resulting bounded operator $H_\delta$.
\end{enumerate}
A Hilbert-module task theory is said to be \emph{spectrally admissible} when
the Hankel operator exists as a bounded operator and is compact, i.e.\ when
conditions~(iv) and~(v) hold.  All singular-value statements below are
asserted only for spectrally admissible task theories.
\end{definition}

If the residual objects are operators, $B$ may be an operator space or a
$C^*$-algebra, and $\phi$ may be a bounded functional selecting the scalar
observable relevant to the impulse response.  The scalar impulse response $h$
is the object seen by the Hankel operator.  Distinct residual objects may
induce the same impulse response; the Hankel theory below depends only on $h$
unless additional domination hypotheses are supplied.

\begin{remark}[Residual Norm versus Hankel Norm]
\label{rem:hilbert-residual-hankel}
The norm used for residual-operator comparisons is the norm of $B$, while the
Hankel gap uses the $\ell^2$ operator norm of $H_\delta$.  No domination
between the $B$-norm and the Hankel operator norm is assumed.
\end{remark}

The model does not assume an unqualified domination hypothesis of the form
\[
\norm{H_\delta-H_{\widehat\delta}}_{\mathrm{op}}
\le
\mathcal L_\infty(\delta,\widehat\delta).
\]
Such an inequality may hold under explicit decay or summability conditions,
but it is not a general axiom.

\begin{definition}[Unrestricted Linear Hankel Relaxation Gap]
\label{def:linear-hankel-gap}
The \textbf{unrestricted linear Hankel $M$-rank relaxation gap} is
\[
\Delta_{\mathbb T}^{\mathrm{unres}}(M)
=
\inf_{\operatorname{rank}\widehat H\le M}
\norm{H_\delta-\widehat H}_{\mathrm{op}}.
\]
The infimum is over bounded operators $\widehat H$ of rank at most $M$, not
necessarily Hankel operators.
\end{definition}

This is an operator-approximation gap.  It is the object to which
Eckart--Young--Mirsky \cite{eckart1936,mirsky1960} applies directly.

\begin{definition}[Hankel-Structured Finite-Rank Gap]
\label{def:hankel-structured-gap}
The \textbf{Hankel-structured finite-rank gap} is
\[
\Delta_{\mathbb T}^{\mathrm{Hank,str}}(M)
=
\inf_{\substack{
\operatorname{rank}\widehat H\le M\\
\widehat H\ \text{Hankel}
}}
\norm{H_\delta-\widehat H}_{\mathrm{op}}.
\]
The infimum is over Hankel operators of rank at most $M$.
\end{definition}

Because the Hankel-structured feasible set is a subset of the unrestricted
feasible set,
\[
\Delta_{\mathbb T}^{\mathrm{Hank,str}}(M)
\ge
\Delta_{\mathbb T}^{\mathrm{unres}}(M).
\]
Equality requires additional structure.

\begin{definition}[Linear Finite-Rank Residual-Cost Gap]
\label{def:lin-finite-rank-gap}
The \textbf{linear finite-rank residual-cost gap} is
\[
\Delta_{\mathbb T}^{\mathrm{lin}}(M)
=
\inf
\left\{
\mathcal L_\infty(\delta,\widehat\delta):
\operatorname{rank}H_{\widehat\delta}\le M
\right\}.
\]
\end{definition}

The residual-cost gap $\Delta_{\mathbb T}^{\mathrm{lin}}(M)$ and the
unrestricted Hankel relaxation gap
$\Delta_{\mathbb T}^{\mathrm{unres}}(M)$ are distinct objects.  The former
measures task loss of finite-rank residual approximants; the latter measures
operator-norm approximation to $H_\delta$.  They are related only under
explicit domination or achievability hypotheses.

\begin{definition}[Weighted Residual Norm]
\label{def:weighted-residual-norm}
For $0<\rho<1$, define
\[
\norm{d}_{\infty,\rho}
=
\sup_{w\in\Sigma^*}
\rho^{-|w|}
|d(w)|.
\]
\end{definition}

\begin{proposition}[Stable-Decay Hankel Domination]
\label{prop:stable-decay}
Suppose
\[
|h_\delta(w)-h_{\widehat\delta}(w)|
\le
\varepsilon\rho^{|w|}
\]
with
\[
0<\rho<\frac1{|\Sigma|}.
\]
Then
\[
\norm{H_\delta-H_{\widehat\delta}}_{\mathrm{op}}
\le
\frac{1}{1-|\Sigma|\rho}
\varepsilon.
\]
Equivalently, if
\[
\norm{h_\delta-h_{\widehat\delta}}_{\infty,\rho}
\le
\varepsilon,
\]
then
\[
\norm{H_\delta-H_{\widehat\delta}}_{\mathrm{op}}
\le
\frac{1}{1-|\Sigma|\rho}
\norm{h_\delta-h_{\widehat\delta}}_{\infty,\rho}.
\]
\end{proposition}

\begin{proof}
Let
\[
A=H_\delta-H_{\widehat\delta}.
\]
The Hankel matrix entries satisfy
\[
|A(u,v)|
=
|h_\delta(uv)-h_{\widehat\delta}(uv)|
\le
\varepsilon\rho^{|u|+|v|}.
\]
By Schur's test,
\[
\sup_u\sum_v |A(u,v)|
\le
\varepsilon
\sup_u
\rho^{|u|}
\sum_{v\in\Sigma^*}
\rho^{|v|}.
\]
Since $\rho^{|u|}\le1$ and the number of words of length $n$ is
$|\Sigma|^n$,
\[
\sum_{v\in\Sigma^*}\rho^{|v|}
=
\sum_{n=0}^\infty
|\Sigma|^n\rho^n
=
\frac{1}{1-|\Sigma|\rho}.
\]
Thus
\[
\sup_u\sum_v |A(u,v)|
\le
\frac{\varepsilon}{1-|\Sigma|\rho}.
\]
The same bound holds for columns.  Schur's test therefore gives
\[
\norm{A}_{\mathrm{op}}
\le
\frac{\varepsilon}{1-|\Sigma|\rho}.
\]
\end{proof}

\begin{remark}[AAK versus EYM]
\label{rem:aak-eym-hilbert}
The default spectral converse for
$\Delta_{\mathbb T}^{\mathrm{unres}}(M)$ is Eckart--Young--Mirsky:
\[
\Delta_{\mathbb T}^{\mathrm{unres}}(M)
=
\sigma_{M+1}(H_\delta).
\]
Adamjan--Arov--Krein applies only after a Hardy-space embedding has been
separately established.  The model does not assume such an embedding by
default.
\end{remark}

\begin{remark}[Infinite support and boundedness]
\label{rem:infinite-support}
Infinite support of an impulse response does not by itself imply that the
associated Hankel operator is unbounded.  Boundedness requires decay or
summability conditions.  For example, scalar impulse responses satisfying
$|h(u)|\le C\rho^{|u|}$ with $0<\rho<1/|\Sigma|$ define bounded Hankel-type
operators by the Schur-test argument above.  Bounded symbolic impulse
responses that fail to decay may or may not define bounded Hankel operators;
boundedness is checked case by case.
\end{remark}

%----------------------------------------------------------------------
\subsection{Stationary Controlled Causal Machines}
\label{subsec:controlled-markov}
%----------------------------------------------------------------------

\begin{definition}[Stationary Controlled Causal Machine]
\label{def:controlled-markov}
A \textbf{stationary controlled causal machine} consists of:
\begin{enumerate}[label=(\arabic*)]
\item a finite input alphabet $\mathcal I$;
\item a standard Borel output space $\mathcal O$;
\item a finite state set $\mathcal S$;
\item a deterministic transition map
\[
\tau:\mathcal S\times\mathcal I\to\mathcal S;
\]
\item emission laws
\[
P_s\in\mathcal P(\mathcal O)
\]
for each $s\in\mathcal S$;
\item a stationary ergodic input process $(X_t)_{t\in\mathbb Z}$;
\item a stationary state process $(S_t)_{t\in\mathbb Z}$ satisfying
\[
S_t=\tau(S_{t-1},X_t)
\quad\text{a.s.};
\]
\item a stationary distribution
\[
\pi_s=\mathbb P(S_t=s);
\]
\item a predictive future variable $Z_t$, by default the next output $Y_t$,
such that $S_t$ is sufficient:
\[
Z_t\perp\!\!\!\perp
\{S_u,X_u,Y_u:u<t\}
\mid S_t.
\]
\end{enumerate}
The predictive distribution of state $s$ is
\[
P_s=\mathcal L(Z_t\mid S_t=s).
\]
The stationary support is
\[
\Splus=\{s\in\mathcal S:\pi_s>0\}.
\]
\end{definition}

\begin{definition}[Stationary Controlled Unifilar Causal Machine]
\label{def:unifilar-machine}
A \textbf{stationary controlled unifilar causal machine} consists of a finite
input alphabet $\mathcal I$, a standard Borel output space $\mathcal O$, a
finite state set $\mathcal S$, emission kernels
\[
P_s^x=\mathcal L\bigl(Y_t\mid S_t=s,\ X_t=x\bigr),
\]
a measurable \emph{unifilar} update
\[
\tau:\mathcal S\times\mathcal I\times\mathcal O\to\mathcal S,
\qquad
S_{t+1}=\tau(S_t,X_t,Y_t)\quad\text{a.s.},
\]
defined on jointly feasible triples, together with a stationary ergodic input
process, a stationary distribution $\pi$, and a declared predictive target
$Z_t$.

An update with $\tau(s,x,y)=\tau_0(s,x)$ for every feasible $y$ is called
\emph{output-independent}; in that case the recursion collapses to
$S_{t+1}=\tau_0(S_t,X_t)$ and, after an index shift, the machine is exactly a
stationary controlled causal machine in the sense of
Definition~\ref{def:controlled-markov}.
\end{definition}

\begin{remark}[The Input-Driven Model Is a Proper Subclass]
\label{rem:unifilar-proper-subclass}
The containment is strict, and the witness is standard.  Take
$|\mathcal I|=1$, $\mathcal O=\{0,1\}$, two states $A,B$, with $A$ emitting
$0$ or $1$ with equal probability and $B$ emitting $1$ almost surely, and
\[
\tau(A,\cdot,0)=A,\qquad \tau(A,\cdot,1)=B,\qquad \tau(B,\cdot,1)=A .
\]
Since the input alphabet is a singleton, the input history of length $t$ is
the unique word of that length, yet the state after two steps is $A$ on the
output history $00$ and $B$ on $01$.  The state is therefore not a function of
the input history, and no map $\tau_0:\mathcal S\times\mathcal I\to\mathcal S$
reproduces this machine.

Consequently Definition~\ref{def:controlled-markov} describes an
\emph{input-driven} transducer, and the descriptions of it as a lumped
$\epsilon$-machine refer to this output-independent subclass rather than to
general controlled $\epsilon$-machines, whose causal states are functions of
joint input--output histories
(Remark~\ref{rem:epsilon-machine-relation}).  The retention theory is developed
for the full unifilar class in
Section~\ref{subsec:unifilar-retention}, with the input-driven model as the
output-independent specialization; the complexity constructions of
Section~\ref{subsec:retention-complexity} are stated for the input-driven model
and transfer to the unifilar class by
Remark~\ref{rem:complexity-transfer}.
\end{remark}

\begin{definition}[$Z$-predictive equivalence]
\label{def:z-predictive-equivalence}
Fix the predictive target $Z_t$.  On the stationary support define
\[
s\sim_Z s'
\iff
\mathcal L(Z_t\mid S_t=s)
=
\mathcal L(Z_t\mid S_t=s').
\]
If $Z_t$ is the entire controlled future, $\sim_Z$ is causal-state
equivalence.  If $Z_t=Y_t$ is only the next output, $\sim_Z$ is one-step
predictive equivalence, which is strictly coarser in general and need not
identify causal states.

Throughout, ``causal equivalence'' means $\sim_Z$ for the declared target
$Z$, and coincides with causal-state equivalence exactly when $Z$ is the full
future.  All retention statements, including the zero-retention threshold,
are to be read with respect to distinct $Z$-predictive states.

For unifilar machines with input-dependent emissions
(Definition~\ref{def:unifilar-machine}) the predictive object attached to a
state is the whole kernel $x\mapsto P_s^x$, and $\sim_Z$ compares kernels on
the input support: $s\sim_Zs'$ if and only if $P_s^x=P_{s'}^x$ for every input
$x$ of positive mass.
\end{definition}

Finite output alphabets are the special case in which all $P_s$ are supported
on a finite set.  Gaussian quadratic retention is obtained by taking
$\mathcal O=\mathbb R^D$ and $P_s$ Gaussian.

\begin{definition}[Synchronized Finite-History Realization]
\label{def:synchronized-realization}
A stationary controlled causal machine has \textbf{synchronization depth}
$L<\infty$ if there exists a surjection
\[
\sigma_L:\mathcal I^L\to\Splus
\]
such that, almost surely,
\[
S_t=\sigma_L(X_{t-L+1},\dots,X_t).
\]
For histories $u$ with $|u|\ge L$, define
\[
\sigma(u)=\sigma_L(\operatorname{suffix}_L(u)).
\]
Then for all sufficiently long $u$,
\[
\sigma(ux)=\tau(\sigma(u),x).
\]
The finite-history trajectory is
\[
\delta(u)=P_{\sigma(u)}.
\]
\end{definition}

Thus the original finite-history/right-congruence syntax is recovered for
definite and synchronizing machines, including all NP-hardness constructions.

\begin{definition}[Causal Equivalence]
\label{def:causal-equivalence}
On the stationary support, causal equivalence is
\[
s\sim_\delta t
\iff
P_s=P_t.
\]
In a synchronized machine, the corresponding history equivalence is
\[
u\sim_\delta v
\iff
\sigma(u)=\sigma(v).
\]
If the predictive distributions $P_s$ are pairwise distinct on $\Splus$, then
causal equivalence coincides with state identity on $\Splus$.
\end{definition}

\begin{definition}[Lumpable Quotient]
\label{def:lumpable-quotient}
A map $\phi:\Splus\to\mathcal K$ is \textbf{lumpable} if there exists a
deterministic quotient transition
\[
\tau_{\mathcal K}:\mathcal K\times\mathcal I\to\mathcal K
\]
such that
\[
\phi(\tau(s,x))
=
\tau_{\mathcal K}(\phi(s),x)
\]
for all $s\in\Splus$ and all inputs $x$ in the support of the input process.
For worst-case aggregation, require this for all $x\in\mathcal I$.
The quotient process is
\[
K_t=\phi(S_t).
\]
\end{definition}

\begin{definition}[Unifilar Lumpability]
\label{def:unifilar-lumpable}
A quotient $\phi:\mathcal S\to\mathcal K$ of a unifilar machine is
\textbf{unifilar-lumpable} if there is
$\tau_{\mathcal K}:\mathcal K\times\mathcal I\times\mathcal O\to\mathcal K$
with
\[
\phi\bigl(\tau(s,x,y)\bigr)=\tau_{\mathcal K}\bigl(\phi(s),x,y\bigr)
\]
for every \emph{jointly feasible} triple $(s,x,y)$, that is, every triple with
$P_s^x(\{y\})>0$ in the discrete case and, in general, every triple in the
support of the joint law.  The singleton quotient is always
unifilar-lumpable.

The quotient has \textbf{block-uniform support} if
\[
\phi(s)=\phi(s')
\implies
\supp P_s^x=\supp P_{s'}^x
\qquad\text{for every }x
\]
in the input support.  More generally, say that $\phi$ has \textbf{connected
support} if for every block $C$ and every input $x$ the graph on $C$ joining
$s$ to $s'$ when $\supp P_s^x\cap\supp P_{s'}^x\neq\varnothing$ is connected;
block-uniform support is the special case in which that graph is complete.
\end{definition}

\begin{remark}[Unifilar Lumpability Does Not Reduce to
Definition~\ref{def:lumpable-quotient}]
\label{rem:unifilar-support-not-automatic}
For an output-independent update Definition~\ref{def:unifilar-lumpable} is
strictly weaker than Definition~\ref{def:lumpable-quotient}, and the two do not
collapse onto each other.  Writing
$\tau(s,x,y)=\tau_0(s,x)$, unifilar lumpability requires
$\phi(\tau_0(s,x))=\phi(\tau_0(s',x))$ only for $s,s'$ in a common block that
share a feasible output $y$ at input $x$; ordinary lumpability requires it for
\emph{all} $s,s'$ in the block.  When within-block emission supports are
disjoint the requirement is vacuous and the unifilar notion is strictly weaker.

An explicit witness has $\mathcal I=\{0,1\}$, $\mathcal O=\{0,1\}$, three
states, and
\[
\tau_0(s_0,0)=s_0,\quad \tau_0(s_0,1)=s_1,\qquad
\tau_0(s_1,0)=s_0,\quad \tau_0(s_1,1)=s_2,
\]
\[
\tau_0(s_2,0)=s_0,\quad \tau_0(s_2,1)=s_0,
\]
with $s_0,s_1$ emitting $0$ almost surely and $s_2$ emitting $1$ almost surely.
The chain is irreducible with strictly positive stationary distribution.  The
partition $\{\{s_0,s_2\},\{s_1\}\}$ is unifilar-lumpable --- $s_0$ and $s_2$
never share a feasible output, so no intertwining is demanded of them --- yet
it is not lumpable, since $\tau_0(s_0,1)=s_1$ and $\tau_0(s_2,1)=s_0$ lie in
different blocks.

Consequently the two feasible sets differ, and the associated retention values
can differ as well: on a three-state machine with
$\tau_0(s_0,\cdot)=(s_1,s_2)$, $\tau_0(s_1,\cdot)=(s_2,s_0)$,
$\tau_0(s_2,\cdot)=(s_0,s_0)$, predictive laws
$P_{s_0}=\delta_2$, $P_{s_2}=\delta_1$, and $P_{s_1}$ interior, the minimum of
the retention cost over unifilar-lumpable quotients with two blocks is strictly
below the minimum over lumpable quotients with two blocks.  The reduction is
therefore genuinely conditional, and the correct statement is
Proposition~\ref{prop:input-driven-specialization}: the two notions agree under
connected support, in particular under block-uniform support, and in particular
whenever all emission laws have full support.
\end{remark}

\begin{proposition}[Right-Congruence Correspondence in Synchronized Machines]
\label{prop:lumpability}
Let the machine have synchronization depth $L$.  A lumpable quotient
$\phi:\Splus\to\mathcal K$ induces a right congruence on sufficiently long
histories by
\[
u\sim_\phi v
\iff
\phi(\sigma(u))=\phi(\sigma(v)).
\]
Conversely, assume that $Z$-predictive equivalence $\sim_Z$ of
Definition~\ref{def:z-predictive-equivalence} is itself a right congruence on
$\Splus$, i.e.\ $s\sim_Zs'$ implies $\tau(s,x)\sim_Z\tau(s',x)$ for every
feasible input $x$.  This holds in particular when $Z$ is the full controlled
future, so that $\sim_Z$ is causal-state equivalence.  Then any finite-index
right congruence coarser than $\sim_Z$ induces a lumpable quotient of
$\Splus$.

The hypothesis cannot be dropped: by Example~\ref{ex:onestep-not-congruence},
one-step predictive equivalence need not be a right congruence, and the
induced partition need not be lumpable.
\end{proposition}

\begin{proof}
If $\phi$ is lumpable and $u\sim_\phi v$, then for sufficiently long $u,v$,
\[
\phi(\sigma(ux))
=
\phi(\tau(\sigma(u),x))
=
\tau_{\mathcal K}(\phi(\sigma(u)),x)
=
\tau_{\mathcal K}(\phi(\sigma(v)),x)
=
\phi(\sigma(vx)),
\]
so $ux\sim_\phi vx$.  Hence $\sim_\phi$ is a right congruence on long
histories.

Conversely, let $\sim$ be a finite-index right congruence coarser than
$\sim_Z$, and assume $\sim_Z$ is a right congruence.  Each $\sim$-class is then
a union of $\sim_Z$-classes.  Define $\phi(s)$ to be the $\sim$-class of any
long history $u$ with $\sigma(u)=s$.  This is well defined: two long histories
with the same $\sim_Z$-state are $\sim_Z$-equivalent, hence $\sim$-equivalent
because $\sim$ is coarser.  The right-congruence property gives
\[
\phi(\tau(s,x))=\tau_{\mathcal K}(\phi(s),x),
\]
so $\phi$ is lumpable.
\end{proof}

\begin{example}[One-Step Predictive Equivalence Need Not Be a Congruence]
\label{ex:onestep-not-congruence}
Let $\mathcal S=\{s_0,s_1,s_2,s_3\}$ with a single input letter $a$,
$\mathcal O=\{0,1\}$, and next-output laws
\[
P_{s_0}=P_{s_1}=\mathrm{Bernoulli}(\tfrac12),
\qquad
P_{s_2}=\delta_0,
\qquad
P_{s_3}=\delta_1 .
\]
Let the transitions be
\[
\tau(s_0,a)=s_2,\quad
\tau(s_1,a)=s_3,\quad
\tau(s_2,a)=s_2,\quad
\tau(s_3,a)=s_3 .
\]
Taking $Z=Y$, the next output, gives $s_0\sim_Ys_1$, since both are fair
coins.  But
\[
\tau(s_0,a)=s_2\not\sim_Y s_3=\tau(s_1,a),
\]
so $\sim_Y$ is not a right congruence, and the partition
$\{\{s_0,s_1\},\{s_2\},\{s_3\}\}$ is not lumpable: the block
$\{s_0,s_1\}$ has no well-defined successor block.

Taking instead $Z$ to be the full future separates $s_0$ from $s_1$, because
their continuations are the constant sequences $0^\infty$ and $1^\infty$.  With
that choice $\sim_Z$ is causal-state equivalence and the converse of
Proposition~\ref{prop:lumpability} applies.

Since input-driven machines are unifilar
(Remark~\ref{rem:unifilar-proper-subclass}), the same example shows that
one-step predictive equivalence need not be a support-relative right
congruence on joint histories in the general model.
\end{example}

%----------------------------------------------------------------------
\subsection{Full-KL Retention Gap}
\label{subsec:full-kl-gap}
%----------------------------------------------------------------------

\begin{definition}[Full-KL Retention Gap]
\label{def:full-kl-retention}
For a stationary controlled causal machine, define
\[
\RetKL(M)
=
\min_{\substack{
\phi:\Splus\to\mathcal K\ \text{lumpable}\\
|\mathcal K|\le M
}}
\sum_{k\in\mathcal K}
\sum_{s:\phi(s)=k}
\pi_s
D_{\mathrm{KL}}(P_s\|\bar P_k),
\]
where
\[
\bar P_k
=
\frac{
\sum_{s:\phi(s)=k}\pi_s P_s
}{
\sum_{s:\phi(s)=k}\pi_s
}.
\]
Blocks of zero stationary mass are ignored.  Because $\Splus$ is finite, the
minimum exists.

For unifilar machines with input-dependent emissions the operative object is
the controlled gap of Definition~\ref{def:controlled-full-kl}, of which the
present definition is the output-independent specialization
(Corollary~\ref{cor:controlled-reduces}).
\end{definition}

In a synchronized machine, lumpable quotients of $\Splus$ correspond to
finite-index right congruences on sufficiently long histories by
Proposition~\ref{prop:lumpability}.  The full-KL centroid $\bar P_k$ is a
mixture distribution.  It is not, in general, an arithmetic mean of
parameters.  The arithmetic-mean centroid appears only in restricted
parametric families, such as the Gaussian quadratic restricted regime.

%----------------------------------------------------------------------
\subsection{Gaussian Quadratic Restricted Retention Gap}
\label{subsec:gaussian-quadratic}
%----------------------------------------------------------------------

\begin{definition}[Gaussian Quadratic Restricted Retention]
\label{def:gaussian-quadratic}
A stationary controlled causal machine is in the \textbf{Gaussian quadratic
restricted regime} if:
\begin{enumerate}[label=(\arabic*)]
\item the output space is the standard Borel space $\mathcal O=\mathbb R^D$;
\item each positive-mass state $s\in\Splus$ has predictive distribution
\[
P_s=\mathcal N(m_s,\tfrac12 I_D);
\]
\item admissible representatives are restricted to the same fixed-covariance
Gaussian family:
\[
Q_c=\mathcal N(c,\tfrac12 I_D);
\]
\item the cost is the KL divergence restricted to this family.
\end{enumerate}
For this family,
\[
D_{\mathrm{KL}}(P_s\|Q_c)
=
\norm{m_s-c}^2.
\]
Equivalently, define Fisher-scaled features
\[
y_s=\sqrt{2}\,m_s,
\qquad
c_y=\sqrt{2}\,c.
\]
Then
\[
D_{\mathrm{KL}}(P_s\|Q_c)
=
\frac12
\norm{y_s-c_y}^2.
\]
The associated quadratic gap is
\[
\RetQuad(M)
=
\min_{\substack{
\phi:\Splus\to\mathcal K\ \text{lumpable}\\
|\mathcal K|\le M
}}
\sum_{k\in\mathcal K}
\sum_{s:\phi(s)=k}
\pi_s
\norm{m_s-\bar m_k}^2,
\]
where
\[
\bar m_k
=
\frac{
\sum_{s:\phi(s)=k}\pi_s m_s
}{
\sum_{s:\phi(s)=k}\pi_s
}.
\]
Equivalently, in Fisher-scaled coordinates,
\[
\RetQuad(M)
=
\min_{\substack{
\phi:\Splus\to\mathcal K\ \text{lumpable}\\
|\mathcal K|\le M
}}
\frac12
\sum_{k\in\mathcal K}
\sum_{s:\phi(s)=k}
\pi_s
\norm{y_s-\bar y_k}^2,
\]
where
\[
\bar y_k
=
\frac{
\sum_{s:\phi(s)=k}\pi_s y_s
}{
\sum_{s:\phi(s)=k}\pi_s
}.
\]
Blocks of zero stationary mass are ignored.
\end{definition}

The restricted Gaussian centroid is the arithmetic mean, either in mean
coordinates $m_s$ or in Fisher-scaled coordinates $y_s$.  The unrestricted
full-KL centroid is the mixture distribution $\bar P_k$.  These are not the
same.

The quadratic gap is a weighted $k$-means objective with a lumpability
constraint.  It is the object used for the spectral retention converse and for
the Gaussian quadratic NP-hardness theorem.  It is not the full-KL retention
gap unless the full-KL problem has been explicitly restricted to the
fixed-covariance Gaussian family.

%======================================================================
\section{Structural Meta-Theorems}
\label{sec:meta}
%======================================================================

The schema supports several structural meta-theorems.  These results are
governed by finite-index right congruences, quotient factorization, and
separated cost functionals.  Their proofs remain regime-specific:
Myhill--Nerode theory, operator approximation, and information processing
supply the analytical engines.

%----------------------------------------------------------------------
\subsection{Master Scope: A Typed Variational Picture}
\label{subsec:meta-scope}
%----------------------------------------------------------------------

The results of this manuscript, across commitment, retention, and grounding,
instantiate one typed variational picture along four pillars.  Stating them
together fixes what is unified and what is not; this subsection is a summary
of scope, not a new derivation.

\begin{enumerate}[label=(\Roman*),leftmargin=2.2em]
\item \textbf{Typed variational schema.}
Every regime gap has the shape
$\Delta_{\mathbb T}(M)=\inf_{\text{budget-}M\text{ feasible }B}
\mathcal L(\delta,B)$,
but the feasible class carries a regime-specific type: a finite right
congruence on histories for commitment (Section~\ref{subsec:det-spec}), a
stationary lumpable quotient for retention
(Section~\ref{subsec:shared-schema-intro}), a deterministic Mealy machine for
the symbolic grounding gap (Definition~\ref{def:symbolic-grounding-gap}), and
a finite-rank bounded operator for the unrestricted linear grounding
relaxation (Theorem~\ref{thm:spectral-grounding}).  The budget $M$ is a
right-congruence index in the first three cases and an operator rank in the
fourth; these are not asserted equivalent
(Remark~\ref{rem:quotient-display-scope}).

\item \textbf{Structural monotonicity and zero thresholds.}
Each regime gap is non-increasing in $M$ and has an exact zero threshold on
its own feasible class: the Myhill--Nerode index for commitment
(Theorem~\ref{thm:commitment-spec}), the stationary support size of distinct
predictive laws for full-KL retention (Theorem~\ref{thm:retention-zero}), and
the operator rank of $H_\nu$ for the linear grounding relaxations
(Proposition~\ref{prop:grounding-structured-zero}).  These thresholds are
proved separately in each regime; no single argument produces all three.

\item \textbf{Conditional response-operator converse.}
Whenever a regime's approximants admit a response operator of bounded rank
dominating the task cost (Definition~\ref{def:response}), the unified
Schatten converse of Theorem~\ref{thm:unified} supplies a spectral-tail lower
bound.  This bridge is conditional on the domination hypothesis and on the
domination modulus, which may itself depend on the budget
(Remark~\ref{rem:bridge-consequence}); it does not by itself make the three
regimes' numerical gaps comparable (Theorem~\ref{thm:schatten-nogo}).

\item \textbf{Regime-specific analytical engines.}
The zero thresholds and quantitative rates are proved by different classical
theories: Myhill--Nerode automata theory and minimax game values for
commitment, information-bottleneck and Fisher-information asymptotics for
retention, and Eckart--Young--Mirsky/Adamjan--Arov--Krein operator
approximation for grounding.  The schema organizes these engines under a
common typed vocabulary; it does not replace them.
\end{enumerate}

Sections~\ref{sec:commitment}--\ref{sec:grounding} develop each regime in
this order, and Section~\ref{sec:joint} states precisely which joint and
comparative statements the typed schema does and does not support across
regimes.

%----------------------------------------------------------------------
\subsection{Monotonicity and Profinite Convergence}
\label{subsec:meta-monotone}
%----------------------------------------------------------------------

\begin{metatheorem}[Monotonicity and Profinite Convergence]
\label{meta:monotone}
Let $\mathbb T$ be a task theory.
\begin{enumerate}[label=(\roman*)]
\item $\Delta_{\mathbb T}(M)$ is non-increasing in $M$.
\item If $\sim_\delta$ has finite index, then, $\mathcal E$ being reflexive by
Definition~\ref{def:task-theory},
\[
\Delta_{\mathbb T}(M)=0
\quad
\forall M\ge\operatorname{index}(\sim_\delta).
\]
\item If $\mathcal E$ is separated and $\{0,\infty\}$-valued, then:
\begin{enumerate}[label=(\alph*)]
\item for worst-case aggregation,
\[
\Delta_{\mathbb T}^{\sup}(M)=0
\text{ for some }M
\iff
\operatorname{index}(\sim_\delta)<\infty;
\]
\item for $\mu$-average aggregation,
\[
\Delta_{\mathbb T}^{\mu}(M)=0
\text{ for some }M
\iff
\kappa_{\mathrm{obs}}(\delta,\mu)<\infty,
\]
with $\kappa_{\mathrm{obs}}$ as in
Definition~\ref{def:observable-support-index}.  If $\mu$ has full support this
reduces to~(a).
\end{enumerate}
\item Suppose $\mathbf R$ is compact metrizable or compact uniform,
$\mathcal E$ is uniformly continuous and separated, the gap is evaluated in
the worst-case/uniform sense, and $\mathcal E$ is \emph{uniformly compatible}
with the uniformity of $\mathbf R$:
\[
\forall\,\text{entourages }U\ \exists\,\eta>0:\quad
\max\{\mathcal E(X,Z),\mathcal E(Y,Z)\}<\eta
\ \Longrightarrow\ (X,Y)\in U .
\]
Conversely, closeness in $\mathbf R$ implies small cost, by uniform continuity
of $\mathcal E$.  On a compact metric space with
$\mathcal E(X,Z)=0\iff X=Z$, joint continuity already implies this
compatibility; stating it explicitly is what makes the passage from ``small
cost to a common representative'' to ``close residual objects'' valid without
a hidden asymmetry.  Then the following are equivalent:
\begin{enumerate}[label=(\alph*)]
\item $\delta$ is uniformly continuous for the profinite uniformity;
\item $\delta$ extends continuously to $\widehat{\mathbf H}$;
\item
\[
\lim_{M\to\infty}\Delta_{\mathbb T}(M)=0.
\]
\end{enumerate}
\end{enumerate}
\end{metatheorem}

\begin{proof}
Parts (i)--(iii) follow from the definitions.  Monotonicity follows because
the feasible set of right congruences of index at most $M$ is nested:
\[
\{\operatorname{index}(\sim)\le M\}
\subseteq
\{\operatorname{index}(\sim)\le M+1\}.
\]
If $\sim_\delta$ has finite index, quotient by $\sim_\delta$ and assign each
class its common residual $\Psi([u])=\delta(u)$, which is well defined since
$u\sim_\delta v$ implies $\delta(u)=\delta(v)$.  Every history then incurs
$\mathcal E(\delta(u),\Psi([u]))=\mathcal E(\delta(u),\delta(u))=0$ by
reflexivity, so the aggregated cost vanishes for any of the three aggregation
schemes.  Reflexivity is exactly what this step consumes; without it the
quotient reproduces the trajectory but need not reproduce it at zero cost.

For the converse in~(iii), suppose $\mathcal E$ is separated and
$\{0,\infty\}$-valued and the worst-case gap is zero at budget $M$.  Then
$\mathcal E(\delta(u),\Psi([u]))=0$ for every $u$, and separatedness upgrades
this to $\delta(u)=\Psi([u])$.  Hence histories in a common class have equal
residuals after every continuation, so the state relation refines
$\sim_\delta$ and $\operatorname{index}(\sim_\delta)\le M$.  Separatedness is
what this direction consumes, and it is assumed in~(iii) but not in~(ii).

For (iv), assume first that $\delta$ is uniformly continuous for the profinite
uniformity.  By the universal property of completion, $\delta$ extends
continuously to a map
\[
\widehat\delta:\widehat{\mathbf H}\to\mathbf R.
\]
Since $\widehat{\mathbf H}$ is compact and $\mathbf R$ is compact uniform,
$\widehat\delta$ is uniformly continuous.  Hence for every cost tolerance
$\varepsilon>0$ there is a finite clopen partition of $\widehat{\mathbf H}$,
equivalently a finite-index right congruence $\sim$, such that
\[
u\sim v
\implies
\mathcal E(\delta(u),\delta(v))<\varepsilon
\]
in the uniform sense.  Choosing representatives for each $\sim$-class gives an
$M$-state approximation with worst-case gap less than $\varepsilon$.

Conversely, suppose $\Delta_{\mathbb T}(M)\to0$ in the worst-case sense.  For
every $\varepsilon>0$, choose a finite quotient $\sim$ and representative map
$\Psi$ such that
\[
\sup_{u\in\mathcal I^*}
\mathcal E(\delta(u),\Psi([u]))
<
\eta,
\]
where $\eta>0$ is chosen small enough, using compactness and separatedness of
$\mathcal E$, to ensure that whenever two histories $u,v$ are assigned the same
representative, their residual objects are $\varepsilon$-close in the uniform
structure of $\mathbf R$.  If $u\sim v$, then $\Psi([u])=\Psi([v])$, so
$\delta(u)$ and $\delta(v)$ are $\varepsilon$-close.  Thus $\delta$ is
uniformly continuous for the profinite uniformity and therefore extends
continuously to $\widehat{\mathbf H}$.
\end{proof}

\begin{remark}[Average-case aggregation]
\label{rem:average-profinite}
Uniform worst-case approximation implies stationary average-case
approximation.  The converse direction, from average-case gap convergence to
profinite uniform continuity, requires additional observability or positivity
assumptions on the stationary process.  The profinite equivalence is therefore
stated for the worst-case/uniform gap unless such assumptions are supplied.
\end{remark}

%----------------------------------------------------------------------
\subsection{Boolean Zero-Error Dichotomy}
\label{subsec:meta-boolean}
%----------------------------------------------------------------------

\begin{definition}[Observable Support Index]
\label{def:observable-support-index}
Let $\mu$ be an average-case history law on $\mathcal I^*$ and let
$S=\operatorname{supp}\mu$.  Such a $\mu$ exists for discounted-prefix
aggregation (Definition~\ref{def:discounted-agg}), where
$\mu(u)=(1-\gamma)\gamma^{|u|}\Pr[\text{prefix }u]$ is a probability measure on
finite histories.  Stationary Ces\`aro aggregation does \emph{not} in general
supply one: Definition~\ref{def:cesaro-agg} averages a cost process along a
stationary trajectory and, when the process has infinite support, the Ces\`aro
averages need not converge to the expectation under any single probability
measure on $\mathcal I^*$.  The Ces\`aro case is treated separately in
Remark~\ref{rem:cesaro-boolean}.
The \textbf{observable support index} $\kappa_{\mathrm{obs}}(\delta,\mu)$ is
the minimum index of a right congruence $\sim$ on $\mathcal I^*$ for which
there exists a representative map $\Psi$ with
\[
\mathcal E\bigl(\delta(u),\Psi([u])\bigr)=0
\qquad
\text{for }\mu\text{-almost every }u .
\]
If $S$ is right-closed and $\mu(u)>0$ for every $u\in S$, define in addition
the \emph{support-relative Nerode relation}
\[
u\sim_{\delta,S}v
\iff
\forall w\in\mathcal I^*\text{ with }uw,vw\in S:\quad
\mathbf R\bigl(\delta(uw),\delta(vw)\bigr)
=
\mathbf R\bigl(\delta(vw),\delta(uw)\bigr)
=
0 .
\]
\end{definition}

\begin{remark}[Restriction to the Support Is Not a Congruence]
\label{rem:support-restriction-scope}
The restriction $\sim_\delta\!\restriction_S$ is not in general a right
congruence.  If $S$ is not closed under right extension, a pair $u,v\in S$ may
have continuations $uw\notin S$ witnessing $u\not\sim_\delta v$ that the
average cost never sees; and the restricted relation need not be compatible
with the transition action.  The operative object is therefore the index of a
global congruence that is observationally exact on $S$, which is what
Definition~\ref{def:observable-support-index} names.  For right-closed $S$ the
two indices are related by the sandwich of
Meta-Theorem~\ref{meta:boolean}(ii); they coincide precisely when the
completion of Lemma~\ref{lem:support-extension} is index-preserving, and
Remark~\ref{rem:support-extension-sharp} exhibits a right-closed support where
they differ.
\end{remark}

\begin{lemma}[Index-Preserving Extension off the Support]
\label{lem:support-extension}
Let $S\subseteq\mathcal I^*$ be right-closed, i.e.\ $u\in S$ implies $uw\in S$
for every $w\in\mathcal I^*$.  Let $\approx$ be an equivalence relation on $S$
that is \emph{stable}: $u\approx v$ and $x\in\mathcal I$ imply $ux\approx vx$.
Then there is a right congruence $\sim$ on all of $\mathcal I^*$ whose
restriction to $S$ is $\approx$, and
\[
\operatorname{index}(\sim)
=
\operatorname{index}(\approx)+c_S ,
\]
where $c_S$ is the number of classes the construction requires off $S$.  One
has $c_S=0$ when $S=\mathcal I^*$, and $c_S\le1$ whenever no history outside
$S$ has an extension inside $S$.  In general $c_S$ may be infinite, in which
case the construction yields a right congruence of infinite index and the
upper bound is vacuous; the sandwich is informative exactly when the
completion has finite index.
\end{lemma}

\begin{proof}
Right-closure gives $u\in S\Rightarrow ux\in S$, so $S$ is forward invariant.
The complement is \emph{not} forward invariant: a history outside $S$ may have
an extension inside $S$.  For $S=\{u:u\text{ contains at least one }\mathtt b\}$,
which is right-closed, one has $\mathtt{aaa}\notin S$ and $\mathtt{aaab}\in S$.
The set $\{u:u\text{ ends in }\mathtt b\}$ would \emph{not} serve here, since
appending $\mathtt a$ leaves it and it is therefore not right-closed.  The
classes off $S$ therefore cannot be merged freely, and the construction below
is stated explicitly.

Define $\sim$ on $\mathcal I^*$ by declaring $u\sim v$ when either
\begin{enumerate}[label=(\alph*)]
\item $u,v\in S$ and $u\approx v$; or
\item $u,v\notin S$ and for every $y\in\mathcal I^*$ one has
$uy\in S\iff vy\in S$, and $uy\in S$ implies $uy\approx vy$.
\end{enumerate}
This is an equivalence relation, and no class meets both $S$ and its
complement.

\emph{Right congruence.}  Let $u\sim v$ and $x\in\mathcal I$.  If $u,v\in S$
then $ux,vx\in S$ by right-closure and $ux\approx vx$ by stability, so
$ux\sim vx$.  Suppose $u,v\notin S$.  For any $y\in\mathcal I^*$, applying~(b)
to the continuation $xy$ gives $uxy\in S\iff vxy\in S$ and, when these hold,
$uxy\approx vxy$.  If $ux\notin S$ then $vx\notin S$ by~(b) with $y=x$, and the
displayed conditions are exactly~(b) for the pair $(ux,vx)$, so $ux\sim vx$.
If $ux\in S$ then~(b) with $y=x$ gives $vx\in S$ and $ux\approx vx$, so
$ux\sim vx$ by~(a).

\emph{Index.}  The classes inside $S$ are exactly the $\approx$-classes, and
those outside are the classes of~(b); writing $c_S$ for their number gives the
stated count.  If no history outside $S$ has an extension inside $S$, then
condition~(b) is vacuous and all such histories form a single class, so
$c_S\le1$.
\end{proof}

\begin{remark}[The Extension Can Strictly Increase the Index]
\label{rem:support-extension-sharp}
The term $c_S$ in Lemma~\ref{lem:support-extension} is not an artifact of the
construction: it can be strictly positive even for right-closed $S$, so the
equality $\kappa_{\mathrm{obs}}(\delta,\mu)=\operatorname{index}(\sim_{\delta,S})$
is false in general.

Let $\mathcal I=\{\mathtt a,\mathtt b\}$ and $S=\mathcal I^{+}$, the set of
nonempty words, which is right-closed.  Let the residual take two values
$A\ne B$ and set
\[
\delta(u)=
\begin{cases}
A, & u\text{ begins with }\mathtt a,\\
B, & u\text{ begins with }\mathtt b,
\end{cases}
\qquad u\in S,
\]
with $\delta(\varepsilon)$ unconstrained since $\varepsilon\notin S$.  For
$u,v\in S$ every continuation $uw,vw$ again lies in $S$ and has the residual
determined by the first letter, so $u\sim_{\delta,S}v$ exactly when $u$ and $v$
begin with the same letter:
\[
\operatorname{index}(\sim_{\delta,S})=2 .
\]

Now let $\sim$ be any right congruence on $\mathcal I^*$ admitting a
representative map exact on $S$.  The class of $\varepsilon$ must send
$\mathtt a$ into the $A$-class and $\mathtt b$ into the $B$-class.  But the
$A$-class sends \emph{both} letters into the $A$-class, and the $B$-class sends
both into the $B$-class, since $u\mathtt a$ and $u\mathtt b$ begin with the same
letter as $u$.  So $\varepsilon$ has a transition pattern matching neither
class and requires a third.  Hence
\[
\kappa_{\mathrm{obs}}(\delta,\mu)=3=\operatorname{index}(\sim_{\delta,S})+1 ,
\]
attaining $c_S=1$.  A three-state witness is $q_\varepsilon$ with
$\tau(q_\varepsilon,\mathtt a)=q_A$, $\tau(q_\varepsilon,\mathtt b)=q_B$ and
$q_A,q_B$ absorbing.

Two consequences follow.  First, the sandwich of
Meta-Theorem~\ref{meta:boolean}(ii) is sharp at both ends.  Second, the
obstruction requires $\mu(\varepsilon)=0$: if $\mu$ is a discounted-prefix law
with $\mu(\varepsilon)>0$, then $\varepsilon\in S$, and right-closure forces
$S=\mathcal I^*$, whence $c_S=0$ and the equality does hold.
\end{remark}

\begin{metatheorem}[Boolean Zero-Error Dichotomy]
\label{meta:boolean}
Let $\mathbb T$ be separated with $\mathcal E$ $\{0,\infty\}$-valued.

\emph{(i) Worst-case aggregation.}  If
\[
\Delta_{\mathbb T}^{\sup}(M)
=
\inf_{\operatorname{index}(\sim)\le M}
\ \sup_{u\in\mathcal I^*}
\mathcal E\bigl(\delta(u),\Psi([u])\bigr),
\]
then
\[
\Delta_{\mathbb T}^{\sup}(M)=0
\iff
\operatorname{index}(\sim_\delta)\le M.
\]

\emph{(ii) Average-case aggregation.}  If the aggregation is discounted-prefix
with history law $\mu$ on $\mathcal I^*$, then
\[
\Delta_{\mathbb T}^{\mu}(M)=0
\iff
\kappa_{\mathrm{obs}}(\delta,\mu)\le M,
\]
with $\kappa_{\mathrm{obs}}$ the observable support index of
Definition~\ref{def:observable-support-index}.  If $\mu$ has full support on
$\mathcal I^*$ then
$\kappa_{\mathrm{obs}}(\delta,\mu)=\operatorname{index}(\sim_\delta)$ and~(ii)
reduces to~(i).  If $S=\operatorname{supp}\mu$ is right-closed with positive
mass, then
\[
\operatorname{index}(\sim_{\delta,S})
\ \le\
\kappa_{\mathrm{obs}}(\delta,\mu)
\ \le\
\operatorname{index}(\sim_{\delta,S})+c_S,
\]
where $c_S$ is the number of additional classes required off $S$ by
Lemma~\ref{lem:support-extension}; in particular $c_S=0$ when
$S=\mathcal I^*$, and equality
$\kappa_{\mathrm{obs}}(\delta,\mu)=\operatorname{index}(\sim_{\delta,S})$ holds
whenever the extension of Lemma~\ref{lem:support-extension} is
index-preserving.
In general the average-case index may be strictly smaller than the worst-case
index, since a machine may err on $\mu$-null histories at no cost.
\end{metatheorem}

\begin{remark}[The Ces\`aro Case]
\label{rem:cesaro-boolean}
Clause~(ii) is stated for discounted-prefix aggregation because that scheme
carries a genuine probability measure on finite histories.  Stationary
Ces\`aro aggregation does not, so the phrase ``$\mu$-almost every history'' has
no direct meaning there and the clause must be reformulated.

Two formulations are available.  On the \emph{stationary state space}, when
the process is a finite-state stationary ergodic chain with positive-mass
state set $\Splus$, the natural statement replaces histories by states: zero
Ces\`aro cost holds at budget $M$ exactly when some right congruence of index
at most $M$ is exact on $\Splus$, which is
Theorem~\ref{thm:retention-zero} in the $\{0,\infty\}$-valued case.  On
\emph{histories}, the correct surrogate for ``$\mu$-almost every'' is
\emph{zero asymptotic frequency of error}: the Ces\`aro cost vanishes exactly
when
\[
\lim_{T\to\infty}\frac1T\,\bigl|\{t<T:\ \mathcal E(\delta(U_t),\Psi([U_t]))>0\}\bigr|
=0
\]
almost surely.  For a finite-state ergodic chain the two agree, since a state
of positive stationary mass is visited with positive asymptotic frequency.
Outside that setting the second is the operative condition, and no formula in
terms of a support-relative Nerode relation is claimed.
\end{remark}

\begin{proof}
\emph{(i)} If $\operatorname{index}(\sim_\delta)\le M$, quotienting by
$\sim_\delta$ and assigning each class its common residual gives an $M$-state
approximation with zero cost at every history.  Conversely, suppose an
$M$-state approximation has zero worst-case cost.  Then for \emph{every}
history $u$ we have $\mathcal E(\delta(u),\Psi([u]))=0$, so by separatedness
$\delta(u)=\Psi([u])$.  Hence any two histories in the same machine state have
equal residuals, and, the state relation being a right congruence, equal
residuals after every continuation; that is, the state relation refines
$\sim_\delta$.  Therefore
$\operatorname{index}(\sim_\delta)\le\operatorname{index}(\sim)\le M$.

\emph{(ii)} Suppose $\Delta_{\mathbb T}^{\mu}(M)=0$.  Then there are a right
congruence $\sim$ of index at most $M$ and a map $\Psi$ with
$\mathcal E(\delta(u),\Psi([u]))=0$ for $\mu$-almost every $u$.  That is
precisely the defining property in
Definition~\ref{def:observable-support-index}, so
$\kappa_{\mathrm{obs}}(\delta,\mu)\le M$.

Conversely, if $\kappa_{\mathrm{obs}}(\delta,\mu)\le M$, take a witnessing
congruence and representative map; the induced $M$-state machine has zero
$\mu$-average cost by construction.

If $\mu$ has full support, ``$\mu$-almost every'' is ``every'', and the
argument of~(i) applies verbatim, giving
$\kappa_{\mathrm{obs}}(\delta,\mu)=\operatorname{index}(\sim_\delta)$.

If $S$ is right-closed with positive mass, then for $u,v\in S$ the constraint
is exactly that all continuations \emph{within $S$} agree, which is
$\sim_{\delta,S}$.  Right-closure makes $\sim_{\delta,S}$ stable under the
transition action in the sense of Lemma~\ref{lem:support-extension}: if
$u\sim_{\delta,S}v$ and $x\in\mathcal I$, then for every $w$ with
$uxw,vxw\in S$ the defining condition applied to the continuation $xw$ gives
$\delta(uxw)=\delta(vxw)$, so $ux\sim_{\delta,S}vx$.

For the lower bound, any global right congruence $\sim$ witnessing
$\kappa_{\mathrm{obs}}$ must, restricted to $S$, refine $\sim_{\delta,S}$: two
$\sim$-equivalent histories in $S$ receive the same representative and hence
the same residual after every continuation remaining in $S$.  Therefore
$\operatorname{index}(\sim_{\delta,S})\le\kappa_{\mathrm{obs}}(\delta,\mu)$.

For the upper bound, apply Lemma~\ref{lem:support-extension} to
$\approx\,=\,\sim_{\delta,S}$ to obtain a global right congruence restricting
to it on $S$, of index $\operatorname{index}(\sim_{\delta,S})+c_S$; assigning
each class inside $S$ its common residual gives zero cost $\mu$-almost
everywhere.  This is the stated sandwich.

The extension step is not vacuous: a history outside $S$ may have an extension
inside $S$, so the classes off $S$ cannot be merged arbitrarily without
destroying the right-congruence property.  Lemma~\ref{lem:support-extension}
supplies the correct construction, and $c_S=0$ exactly when it is
index-preserving.
\end{proof}

\begin{remark}[Why the Qualifier Is Necessary]
\label{rem:boolean-aggregation-qualifier}
The two clauses differ.  Let $\delta$ be a residual trajectory
requiring $\operatorname{index}(\sim_\delta)=k$ on all of $\mathcal I^*$, but
suppose the input process never visits a set of histories whose removal
collapses $\sim_\delta$ to $k'<k$ classes.  Then a $k'$-state machine has zero
average cost while failing to realize $\delta$ exactly, so the unqualified
biconditional would be false.  Statements elsewhere in this manuscript that
invoke ``zero error'' should be read with clause~(i) unless a stationary
support is explicitly named, as in Theorem~\ref{thm:retention-zero}, which is
a clause~(ii) statement with $\operatorname{supp}\mu=\Splus$, so the
relevant quantity there is $\kappa_{\mathrm{obs}}$ rather than the restriction
of $\sim_\delta$ to the support.
\end{remark}

\begin{corollary}[Normalized $\{0,1\}$ Commitment Cost]
\label{cor:boolean-01}
If
\[
\mathcal E(f,g)=
\begin{cases}
0,&f=g,\\
1,&f\neq g,
\end{cases}
\]
and the aggregation is worst-case, then
\[
\Com(M)=
\begin{cases}
0,&M\ge\operatorname{index}(\sim_\delta),\\
1,&M<\operatorname{index}(\sim_\delta).
\end{cases}
\]
\end{corollary}

This corollary separates the structural Boolean valuation $\{0,\infty\}$ from
the normalized operational valuation $\{0,1\}$.  The threshold is the same;
only the positive mismatch value changes.

%----------------------------------------------------------------------
\subsection{Exact Spectral Converse for Linear Hankel Gaps}
\label{subsec:meta-spectral}
%----------------------------------------------------------------------

\begin{metatheorem}[Eckart--Young--Mirsky Spectral Converse]
\label{meta:spectral}
Let $\mathbb T$ be a Hilbert-module task theory with compact Hankel operator
$H_\delta$.  Define the unrestricted linear rank-$M$ gap by
\[
\Delta_{\mathbb T}^{\mathrm{unres}}(M)
=
\inf_{\operatorname{rank}B\le M}
\norm{H_\delta-B}_{\mathrm{op}}.
\]
Then
\[
\Delta_{\mathbb T}^{\mathrm{unres}}(M)
=
\sigma_{M+1}(H_\delta).
\]
Define the Hankel-structured finite-rank gap by
\[
\Delta_{\mathbb T}^{\mathrm{Hank,str}}(M)
=
\inf_{\substack{
\operatorname{rank}B\le M\\
B\ \text{Hankel}
}}
\norm{H_\delta-B}_{\mathrm{op}}.
\]
Then
\[
\Delta_{\mathbb T}^{\mathrm{Hank,str}}(M)
\ge
\sigma_{M+1}(H_\delta),
\]
with equality if an optimal rank-$M$ truncation can be chosen Hankel, or under
a valid Hardy-space Adamjan--Arov--Krein embedding satisfying its hypotheses.
\end{metatheorem}

\begin{proof}
The unrestricted equality is the Eckart--Young--Mirsky theorem for compact
operators on Hilbert space:
\[
\inf_{\operatorname{rank}B\le M}
\norm{H_\delta-B}_{\mathrm{op}}
=
\sigma_{M+1}(H_\delta).
\]
The Hankel-structured feasible set is a subset of the unrestricted feasible
set, so restricting the infimum cannot decrease the value.  Hence
\[
\Delta_{\mathbb T}^{\mathrm{Hank,str}}(M)
\ge
\Delta_{\mathbb T}^{\mathrm{unres}}(M)
=
\sigma_{M+1}(H_\delta).
\]
Equality requires that an unrestricted optimal rank-$M$ approximant can be
chosen Hankel, or that a separate Hankel approximation theorem applies.
\end{proof}

\begin{remark}[AAK versus EYM]
\label{rem:aak-eym-meta}
The theorem above uses Eckart--Young--Mirsky \cite{eckart1936,mirsky1960}.  As
already noted in Remark~\ref{rem:aak-eym-hilbert} for the Hilbert-module
instantiation, the Adamjan--Arov--Krein theorem \cite{aak1971} applies only
after a Hardy-space embedding has been separately established, and the model
does not assume one by default.
\end{remark}

\begin{corollary}[Stable Residual-Norm Spectral Bound]
\label{cor:stable-residual-spectral}
Suppose a residual-cost functional satisfies
\[
\norm{h_\delta-h_{\widehat\delta}}_{\infty,\rho}
\le
\mathcal L_{\infty,\rho}(\delta,\widehat\delta)
\]
with $0<\rho<1/|\Sigma|$.  Then for every linear $M$-rank approximation
$\widehat\delta$ with $\operatorname{rank}H_{\widehat\delta}\le M$,
\[
\mathcal L_{\infty,\rho}(\delta,\widehat\delta)
\ge
(1-|\Sigma|\rho)
\sigma_{M+1}(H_\delta).
\]
\end{corollary}

\begin{proof}
By Proposition~\ref{prop:stable-decay},
\[
\norm{H_\delta-H_{\widehat\delta}}_{\mathrm{op}}
\le
\frac{1}{1-|\Sigma|\rho}
\norm{h_\delta-h_{\widehat\delta}}_{\infty,\rho}.
\]
By Eckart--Young--Mirsky, since
$\operatorname{rank}H_{\widehat\delta}\le M$,
\[
\norm{H_\delta-H_{\widehat\delta}}_{\mathrm{op}}
\ge
\sigma_{M+1}(H_\delta).
\]
Combining the two inequalities gives
\[
\sigma_{M+1}(H_\delta)
\le
\frac{1}{1-|\Sigma|\rho}
\norm{h_\delta-h_{\widehat\delta}}_{\infty,\rho}.
\]
Using the assumed domination of the weighted residual norm by the cost gives
\[
\sigma_{M+1}(H_\delta)
\le
\frac{1}{1-|\Sigma|\rho}
\mathcal L_{\infty,\rho}(\delta,\widehat\delta),
\]
which rearranges to the stated bound.
\end{proof}

This corollary is a domination-bridge result.  It does not assert that the
residual-cost linear finite-rank gap equals the Hankel spectral tail unless
the task cost is exactly the relevant operator or weighted residual norm and
the feasible set is unrestricted.

%----------------------------------------------------------------------
\subsection{Information-Bottleneck Meta-Theorem for Full KL}
\label{subsec:meta-ib}
%----------------------------------------------------------------------

\begin{metatheorem}[Finite-State Information Bottleneck]
\label{meta:ib}
Let $S=S_t$, let $K=K_t=\phi(S_t)$ be a lumpable quotient, and let $Z=Z_t$ be
the predictive future, by default the next output.  For every lumpable quotient
$\phi$,
\[
\RetKL(\phi)
=
I(S;Z\mid K).
\]
Consequently,
\[
\RetKL(\phi)
=
I(S;Z)-I(K;Z).
\]
Therefore,
\[
\RetKL(M)
=
I(S;Z)
-
\max_{\substack{
\phi\ \text{lumpable}\\
|\mathcal K|\le M
}}
I(K;Z).
\]
\end{metatheorem}

\begin{proof}
Since $K=\phi(S)$, $K$ is a function of $S$.  Since $S$ is sufficient for $Z$,
\[
Z\perp\!\!\!\perp K\mid S.
\]
Therefore,
\[
I(K;Z\mid S)=0.
\]
For a fixed quotient $K$, the optimal representative in each block $k$ is the
conditional law $P_{Z\mid K=k}$, which equals the mixture centroid
$\bar P_k$.  Hence
\[
\RetKL(\phi)
=
\mathbb E
D_{\mathrm{KL}}
\bigl(
P_{Z\mid S}
\|
P_{Z\mid K}
\bigr)
=
I(S;Z\mid K).
\]
By the chain rule,
\[
I(S,K;Z)
=
I(K;Z)+I(S;Z\mid K).
\]
But $K$ is a function of $S$, so
\[
I(S,K;Z)=I(S;Z).
\]
The first identity $\RetKL(\phi)=I(S;Z\mid K)$ holds in $[0,\infty]$ without
restriction, regular conditional probabilities existing because $Z$ is standard
Borel.  The subtracted form
\[
I(S;Z\mid K)
=
I(S;Z)-I(K;Z)
\]
requires $I(S;Z)<\infty$, since otherwise the right-hand side is the undefined
difference $\infty-\infty$.  In the present finite-state setting that condition
is automatic: $S$ takes finitely many values, so
\[
I(S;Z)\le H(S)<\infty ,
\]
and $I(K;Z)\le I(S;Z)<\infty$ by data processing, so the rearrangement is
always legitimate here.  The statement carries the hypothesis explicitly
because the identity is applied in the sequel to state sets that are finite by
assumption rather than by construction.  All variational
statements below are to be read in the conditional-mutual-information form,
which is unconditionally defined.
Minimizing over lumpable quotients of size at most $M$ is equivalent to
maximizing $I(K;Z)$ over the same feasible set.
\end{proof}

The identity is stated for the input-marginal target; its controlled analogue,
conditional on the input and valid for general unifilar machines, is
Theorem~\ref{thm:controlled-ib}.

\begin{corollary}[Lumpable Partition Form]
\label{cor:ib-lumpable}
If $\phi:\Splus\to\mathcal K$ is a lumpable quotient with blocks
$C_k=\{s\in\Splus:\phi(s)=k\}$, then
\[
\RetKL(M)
=
\min_{\substack{
\phi:\Splus\to\mathcal K\ \text{lumpable}\\
|\mathcal K|\le M
}}
\sum_{k\in\mathcal K}
\sum_{s\in C_k}
\pi_s
D_{\mathrm{KL}}(P_s\|\bar P_k),
\]
where
\[
\bar P_k
=
\frac{
\sum_{s\in C_k}\pi_s P_s
}{
\sum_{s\in C_k}\pi_s
}.
\]
\end{corollary}

\begin{remark}[Role of the Right-Congruence Constraint]
\label{rem:ib-what-buys}
The identity in Meta-Theorem~\ref{meta:ib} is the information bottleneck \cite{tishby1999,bialek2001} in
native form.  The schema constrains the optimizer: the maximum runs over
finite-index right congruences in synchronized machines, or equivalently over
lumpable deterministic quotients of the stationary support, rather than over
arbitrary Markov encoders.
\end{remark}

\begin{remark}[Full KL versus Quadratic Surrogate]
\label{rem:ib-vs-quad}
The information-bottleneck identity is exact for the full-KL retention gap.
The spectral covariance-tail bound proved later applies to the Gaussian
quadratic restricted gap $\RetQuad(M)$, and locally to full KL under
additional regularity assumptions.  The full-KL centroid is the mixture
distribution $\bar P_k$; the arithmetic mean centroid is a restricted
Gaussian/Fisher phenomenon.
\end{remark}

%----------------------------------------------------------------------
\subsection{The Typed Principle}
\label{subsec:typed-principle}
%----------------------------------------------------------------------

\begin{metatheorem}[Typed Finite-State Rate--Distortion Principle]
\label{meta:typed-rate-distortion}
Fix a regime $\mathsf r$ specified by a rooted history right action
$H_{\mathsf r}$ over an event alphabet $\mathcal A_{\mathsf r}$, a residual map
$\delta_{\mathsf r}:H_{\mathsf r}\to\mathbf R_{\mathsf r}$, a reflexive cost
$\mathcal E_{\mathsf r}$, a declared aggregation and support, and a class
$\mathfrak Q_{\mathsf r}(M)$ of admissible finite-state factors of index at
most $M$, with
\[
\Delta_{\mathsf r}(M)
=
\inf_{(Q,\Psi)\in\mathfrak Q_{\mathsf r}(M)}
\mathcal L_{\mathsf r}\bigl(\delta_{\mathsf r},\Psi\circ Q\bigr).
\]
Then:
\begin{enumerate}[label=(\alph*)]
\item $\Delta_{\mathsf r}(M)$ is nonincreasing in $M$
(Meta-Theorem~\ref{meta:monotone}(i));
\item exact zero distortion is governed by the least admissible factor
through which $\delta_{\mathsf r}$ factors on the declared support: for
separated $\{0,\infty\}$-valued costs this is
Meta-Theorem~\ref{meta:monotone}(iii), with the index taken relative to the
support in the average-aggregation case;
\item every quantitative converse beyond~(b) requires a regime-specific
analytical bridge; no such bound follows from the factorization syntax alone;
\item if the regime admits a Schatten-$p$ response representation with
response operator $A_{\delta_{\mathsf r}}$, effective rank budget
$r_{\mathsf r}(M)$ and modulus $c_{\mathsf r}$, then
\[
\Delta_{\mathsf r}(M)
\ \ge\
c_{\mathsf r}^{-1}\,
\Phi^{(p_{\mathsf r})}_{r_{\mathsf r}(M)}\bigl(A_{\delta_{\mathsf r}}\bigr)
\]
(Theorem~\ref{thm:unified});
\item no comparison between two regimes follows without an explicit map
identifying their history actions, admissible factors, costs and response
operators; in the three principal regimes no such map exists
(Theorem~\ref{thm:schatten-nogo}).
\end{enumerate}

In the three principal regimes the instantiation is:

\begin{center}
\begin{tabular}{llll}
\toprule
regime & admissible factor & exact invariant & analytical bridge \\
\midrule
commitment &
deterministic right congruence &
Myhill--Nerode index &
Boolean realizability \\
retention &
lumpable quotient &
$Z$-predictive support index &
bottleneck / covariance \\
grounding &
linear rank factor &
Hankel rank (unrestricted gap) &
Eckart--Young--Mirsky \\
\bottomrule
\end{tabular}
\end{center}

\noindent
The grounding entry is stated for the \emph{unrestricted} gap $\Dunres$, where
$\operatorname{rank}(H_\nu)=r<\infty$ gives $\Dunres(M)=0$ for $M\ge r$ by
Corollary~\ref{cor:grounding-finite-rank}.  For the Hankel-structured gap
$\DHankstr$ only the inequality $\DHankstr(M)\ge\sigma_{M+1}(H_\nu)$ holds in
general, and the corresponding exact statement needs the additional hypotheses
of Theorems~\ref{thm:aak-equality} and~\ref{thm:aak-multiletter}.
\end{metatheorem}

\begin{remark}[What the Principle Does and Does Not Unify]
\label{rem:typed-principle-scope}
The principle is a statement about shared \emph{syntax} together with an
explicit refusal to share \emph{analysis}.  Clauses~(a) and~(b) are genuinely
regime-independent: monotonicity and the zero-distortion threshold follow from
the factorization structure alone.  Clause~(d) is conditional on a
representation being supplied, and the three regimes supply different operators
--- $H_\nu$, $\Sigma_\pi$, $N_\Lambda$ --- with different effective rank
budgets $M$, $M-1$, $M$ and unrelated moduli, which is precisely the content of
clause~(e).

Two things the principle deliberately does not assert.  It does not claim a
common history alphabet: commitment acts on $\mathcal I^*$ while retention in
the unifilar generality of Definition~\ref{def:unifilar-machine} and grounding
act on $(\mathcal I\times\mathcal O)^*$, and the input-driven retention model
of Definition~\ref{def:controlled-markov} is the output-independent case in
which the action factors through the input projection
(Remark~\ref{rem:unifilar-proper-subclass}).  And it does not assert a
per-symbol rate: the resource is the quotient index, whose description length
is $R(M)=\log M$ by Definition~\ref{def:state-rate}.
\end{remark}

%======================================================================
\section{Retention}
\label{sec:retention}
%======================================================================

The retention regime approximates a stochastic predictive process by a
finite-state stochastic predictor.  The machine type is not a deterministic
Mealy BSG.  In general it is a stationary controlled unifilar predictor, and in
the input-driven synchronized subclass a deterministic-transition,
stochastic-emission lumped $\epsilon$-machine.  The cost is KL divergence
between predictive distributions, and the structural constraint is lumpability
of the stationary causal-state process, in the output-aware sense of
Definition~\ref{def:unifilar-lumpable} for unifilar machines.

%----------------------------------------------------------------------
\subsection{Retention Instantiation}
\label{subsec:retention-instantiation}
%----------------------------------------------------------------------

In the retention instantiation, residual objects are probability distributions
on standard Borel output spaces.  In the synchronized finite-history subclass,
the true trajectory is
\[
\delta(u)=P_{\sigma(u)}
\]
for sufficiently long histories $u$, where $\sigma(u)$ is the causal state
associated with $u$.  In the general stationary formulation, the residual at
time $t$ is the predictive distribution
\[
P_{S_t}
=
\mathcal L(Z_t\mid S_t).
\]

The cost profunctor is KL divergence:
\[
\mathcal E(P,Q)=D_{\mathrm{KL}}(P\|Q).
\]

The default aggregation is stationary average cost, either Ces\`aro or
discounted.  Finite prefixes do not define a probability measure over all
lengths.

An $M$-state retention approximant is a stationary lumpable stochastic
finite-state predictor.  The machine type is a stationary controlled unifilar
predictor (Definition~\ref{def:unifilar-machine}); the input-driven
synchronized subclass is the deterministic-transition case, and in it the
approximant's states are lumped causal states.  The transition structure is deterministic and must be lumpable
with respect to the causal-state transition map $\tau$.  The emission
associated with each machine state is a predictive distribution chosen to
minimize stationary KL cost.

The full-KL retention gap is
\[
\RetKL(M)
=
\min_{\substack{
\phi:\Splus\to\mathcal K\ \text{lumpable}\\
|\mathcal K|\le M
}}
\sum_{k\in\mathcal K}
\sum_{s:\phi(s)=k}
\pi_s
D_{\mathrm{KL}}(P_s\|\bar P_k),
\]
where
\[
\bar P_k
=
\frac{
\sum_{s:\phi(s)=k}\pi_s P_s
}{
\sum_{s:\phi(s)=k}\pi_s
}.
\]

When the quotient is synchronized, lumpable quotients of $\Splus$ correspond
to finite-index right congruences on sufficiently long histories.  The gap is
then a constrained information-bottleneck objective over lumpable partitions
of causal states.

%----------------------------------------------------------------------
\subsection{Full-KL Retention}
\label{subsec:full-kl-retention}
%----------------------------------------------------------------------

\begin{lemma}[Mixture-Centroid Optimality]
\label{lem:mixture-centroid}
Let $P_1,\dots,P_m$ be probability measures on a standard Borel space, and let
$w_1,\dots,w_m>0$ with $\sum_i w_i=1$.  Define the mixture
\[
\bar P=\sum_{i=1}^m w_i P_i.
\]
For any probability measure $Q$ such that $\bar P\ll Q$,
\[
\sum_{i=1}^m w_i D_{\mathrm{KL}}(P_i\|Q)
=
\sum_{i=1}^m w_i D_{\mathrm{KL}}(P_i\|\bar P)
+
D_{\mathrm{KL}}(\bar P\|Q).
\]
Consequently,
\[
\arg\min_Q
\sum_{i=1}^m w_i D_{\mathrm{KL}}(P_i\|Q)
=
\bar P,
\]
and the minimum value is
\[
\sum_{i=1}^m w_i D_{\mathrm{KL}}(P_i\|\bar P).
\]
\end{lemma}

\begin{proof}
For each $i$,
\[
D_{\mathrm{KL}}(P_i\|Q)
=
D_{\mathrm{KL}}(P_i\|\bar P)
+
\int
\log\frac{d\bar P}{dQ}\,dP_i.
\]
Multiplying by $w_i$ and summing gives
\[
\sum_i w_i D_{\mathrm{KL}}(P_i\|Q)
=
\sum_i w_i D_{\mathrm{KL}}(P_i\|\bar P)
+
\int
\log\frac{d\bar P}{dQ}
\,d\!\left(\sum_i w_i P_i\right).
\]
The last integral is $D_{\mathrm{KL}}(\bar P\|Q)$.  Since KL divergence is
nonnegative and equals zero iff $\bar P=Q$, the unique minimizer is
$Q=\bar P$.
\end{proof}

\begin{theorem}[Zero Retention on Stationary Support]
\label{thm:retention-zero}
Assume the predictive distributions $\{P_s:s\in\Splus\}$ are pairwise distinct
on the stationary support.  Then
\[
\RetKL(M)=0
\iff
M\ge |\Splus|.
\]
For worst-case aggregation over all reachable states, replace $|\Splus|$ by
the number of reachable states with pairwise distinct predictive
distributions.
\end{theorem}

\begin{proof}
Zero KL cost requires that, for every block $k$,
\[
P_s=\bar P_k
\]
for every positive-mass state $s$ in that block.  If $P_s\neq P_t$ and
$\pi_s,\pi_t>0$, then merging $s$ and $t$ gives strictly positive KL cost by
strict convexity of KL divergence in its second argument.
Therefore zero cost requires every positive-mass predictive state to occupy
its own block.  The singleton partition is lumpable, so $M=|\Splus|$
suffices.  Conversely, if $M<|\Splus|$, at least two distinct positive-mass
predictive states are merged, giving positive cost.
\end{proof}

\begin{example}[Parity]
\label{ex:parity}
Let $\mathcal I=\{0,1\}$ with fair i.i.d.\ inputs, and let the output emit
running input parity.  There are two causal states with predictive
distributions
\[
P_0=\delta_0,
\qquad
P_1=\delta_1.
\]
Then
\[
\RetKL(2)=0.
\]
The best one-state stochastic-emission approximation emits
$\mathrm{Bernoulli}(q)$.  Its stationary KL cost is
\[
\frac12D_{\mathrm{KL}}(\delta_0\|\mathrm{Bernoulli}(q))
+
\frac12D_{\mathrm{KL}}(\delta_1\|\mathrm{Bernoulli}(q))
=
\frac12[-\log q-\log(1-q)].
\]
The minimum occurs at $q=1/2$, giving
\[
\RetKL(1)=\log 2.
\]
If one restricts the model to deterministic point-mass emissions, the
one-state cost is infinite.  That deterministic-emission model is not the
stochastic-emission retention model studied here.
\end{example}

\begin{lemma}[Coarsening Lemma]
\label{lem:coarsening}
Let $\phi:\Splus\to\mathcal K$ be a lumpable quotient with blocks
\[
C_k=\{s\in\Splus:\phi(s)=k\}.
\]
Then
\[
\RetKL(\phi)
=
\sum_{k\in\mathcal K}
\sum_{s\in C_k}
\pi_s
D_{\mathrm{KL}}(P_s\|\bar P_k),
\]
where
\[
\bar P_k
=
\frac{
\sum_{s\in C_k}\pi_s P_s
}{
\sum_{s\in C_k}\pi_s
}.
\]
Blocks of zero stationary mass are ignored.
\end{lemma}

\begin{proof}
For a fixed quotient, the stationary cost is
\[
\sum_{k\in\mathcal K}
\sum_{s\in C_k}
\pi_s
D_{\mathrm{KL}}(P_s\|Q_k),
\]
where $Q_k$ is the emission distribution assigned to block $C_k$.  By
Lemma~\ref{lem:mixture-centroid}, the minimizer over $Q_k$ is the weighted
mixture centroid $\bar P_k$.  Substituting $\bar P_k$ gives the stated formula.
\end{proof}

\begin{theorem}[Predictive-Information Decomposition]
\label{thm:predictive-info}
For every lumpable quotient $\phi$,
\[
\RetKL(\phi)
=
I(S;Z\mid K_\phi)
=
I(S;Z)-I(K_\phi;Z).
\]
Consequently,
\[
\RetKL(M)
=
I(S;Z)
-
\max_{\substack{
\phi\ \text{lumpable}\\
|\mathcal K|\le M
}}
I(K_\phi;Z).
\]
\end{theorem}

\begin{proof}
This is Meta-Theorem~\ref{meta:ib} specialized to the stationary controlled
causal machine of Definition~\ref{def:controlled-markov}.
\end{proof}

\begin{corollary}[Range, One-State Value, and Refinement Monotonicity]
\label{cor:retention-elementary}
For every $M\ge1$,
\[
0\le \RetKL(M)\le I(S;Z).
\]
Moreover,
\[
\RetKL(1)=I(S;Z).
\]
If $\phi$ and $\psi$ are lumpable quotients and $\phi$ refines $\psi$, then
\[
\RetKL(\phi)\le \RetKL(\psi).
\]
Consequently, the optimized gap $\RetKL(M)$ is nonincreasing in $M$.
\end{corollary}

\begin{proof}
Since $K_\phi$ is a function of $S$,
\[
0\le I(K_\phi;Z)\le I(S;Z),
\]
and the range follows from Theorem~\ref{thm:predictive-info}.  For $M=1$, the
only lumpable quotient is constant, so $I(K_\phi;Z)=0$, giving
$\RetKL(1)=I(S;Z)$.  If $\phi$ refines $\psi$, then $K_\psi$ is a function of
$K_\phi$, so the data-processing inequality gives
\[
I(K_\psi;Z)\le I(K_\phi;Z),
\]
and therefore $\RetKL(\phi)\le\RetKL(\psi)$.  Monotonicity in $M$ follows
because the feasible set of quotients with at most $M$ blocks is nested in
$M$.
\end{proof}


\begin{remark}[Role of the Right-Congruence Constraint in Retention]
\label{rem:retention-what-buys}
Theorem~\ref{thm:predictive-info} is the information bottleneck
\cite{tishby1999,bialek2001} in native form, and what the schema contributes is
the constraint on the optimizer, exactly as recorded in
Remark~\ref{rem:ib-what-buys}: the maximum runs over finite-index right
congruences in synchronized machines, equivalently over lumpable deterministic
quotients of the stationary support, rather than over arbitrary Markov
encoders.
\end{remark}

%----------------------------------------------------------------------
\subsection{Quadratic Restricted Spectral Converse}
\label{subsec:quad-spectral}
%----------------------------------------------------------------------

The spectral retention converse applies to the Gaussian quadratic restricted
regime of Definition~\ref{def:gaussian-quadratic}.  It is not a global theorem
for unrestricted full KL.

Let $y_s$ denote Fisher-scaled predictive features for positive-mass states
$s\in\Splus$.  In the common-covariance Gaussian family with covariance
$\frac12 I$, one may take
\[
y_s=\sqrt{2}\,m_s,
\]
so that
\[
D_{\mathrm{KL}}(P_s\|Q_c)
=
\frac12
\norm{y_s-c_y}^2.
\]

Define the stationary Fisher covariance
\[
\Sigma_\pi
=
\sum_{s\in\Splus}
\pi_s
(y_s-\bar y)(y_s-\bar y)^\top,
\qquad
\bar y=\sum_{s\in\Splus}\pi_s y_s.
\]

In mean coordinates $m_s$, the same covariance is scaled by a factor of $2$.
The quadratic gap in mean coordinates is
\[
\sum_k\sum_{s:\phi(s)=k}\pi_s
\norm{m_s-\bar m_k}^2,
\]
while in Fisher coordinates it is
\[
\frac12
\sum_k\sum_{s:\phi(s)=k}\pi_s
\norm{y_s-\bar y_k}^2.
\]
The spectral statement below uses the Fisher-coordinate covariance
$\Sigma_\pi$.

\begin{theorem}[Spectral Converse for Quadratic Restricted Retention]
\label{thm:spectral-retention}
In the Gaussian quadratic restricted regime, for every $M\ge1$,
\[
\RetQuad(M)
\ge
\frac12
\sum_{i\ge M}
\lambda_i(\Sigma_\pi),
\]
where the eigenvalues are ordered nonincreasingly.  Equivalently,
\[
\RetQuad(M)
\ge
\frac12
\Phi^{(1)}_{M-1}(\Sigma_\pi).
\]
\end{theorem}

\begin{proof}
In Fisher coordinates, the quadratic gap is weighted $k$-means:
\[
\RetQuad(M)
=
\min_{\phi}
\frac12
\sum_{k}
\sum_{s:\phi(s)=k}
\pi_s
\norm{y_s-c_k}^2.
\]
For any lumpable quotient $\phi$ with $K\le M$ cells, let
\[
\Pi_k=\sum_{s:\phi(s)=k}\pi_s,
\qquad
m_k=\Pi_k^{-1}\sum_{s:\phi(s)=k}\pi_s y_s.
\]
Completing the square inside each cluster gives
\[
\sum_{s:\phi(s)=k}
\pi_s
\norm{y_s-c_k}^2
=
\sum_{s:\phi(s)=k}
\pi_s
\norm{y_s-m_k}^2
+
\Pi_k
\norm{m_k-c_k}^2.
\]
Therefore
\[
\frac12
\sum_{k}
\sum_{s:\phi(s)=k}
\pi_s
\norm{y_s-c_k}^2
\ge
\frac12
\sum_{k}
\sum_{s:\phi(s)=k}
\pi_s
\norm{y_s-m_k}^2.
\]

Define the within-cluster covariance
\[
W
=
\sum_{k}
\sum_{s:\phi(s)=k}
\pi_s
(y_s-m_k)(y_s-m_k)^\top
\]
and the between-cluster covariance
\[
B
=
\sum_{k}
\Pi_k
(m_k-\bar y)(m_k-\bar y)^\top.
\]
The total covariance decomposes as
\[
\Sigma_\pi=W+B.
\]
Hence
\[
\frac12
\sum_{k}
\sum_{s:\phi(s)=k}
\pi_s
\norm{y_s-m_k}^2
=
\frac12\tr(W)
=
\frac12\tr(\Sigma_\pi-B).
\]

Because
\[
\sum_k \Pi_k(m_k-\bar y)=0,
\]
the centered cluster-mean vectors satisfy one linear relation, so
\[
\rank(B)\le K-1\le M-1.
\]
By Ky Fan's maximum principle \cite{kyfan1951},
\[
\tr(B)
\le
\sum_{i=1}^{M-1}
\lambda_i(\Sigma_\pi).
\]
Therefore
\[
\RetQuad(M)
\ge
\frac12\tr(\Sigma_\pi)
-
\frac12
\sum_{i=1}^{M-1}
\lambda_i(\Sigma_\pi)
=
\frac12
\sum_{i\ge M}
\lambda_i(\Sigma_\pi).
\]
\end{proof}

\begin{corollary}[Trace Form and Equality Condition]
\label{cor:quad-trace-equality}
For any lumpable quotient $\phi$,
\[
\RetQuad(\phi)
=
\frac12\tr(\Sigma_\pi)
-
\frac12\tr(B_\phi),
\]
where
\[
B_\phi
=
\sum_{k\in\mathcal K}
\Pi_k
(m_k-\bar y)(m_k-\bar y)^\top,
\]
with
\[
\Pi_k=\sum_{s:\phi(s)=k}\pi_s,
\qquad
m_k=\Pi_k^{-1}\sum_{s:\phi(s)=k}\pi_s y_s.
\]
Consequently,
\[
\RetQuad(M)
=
\frac12\tr(\Sigma_\pi)
-
\frac12
\max_{\substack{
\phi\ \text{lumpable}\\
|\mathcal K|\le M
}}
\tr(B_\phi).
\]

The lower bound in Theorem~\ref{thm:spectral-retention} is tight iff there
exists a lumpable quotient $\phi$ with $|\mathcal K|\le M$ such that
\[
\tr(B_\phi)
=
\sum_{i=1}^{M-1}
\lambda_i(\Sigma_\pi).
\]
Equivalently, the lumpable quotient realizes the constrained Ky Fan optimum
capturing the top $M-1$ Fisher covariance directions.
\end{corollary}

\begin{proof}
The trace identity is the ANOVA decomposition used in the proof of
Theorem~\ref{thm:spectral-retention}.  Maximizing $\tr(B_\phi)$ over lumpable
quotients gives the stated formula for $\RetQuad(M)$.  The equality condition
is exactly the condition that the constrained maximizer attains the Ky Fan
upper bound.
\end{proof}


\begin{remark}[Scope of the Spectral Theorem]
\label{rem:spectral-not-full-kl}
The theorem is a statement about $\RetQuad(M)$, not about $\RetKL(M)$.  For
unrestricted full KL, the centroid is the mixture distribution $\bar P_k$, and
the quadratic spectral bound can be loose or can exceed the total predictive
information if interpreted as a full-KL bound.
\end{remark}

\begin{remark}[Tightness]
\label{rem:quad-tightness}
The bound is tight for the quadratic restricted gap whenever an optimal
$M$-means solution splits along the top $M-1$ eigenvectors of
$\Sigma_\pi$.  No tightness claim is made for full KL.
\end{remark}

\begin{lemma}[Uniform Second-Order Expansion of Mixture KL]
\label{lem:fisher-uniform-expansion}
Let $\{P_s^{(r)}\}_{s\in\mathcal S}$ be a finite family of distributions on a
common measurable space, indexed by $r>0$.  Suppose there exist a base
probability measure $P_0$, functions $f_s\in L^2(P_0)$, and constants
$C,C_3<\infty$ such that, uniformly in $s$,
\[
\frac{dP_s^{(r)}}{dP_0}
=
1+r f_s+O(r^2),
\]
with
\[
\int f_s\,dP_0=0,
\qquad
\|f_s\|_{L^3(P_0)}\le C,
\]
and third-order remainders bounded by $C_3 r^3$ in KL expansion.  For any
block $C\subseteq\mathcal S$ and weights
\[
w_s=\frac{\pi_s}{\sum_{t\in C}\pi_t},
\]
define the block mixture
\[
\bar P_C^{(r)}
=
\sum_{s\in C}w_s P_s^{(r)}.
\]
Then uniformly over all blocks,
\[
\sum_{s\in C}w_s
D_{\mathrm{KL}}(P_s^{(r)}\|\bar P_C^{(r)})
=
\frac{r^2}{2}
\sum_{s\in C}
w_s
\|f_s-\bar f_C\|_{L^2(P_0)}^2
+
O(r^3),
\]
where
\[
\bar f_C=\sum_{s\in C}w_s f_s.
\]
If the family is a smooth exponential family with natural-parameter
perturbations
\[
\eta_s^{(r)}=\eta_0+r\theta_s+o(r),
\]
then
\[
\|f_s-f_t\|_{L^2(P_0)}^2
=
\|\theta_s-\theta_t\|_{I_0}^2
+
o(1),
\]
where $I_0$ is the Fisher information matrix at $\eta_0$.
\end{lemma}

\begin{proof}
Write densities
\[
q_s^{(r)}=1+r f_s+O(r^2),
\qquad
q_C^{(r)}=1+r\bar f_C+O(r^2).
\]
Using
\[
\log(1+x)=x-\frac{x^2}{2}+O(x^3),
\]
we obtain
\[
\log\frac{q_s^{(r)}}{q_C^{(r)}}
=
r(f_s-\bar f_C)
-
\frac{r^2}{2}(f_s^2-\bar f_C^2)
+
O(r^3).
\]
Multiplying by $q_s^{(r)}=1+r f_s+O(r^2)$ gives
\[
q_s^{(r)}
\log\frac{q_s^{(r)}}{q_C^{(r)}}
=
r(f_s-\bar f_C)
+
r^2
\left[
f_s(f_s-\bar f_C)
-
\frac12(f_s^2-\bar f_C^2)
\right]
+
O(r^3).
\]
Integrating against $P_0$, the first-order term vanishes because
\[
\int(f_s-\bar f_C)\,dP_0=0.
\]
The second-order term simplifies:
\[
\int
\left[
f_s(f_s-\bar f_C)
-
\frac12(f_s^2-\bar f_C^2)
\right]
dP_0
=
\frac12
\int
(f_s-\bar f_C)^2\,dP_0.
\]
Therefore
\[
D_{\mathrm{KL}}(P_s^{(r)}\|\bar P_C^{(r)})
=
\frac{r^2}{2}
\|f_s-\bar f_C\|_{L^2(P_0)}^2
+
O(r^3).
\]
Multiplying by $w_s$ and summing over $s\in C$ gives the block expansion.

Because $\mathcal S$ is finite and $M$ is fixed, the number of lumpable
partitions is finite.  Hence the $O(r^3)$ remainder is uniform over all
lumpable partitions of size at most $M$.

For a smooth exponential family, the score directions associated with
natural-parameter perturbations $\theta_s$ have $L^2(P_0)$ inner product
\[
\langle \theta_s,\theta_t\rangle_{I_0}
=
\theta_s^\top I_0\theta_t.
\]
Thus the $L^2(P_0)$ variance equals the Fisher-coordinate variance.
\end{proof}

\begin{corollary}[Uniform Fisher Remainder]
\label{cor:fisher-uniform-remainder}
Under the hypotheses of Theorem~\ref{thm:local-full-kl}, there exists a
constant $C<\infty$ such that, for all sufficiently small $r>0$ and all
lumpable quotients $\phi$ with $|\mathcal K|\le M$,
\[
\left|
\RetKLr{r}(\phi)
-
\frac{r^2}{2}
\operatorname{tr}(\Sigma_\pi-B_\phi)
\right|
\le
C r^3,
\]
where $\RetKLr{r}(\phi)$ denotes the full-KL cost of quotient $\phi$ in the
perturbed family.  Consequently,
\[
\RetKLr{r}(M)
=
\frac{r^2}{2}
\min_{\substack{
\phi\ \text{lumpable}\\
|\mathcal K|\le M
}}
\operatorname{tr}(\Sigma_\pi-B_\phi)
+
O(r^3).
\]
\end{corollary}

\begin{proof}
This follows from Lemma~\ref{lem:fisher-uniform-expansion} and the finiteness
of the set of lumpable quotients with at most $M$ blocks.  Uniformity over
that finite set converts the blockwise $O(r^3)$ remainders into a uniform
$O(r^3)$ remainder for the optimized gap.
\end{proof}

\begin{theorem}[Local Fisher Asymptotic Equality]
\label{thm:local-full-kl}
Let $\{P_s^{(r)}\}_{s\in\mathcal S}$ be a family of distributions in a smooth
finite-dimensional exponential family, indexed by $r>0$.  Suppose:
\begin{enumerate}[label=(\arabic*)]
\item the natural parameters satisfy
\[
\eta_s^{(r)}
=
\eta_0+r\theta_s+o(r)
\]
uniformly in $s$;
\item the Fisher information matrix
\[
I_0=I(\eta_0)
\]
is positive definite;
\item the stationary weights satisfy
\[
\pi_s>0,
\qquad
\sum_s\pi_s=1;
\]
\item third derivatives of the log-partition function are uniformly bounded in
a neighborhood of $\eta_0$;
\item the state set $\mathcal S$ and budget $M$ are fixed.
\end{enumerate}
Define Fisher-scaled local features
\[
y_s=I_0^{1/2}\theta_s,
\]
and the stationary local covariance
\[
\Sigma_\pi
=
\sum_s
\pi_s
(y_s-\bar y)(y_s-\bar y)^\top,
\qquad
\bar y=\sum_s\pi_s y_s.
\]
Then, as $r\downarrow0$,
\[
\RetKLr{r}(M)
=
\frac{r^2}{2}
\min_{\substack{
\phi\ \text{lumpable}\\
|\mathcal K|\le M
}}
\operatorname{tr}(\Sigma_\pi-B_\phi)
+
o(r^2),
\]
where
\[
B_\phi
=
\sum_{k\in\mathcal K}
\Pi_k
(m_k-\bar y)(m_k-\bar y)^\top,
\]
with
\[
\Pi_k=\sum_{s:\phi(s)=k}\pi_s,
\qquad
m_k=\Pi_k^{-1}\sum_{s:\phi(s)=k}\pi_s y_s.
\]
Consequently,
\[
\RetKLr{r}(M)
\ge
\frac{r^2}{2}
\sum_{i\ge M}
\lambda_i(\Sigma_\pi)
+
o(r^2).
\]
If a lumpable partition attains the constrained PCA optimum capturing the top
$M-1$ eigenvalues, then equality holds in the leading term:
\[
\RetKLr{r}(M)
=
\frac{r^2}{2}
\sum_{i\ge M}
\lambda_i(\Sigma_\pi)
+
o(r^2).
\]
\end{theorem}

\begin{proof}
KL divergence agrees with the Fisher quadratic form to second order:
\[
D_{\mathrm{KL}}(P_s^{(r)}\|P_t^{(r)})
=
\frac12
\norm{\eta_s-\eta_t}_{I_0}^2
+
O(r^3),
\]
where $\eta_s$ are local Fisher coordinates and $I_0$ is the Fisher information
at the base point.  The assumptions give bounded third derivatives in a
neighborhood of the base parameter.

Because the state set and the budget $M$ are fixed, there are only finitely
many lumpable quotients with at most $M$ blocks.  The $O(r^3)$ remainder is
therefore uniform over all such quotients.  Hence
\[
\RetKLr{r}(M)
=
\frac{r^2}{2}
\min_{\phi}
\operatorname{tr}(\Sigma_\pi-B_\phi)
+
o(r^2).
\]
Since $\rank(B_\phi)\le M-1$, Ky Fan's maximum principle \cite{kyfan1951} gives
\[
\operatorname{tr}(B_\phi)
\le
\sum_{i=1}^{M-1}
\lambda_i(\Sigma_\pi),
\]
and therefore
\[
\operatorname{tr}(\Sigma_\pi-B_\phi)
\ge
\sum_{i\ge M}
\lambda_i(\Sigma_\pi).
\]
This yields the spectral-tail lower bound.  If a lumpable quotient attains the
constrained PCA optimum, equality holds in the leading term.
\end{proof}

\begin{remark}[Scope of the Local Fisher Theorem]
\label{rem:local-fisher-scope}
Theorem~\ref{thm:local-full-kl} is a small-radius asymptotic statement about
the family $\{P_s^{(r)}\}$ as $r\downarrow0$.  It is not a global spectral
equality for an arbitrary fixed full-KL retention instance.
\end{remark}

\begin{theorem}[Global Full-KL Spectral Converse in Probability Coordinates]
\label{thm:global-kl-simplex}
Let $\mathcal O$ be finite and let each state $s\in\Splus$ carry the predictive
probability vector $p_s\in\Delta^{|\mathcal O|-1}$.  Put
\[
\bar p=\sum_{s\in\Splus}\pi_sp_s,
\qquad
\Sigma_p
=
\sum_{s\in\Splus}\pi_s(p_s-\bar p)(p_s-\bar p)^{\!\top}.
\]
Then for every lumpable quotient $\phi$ with at most $M$ blocks,
\[
\RetKL(\phi)
\ \ge\
\sum_{i\ge M}\lambda_i(\Sigma_p),
\]
and consequently
\[
\RetKL(M)
\ \ge\
\sum_{i\ge M}\lambda_i(\Sigma_p).
\]
The bound is global: no small-radius, local-asymptotic, or regularity
hypothesis is used, and $\Sigma_p$ is the covariance of the predictive
\emph{probability vectors}, not the Fisher covariance.  The constant $1$ is
the largest possible: by Proposition~\ref{prop:kl-simplex-sharp} there are
instances on which the ratio of the two sides tends to $1$.
\end{theorem}

\begin{proof}
\emph{Step 1: a quadratic minorant for KL, with the sharp constant.}
Pinsker's inequality gives
$D_{\mathrm{KL}}(p\|q)\ge\tfrac12\|p-q\|_1^2$.  The crude bound
$\|\cdot\|_2\le\|\cdot\|_1$ would lose a factor of two here; on the simplex a
strictly better comparison is available, because the difference of two
probability vectors is a \emph{centred} vector.

\begin{quote}
\emph{Claim.}  If $p,q\in\Delta^{|\mathcal O|-1}$ and $\delta=p-q$, then
$\tfrac12\|\delta\|_1^2\ge\|\delta\|_2^2$.
\end{quote}

\noindent
Indeed $\sum_i\delta_i=0$, so the positive and negative parts of $\delta$ carry
equal mass: writing
$a=\sum_i(\delta_i)_+=\sum_i(\delta_i)_-$ one has $\|\delta\|_1=2a$.  Moreover
$\|\delta\|_\infty\le a$, since a single coordinate cannot exceed the total
mass of its own sign class.  Hence
\[
\|\delta\|_2^2
\ \le\
\|\delta\|_\infty\|\delta\|_1
\ \le\
a\cdot 2a
\ =\
2a^2
\ =\
\tfrac12\|\delta\|_1^2 ,
\]
which is the claim.  Combining it with Pinsker,
\[
D_{\mathrm{KL}}(p\|q)\ \ge\ \tfrac12\|p-q\|_1^2\ \ge\ \|p-q\|_2^2
\tag{$*$}
\]
for all $p,q\in\Delta^{|\mathcal O|-1}$, including boundary points, where both
sides are interpreted in $[0,\infty]$.

Only~$(*)$ is used below.  Two remarks on why it is stated in this form.
First, the centring step is what produces the constant $1$ rather than
$\tfrac12$, and by Proposition~\ref{prop:kl-simplex-sharp} the constant $1$
cannot be improved.  Second, on the \emph{interior} of the simplex the weaker
inequality $D_{\mathrm{KL}}(p\|q)\ge\tfrac12\|p-q\|_2^2$ is the statement that
the negative entropy $\Phi(p)=\sum_ip_i\log p_i$, whose Bregman divergence is
$D_{\mathrm{KL}}$, has Hessian $\operatorname{diag}(1/p_i)\succeq I$ and is
therefore $1$-strongly convex in $\|\cdot\|_2$.  That formulation is not
available on the boundary, where the Hessian is undefined, and it is in any
case lossy by the factor just recovered; this is why the proof is routed
through Pinsker together with the centring claim rather than through strong
convexity.

\emph{Step 2: mixture centroids.}
By Lemma~\ref{lem:mixture-centroid} the block representative minimizing the
$\pi$-weighted KL cost of a block $C$ is the mixture centroid
$\bar p_C=\bigl(\sum_{s\in C}\pi_sp_s\bigr)/\pi(C)$.  Hence, writing $\phi$ for
the partition into blocks $C_1,\dots,C_M$,
\[
\RetKL(\phi)
=
\sum_{k}\sum_{s\in C_k}\pi_s\,D_{\mathrm{KL}}(p_s\|\bar p_{C_k})
\ \overset{(*)}{\ge}\
\sum_k\sum_{s\in C_k}\pi_s\|p_s-\bar p_{C_k}\|_2^2
=
\tr(W_\phi),
\]
where $W_\phi=\sum_k\sum_{s\in C_k}\pi_s(p_s-\bar p_{C_k})(p_s-\bar p_{C_k})^{\!\top}$
is the within-block covariance.

\emph{Step 3: ANOVA and Ky Fan.}
The usual variance decomposition gives $\Sigma_p=W_\phi+B_\phi$ with
$B_\phi=\sum_k\pi(C_k)(\bar p_{C_k}-\bar p)(\bar p_{C_k}-\bar p)^{\!\top}$ the
between-block covariance.  The $M$ centroids satisfy the single affine relation
$\sum_k\pi(C_k)(\bar p_{C_k}-\bar p)=0$, so $\rank(B_\phi)\le M-1$.  Both
$\Sigma_p$ and $B_\phi$ are positive semidefinite and symmetric, so Ky Fan's
maximum principle \cite{kyfan1951} applies to eigenvalues and gives
\[
\tr(B_\phi)\ \le\ \sum_{i=1}^{M-1}\lambda_i(\Sigma_p).
\]
Therefore
\[
\tr(W_\phi)=\tr(\Sigma_p)-\tr(B_\phi)\ \ge\ \sum_{i\ge M}\lambda_i(\Sigma_p),
\]
and combining with Step~2 yields the claim.  Taking the infimum over quotients
with at most $M$ blocks gives the statement for $\RetKL(M)$.
\end{proof}

\begin{proposition}[Optimality of the Constant]
\label{prop:kl-simplex-sharp}
The constant $1$ in Theorem~\ref{thm:global-kl-simplex} is the largest
universal constant.  Explicitly, let $\mathcal O=\{0,1\}$, let two predictive
states carry weight $\tfrac12$ each with
\[
p_\pm=\Bigl(\tfrac12\pm\varepsilon,\ \tfrac12\mp\varepsilon\Bigr),
\qquad
0<\varepsilon<\tfrac12,
\]
and take $M=1$.  Then
\[
\sum_{i\ge1}\lambda_i(\Sigma_p)=\tr(\Sigma_p)=2\varepsilon^2,
\qquad
\RetKL(1)=2\varepsilon^2+O(\varepsilon^4),
\]
so
\[
\frac{\RetKL(1)}{\sum_{i\ge1}\lambda_i(\Sigma_p)}
\ \longrightarrow\ 1
\qquad(\varepsilon\downarrow0).
\]
Consequently no constant $c>1$ satisfies
$\RetKL(M)\ge c\sum_{i\ge M}\lambda_i(\Sigma_p)$ for all instances.
\end{proposition}

\begin{proof}
The stationary mixture is $\bar p=(\tfrac12,\tfrac12)$, and
$p_\pm-\bar p=\pm(\varepsilon,-\varepsilon)$, so
\[
\Sigma_p
=
\tfrac12\begin{pmatrix}\varepsilon\\-\varepsilon\end{pmatrix}
\begin{pmatrix}\varepsilon&-\varepsilon\end{pmatrix}
+
\tfrac12\begin{pmatrix}-\varepsilon\\\varepsilon\end{pmatrix}
\begin{pmatrix}-\varepsilon&\varepsilon\end{pmatrix}
=
\begin{pmatrix}\varepsilon^2&-\varepsilon^2\\-\varepsilon^2&\varepsilon^2\end{pmatrix},
\]
whose trace is $2\varepsilon^2$.  With $M=1$ the only quotient is the trivial
one, whose representative is the mixture centroid $\bar p$
(Lemma~\ref{lem:mixture-centroid}), so
\[
\RetKL(1)
=
\tfrac12D_{\mathrm{KL}}(p_+\|\bar p)+\tfrac12D_{\mathrm{KL}}(p_-\|\bar p)
=
\log2-h\bigl(\tfrac12+\varepsilon\bigr),
\]
$h$ denoting the binary entropy in nats.  Expanding
$h(\tfrac12+\varepsilon)=\log2-2\varepsilon^2-\tfrac43\varepsilon^4+O(\varepsilon^6)$
gives $\RetKL(1)=2\varepsilon^2+\tfrac43\varepsilon^4+O(\varepsilon^6)$, and the
stated ratio follows.
\end{proof}

\begin{remark}[Where the Factor of Two Was Lost]
\label{rem:kl-simplex-constant}
The constant is governed entirely by the quadratic minorant~$(*)$.  Bounding
$\|\delta\|_2\le\|\delta\|_1$ is valid on all of $\mathbb R^{|\mathcal O|}$ but
ignores that differences of probability vectors are centred; the centring
argument in Step~1 of Theorem~\ref{thm:global-kl-simplex} recovers the missing
factor, and Proposition~\ref{prop:kl-simplex-sharp} shows that nothing further
is available.  The extremal configuration is the one that makes both
inequalities in~$(*)$ asymptotically tight at once: two antipodal binary
perturbations of the uniform law, for which Pinsker is tight to leading order
and $\delta$ has exactly two nonzero coordinates of equal magnitude.
\end{remark}

\begin{remark}[Why This Does Not Contradict the Fisher No-Go]
\label{rem:simplex-vs-fisher}
Theorem~\ref{thm:global-kl-simplex} is stated in the covariance $\Sigma_p$ of
the predictive \emph{probability vectors}.  The Fisher covariance
$\Sigma_\pi$ of Theorem~\ref{thm:local-full-kl} is a different object, computed
in natural parameters, and the two are comparable only away from the simplex
boundary.

A bi-Lipschitz reparameterization does not by itself give eigenvalue-by-
eigenvalue comparability of the two covariance matrices, so the transfer is not
immediate.  What does hold is a converse stated directly in a fixed
natural-parameter chart, by strong convexity of the log-partition function
rather than by comparing covariances; this is
Theorem~\ref{thm:global-interior-fisher}.

Note also that ``the natural-parameter Fisher covariance'' is not canonical for
a global family: it depends on a choice of minimal chart and, if one wishes a
Fisher-weighted form, on a reference point.  Fixing the natural parameter of
the stationary mixture $\bar p$ and setting
$\Sigma_F=I(\eta_{\bar p})^{1/2}\Sigma_\eta I(\eta_{\bar p})^{1/2}$ is one such
choice.  Theorem~\ref{thm:no-global-fisher-converse} is stated against a fixed
chart in this sense, and rules out a universal constant for either form.
\end{remark}

\begin{theorem}[Global Interior Fisher-Coordinate Converse]
\label{thm:global-interior-fisher}
Let $\mathcal O$ be finite and write the strictly positive laws on
$\mathcal O$ in a fixed minimal natural-parameter chart
\[
p_\eta(y)=\exp\bigl\{\langle\eta,T(y)\rangle-A(\eta)\bigr\}.
\]
Let $K$ be a compact convex set of natural parameters containing the
parameters $\eta_s$ of all predictive laws \emph{and the parameters of all
their mixture centroids}, and suppose
\[
\nabla^2A(\eta)\ \succeq\ m_K I
\qquad\text{for all }\eta\in K,
\]
for some $m_K>0$.  Put $\bar\eta=\sum_s\pi_s\eta_s$ and
\[
\Sigma_\eta=\sum_s\pi_s(\eta_s-\bar\eta)(\eta_s-\bar\eta)^{\!\top}.
\]
Then
\[
\RetKL(M)\ \ge\ \frac{m_K}{2}\sum_{i\ge M}\lambda_i(\Sigma_\eta).
\]
\end{theorem}

\begin{proof}
Fix a block $C$ with weights $w_s=\pi_s/\pi(C)$, mixture centroid $\bar p_C$,
and let $\zeta_C\in K$ be its natural parameter.  For an exponential family the
Kullback--Leibler divergence is the Bregman divergence of $A$,
\[
D_{\mathrm{KL}}(p_{\eta_s}\|p_{\zeta_C})=B_A(\zeta_C,\eta_s),
\]
so $m_K$-strong convexity of $A$ on $K$ gives
\[
D_{\mathrm{KL}}(p_{\eta_s}\|p_{\zeta_C})\ \ge\ \frac{m_K}{2}\|\eta_s-\zeta_C\|_2^2 .
\]
The hypothesis that $K$ contains the centroid parameters is what places the
segment $[\eta_s,\zeta_C]$ inside $K$; it is not implied by $K\supseteq
\{\eta_s\}$, since $\eta\mapsto p_\eta$ is nonlinear and a mixture centroid
need not have its parameter in the convex hull of the $\eta_s$.

Writing $\bar\eta_C=\sum_{s\in C}w_s\eta_s$ for the \emph{parameter} mean and
using that $\zeta\mapsto\sum_{s\in C}w_s\|\eta_s-\zeta\|_2^2$ is minimized at
$\zeta=\bar\eta_C$,
\[
\sum_{s\in C}w_s D_{\mathrm{KL}}(p_{\eta_s}\|\bar p_C)
\ \ge\
\frac{m_K}{2}\sum_{s\in C}w_s\|\eta_s-\zeta_C\|_2^2
\ \ge\
\frac{m_K}{2}\sum_{s\in C}w_s\|\eta_s-\bar\eta_C\|_2^2 .
\]
Multiplying by $\pi(C)$ and summing over blocks gives
$\RetKL(\phi)\ge\frac{m_K}{2}\tr(W_\phi)$ with $W_\phi$ the within-block
covariance of the $\eta_s$.  The ANOVA decomposition
$\Sigma_\eta=W_\phi+B_\phi$ with $\rank(B_\phi)\le M-1$ and Ky Fan's principle
\cite{kyfan1951} give $\tr(W_\phi)\ge\sum_{i\ge M}\lambda_i(\Sigma_\eta)$,
exactly as in Theorem~\ref{thm:global-kl-simplex}.  Taking the infimum over
quotients with at most $M$ blocks completes the proof.
\end{proof}

\begin{remark}[Degeneration at the Boundary]
\label{rem:interior-fisher-boundary}
The constant $m_K$ is the operative quantity and it degrades toward the
boundary: for the binary family, $\nabla^2A(\eta)=p(1-p)\to0$ as
$p\to\{0,1\}$.  This is the same mechanism recorded in
Proposition~\ref{prop:bernoulli-fisher-scales}, and it is why the theorem is
stated on a compact interior set rather than globally.
\end{remark}

\begin{proposition}[The Two-Point Bernoulli Family Does Not Separate the Scales]
\label{prop:bernoulli-fisher-scales}
Let $\pi_\pm=\tfrac12$, let the output be binary, and set
\[
P_+=\mathrm{Bernoulli}(1-\varepsilon),
\qquad
P_-=\mathrm{Bernoulli}(1-2\varepsilon),
\qquad
\varepsilon\downarrow0 .
\]
Then, with $\Sigma_\pi$ the Fisher covariance of
Theorem~\ref{thm:local-full-kl} computed in the natural parameter,
\[
\RetKL(1)=\Theta(\varepsilon),
\qquad
\tr(\Sigma_\pi)=\Theta(\varepsilon),
\qquad
\frac{\RetKL(1)}{\tr(\Sigma_\pi)}
\longrightarrow
\frac{\log2-\tfrac32\log\tfrac32}{\sqrt2\,(\tfrac12\log2)^2}
=
0.50009\ldots
\]
In particular the ratio is bounded away from $0$, and this family does
\emph{not} witness failure of a global Fisher-covariance converse.
\end{proposition}

\begin{proof}
\emph{The information term.}  With $M=1$ the only quotient is trivial, so
$\RetKL(1)=I(S;Y)=H(\bar P)-\tfrac12H(P_+)-\tfrac12H(P_-)$ with
$\bar P=\mathrm{Bernoulli}(1-\tfrac32\varepsilon)$.  Using
\[
H\bigl(\mathrm{Bernoulli}(1-a\varepsilon)\bigr)
=
a\varepsilon\bigl(-\log\varepsilon+1-\log a\bigr)+o(\varepsilon)
\]
for fixed $a>0$, the $-\log\varepsilon$ terms cancel because
$\tfrac32=\tfrac12\cdot1+\tfrac12\cdot2$, leaving
\[
\RetKL(1)
=
\Bigl(\log2-\tfrac32\log\tfrac32\Bigr)\varepsilon+o(\varepsilon),
\qquad
\log2-\tfrac32\log\tfrac32=0.08494\ldots>0 .
\]

\emph{The Fisher term.}  Write the family in natural-parameter form
$P_\eta(y)=\exp(\eta y-A(\eta))$ with $A(\eta)=\log(1+e^\eta)$ and
$\eta(p)=\log\frac{p}{1-p}$.  The Fisher information in the natural parameter
is the second derivative of the log-partition function,
\[
I(\eta)=A''(\eta)=p(1-p),
\]
\emph{not} $[p(1-p)]^{-1}$; the latter is the Fisher information in the mean
parameter $p$, and the two are reciprocal because $dp/d\eta=p(1-p)$.  Setting
$\eta_\pm=\eta(p_\pm)$ and $\eta_0=\tfrac12(\eta_++\eta_-)$, equal weights give
\[
\tr(\Sigma_\pi)
=
I(\eta_0)\left(\frac{\eta_+-\eta_-}{2}\right)^{2}.
\]
As $\varepsilon\downarrow0$ one has $\eta_+-\eta_-\to\log2$ and
$p_0=1-\sqrt2\,\varepsilon+o(\varepsilon)$, hence
$I(\eta_0)=p_0(1-p_0)\sim\sqrt2\,\varepsilon$ and
\[
\tr(\Sigma_\pi)
=
\sqrt2\,\varepsilon\left(\tfrac12\log2\right)^{2}+o(\varepsilon).
\]

\emph{Conclusion.}  Both quantities are $\Theta(\varepsilon)$ and the ratio
converges to the stated positive constant.
\end{proof}

\begin{theorem}[No Global Fisher-Coordinate Converse]
\label{thm:no-global-fisher-converse}
Fix a minimal natural-parameter chart for the strictly positive laws on a
finite $\mathcal O$, write $\eta_s$ for the parameters of the predictive laws,
and let $\Sigma_\eta$ be their $\pi$-covariance in that chart; for the
Fisher-weighted form fix in addition the reference point $\eta_{\bar p}$ of the
stationary mixture and set
$\Sigma_F=I(\eta_{\bar p})^{1/2}\Sigma_\eta I(\eta_{\bar p})^{1/2}$.

There is \emph{no} universal constant $c>0$, independent of the instance, with
\[
\RetKL(M)\ \ge\ c\sum_{i\ge M}\lambda_i(\Sigma_\eta)
\qquad\text{or}\qquad
\RetKL(M)\ \ge\ c\sum_{i\ge M}\lambda_i(\Sigma_F)
\]
for all instances.  The failure occurs already for two binary predictive states
and $M=1$.
\end{theorem}

\begin{proof}
Let $\mathcal O=\{0,1\}$, let two predictive states carry weight $\tfrac12$
each, and for $\varepsilon\in(0,\tfrac12)$ set
\[
P_-=\mathrm{Bernoulli}(\varepsilon),
\qquad
P_+=\mathrm{Bernoulli}(1-\varepsilon).
\]

\emph{The divergence is bounded.}  The stationary mixture is
$\bar p=\tfrac12(P_-+P_+)=\mathrm{Bernoulli}(\tfrac12)$.  With $M=1$ the only
quotient is trivial and its representative is the mixture centroid
(Lemma~\ref{lem:mixture-centroid}), so
\[
\RetKL(1)
=
\tfrac12D_{\mathrm{KL}}(P_+\|\bar p)+\tfrac12D_{\mathrm{KL}}(P_-\|\bar p)
=
\log2-h(\varepsilon)
\ \le\
\log2,
\]
where $h$ is the binary entropy in nats.  In particular $\RetKL(1)$ stays
bounded, and indeed increases to $\log2$, as $\varepsilon\downarrow0$.

\emph{The parameter covariance diverges.}  In the minimal natural parameter
$\eta(p)=\log\frac{p}{1-p}$ the two states sit at
\[
\eta_\pm=\pm L_\varepsilon,
\qquad
L_\varepsilon=\log\frac{1-\varepsilon}{\varepsilon}.
\]
Equal weights give $\bar\eta=0$ and hence the scalar covariance
\[
\Sigma_\eta=L_\varepsilon^2\ \longrightarrow\ \infty
\qquad(\varepsilon\downarrow0).
\]

\emph{Conclusion for the unweighted form.}  Therefore
\[
\frac{\RetKL(1)}{\lambda_1(\Sigma_\eta)}
\ \le\
\frac{\log2}{L_\varepsilon^2}
\ \longrightarrow\ 0 ,
\]
so no positive $c$ can work.

\emph{Conclusion for the Fisher-weighted form.}  The reference point is
$\bar p=\tfrac12$, at which the natural-parameter Fisher information is
$I(\eta_{\bar p})=A''(\eta_{\bar p})=\bar p(1-\bar p)=\tfrac14$
(Proposition~\ref{prop:bernoulli-fisher-scales} records this normalization).
Hence $\Sigma_F=\tfrac14L_\varepsilon^2\to\infty$ as well, and
$\RetKL(1)/\lambda_1(\Sigma_F)\le 4\log2/L_\varepsilon^2\to0$.  Weighting by the
Fisher information at the mixture therefore does not repair the bound: the
reference point stays in the interior while the states escape to the boundary.
\end{proof}

\begin{remark}[Reading the Two Fisher Statements Together]
\label{rem:fisher-nogo-reading}
Theorem~\ref{thm:no-global-fisher-converse} and
Theorem~\ref{thm:global-interior-fisher} are complementary rather than
competing, and together they settle the natural-parameter question completely.

The interior theorem gives $\RetKL(M)\ge\frac{m_K}{2}\sum_{i\ge M}
\lambda_i(\Sigma_\eta)$ with $m_K=\inf_{\eta\in K}\lambda_{\min}(\nabla^2A(\eta))$
over a compact $K$ containing the state parameters and their mixture
centroids.  The no-go theorem shows that the dependence of the constant on $K$
is not an artefact of the proof: $m_K\downarrow0$ as $K$ is enlarged toward the
boundary, and no boundary-independent constant survives.  The witness is not
the family of Proposition~\ref{prop:bernoulli-fisher-scales}, in which the two
states approach the boundary at comparable rates and the ratio tends to
$0.50009\ldots$; it is the family in which the two states approach
\emph{opposite} faces of the simplex, so that the parameter separation
$L_\varepsilon$ diverges while the divergence saturates at $\log2$.

The contrast with probability coordinates is the substance of the matter.  In
those coordinates Theorem~\ref{thm:global-kl-simplex} holds globally with the
sharp constant $1$, because the simplex is bounded and $\|p_s-\bar p\|_2$
cannot exceed $\sqrt2$.  The natural-parameter chart is unbounded, and it is
exactly this unboundedness that the counterexample exploits.
\end{remark}

\begin{remark}[What Survives, and Why]
\label{rem:global-kl-what-survives}
The delicate regime is the boundary of the simplex.  In the natural parameter
the Fisher information of a Bernoulli law is $p(1-p)$, which \emph{vanishes} at
the boundary, while $\RetKL$ vanishes there as well; the two scales must
therefore be compared rather than assumed to separate, and
Proposition~\ref{prop:bernoulli-fisher-scales} shows that on the natural
two-point family they do not.  The local theorem
(Theorem~\ref{thm:local-full-kl}) is unaffected in any case, since there the
family is fixed and only the perturbation scale $r$ tends to zero.

The hypotheses under which a global Fisher-covariance converse is known are
those that make the parameter-to-probability map bi-Lipschitz: uniform
interiority of the predictive laws, or equivalently a fixed compact
exponential-family parameter set bounded away from the boundary.  Under any of
them Theorem~\ref{thm:global-kl-simplex} transfers, by
Remark~\ref{rem:simplex-vs-fisher}.  Without them no such converse exists at
all: Theorem~\ref{thm:no-global-fisher-converse} exhibits binary families on
which $\RetKL$ stays bounded while the parameter covariance diverges.  The
interiority hypothesis is therefore necessary, not merely convenient.
\end{remark}


\begin{remark}[Machine-Checked Fragments]
\label{rem:lean-formalization}
Several of the inequalities on which the results below depend have been
formalized in Lean~4 against Mathlib and are machine-checked: the centring step of
Theorem~\ref{thm:global-kl-simplex}, namely that $\sum_id_i=0$ implies
$\Vert d\Vert_2^2\le\tfrac12\Vert d\Vert_1^2$, together with its two
supporting lemmas; the halving step
$c_2\le c_1$ and $c_1+c_2\le n$ imply $2c_2\le n$ of
Proposition~\ref{prop:esyncsi-log}, in a form making its independence of
$|\mathcal O|$ explicit; the comparison
$\min\{a(M),b(M)\}\le a(N)+b(N)$ underlying
Lemma~\ref{lem:discrete-bv-sandwich} and the mediant inequality underlying
Lemma~\ref{lem:kappa-bounded}; the parallel-axis identity behind the ANOVA
split and the minimality of the weighted mean; and the counting core of
Proposition~\ref{prop:lsyncu-single-input}, that a bounded strictly increasing
block count admits at most $M-1$ increases.  Fifteen statements are checked in
total, with no appeal to \texttt{sorry} and no axioms beyond Lean's standard
three.

The scope of this is deliberately narrow and should not be overstated.  What is
verified is the arithmetic skeleton: the inequalities, identities and counting
arguments.  Pinsker's inequality is taken as a hypothesis rather than imported;
Ky Fan's maximum principle, the Kullback--Leibler functional itself, and all
automata-theoretic constructions --- Mealy machines, version spaces, the
adaptive games --- are not formalized, so the surrounding theorems are not
verified end to end.

\end{remark}

\begin{remark}[Sharpness of the Quadratic Minorant]
\label{rem:kl-minorant-sharp}
The constant in the minorant $(*)$ of Theorem~\ref{thm:global-kl-simplex} is
exactly $1$, and both inequalities composing it are individually sharp.

For the centring step, $\sup\{\|\delta\|_2^2/(\tfrac12\|\delta\|_1^2)\}$ over
nonzero centred $\delta\in\mathbb R^k$ equals $1$ for every $k\ge2$, attained
precisely when $\delta$ has exactly two nonzero coordinates of equal magnitude
and opposite sign; direct maximization over $k=2,\dots,12$ returns
$1.000000000000$.  For the composite, $\inf\{D_{\mathrm{KL}}(p\|q)/
\|p-q\|_2^2\}$ over distinct $p,q\in\Delta^{k-1}$ equals $1$ and is not
attained: over $58{,}186$ deliberately near-degenerate pairs evaluated in
$60$-digit arithmetic the minimum observed ratio is $1.0000000053$.  Along the
extremal family $p=(\tfrac12+t,\tfrac12-t)$, $q=(\tfrac12-t,\tfrac12+t)$ one
has $D_{\mathrm{KL}}(p\|q)/\|p-q\|_2^2=1+\tfrac43t^2+O(t^4)$, decreasing to $1$.

A caution on floating-point evaluation: at separations $\|p-q\|_2^2\sim10^{-12}$
the ratio computed in double precision can read below $1$ purely through
cancellation in the logarithms.  Re-evaluating the same pairs in extended
precision restores the inequality.
\end{remark}

\begin{remark}[Numerical Verification]
\label{rem:retention-numerical}
The predictive-information identity of Theorem~\ref{thm:predictive-info} is an
exact identity and can be checked numerically on finite instances: over
stationary chains with $|\Splus|\le5$ and $|\mathcal O|\le4$, enumerating
\emph{all} set partitions, the three quantities $\RetKL(\phi)$,
$I(S;Z\mid K_\phi)$ and $I(S;Z)-I(K_\phi;Z)$ agree to $7.6\times10^{-16}$.

Theorem~\ref{thm:global-kl-simplex} admits an end-to-end check of the same
kind, and each step of its proof can be isolated.  Over $2{,}434$
(instance, partition) pairs with $|\Splus|\le5$, $|\mathcal O|\le4$, again
enumerating all set partitions, the maximum violation is
$7.5\times10^{-17}$ for the minorant $\RetKL(\phi)\ge\tr(W_\phi)$,
$1.1\times10^{-16}$ for the ANOVA decomposition
$\Sigma_p=W_\phi+B_\phi$, $1.7\times10^{-16}$ for the Ky Fan bound
$\tr(B_\phi)\le\sum_{i\le M-1}\lambda_i$, and $7.5\times10^{-17}$ for the
composite conclusion $\RetKL(\phi)\ge\sum_{i\ge M}\lambda_i(\Sigma_p)$.  The
smallest ratio of the two sides over those pairs is $1.0015$, so the bound is
not vacuous on generic instances.

Theorem~\ref{thm:global-interior-fisher} requires more care to test, since
$m_K$ is the least eigenvalue of $\nabla^2A$ over a convex set containing the
state parameters \emph{and} the centroid parameters, not a coarser proxy such
as the least predictive probability.  Computing $m_K$ by minimizing
$\lambda_{\min}(\nabla^2A)$ over the convex hull gives $0$ violations on
$1{,}664$ (instance, partition) pairs, with least ratio $1.0011$; substituting
a proxy for $m_K$ produces spurious violations and tests a different statement.

The quadratic-surrogate spectral bound can be checked directly on finite
weighted point clouds.  The local full-KL Fisher expansion can be checked on
small-radius exponential-family instances by verifying the predicted $r^2$
leading term and an $O(r^3)$ remainder.  These computations illustrate the
theorems; they do not establish their hypotheses.
\end{remark}

\begin{proposition}[Relation between Full-KL and Quadratic Restricted Gaps]
\label{prop:retkl-le-retquad}
In the Gaussian quadratic restricted regime, for every $M\ge1$,
\[
\RetKL(M)\le \RetQuad(M).
\]
\end{proposition}

\begin{proof}
Fix a lumpable quotient $\phi$.  In the full-KL problem, the block
representative is unrestricted, and the optimal representative is the mixture
centroid $\bar P_k$.  In the quadratic restricted problem, the representative
is constrained to the fixed-covariance Gaussian family
$Q_c=\mathcal N(c,\tfrac12 I)$.  A constrained minimization problem has value
no smaller than the corresponding unconstrained problem.  Hence, for each
$\phi$,
\[
\RetKL(\phi)\le \RetQuad(\phi).
\]
Minimizing over lumpable quotients gives $\RetKL(M)\le \RetQuad(M)$.
\end{proof}

%----------------------------------------------------------------------
\subsection{Unifilar Retention and the Controlled Information Bottleneck}
\label{subsec:unifilar-retention}
%----------------------------------------------------------------------

The retention theory extends from the input-driven model to general unifilar
machines without new analytical mechanism: the information-theoretic identities
depend only on the sufficiency structure $K=\phi(S)$, and the spectral
converses apply fiberwise in the input.  What does change is the feasible set,
which is now the unifilar-lumpable quotients of
Definition~\ref{def:unifilar-lumpable}, and with it the zero-retention
threshold.

\begin{definition}[Controlled Full-KL Retention Gap]
\label{def:controlled-full-kl}
Let the input process be i.i.d.\ with law $p$ and independent of the state
process.  For a unifilar-lumpable quotient $\phi:\Splus\to\mathcal K$ with
blocks $C_k$ and representatives $\{Q_k^x\}$, set
\[
\RetKLc(\phi)
=
\sum_xp(x)\sum_k\sum_{s\in C_k}\pi_s\,
D_{\mathrm{KL}}\bigl(P_s^x\,\big\|\,Q_k^x\bigr),
\]
the representatives being optimized for each $\phi$, and
\[
\RetKLc(M)
=
\min_{\substack{\phi\ \text{unifilar-lumpable}\\ |\mathcal K|\le M}}
\RetKLc(\phi).
\]
Blocks of zero stationary mass are ignored.
\end{definition}

The independence of the input from the state is a genuine hypothesis and not a
normalization: it is what makes $\mathbb P(S_t=s\mid K_t=k,X_t=x)$ equal to
$\pi_s/\pi(C_k)$, and hence what identifies the block centroid with a
conditional law.  Remark~\ref{rem:controlled-ib-independence} records what
fails without it.

\begin{theorem}[Controlled Information-Bottleneck Identity]
\label{thm:controlled-ib}
Under the hypotheses of Definition~\ref{def:controlled-full-kl}, for every
unifilar-lumpable $\phi$,
\[
\RetKLc(\phi)
=
I(S;Y\mid K_\phi,X)
=
I(S;Y\mid X)-I(K_\phi;Y\mid X),
\]
and consequently
\[
\RetKLc(M)
=
I(S;Y\mid X)
-
\max_{\substack{\phi\ \text{unifilar-lumpable}\\|\mathcal K|\le M}}
I(K_\phi;Y\mid X).
\]
The optimal representative is the controlled mixture centroid
\[
\bar P_k^x
=
\frac{\sum_{s\in C_k}\pi_sP_s^x}{\sum_{s\in C_k}\pi_s}.
\]
All the mutual informations are finite, since $S$ is finite and
$I(S;Y\mid X)\le H(S)$.
\end{theorem}

\begin{proof}
Fix an input $x$.  Conditionally on $X_t=x$, the block cost
$\sum_{s\in C_k}\pi_sD_{\mathrm{KL}}(P_s^x\|Q_k^x)$ is minimized over $Q_k^x$
at the mixture centroid $\bar P_k^x$, by Lemma~\ref{lem:mixture-centroid}
applied to the family $\{P_s^x:s\in C_k\}$ with weights proportional to
$\pi_s$.  Substituting the centroids and summing against $p(x)$,
\[
\RetKLc(\phi)
=
\mathbb E\,D_{\mathrm{KL}}
\bigl(P_{Y\mid S,X}\,\big\|\,P_{Y\mid K,X}\bigr),
\]
where the identification $P_{Y\mid K,X}(\cdot\mid k,x)=\bar P_k^x$ uses
independence of $X_t$ from $S_t$: it gives
$\mathbb P(S_t=s\mid K_t=k,X_t=x)=\pi_s/\pi(C_k)$ for $s\in C_k$, so the
conditional law of $Y_t$ given $(K_t,X_t)=(k,x)$ is precisely the
$\pi$-weighted mixture $\bar P_k^x$.  Since $K=\phi(S)$ is a function of $S$,
we have $P_{Y\mid S,K,X}=P_{Y\mid S,X}$, and the displayed expectation is by
definition $I(S;Y\mid K,X)$.

The chain rule for conditional mutual information gives
\[
I(S;Y\mid X)=I(K;Y\mid X)+I(S;Y\mid K,X),
\]
valid because $K$ is a function of $S$, whence the subtracted form.  The
rearrangement is legitimate because $I(S;Y\mid X)\le H(S)<\infty$, so no
difference of infinities arises.  Minimizing over unifilar-lumpable quotients
with at most $M$ blocks turns the identity into the stated variational form.
\end{proof}

\begin{remark}[Role of the Independence Hypothesis]
\label{rem:controlled-ib-independence}
Definition~\ref{def:unifilar-machine} requires only that the input process be
stationary and ergodic, which permits $X_t$ to be correlated with $S_t$.  Under
such correlation the two sides of Theorem~\ref{thm:controlled-ib} differ: the
conditional weight of $s$ within its block given $X_t=x$ is not
$\pi_s/\pi(C_k)$, so the centroid $\bar P_k^x$ is not a conditional law.
Already with two states, binary inputs and a ternary output alphabet, a
non-product stationary law of $(S,X)$ makes the quantity of
Definition~\ref{def:controlled-full-kl} exceed $I(S;Y\mid X)-I(K;Y\mid X)$ by
more than $0.13$ nats on the one-block quotient.

The hypothesis is a property of the cost functional rather than of the
machine, and it can be removed by measuring the cost against the fiberwise
conditional weights.  This is carried out in
Theorem~\ref{thm:controlled-ib-general} below, of which
Theorem~\ref{thm:controlled-ib} is the independent case.
\end{remark}

\begin{definition}[Fiberwise Controlled Retention Gap]
\label{def:controlled-full-kl-general}
Let $(S_t,X_t)$ have an arbitrary stationary joint law, with input marginal
$p(x)=\mathbb P(X_t=x)$ and fiberwise conditional state weights
\[
\pi_s^x=\mathbb P(S_t=s\mid X_t=x),
\qquad
x\in\supp p .
\]
For a unifilar-lumpable quotient $\phi$ with blocks $C_k$ and representatives
$\{Q_k^x\}$, define
\[
\RetKLg(\phi)
=
\sum_xp(x)\sum_k\sum_{s\in C_k}\pi_s^x\,
D_{\mathrm{KL}}\bigl(P_s^x\,\big\|\,Q_k^x\bigr),
\]
the representatives being optimized for each $\phi$, and
\[
\RetKLg(M)
=
\min_{\substack{\phi\ \text{unifilar-lumpable}\\|\mathcal K|\le M}}
\RetKLg(\phi).
\]
Fibers and blocks of zero mass are ignored.  When $X_t$ is independent of
$S_t$ one has $\pi_s^x=\pi_s$ and
Definition~\ref{def:controlled-full-kl} is recovered.
\end{definition}

\begin{theorem}[Controlled Information-Bottleneck Identity, General Inputs]
\label{thm:controlled-ib-general}
Let $(S_t,X_t)$ be stationary with arbitrary dependence, and let $\phi$ be
unifilar-lumpable.  Then
\[
\RetKLg(\phi)
=
I(S;Y\mid K_\phi,X)
=
I(S;Y\mid X)-I(K_\phi;Y\mid X),
\]
and consequently
\[
\RetKLg(M)
=
I(S;Y\mid X)
-
\max_{\substack{\phi\ \text{unifilar-lumpable}\\|\mathcal K|\le M}}
I(K_\phi;Y\mid X).
\]
The optimal representative is the fiberwise conditional centroid
\[
\bar P_k^x
=
\frac{\sum_{s\in C_k}\pi_s^xP_s^x}{\sum_{s\in C_k}\pi_s^x}
=
\mathcal L\bigl(Y_t\mid K_t=k,\ X_t=x\bigr).
\]
\end{theorem}

\begin{proof}
Fix $x\in\supp p$ and a block $C_k$ of positive conditional mass.  Conditionally
on $X_t=x$, the inner sum $\sum_{s\in C_k}\pi_s^xD_{\mathrm{KL}}(P_s^x\|Q_k^x)$
is minimized over $Q_k^x$ at the $\pi^x$-weighted mixture of the $P_s^x$, by
Lemma~\ref{lem:mixture-centroid} applied with weights proportional to
$\pi_s^x$.  That mixture is exactly
$\mathcal L(Y_t\mid K_t=k,X_t=x)$: conditioning on $X_t=x$ first and on the
block second,
\[
\mathbb P(S_t=s\mid K_t=k,X_t=x)
=
\frac{\pi_s^x}{\sum_{s'\in C_k}\pi_{s'}^x}
\qquad(s\in C_k),
\]
which holds for \emph{any} joint law of $(S_t,X_t)$, no independence being
used; independence was needed in Theorem~\ref{thm:controlled-ib} only to
replace $\pi_s^x$ by $\pi_s$.  Substituting the centroids,
\[
\RetKLg(\phi)
=
\mathbb E\,D_{\mathrm{KL}}
\bigl(P_{Y\mid S,X}\,\big\|\,P_{Y\mid K,X}\bigr)
=
I(S;Y\mid K,X),
\]
the last equality being the definition of conditional mutual information,
together with $P_{Y\mid S,K,X}=P_{Y\mid S,X}$ because $K=\phi(S)$.  The chain
rule $I(S;Y\mid X)=I(K;Y\mid X)+I(S;Y\mid K,X)$ gives the subtracted form, and
is legitimate since $I(S;Y\mid X)\le H(S)<\infty$.  Minimizing over
unifilar-lumpable quotients with at most $M$ blocks yields the variational
statement.
\end{proof}

\begin{corollary}[Elementary Properties, General Inputs]
\label{cor:controlled-elementary-general}
$0\le\RetKLg(M)\le I(S;Y\mid X)$ with $\RetKLg(1)=I(S;Y\mid X)$; the gap is
nonincreasing in $M$; and $\RetKLg(\phi)\le\RetKLg(\psi)$ whenever the
unifilar-lumpable $\phi$ refines $\psi$.
\end{corollary}

\begin{proof}
As in Corollary~\ref{cor:controlled-elementary}, every step being conditional
on $X$ and using Theorem~\ref{thm:controlled-ib-general} in place of
Theorem~\ref{thm:controlled-ib}.
\end{proof}

\begin{corollary}[Fiberwise Spectral Converse, General Inputs]
\label{cor:controlled-simplex-general}
For finite $\mathcal O$, let $\Sigma_p^{x}$ denote the covariance of the
predictive probability vectors $\{p_s^x\}$ under the \emph{conditional} weights
$\{\pi_s^x\}$.  Then
\[
\RetKLg(M)
\ \ge\
\sum_xp(x)\sum_{i\ge M}\lambda_i\bigl(\Sigma_p^{x}\bigr).
\]
\end{corollary}

\begin{proof}
On each fiber the family $\{p_s^x\}$ carries the probability weights
$\{\pi_s^x\}$ and the quotient is the same for every fiber, the lumpability
constraint being $x$-independent.  Theorem~\ref{thm:global-kl-simplex} applies
verbatim on the fiber with those weights; multiplying by $p(x)$ and summing
gives the bound.
\end{proof}

\begin{remark}[The Reweighting Is Not Cosmetic]
\label{rem:controlled-general-reweighting}
In Corollary~\ref{cor:controlled-simplex-general} the fiber covariance must be
formed with the conditional weights $\pi_s^x$.  Substituting the unconditional
weights $\pi_s$, which is the form that Corollary~\ref{cor:controlled-simplex-spectral}
inherits from the independent case, produces a quantity that is not a lower
bound: among randomly generated instances with correlated $(S,X)$ and all
budgets $M$, the unconditional-weight expression exceeds $\RetKLg(M)$ in
several per cent of cases, while the conditional-weight bound of
Corollary~\ref{cor:controlled-simplex-general} is never violated.  Dependence
between input and state therefore changes the operative covariance, not only
the block weights.
\end{remark}

\begin{corollary}[Range, One-State Value, and Refinement Monotonicity, Controlled]
\label{cor:controlled-elementary}
$0\le\RetKLc(M)\le I(S;Y\mid X)$, with
$\RetKLc(1)=I(S;Y\mid X)$; if $\phi$ refines $\psi$ and both are
unifilar-lumpable then
$\RetKLc(\phi)\le\RetKLc(\psi)$; and
$\RetKLc(M)$ is nonincreasing in $M$.
\end{corollary}

\begin{proof}
This is Corollary~\ref{cor:retention-elementary} with every quantity
conditioned on $X$.  A single block makes $K$ constant, so $I(K;Y\mid X)=0$;
refinement is data processing conditional on $X$, since $K_\psi$ is then a
function of $K_\phi$; and the feasible sets are nested in $M$.  Feasibility of
the one-block quotient is Definition~\ref{def:unifilar-lumpable}.
\end{proof}

The zero-retention threshold does \emph{not} transfer verbatim.  In the
input-driven theory the singleton quotient is lumpable and the predictive
partition can always be realized, so the threshold is the number of distinct
predictive laws (Theorem~\ref{thm:retention-zero}).  For unifilar machines the
partition by predictive kernel need not be unifilar-lumpable, and the correct
threshold is the index of its coarsest unifilar-lumpable refinement.

Corollary~\ref{cor:controlled-elementary} is the controlled counterpart of the
elementary facts recorded in Corollary~\ref{cor:retention-elementary}, and the
boundary behaviour of the natural-parameter converses is unchanged under
control, as Remark~\ref{rem:controlled-fisher-nogo} records.

\begin{definition}[Kernel Partition and Its Stable Refinement]
\label{def:kernel-refinement}
Let $\phi_{\ker}$ be the partition of $\Splus$ by predictive kernel, so
$\phi_{\ker}(s)=\phi_{\ker}(s')$ if and only if $P_s^x=P_{s'}^x$ for every $x$
in the input support.  Define $\phi_{\ker}^{\,\ast}$ to be the coarsest
unifilar-lumpable quotient refining $\phi_{\ker}$, and write
$N^{\ast}=|\phi_{\ker}^{\,\ast}(\Splus)|$ for its index.
\end{definition}

\begin{proposition}[Existence of the Stable Refinement]
\label{prop:kernel-refinement-exists}
$\phi_{\ker}^{\,\ast}$ exists and is unique, and is computed by the refinement
recursion
\[
\phi^{(0)}=\phi_{\ker},
\qquad
\phi^{(m+1)}(s)
=
\Bigl(\phi^{(m)}(s),\ \bigl(\phi^{(m)}(\tau(s,x,y))\bigr)_{(x,y)\ \text{feasible from }s}\Bigr),
\]
which stabilizes after at most $|\Splus|$ steps.
\end{proposition}

\begin{proof}
Call a partition \emph{admissible} if it is unifilar-lumpable and refines
$\phi_{\ker}$.  The singleton partition is admissible, so the family is
nonempty.  The recursion produces a decreasing chain of partitions, each
refining $\phi_{\ker}$, in a set of size $|\Splus|$; the number of blocks is
nondecreasing and bounded by $|\Splus|$, so the chain stabilizes after at most
$|\Splus|$ steps, at a partition $\phi^{\infty}$.  At the fixed point,
$\phi^{\infty}(s)=\phi^{\infty}(s')$ implies
$\phi^{\infty}(\tau(s,x,y))=\phi^{\infty}(\tau(s',x,y))$ for every commonly
feasible $(x,y)$, which is unifilar lumpability with
$\tau_{\mathcal K}(\phi^{\infty}(s),x,y)=\phi^{\infty}(\tau(s,x,y))$; and
$\phi^{\infty}$ refines $\phi_{\ker}=\phi^{(0)}$ by construction.  So
$\phi^{\infty}$ is admissible.

For maximality, let $\psi$ be any admissible partition.  We show by induction
that $\psi$ refines $\phi^{(m)}$ for every $m$.  For $m=0$ this is the
hypothesis that $\psi$ refines $\phi_{\ker}$.  If $\psi$ refines $\phi^{(m)}$
and $\psi(s)=\psi(s')$, then $\phi^{(m)}(s)=\phi^{(m)}(s')$, and unifilar
lumpability of $\psi$ gives $\psi(\tau(s,x,y))=\psi(\tau(s',x,y))$ at every
commonly feasible event, hence
$\phi^{(m)}(\tau(s,x,y))=\phi^{(m)}(\tau(s',x,y))$ by the inductive
hypothesis; the two defining data of $\phi^{(m+1)}$ therefore agree, so
$\phi^{(m+1)}(s)=\phi^{(m+1)}(s')$.  Taking $m$ at the stabilization index
shows $\psi$ refines $\phi^{\infty}$.  Thus $\phi^{\infty}$ is the coarsest
admissible partition, and uniqueness follows.
\end{proof}

\begin{theorem}[Zero Controlled Retention]
\label{thm:controlled-zero}
$\RetKLc(M)=0$ if and only if $M\ge N^{\ast}$, with $N^{\ast}$
the index of Definition~\ref{def:kernel-refinement}.
\end{theorem}

\begin{proof}
Suppose $\RetKLc(\phi)=0$ for a unifilar-lumpable $\phi$.  The
cost is a sum of nonnegative terms $p(x)\pi_sD_{\mathrm{KL}}(P_s^x\|\bar P_k^x)$,
so every term vanishes; by the equality condition for KL divergence,
$P_s^x=\bar P_{\phi(s)}^x$ for every positive-mass $s$ and every $x$ in the
input support.  Two states in a common block therefore have equal kernels, so
$\phi$ refines $\phi_{\ker}$, and being unifilar-lumpable it is admissible in
the sense of Proposition~\ref{prop:kernel-refinement-exists}.  Hence $\phi$
refines $\phi_{\ker}^{\,\ast}$ and has at least $N^{\ast}$ blocks.

Conversely $\phi_{\ker}^{\,\ast}$ is itself unifilar-lumpable and refines
$\phi_{\ker}$, so each of its blocks carries a single kernel, every block
centroid equals that kernel, and the cost is zero.  It has $N^{\ast}$ blocks,
so $M\ge N^{\ast}$ suffices.
\end{proof}

\begin{remark}[The Threshold Is Not the Number of Distinct Kernels]
\label{rem:controlled-zero-not-kernels}
In the input-driven setting of Theorem~\ref{thm:retention-zero} the threshold
is the number of distinct predictive laws, because the predictive partition is
automatically lumpable there.  In the unifilar class the two quantities
separate, and the gap $N^{\ast}-|\phi_{\ker}(\Splus)|$ can be positive.  Take $\mathcal I$ a singleton, $\mathcal O=\{0,1\}$, three states
$A,B,C$, kernels
\[
P_A=P_B=\mathrm{Bernoulli}(\tfrac12),
\qquad
P_C=\mathrm{Bernoulli}(\tfrac1{10}),
\]
and the unifilar update
\[
\tau(A,0)=A,\quad\tau(A,1)=C,\qquad
\tau(B,0)=C,\quad\tau(B,1)=B,\qquad
\tau(C,0)=A,\quad\tau(C,1)=B .
\]
The stationary distribution is strictly positive.  There are two distinct
kernels, but the partition $\{\{A,B\},\{C\}\}$ is not unifilar-lumpable: $A$ and
$B$ are in one block and the output $0$ is feasible from both, yet
$\tau(A,0)=A$ and $\tau(B,0)=C$ lie in different blocks.  The refinement of
Proposition~\ref{prop:kernel-refinement-exists} separates $A$ from $B$, so
$N^{\ast}=3$ while $|\phi_{\ker}(\Splus)|=2$, and indeed
$\RetKLc(2)>0$ whereas $\RetKLc(3)=0$.  The
separation is not exceptional: over randomly generated unifilar machines with
at most five states and at least one coincidence of predictive kernels, the
kernel partition fails to be unifilar-lumpable in approximately fifty-five per
cent of instances.

This is exactly the distinction between predictive equivalence and causal-state
equivalence recorded in Definition~\ref{def:z-predictive-equivalence}: the
refinement $\phi_{\ker}^{\,\ast}$ is the causal-state partition for the full
controlled future, and $\phi_{\ker}$ is its one-step coarsening.
\end{remark}

The separation of Remark~\ref{rem:controlled-zero-not-kernels} is as large as
the state set permits, and the refinement recursion is essentially as slow as
Proposition~\ref{prop:kernel-refinement-exists} allows.  Both are witnessed by
one family.

\begin{definition}[Counter Family $\mathcal C_M$]
\label{def:counter-family}
For $M\ge3$ and $\beta,\gamma\in(0,1)$ with $\beta\neq\gamma$, let
$\mathcal C_M$ be the unifilar machine with a single input letter, output
alphabet $\mathcal O=\{0,1\}$, state set $\{0,1,\dots,M-1\}$, update
\[
\tau(s,0)=\min\{s+1,\,M-1\},
\qquad
\tau(s,1)=0,
\]
and emission laws
\[
P_s=\mathrm{Bernoulli}(\beta)\ \ (0\le s\le M-2),
\qquad
P_{M-1}=\mathrm{Bernoulli}(\gamma).
\]
Output $0$ advances the counter and saturates at $M-1$; output $1$ resets it.
The restriction to the open interval is required for
Proposition~\ref{prop:refinement-extremal}(i): at $\gamma=0$ the state $M-1$
emits $0$ almost surely and, being saturated, is absorbing, while at $\beta=1$
every state below $M-1$ resets almost surely and $M-1$ is never reached.  In
either case the chain fails to be irreducible and some state leaves the
stationary support.  The input alphabet being a singleton, $X_t$ is
deterministic and hence independent of $S_t$, so
Theorem~\ref{thm:controlled-ib} applies and $\RetKLg$ coincides with
$\RetKLc$ on this family.
\end{definition}

\begin{proposition}[The Refinement Gap Is Maximal and the Bound Is Tight]
\label{prop:refinement-extremal}
Let $\beta,\gamma\in(0,1)$ with $\beta\neq\gamma$.  For $\mathcal C_M$:
\begin{enumerate}[label=(\roman*)]
\item the chain is irreducible with strictly positive stationary
distribution, so $\Splus$ is the whole state set;
\item the kernel partition has exactly two blocks,
$|\phi_{\ker}(\Splus)|=2$;
\item the stable refinement is the discrete partition, $N^{\ast}=M$, so the
gap $N^{\ast}-|\phi_{\ker}(\Splus)|=M-2$ is the largest possible on $M$
states;
\item the refinement recursion of
Proposition~\ref{prop:kernel-refinement-exists} stabilizes after exactly
$M-1$ steps, so its bound of $|\Splus|$ steps is tight to within one;
\item consequently $\RetKLc(M-1)>0$ while $\RetKLc(M)=0$.
\end{enumerate}
\end{proposition}

\begin{proof}
\emph{(i)} Every state emits both outputs with positive probability: states
$0,\dots,M-2$ because $\beta\in(0,1)$, and state $M-1$ because
$\gamma\in(0,1)$.  Consequently output $1$ carries any state to $0$, and a run
of $M-1$ outputs $0$ carries $0$ to $M-1$, so every state reaches every other.
The chain is therefore irreducible on a finite set, and its stationary
distribution is strictly positive, so $\Splus=\{0,\dots,M-1\}$.

\emph{(ii)} States $0,\dots,M-2$ carry the law $\mathrm{Bernoulli}(\beta)$ and
state $M-1$ carries $\mathrm{Bernoulli}(\gamma)\neq\mathrm{Bernoulli}(\beta)$,
giving exactly two kernel classes.

\emph{(iii)} and \emph{(iv)} Write $\phi^{(m)}$ for the $m$-th iterate, with
$\phi^{(0)}=\phi_{\ker}$. We claim, by induction on $m\ge0$, that $\phi^{(m)}$
separates exactly the top $m+1$ states, that is
\[
\phi^{(m)}(s)=\phi^{(m)}(s')
\iff
s=s'
\ \text{ or }\
s,s'\le M-2-m .
\]
For $m=0$ this is~(ii). Assume it for $m$ and let $s<s'\le M-2-m$. The
recursion refines by the pair of successor classes
$\bigl(\phi^{(m)}(\tau(s,0)),\phi^{(m)}(\tau(s,1))\bigr)$. Since
$s,s'\le M-2-m\le M-2$, neither is saturated, so $\tau(s,0)=s+1$ and
$\tau(s',0)=s'+1$, while $\tau(s,1)=\tau(s',1)=0$. By the inductive hypothesis
$s+1$ and $s'+1$ are separated precisely when at least one of them exceeds
$M-2-m$, that is precisely when $s'=M-2-m$. Hence $\phi^{(m+1)}$ splits
$M-2-m$ off from the states below it and leaves $0,\dots,M-3-m$ together,
which is the claim for $m+1$.

At $m=M-2$ the block $\{s:s\le M-2-m\}=\{0\}$ is a singleton, so
$\phi^{(M-2)}$ is the discrete partition and $N^{\ast}=M$, giving~(iii). The
recursion detects stabilization one step later, at $m=M-1$, since
$\phi^{(M-2)}$ and $\phi^{(M-1)}$ have the same number of blocks while
$\phi^{(M-3)}$ has fewer; the number of iterations is therefore exactly $M-1$,
which is~(iv). The bound $N^{\ast}\le|\Splus|$ of
Proposition~\ref{prop:kernel-refinement-exists} is attained, and no partition
of an $M$-element set has more than $M$ blocks, so the gap $M-2$ is maximal.

\emph{(v)} Immediate from Theorem~\ref{thm:controlled-zero} and $N^{\ast}=M$.
\end{proof}

\begin{remark}[Scope of the Extremal Family]
\label{rem:refinement-extremal-scope}
Proposition~\ref{prop:refinement-extremal} settles the worst case, not the
typical one. It shows that no bound of the form
$N^{\ast}\le f(|\phi_{\ker}(\Splus)|)$ can hold, since $\mathcal C_M$ has two
kernel classes for every $M$, and that the round count of
Proposition~\ref{prop:kernel-refinement-exists} cannot be improved below
$|\Splus|-1$. It uses a single input letter, so the phenomenon is driven
entirely by the output-dependence of the update and is invisible to any
input-driven analysis. For $M=3,\dots,7$ the values $\RetKLc(M-1)$ are
$0.0481$, $0.0321$, $0.0192$, $0.0107$ and $0.0057$ nats at
$\beta=\tfrac12$, $\gamma=\tfrac1{10}$, decreasing in $M$ as the extra
resolved state carries less stationary mass.
\end{remark}

\begin{corollary}[Recovery of the Input-Driven Theory]
\label{cor:controlled-reduces}
Let the emissions be output-independent, $P_s^x=P_s$ for every $x$, as in every
construction of Section~\ref{subsec:retention-complexity}.  Then
$\RetKLc(\phi)=\RetKL(\phi)$ for every quotient $\phi$ feasible
in both theories.  If moreover the machine is input-driven and every quotient
under consideration has connected support --- in particular if every $P_s$ has
full support --- the two feasible sets coincide by
Proposition~\ref{prop:input-driven-specialization}(i), and therefore
$\RetKLc(M)=\RetKL(M)$ for every $M$, the controlled identity
reducing to Meta-Theorem~\ref{meta:ib}.
\end{corollary}

\begin{proof}
The first claim is immediate: with $P_s^x=P_s$ the centroids $\bar P_k^x$ are
independent of $x$ and $\sum_xp(x)=1$.  The second follows because equality of
the two minima requires equality of the feasible sets over which they are
taken; Remark~\ref{rem:controlled-reduces-sharp} shows that this hypothesis
cannot be dropped.
\end{proof}

\begin{remark}[Why the Support Hypothesis Cannot Be Dropped Here]
\label{rem:controlled-reduces-sharp}
Without a support hypothesis the two optimized gaps genuinely differ, even for
output-independent emissions, because the unifilar feasible set is strictly
larger.  The three-state witness of
Remark~\ref{rem:unifilar-support-not-automatic} has
$\RetKLc(2)<\RetKL(2)$, the two values differing by more than
$0.15$ nats.  Corollary~\ref{cor:controlled-reduces} is therefore stated with
the hypothesis, and the complexity transfer of
Remark~\ref{rem:complexity-transfer} is argued from full supports rather than
from the general reduction.
\end{remark}

\begin{corollary}[Controlled Quadratic Spectral Converse]
\label{cor:controlled-quad-spectral}
In the controlled Gaussian quadratic regime $P_s^x=\mathcal N(m_s^x,\tfrac12I_D)$,
with Fisher features $y_s^x=\sqrt2\,m_s^x$ and per-input covariances
\[
\Sigma_\pi^x=\sum_s\pi_s(y_s^x-\bar y^x)(y_s^x-\bar y^x)^{\!\top},
\qquad
\bar y^x=\sum_s\pi_sy_s^x,
\]
the controlled quadratic gap satisfies
\[
\RetQuadc(M)
\ \ge\
\frac12\sum_xp(x)\sum_{i\ge M}\lambda_i(\Sigma_\pi^x).
\]
\end{corollary}

\begin{proof}
The lumpability constraint does not depend on $x$, so a single quotient is used
on every fiber.  Fix $x$: the inner cost is the Gaussian quadratic restricted
gap of the family $\{y_s^x\}$ under that quotient, and the ANOVA-plus-Ky-Fan
argument of Theorem~\ref{thm:spectral-retention} bounds it below by
$\tfrac12\sum_{i\ge M}\lambda_i(\Sigma_\pi^x)$, the effective rank of the
between-block covariance being $M-1$ on each fiber.  Multiplying by $p(x)$ and
summing gives the bound, and taking the minimum over quotients preserves it
since the bound is quotient-free.
\end{proof}

\begin{corollary}[Controlled Probability-Coordinate Spectral Converse]
\label{cor:controlled-simplex-spectral}
For finite $\mathcal O$, let $p_s^x$ be the predictive probability vectors and
$\Sigma_p^x$ their $\pi$-covariance on fiber $x$.  Then, globally and with no
regularity hypothesis,
\[
\RetKLc(M)
\ \ge\
\sum_xp(x)\sum_{i\ge M}\lambda_i(\Sigma_p^x),
\]
and the constant $1$ is sharp on any fiber of positive mass.
\end{corollary}

\begin{proof}
Apply Theorem~\ref{thm:global-kl-simplex} fiberwise, exactly as in
Corollary~\ref{cor:controlled-quad-spectral}; the inequality $(*)$ of that
proof is valid at boundary points, so no interiority is needed.  Sharpness
follows by placing the two-point family of
Proposition~\ref{prop:kl-simplex-sharp} on a single fiber with $p(x)>0$ and
making the emissions constant in $x$ elsewhere.
\end{proof}

\begin{corollary}[Controlled Interior Fisher Converse]
\label{cor:controlled-fisher}
Suppose that for each $x$ the strictly positive laws $\{p_s^x\}$ lie in a fixed
minimal natural-parameter chart with a compact convex parameter set $K^x$
containing the state parameters \emph{and} their mixture-centroid parameters,
and that $\nabla^2A(\eta)\succeq m_{K}^{x}I$ on $K^x$.  Then
\[
\RetKLc(M)
\ \ge\
\frac12\Bigl(\min_xm_{K}^{x}\Bigr)\sum_xp(x)\sum_{i\ge M}\lambda_i(\Sigma_\eta^x),
\]
where $\Sigma_\eta^x$ is the $\pi$-covariance of the natural parameters on
fiber $x$.
\end{corollary}

\begin{proof}
Theorem~\ref{thm:global-interior-fisher} applied on each fiber, with the worst
fiber constant taken out of the $p$-weighted sum.  As there, the hypothesis
that $K^x$ contains the centroid parameters is what places the segments
$[\eta_s,\zeta_C]$ inside $K^x$.
\end{proof}

\begin{remark}[The Fisher No-Go Persists Under Control]
\label{rem:controlled-fisher-nogo}
The counterexample of Theorem~\ref{thm:no-global-fisher-converse} lives on a
single fiber, so it is a controlled instance with constant emissions.  No
universal constant exists for a natural-parameter converse in the controlled
theory either: interiority is necessary exactly as in the unconditional case,
and Corollary~\ref{cor:controlled-simplex-spectral} is the boundary-safe
statement.
\end{remark}

\begin{remark}[Complexity Results Transfer to the Unifilar Class]
\label{rem:complexity-transfer}
Every machine appearing in Theorems~\ref{thm:retention-reset-np},
\ref{thm:retention-definite-np}, \ref{thm:full-kl-promise-np} and
Corollary~\ref{cor:full-kl-apx} is input-driven with output-independent
emissions, hence a valid instance of the unifilar class.  Transfer of hardness
requires an argument, since enlarging the feasible set of quotients can only
lower the optimum and could in principle convert a NO instance into a YES
instance.  It is valid here because in each construction the two feasible sets
coincide.

In the reset construction of Theorem~\ref{thm:retention-reset-np}, and in the
reset machine reused in Theorem~\ref{thm:full-kl-promise-np}, \emph{every}
partition of the state set is already lumpable, since states in a common block
have identical successors $\tau(s_i,\ell)=s_\ell$; the unifilar feasible set,
being a subset of the set of all partitions, therefore coincides with the
lumpable one.  In the depth-two construction of
Theorem~\ref{thm:retention-definite-np} not every partition is lumpable, but
the NO-case bound is proved against an arbitrary machine-state variable $K$
with $|K|\le k$ through the fractional $k$-means Lemma~\ref{lem:fractional-kmeans};
a fractional assignment dominates every partition, lumpable or not, so the
lower bound already covers the unifilar feasible set.  In
Theorem~\ref{thm:full-kl-promise-np} the predictive laws $p_i=u+\delta z_i$
have all coordinates strictly positive, so supports are full and block-uniform
for every partition, and Proposition~\ref{prop:input-driven-specialization}(i)
applies directly.

The NP-completeness, promise NP-hardness, and APX-hardness statements therefore
hold verbatim when the retention decision problem ranges over unifilar
machines; certificates and membership are unchanged.
\end{remark}

%----------------------------------------------------------------------
\subsection{Computational Complexity of Retention}
\label{subsec:retention-complexity}
%----------------------------------------------------------------------

The complexity results below are theorems of the Gaussian quadratic restricted
retention regime.  They do not establish NP-hardness for unrestricted full-KL
retention.

\begin{lemma}[Polynomial Gap Amplification for Rational $k$-Means]
\label{lem:kmeans-gap}
Given a rational Euclidean $k$-means instance
\[
v_1,\dots,v_n\in\mathbb Q^d,
\]
integer $k$, and rational threshold $\tau$, there is a polynomial-time
transformation to integer points
\[
a_1,\dots,a_n\in\mathbb Z^D
\]
and a rational threshold $T$ such that exactly one of the following holds:
\[
\mathrm{OPT}_a\le T,
\]
or
\[
\mathrm{OPT}_a\ge T+2n.
\]
The bit complexity of the transformed instance is polynomial in the original
bit complexity.
\end{lemma}

\begin{proof}
Clear denominators to obtain integer points $b_i$ and rational threshold
$\tau'$.  Let
\[
D_0=\operatorname{lcm}(1,2,\dots,n)\cdot \operatorname{denom}(\tau').
\]
For integer points, the ordinary $k$-means objective of a cluster $C$ of size
$m$ is
\[
W(C)
=
\sum_{i\in C}\|b_i\|^2
-
\frac1m
\left\|
\sum_{i\in C}b_i
\right\|^2.
\]
The first term is an integer.  The second term has denominator dividing $m$.
Therefore $mW(C)\in\mathbb Z$, and since $D_0$ is a multiple of every
$m\le n$,
\[
D_0 W(C)\in\mathbb Z.
\]
Summing over clusters shows that every $k$-means objective value lies in
\[
\frac1{D_0}\mathbb Z.
\]
Let
\[
q=\lfloor D_0\tau'\rfloor.
\]
Because all objective values lie on the grid $\frac1{D_0}\mathbb Z$,
\[
\mathrm{OPT}_b\le\tau'
\implies
\mathrm{OPT}_b\le \frac{q}{D_0},
\]
while
\[
\mathrm{OPT}_b>\tau'
\implies
\mathrm{OPT}_b\ge \frac{q+1}{D_0}.
\]
Now scale all coordinates by an integer
\[
L=\left\lceil \sqrt{2nD_0}\right\rceil,
\]
and set
\[
a_i=Lb_i.
\]
Every $k$-means objective scales by $L^2$.  Define
\[
T=\frac{L^2 q}{D_0}.
\]
Then in the YES case,
\[
\mathrm{OPT}_a
\le
\frac{L^2 q}{D_0}
=
T,
\]
while in the NO case,
\[
\mathrm{OPT}_a
\ge
\frac{L^2(q+1)}{D_0}
=
T+\frac{L^2}{D_0}
\ge
T+2n.
\]
The bit length of $D_0$ is polynomial because
\[
\log D_0
=
O(\log\operatorname{denom}(\tau'))+O(\log\operatorname{lcm}(1,\dots,n)),
\]
and $\log\operatorname{lcm}(1,\dots,n)=O(n)$.  Therefore $\log L$ is
polynomial, and the scaled instance has polynomial bit complexity.
\end{proof}

\begin{lemma}[Fractional Assignment Cannot Beat Hard $k$-Means]
\label{lem:fractional-kmeans}
Let $a_1,\dots,a_n$ be points in Euclidean space.  For any fractional
assignment weights
\[
w_{jr}\ge0,
\qquad
\sum_{r=1}^k w_{jr}=1,
\]
and any centers $c_1,\dots,c_k$, define the fractional objective
\[
F
=
\frac1n
\sum_{j=1}^n
\sum_{r=1}^k
w_{jr}
\|a_j-c_r\|^2.
\]
Let $\mathrm{OPT}$ be the optimal hard $k$-means objective.  Then
\[
F\ge \mathrm{OPT}.
\]
\end{lemma}

\begin{proof}
For fixed centers $c_1,\dots,c_k$, define the hard nearest-center assignment
\[
r(j)\in\argmin_r\|a_j-c_r\|^2.
\]
Then
\[
\frac1n
\sum_{j=1}^n
\|a_j-c_{r(j)}\|^2
\le
F.
\]
Recomputing the centers as the means of the assigned clusters cannot increase
the hard $k$-means objective.  Therefore
\[
\mathrm{OPT}
\le
\frac1n
\sum_{j=1}^n
\|a_j-c_{r(j)}\|^2
\le
F.
\]
\end{proof}

\begin{theorem}[Reset-Vacuity NP-Completeness]
\label{thm:retention-reset-np}
In the Gaussian quadratic restricted retention regime with common covariance
$\frac12 I$, the decision problem
\[
\RetQuad(M)<\theta
\]
is NP-complete for rational data.  NP-hardness holds even for minimal definite
machines of synchronization depth one.
\end{theorem}

\begin{proof}
Reduce from Euclidean $k$-means \cite{aloise2009}.  By Lemma~\ref{lem:kmeans-gap}, assume we are
given integer points
\[
a_1,\dots,a_n\in\mathbb Z^d,
\]
integer $k$, and rational threshold $T$ such that either
\[
\mathrm{OPT}_a\le T
\]
or
\[
\mathrm{OPT}_a\ge T+2n.
\]

Construct a retention instance as follows.  Let
\[
\mathcal S=\{s_1,\dots,s_n\},
\qquad
\mathcal I=\{1,\dots,n\}.
\]
Use uniform i.i.d.\ inputs.  Define reset transitions
\[
\tau(s_i,\ell)=s_\ell.
\]
After one input, the state equals the input symbol.  Hence the synchronization
depth is one.

Let $e_i$ be the $i$-th standard basis vector in $\mathbb R^n$, and set
\[
\eta=\frac18.
\]
Define Gaussian means
\[
m_i=(a_i,\eta e_i)\in\mathbb R^{d+n}.
\]
Let emissions be
\[
P_i=\mathcal N(m_i,\tfrac12 I).
\]
The stationary distribution is uniform on $\mathcal S$.

Every partition of $\mathcal S$ is lumpable.  Indeed, if $s_i,s_j$ lie in the
same block, then for every input $\ell$,
\[
\tau(s_i,\ell)=s_\ell=\tau(s_j,\ell),
\]
so the successors are identical and therefore lie in the same block.

For a partition $\mathcal C=\{C_1,\dots,C_K\}$ with $K\le k$, the quadratic
retention cost is
\[
\RetQuad(\mathcal C)
=
\frac1n
\sum_{r=1}^K
\sum_{i\in C_r}
\|m_i-\bar m_r\|^2.
\]
Decompose into main and tag coordinates.

The main-coordinate cost is
\[
\frac1n
\sum_{r=1}^K
\sum_{i\in C_r}
\|a_i-\bar a_r\|^2,
\]
which is exactly the ordinary $k$-means objective divided by $n$.

For a block of size $m$, the tag-coordinate cost is
\[
\frac1n
\sum_{i\in C_r}
\left\|
\eta e_i-\frac{\eta}{m}\sum_{j\in C_r}e_j
\right\|^2
=
\frac{\eta^2}{n}(m-1).
\]
Summing over blocks gives
\[
\text{tag cost}
=
\eta^2\frac{n-K}{n}
\le
\eta^2
=
\frac1{64}.
\]

If $\mathrm{OPT}_a\le T$, choose an optimal $k$-clustering.  Then
\[
\RetQuad(k)
\le
\frac{T}{n}+\frac1{64}.
\]
If $\mathrm{OPT}_a\ge T+2n$, then for every partition into at most $k$
blocks, the main-coordinate cost is at least
\[
\frac{\mathrm{OPT}_a}{n}
\ge
\frac{T+2n}{n}
=
\frac{T}{n}+2.
\]
The tag cost is nonnegative, so
\[
\RetQuad(k)
\ge
\frac{T}{n}+2.
\]
Set
\[
\theta=\frac{T}{n}+1.
\]
Then
\[
\mathrm{OPT}_a\le T
\implies
\RetQuad(k)<\theta,
\]
while
\[
\mathrm{OPT}_a\ge T+2n
\implies
\RetQuad(k)>\theta.
\]
Thus the strict decision problem is NP-hard.

Membership in NP follows because a certificate is a partition of the $n$
states.  For each block, the centroid is a rational vector whose coordinates
have polynomial bit length.  The quadratic cost is a sum of squared Euclidean
distances with rational coefficients.  The tag contribution is
\[
\eta^2\frac{n-K}{n},
\qquad
\eta=\frac18,
\]
which is rational with constant denominator.  The total cost can therefore be
computed exactly by rational arithmetic in time polynomial in $n$, $d$, and
the input bit length, and compared with the rational threshold $\theta$.
Hence the problem is NP-complete.
\end{proof}

\begin{theorem}[NP-Hardness for Minimal Definite Depth-$2$ $\epsilon$-Machines]
\label{thm:retention-definite-np}
In the Gaussian quadratic restricted retention regime with common covariance
$\frac12 I$, the decision problem
\[
\RetQuad(M)<\theta
\]
is NP-hard even when restricted to minimal definite $\epsilon$-machines of
synchronization depth $2$.  With rational means and rational threshold, it is
NP-complete.
\end{theorem}

\begin{proof}
Use the tagged depth-$2$ construction.  By Lemma~\ref{lem:kmeans-gap}, assume
we are given integer points $a_1,\dots,a_n$, integer $k$, and rational
threshold $T$ with either
\[
\mathrm{OPT}_a\le T
\]
or
\[
\mathrm{OPT}_a\ge T+2n.
\]

Let the input alphabet be
\[
\mathcal I=\{1,\dots,n\},
\]
with uniform i.i.d.\ inputs.  Create states
\[
s_{ij},
\qquad
i,j\in\{1,\dots,n\}.
\]
Define the transition
\[
\tau(s_{ij},\ell)=s_{j\ell}.
\]
After two inputs, the state is determined by the last two inputs, independently
of the initial state.  Hence the synchronization depth is $2$, and the machine
is definite of depth $2$.

Let $e_{ij}$ be the standard basis vector indexed by $(i,j)$ in
$\mathbb R^{n^2}$, and set
\[
\eta=\frac18.
\]
Define means
\[
m_{ij}=(a_j,\eta e_{ij})
\in
\mathbb R^{d+n^2}.
\]
Let the emission be
\[
P_{ij}=\mathcal N(m_{ij},\tfrac12 I).
\]
Because the covariance is $\frac12 I$,
\[
D_{\mathrm{KL}}(P_{ij}\|P_c)
=
\norm{m_{ij}-c}^2
\]
for Gaussian representatives $P_c=\mathcal N(c,\frac12 I)$.

The means $m_{ij}$ are all distinct because the tag coordinates
$\eta e_{ij}$ are distinct.  Therefore the one-step predictive distributions
are all distinct, and the minimal $\epsilon$-machine has $n^2$ causal states.

Suppose first that $\mathrm{OPT}_a\le T$.  Let $C_1,\dots,C_k$ be a
$k$-clustering of $\{1,\dots,n\}$ achieving objective at most $T$.  Define
\[
D_r=\{s_{ij}:j\in C_r\}.
\]
This partition is lumpable: if $s_{ij},s_{i'j'}\in D_r$, then
$j,j'\in C_r$, and for every input $\ell$,
\[
\tau(s_{ij},\ell)=s_{j\ell},
\qquad
\tau(s_{i'j'},\ell)=s_{j'\ell},
\]
both of which have second coordinate $\ell$, hence lie in the same block
$D_{c(\ell)}$, where $c(\ell)$ is the cluster containing $\ell$.

The stationary distribution on pairs $(i,j)$ is uniform.  The main-coordinate
cost is therefore
\[
\frac1n
\sum_{r=1}^k
\sum_{j\in C_r}
\norm{a_j-\bar a_r}^2
\le
\frac{T}{n}.
\]
The tag-coordinate cost is at most $1/64$.  Hence
\[
\RetQuad(k)
\le
\frac{T}{n}+\frac1{64}.
\]

Now suppose that $\mathrm{OPT}_a\ge T+2n$.  Consider any $k$-state
approximation.  It induces a machine-state variable $K$ with $|K|\le k$.
Decompose the cost into main and tag coordinates.  The tag cost is
nonnegative.

For the main coordinate, let $J$ be the second state index.  Under the uniform
stationary distribution, $J$ is uniform on $\{1,\dots,n\}$.  Let
\[
w_{jr}=\mathbb P(K=r\mid J=j).
\]
The main cost is
\[
\frac1n
\sum_{j=1}^n
\sum_{r=1}^k
w_{jr}
\norm{a_j-c_r^{\mathrm{main}}}^2.
\]
This is a fractional $k$-means objective on the points $a_j$, multiplied by
$1/n$.  By Lemma~\ref{lem:fractional-kmeans}, fractional assignment cannot
beat the optimal hard $k$-means objective.  Therefore the main cost is at least
\[
\frac{\mathrm{OPT}_a}{n}
\ge
\frac{T+2n}{n}
=
\frac{T}{n}+2.
\]
Set
\[
\theta=\frac{T}{n}+1.
\]
Then YES implies
\[
\RetQuad(k)<\theta,
\]
while NO implies
\[
\RetQuad(k)>\theta.
\]
Thus the retention decision problem is NP-hard.  Rational parameters give
NP-completeness by the same certificate argument as in
Theorem~\ref{thm:retention-reset-np}.
\end{proof}

\begin{theorem}[Promise NP-Hardness of Finite-Alphabet Full-KL Retention]
\label{thm:full-kl-promise-np}
Finite-alphabet full-KL retention is NP-hard under a promise, already for
synchronization-depth-one reset machines in which every state partition is
lumpable.  Precisely, there are polynomial-time computable rationals
$\theta_{\mathrm{yes}}<\theta_{\mathrm{no}}$ such that deciding
\[
\RetKL(k)\le\theta_{\mathrm{yes}}
\quad\text{versus}\quad
\RetKL(k)\ge\theta_{\mathrm{no}}
\]
is NP-hard, the promise being that one of the two holds.
\end{theorem}

\begin{proof}
Reduce from Euclidean $k$-means.  By Lemma~\ref{lem:kmeans-gap} we may assume
integer points $a_1,\dots,a_n\in\mathbb Z^{d}$, an integer $k$, and a rational
$T$ with either $\mathrm{OPT}_a\le T$ or $\mathrm{OPT}_a\ge T+2n$, all of
polynomial bit complexity.

\emph{The machine.}  Take $\mathcal S=\{s_1,\dots,s_n\}$ and
$\mathcal I=\{1,\dots,n\}$ with reset transitions $\tau(s_i,\ell)=s_\ell$ and
uniform i.i.d.\ inputs, exactly as in Theorem~\ref{thm:retention-reset-np}.
The synchronization depth is one, the stationary distribution is uniform, and
every partition is lumpable, since $s_i,s_j$ in a common block have identical
successors $\tau(s_i,\ell)=s_\ell=\tau(s_j,\ell)$.

\emph{The output alphabet and the tangent embedding.}  Let
$\mathcal O=\{1,\dots,2d\}$, so predictive laws are points of
$\Delta^{2d-1}$.  Embed each $a_i$ in the tangent space of the simplex at the
uniform law by doubling coordinates with opposite signs:
\[
z_i=\bigl(a_{i1},-a_{i1},\ a_{i2},-a_{i2},\ \dots,\ a_{id},-a_{id}\bigr)
\in\mathbb Z^{2d},
\qquad
\sum_{j}(z_i)_j=0 .
\]
The doubling is what makes the embedding \emph{centred}, hence tangent to the
simplex, and it preserves the Euclidean geometry up to a factor of two:
\[
\|z_i-z_{i'}\|_2^2=2\|a_i-a_{i'}\|_2^2 .
\]
Let $u=(1/2d,\dots,1/2d)$ and set
\[
p_i=u+\delta z_i,
\qquad
\delta=2^{-\Lambda},
\]
with $\Lambda$ a positive integer fixed below.  Each $p_i$ is a rational
probability vector, and for $\delta<\bigl(2d\max_{i,j}|(z_i)_j|\bigr)^{-1}$ all
coordinates are strictly positive, so every $p_i$ lies in a fixed compact
subset of the interior.  Assign $P_i=p_i$ as the predictive law of state $s_i$.

\emph{The objective is a Jensen--Shannon cost.}  For a block $C$ with uniform
weights $w_i=1/|C|$, Lemma~\ref{lem:mixture-centroid} makes the optimal
representative the mixture centroid $\bar p_C$, so the block cost is
\[
J_C=\sum_{i\in C}w_i\,D_{\mathrm{KL}}(p_i\|\bar p_C),
\]
the Jensen--Shannon divergence of the block, and
$\RetKL(k)=\min_{|\mathcal C|\le k}\sum_C\frac{|C|}{n}J_C$.

\emph{Uniform quadratic expansion.}  Write $\bar z_C=\sum_{i\in C}w_iz_i$, so
$\bar p_C=u+\delta\bar z_C$.  Since all points lie in a compact interior set,
a third-order Taylor expansion of $D_{\mathrm{KL}}$ about $u$, with the linear
terms vanishing because $\sum_j(z_i-\bar z_C)_j=0$, gives
\[
J_C
=
d\,\delta^2\sum_{i\in C}w_i\|z_i-\bar z_C\|_2^2
+
\rho_C,
\qquad
|\rho_C|\ \le\ C_0\,d^{2}\delta^{3}Z^{3},
\]
where $Z=\max_i\|z_i\|_1$ and $C_0$ is absolute; the Hessian of
$p\mapsto\sum_jp_j\log p_j$ at $u$ is $2d\cdot I$ on the tangent space, which
supplies the factor $d$ after the $\tfrac12$ of the quadratic term.  The
remainder bound is uniform over blocks and partitions because it depends only
on $d$, $\delta$ and $Z$.  By the geometry of the embedding,
\[
\sum_{i\in C}w_i\|z_i-\bar z_C\|_2^2
=
2\sum_{i\in C}w_i\|a_i-\bar a_C\|_2^2 ,
\]
so the leading term is $2d\delta^2$ times the $k$-means cost of $C$.
Summing over blocks with weights $|C|/n$,
\[
\Bigl|\ \RetKL(k)-\frac{2d\delta^2}{n}\,\mathrm{OPT}_a\ \Bigr|
\ \le\
C_0\,d^{2}\delta^{3}Z^{3},
\]
the same bound applying to the value of every individual partition, so that the
optimizing partitions of the two objectives need not coincide for the estimate
to hold.

\emph{Choice of $\delta$ and the thresholds.}  The promise gap in the
$k$-means objective is $2n$, so the gap in the leading term is
$2d\delta^2\cdot 2n/n=4d\delta^2$.  Choose $\Lambda$ polynomial in the input
bit length so that
\[
C_0\,d^{2}\delta^{3}Z^{3}\ <\ d\,\delta^{2},
\]
which holds as soon as $\delta<\bigl(C_0dZ^3\bigr)^{-1}$; since $Z$ has
polynomially many bits, such a $\Lambda$ is polynomial and $\delta=2^{-\Lambda}$
is computable in polynomial time.  Set
\[
\theta_{\mathrm{yes}}=\frac{2d\delta^2}{n}T+d\delta^2,
\qquad
\theta_{\mathrm{no}}=\frac{2d\delta^2}{n}(T+2n)-d\delta^2
=
\theta_{\mathrm{yes}}+2d\delta^{2},
\]
both rational with polynomially many bits and
$\theta_{\mathrm{yes}}<\theta_{\mathrm{no}}$.  In the YES case
$\mathrm{OPT}_a\le T$ gives $\RetKL(k)\le\theta_{\mathrm{yes}}$; in the NO case
$\mathrm{OPT}_a\ge T+2n$ gives $\RetKL(k)\ge\theta_{\mathrm{no}}$.  The
reduction is polynomial, and the promise is met in both cases.
\end{proof}

\begin{corollary}[APX-Hardness of Finite-Alphabet Full-KL Retention]
\label{cor:full-kl-apx}
There is a constant $\varepsilon_0>0$ such that approximating $\RetKL(k)$ to
within a factor $1+\varepsilon_0$ is NP-hard, for finite output alphabets and
already for synchronization-depth-one reset machines in which every partition
is lumpable.  In particular the problem admits no PTAS unless $\mathsf P=\mathsf{NP}$.
\end{corollary}

\begin{proof}
Euclidean $k$-means is APX-hard: there is $\varepsilon_1>0$ such that
approximating it within $1+\varepsilon_1$ is NP-hard \cite{awasthi2015}, with
the constant subsequently improved \cite{lee2017}.  We transfer this through the
embedding of Theorem~\ref{thm:full-kl-promise-np}, for which it suffices to
control the objective \emph{relatively} rather than only across a promise gap.

Fix an instance $a_1,\dots,a_n\in\mathbb Z^d$ and let $\Xi$ denote its
$k$-means objective, so that for a partition $\mathcal C$ with at most $k$
blocks the embedded cost obeys, uniformly in $\mathcal C$,
\[
\RetKL(\mathcal C)
=
\frac{2d\delta^2}{n}\,\Xi(\mathcal C)\;+\;\rho(\mathcal C),
\qquad
|\rho(\mathcal C)|\ \le\ C_0d^2\delta^3Z^3 ,
\]
by the uniform third-order expansion established in the proof of
Theorem~\ref{thm:full-kl-promise-np}.

Two normalizations make the transfer relative.  First, we may assume
$\Xi(\mathcal C)\ge1$ for every partition that is not cost-zero: the points are
integral, so by the granularity argument of Lemma~\ref{lem:kmeans-gap} every
nonzero value of $\Xi$ is a rational with denominator dividing
$\operatorname{lcm}(1,\dots,n)$, and rescaling the input by that integer --- a
polynomial-time operation multiplying $\Xi$ by its square --- makes all nonzero
values at least $1$.  Second, choose $\delta=2^{-\Lambda}$ with $\Lambda$
polynomial in the input size so that
\[
C_0d^2\delta^3Z^3\ \le\ \frac{\varepsilon_1}{4}\cdot\frac{2d\delta^2}{n},
\]
which holds as soon as $\delta\le\varepsilon_1/(2C_0\,d\,n\,Z^3)$ and therefore
costs only polynomially many bits.

Write $\beta=2d\delta^2/n$.  For every partition,
\[
\beta\,\Xi(\mathcal C)-\tfrac{\varepsilon_1}{4}\beta
\ \le\
\RetKL(\mathcal C)
\ \le\
\beta\,\Xi(\mathcal C)+\tfrac{\varepsilon_1}{4}\beta ,
\]
so, using $\Xi\ge1$ on the relevant partitions,
\[
\bigl(1-\tfrac{\varepsilon_1}{4}\bigr)\beta\,\Xi(\mathcal C)
\ \le\
\RetKL(\mathcal C)
\ \le\
\bigl(1+\tfrac{\varepsilon_1}{4}\bigr)\beta\,\Xi(\mathcal C).
\]
Hence $\RetKL$ and $\beta\Xi$ agree to within the factor
$(1+\tfrac{\varepsilon_1}{4})/(1-\tfrac{\varepsilon_1}{4})$ \emph{simultaneously
for all partitions}, and the same therefore holds for their minima.

The uniformity is the operative point and is visible numerically.  Enumerating
\emph{all} partitions with at most $k$ blocks for $d\in\{2,3,4\}$ and
$n\in\{5,6\}$, the ratio $\RetKL(\mathcal C)/(\beta\,\Xi(\mathcal C))$ has
spread at most $9.0\times10^{-5}$ at $\delta=10^{-3}$, $9.0\times10^{-9}$ at
$\delta=10^{-5}$, and $9.0\times10^{-13}$ at $\delta=10^{-7}$: the spread
falls as $\delta^2$, so a polynomial number of bits in $\delta$ makes the
two-sided bound hold for every partition at once.  In particular the minimizing
partitions of the two objectives coincide on the instances tested.

Suppose an algorithm returned a partition $\widehat{\mathcal C}$ with
$\RetKL(\widehat{\mathcal C})\le(1+\varepsilon_0)\min_{\mathcal C}
\RetKL(\mathcal C)$.  Then
\[
\Xi(\widehat{\mathcal C})
\ \le\
\frac{1+\tfrac{\varepsilon_1}{4}}{1-\tfrac{\varepsilon_1}{4}}
\,(1+\varepsilon_0)\,
\min_{\mathcal C}\Xi(\mathcal C),
\]
and choosing $\varepsilon_0>0$ small enough that this prefactor is at most
$1+\varepsilon_1$ --- for instance $\varepsilon_0=\varepsilon_1/4$, for
$\varepsilon_1\le1$ --- yields a $(1+\varepsilon_1)$-approximation to
$k$-means, which is NP-hard.  The machine construction, lumpability of all
partitions, and depth-one synchronization are exactly as in
Theorem~\ref{thm:full-kl-promise-np}.
\end{proof}

\begin{remark}[Why Only a Promise, and Only Hardness]
\label{rem:full-kl-promise-scope}
Two limitations of Theorem~\ref{thm:full-kl-promise-np} are intrinsic to the
construction and should not be read as slack.

First, the statement is a promise problem.  The reduction controls the objective
only up to the cubic remainder $\rho_C$, so it separates instances whose values
differ by at least $2d\delta^2$ but says nothing about instances whose values
straddle a threshold more closely.  Exact comparison of $\RetKL(k)$ with a
rational threshold is a different question, and the promise formulation is the
honest one.

Second, no NP-\emph{completeness} is claimed.  Membership in NP would require
verifying, in polynomial time, an inequality between a rational threshold and a
sum of terms of the form $p\log(p/q)$ with rational $p,q$.  Deciding such
inequalities is not known to be in P, and the general problem of comparing sums
of logarithms of rationals to a rational is a well-known obstacle to exact
membership.  Proposition~\ref{prop:full-kl-algebraic-form} below reduces that
sum to a single comparison $R\ge e^{s}$ between a factored rational and an
exponential, which is the sharp form of the obstruction and is discussed in
Remarks~\ref{rem:full-kl-slp} and~\ref{rem:full-kl-exp-membership}.  The Gaussian quadratic regime of
Theorem~\ref{thm:retention-reset-np} avoids this entirely, since its objective
is a rational function of the data; that is precisely why completeness is
available there and only hardness here.
\end{remark}

The obstruction just described can be located exactly.  The sum of logarithms
collapses to a single comparison between a rational and an exponential, so
neither the number of logarithmic terms nor their coefficients is the
difficulty.

\begin{proposition}[Algebraic Form of the Full-KL Verification Problem]
\label{prop:full-kl-algebraic-form}
Let a full-KL retention instance have rational stationary weights $\pi_s$ and
rational, strictly positive predictive laws $P_s$ on a finite alphabet, and fix
a partition $\phi$.  Then there are a positive integer $D$ and a positive
rational $R$ with
\[
\RetKL(\phi)=\frac1D\log R ,
\]
where $D$ and a \emph{factored} representation
$R=\prod_jb_j^{\,m_j}$, with the $b_j$ rational and the $m_j$ integers, are
computable in time polynomial in the instance size.  Consequently, comparing
$\RetKL(\phi)$ with a rational threshold $\tau$ reduces in polynomial time to
the problem
\[
\textsc{RationalExpCompare}:\quad
\text{given a factored positive rational }R\text{ and a rational }s,
\ \text{decide whether }R\ge e^{s}.
\]
\end{proposition}

\begin{proof}
For the fixed partition, each block centroid
$\bar P_k=\bigl(\sum_{s\in C_k}\pi_sP_s\bigr)/\bigl(\sum_{s\in C_k}\pi_s\bigr)$
is a rational probability vector computable by rational arithmetic.  Expanding
the definition,
\[
\RetKL(\phi)
=
\sum_k\sum_{s\in C_k}\sum_y\pi_sP_s(y)\log\frac{P_s(y)}{\bar P_k(y)}
=
\sum_ja_j\log b_j ,
\]
a sum of at most $|\Splus||\mathcal O|$ terms with all $a_j$ and $b_j$ rational
and $b_j>0$.  Let $D$ be a common denominator of the $a_j$ and write
$a_j=m_j/D$ with $m_j\in\mathbb Z$.  Then
\[
\RetKL(\phi)=\frac1D\sum_jm_j\log b_j=\frac1D\log\Bigl(\prod_jb_j^{\,m_j}\Bigr),
\]
so $R=\prod_jb_j^{\,m_j}$, negative exponents contributing reciprocals.  Both
$D$ and the list of pairs $(b_j,m_j)$ have polynomial bit size.  Finally
$\RetKL(\phi)\le\tau$ if and only if $R\le e^{D\tau}$, and $s=D\tau$ is
rational of polynomial bit size.
\end{proof}

\begin{remark}[The Factored Representation Is Necessary]
\label{rem:full-kl-slp}
Proposition~\ref{prop:full-kl-algebraic-form} asserts polynomial computability
of $R$ only in factored form, and the qualification is not a formality.  The
exponents $m_j=a_jD$ are integers of polynomial bit \emph{length} but
exponential \emph{magnitude}, so the expanded numerator and denominator of
$R=\prod_jb_j^{m_j}$ have bit size $\sum_j|m_j|\log b_j$, which grows
exponentially in the instance size: on rational instances with $|\Splus|\le6$
and denominators below $32$, inputs of a few hundred bits already produce
expanded representations of billions of bits, while the factored
representation stays in the low thousands.  \textsc{RationalExpCompare} must
therefore be posed for $R$ given as a product of powers, that is by a straight-line
program, and it is in that form that it is a \textsc{PosSLP}-type question.
\end{remark}

\begin{remark}[Exact Status of Membership]
\label{rem:full-kl-exp-membership}
Proposition~\ref{prop:full-kl-algebraic-form} isolates the obstruction to
membership in NP: it is \textsc{RationalExpCompare}, and not the number of
logarithmic terms, which collapse into one.  That problem is decidable.  For
$s=0$ one compares the rational $R$ with $1$, which is exact arithmetic on the
factored form.  For $s\neq0$ the Lindemann--Weierstrass theorem makes $e^{s}$
transcendental while $R$ is rational, hence $R\neq e^{s}$, so the comparison is
strict and is resolved by evaluating $e^{s}$ to sufficient precision.  What is
missing is a polynomial bound on the precision required, equivalently a
polynomial lower bound on $|R-e^{s}|$ in terms of the bit sizes of the factored
$R$ and of $s$.

Accordingly, finite-alphabet full-KL retention lies in NP whenever
\textsc{RationalExpCompare} admits polynomial-time verifiable certificates.
This is the precise sense in which the Gaussian quadratic regime of
Theorem~\ref{thm:retention-reset-np}, whose objective is a rational function of
the data and so never meets a transcendental comparison, admits
NP-completeness, while the full-KL regime of
Theorem~\ref{thm:full-kl-promise-np} admits promise NP-hardness only.
\end{remark}

\begin{remark}[The Output Alphabet of the Embedding]
\label{rem:output-alphabet-2d}
The tangent-space embedding of Theorem~\ref{thm:full-kl-promise-np} needs a map
$T:\mathbb R^{d}\to\mathbb R^{|\mathcal O|}$ subject to three constraints: its
image must lie in the zero-sum subspace, so that $u+\delta Tz$ is a probability
vector; it must be a similarity, $T^{\!\top}T=cI$, so that the leading
Jensen--Shannon term is a constant multiple of the $k$-means objective; and it
must be rational, so that integer inputs give rational laws and the promise
structure survives.  The doubling map
$a\mapsto(a_1,-a_1,\dots,a_d,-a_d)$ meets all three with $|\mathcal O|=2d$ and
$c=2$.

The zero-sum subspace of $\mathbb R^{|\mathcal O|}$ has dimension
$|\mathcal O|-1$, so $|\mathcal O|\ge d+1$ is forced.  Whether the minimum is
attained is a rational-equivalence question, and the answer depends on $d$
rather than on a universal factor.

For $d=2$ the minimum is \emph{not} attained.  In the basis $(1,-1,0)$,
$(1,1,-2)$ the zero-sum subspace of $\mathbb Q^{3}$ carries the Gram matrix
$\operatorname{diag}(2,6)$, of discriminant $12$, whose class modulo rational
squares is $3$; an orthogonal rational basis of equal norms would give
$\operatorname{diag}(c,c)$, of discriminant $c^{2}$ and class $1$.  Since
$3\not\equiv1$ modulo squares, no rational centred similarity
$\mathbb R^{2}\to\mathbb R^{3}$ exists, and $|\mathcal O|=4=2d$ is optimal.

For $d=3$, by contrast, the minimum \emph{is} attained.  The three non-constant
rows of a Hadamard matrix of order $4$,
\[
(-1,-1,1,1),\qquad(-1,1,-1,1),\qquad(-1,1,1,-1),
\]
are zero-sum, pairwise orthogonal and of common norm $4$, and therefore define
a rational centred similarity $\mathbb R^{3}\to\mathbb R^{4}$ with $c=4$.  The
alphabet may thus be taken of size $d+1=4$ rather than $2d=6$.  The same
construction applies whenever a Hadamard matrix of order $d+1$ exists, so
$|\mathcal O|=d+1$ suffices for $d=3,7,15,\dots$, and there the factor two is
slack.

The operative statement is therefore that $|\mathcal O|$ may be taken as small
as the least $n\ge d+1$ for which the zero-sum subspace of $\mathbb Q^{n}$
contains $d$ pairwise orthogonal rational vectors of equal norm, a condition of
Hadamard type; $n=2d$ always works and $n=d+1$ works exactly when the
corresponding quadratic form is rationally equivalent to a scalar matrix.  In
every case $|\mathcal O|$ grows with $d$, since the zero-sum subspace must have
dimension at least $d$.  Hardness of full-KL retention for a \emph{fixed}
output alphabet therefore remains open and would require a reduction from a
source problem hard in bounded dimension, Euclidean $k$-means in fixed
dimension not being known to be so.
\end{remark}

\begin{remark}[Scope of the Complexity Theorems]
\label{rem:complexity-scope}
Theorems~\ref{thm:retention-reset-np}
and~\ref{thm:retention-definite-np} are theorems of the Gaussian quadratic
restricted retention regime, where the objective is rational in the data and
NP-completeness is available.  The reset construction shows that the
right-congruence constraint can be made vacuous while preserving reachability
and minimality.  The depth-$2$ construction uses tagged means to make all
$n^2$ predictive distributions distinct.

For unrestricted full-KL retention over a finite output alphabet,
Theorem~\ref{thm:full-kl-promise-np} gives NP-hardness under a promise, by
embedding $k$-means into a small interior neighborhood of the uniform law so
that the Jensen--Shannon block cost is a controlled perturbation of the
Euclidean one.  What is not established for that regime is exact
NP-completeness, for the arithmetic reason recorded in
Remark~\ref{rem:full-kl-promise-scope}.
\end{remark}

%======================================================================
\section{Commitment}
\label{sec:commitment}
%======================================================================

The commitment regime concerns deterministic finite-state realization and
enforcement.  Its exact zero-error theory is governed by Myhill--Nerode
residual equivalence.  Its quantitative theory is governed by dynamic
programming and, in simultaneous-move settings, by strategic-spread lower
bounds.  The exponent correspondence draws an \emph{analogy} between commitment and
the Booleanized symmetrized $\alpha=0$ support vertex: both are zero-set /
support-type objects rather than magnitude-type objects.  This is a structural
analogy, not an identification; commitment is not raw one-sided R\'enyi order
$0$, and it is not ordinary Schatten rank.  Section~\ref{subsec:commitment-rd}
below supplies a genuinely quantitative, non-threshold commitment gap for
readers who want a magnitude-type commitment object.

%----------------------------------------------------------------------
\subsection{Deterministic Specifications}
\label{subsec:det-spec}
%----------------------------------------------------------------------

\begin{definition}[Causal Length-Preserving Specification]
\label{def:causal-lp-spec}
A \textbf{deterministic specification} is a map
\[
F:\mathcal I^*\to\mathcal O^*
\]
that is \textbf{length-preserving}, $|F(u)|=|u|$ for every $u\in\mathcal I^*$,
and \textbf{causal}: $F(u)$ is a prefix of $F(uw)$ for every
$u,w\in\mathcal I^*$.  Causality and length-preservation make $F$ a
well-typed finite-word transduction, and by the standard inverse-limit
argument for prefix-consistent finite-word maps, $F$ extends uniquely to a
continuous causal map $F^\infty:\mathcal I^{\mathbb N}\to\mathcal O^{\mathbb
N}$ on infinite input streams, recovering the BSG map type of
Section~\ref{subsec:bsg-intro}.  All commitment quantities below are stated
for finite $u,w\in\mathcal I^*$, for which $uw\in\mathcal I^*$ and every
displayed expression is well typed.
\end{definition}

For each finite input history $u$, causality lets us define the residual
\[
F_u(w)=F(uw)_{|u|+1:|u|+|w|},
\qquad w\in\mathcal I^*,
\]
the unique word with $F(uw)=F(u)F_u(w)$.  Both $u$ and $w$ range over finite
words, so $uw\in\mathcal I^*$ remains in the domain of $F$, and $F_u$ is the
continuation behavior required after the history $u$.

Define the Myhill--Nerode equivalence
\[
u\sim_F v
\iff
F_u=F_v,
\]
and let
\[
\kappa_{\det}(F)=\operatorname{index}(\sim_F).
\]

\begin{theorem}[Exact Commitment Bound]
\label{thm:commitment-spec}
A deterministic BSG with $M$ states implements $F$ iff
\[
M\ge\kappa_{\det}(F).
\]
Under normalized worst-case mismatch cost,
\[
\Com(M)=
\begin{cases}
0,&M\ge\kappa_{\det}(F),\\
1,&M<\kappa_{\det}(F).
\end{cases}
\]
\end{theorem}

\begin{proof}
Necessity is the Myhill--Nerode residual lower bound: if two histories induce
distinct residuals, no deterministic machine can assign them the same state and
still implement $F$ on all continuations.  Sufficiency is realized by the
residual automaton whose states are the equivalence classes of $\sim_F$, with
transition
\[
[u]\cdot x=[ux]
\]
and output determined by the residual class.
\end{proof}

The structural Boolean version uses the $\{0,\infty\}$ valuation:
\[
\mathcal E(f,g)=
\begin{cases}
0,&f=g,\\
\infty,&f\neq g.
\end{cases}
\]
The normalized operational version uses the $\{0,1\}$ valuation.  The threshold
is the same; only the positive mismatch value changes.

%----------------------------------------------------------------------
\subsection{Quantitative Commitment: A Distributional Rate--Distortion Gap}
\label{subsec:commitment-rd}
%----------------------------------------------------------------------

The exact commitment gap $\Comex(M)$ of Theorem~\ref{thm:commitment-spec} is
a threshold: it is $0$ once $M\ge\kappa_{\det}(F)$ and constant and positive
otherwise, and it carries no information about how a sub-threshold budget
should be spent.  We now supply a genuinely quantitative, rate-distortion-typed
commitment gap over right congruences, parallel in shape to the retention
information-bottleneck gap $\RetKL(M)$ and to the grounding tracking deficit
(Definition~\ref{def:tracking-deficit}), with its own exact fixed-budget
formula and its own support-sensitive zero threshold.

Fix a discounted-prefix history law $\mu$ on $\mathcal I^*$,
\[
\mu(u)=(1-\gamma)\gamma^{|u|}\Pr[\text{prefix }u],
\qquad \gamma\in(0,1),
\]
as in Definition~\ref{def:observable-support-index}; this is a genuine
probability measure on all of $\mathcal I^*$, since finite-length prefix
masses are mixed geometrically across lengths.  Let $F$ be a causal
length-preserving specification (Definition~\ref{def:causal-lp-spec}).

\begin{definition}[Implemented Machine and One-Step Mismatch]
\label{def:com-rd-machine}
For a right congruence $\sim$ of index at most $M$
(Definition~\ref{def:right-cong}) and an output rule
\[
\rho:\bigl(\mathcal I^*/\!\sim\bigr)\times\mathcal I\to\mathcal O,
\]
the pair $(\sim,\rho)$ realizes the Mealy machine that, having read a
nonempty history $u$ of length $t$, outputs
$\rho\bigl([u_{1:t-1}]_\sim,\,u_t\bigr)$
at position $t$.  This is a BSG output map of type
$\lambda:(\mathcal I^*/\!\sim)\times\mathcal I\to\mathcal O$
(Section~\ref{subsec:bsg-intro}), so $(\sim,\rho)$ is implementable with
$\operatorname{index}(\sim)$ states.  Its one-step mismatch indicator at $u$
is
\[
\mathrm e_{\sim,\rho}(u)
=
\mathbf 1
\Bigl\{
\rho\bigl([u_{1:t-1}]_\sim,\,u_t\bigr)
\ne
F(u)_t
\Bigr\},
\qquad t=|u|\ge1.
\]
\end{definition}

\begin{definition}[Distributional Commitment Rate--Distortion Gap]
\label{def:com-rd-gap}
\[
\ComRD(M)
=
\min_{\substack{\sim\ \text{right congruence}\\ \operatorname{index}(\sim)\le M}}
\ \min_{\rho}
\ \sum_{u\in\mathcal I^*\setminus\{\varepsilon\}}
\mu(u)\,
\mathrm e_{\sim,\rho}(u).
\]
\end{definition}

For a class $k$ of $\sim$ and a symbol $a\in\mathcal I$, let
\[
\mu(k,a)
=
\sum_{\substack{u\ne\varepsilon\\ [u_{1:|u|-1}]_\sim=k,\ u_{|u|}=a}}
\mu(u),
\qquad
q_{k,a}(o)
=
\frac1{\mu(k,a)}
\sum_{\substack{u\ne\varepsilon\\ [u_{1:|u|-1}]_\sim=k,\ u_{|u|}=a\\ F(u)_{|u|}=o}}
\mu(u)
\quad(\mu(k,a)>0),
\]
the $\mu$-conditional law, given class $k$ and next symbol $a$, of the true
next output demanded by $F$.

\begin{definition}[One-Step Determination Index]
\label{def:pair-determination-index}
The \textbf{one-step determination index} $\kappa_{\mathrm{pair}}(F,\mu)$ is
the minimum index of a right congruence $\sim$ on $\mathcal I^*$ for which
$F(u)_{|u|}$ is $\mu$-almost-surely determined by the pair
$\bigl([u_{1:|u|-1}]_\sim,\,u_{|u|}\bigr)$: explicitly, there exists a map
$\rho:\bigl(\mathcal I^*/\!\sim\bigr)\times\mathcal I\to\mathcal O$ with
$\rho\bigl([u_{1:|u|-1}]_\sim,u_{|u|}\bigr)=F(u)_{|u|}$ for $\mu$-almost
every nonempty $u$.  Only nonempty $u$ carry a one-step output, so the empty
word --- which has positive mass, $\mu(\varepsilon)=1-\gamma$, under a
discounted law (Definition~\ref{def:discounted-agg}) --- does not enter the
condition; the index depends on $\mu$ only through its support.
\end{definition}

\begin{theorem}[Exact Fixed-Budget Commitment Rate--Distortion Formula]
\label{thm:com-rd-formula}
For every right congruence $\sim$,
\[
\min_\rho
\sum_{u\ne\varepsilon}
\mu(u)\,\mathrm e_{\sim,\rho}(u)
=
\sum_{\substack{k,a\\ \mu(k,a)>0}}
\mu(k,a)
\Bigl(1-\max_{o\in\mathcal O}q_{k,a}(o)\Bigr),
\]
attained by the greedy rule $\rho(k,a)\in\arg\max_oq_{k,a}(o)$.  Consequently
\[
\ComRD(M)
=
\min_{\operatorname{index}(\sim)\le M}
\ \sum_{\substack{k,a\\ \mu(k,a)>0}}
\mu(k,a)
\Bigl(1-\max_{o}q_{k,a}(o)\Bigr).
\]
Moreover $\ComRD(M)=0$ if and only if
$M\ge\kappa_{\mathrm{pair}}(F,\mu)$, the one-step determination index of
Definition~\ref{def:pair-determination-index}.  In particular
$\kappa_{\mathrm{pair}}(F,\mu)\le\kappa_{\det}(F)$, the threshold of
Theorem~\ref{thm:commitment-spec}: the two indices coincide when $\mu$ has
full support, and the inequality is strict under restricted supports
(Remark~\ref{rem:pair-vs-class}).
\end{theorem}

\begin{proof}
Fix $\sim$.  Group the sum defining $\ComRD$ by the pair
$(k,a)=([u_{1:|u|-1}]_\sim,u_{|u|})$: for fixed $(k,a)$ with $\mu(k,a)>0$, the
rule $\rho$ assigns a single value $\rho(k,a)\in\mathcal O$ to every $u$ in
that group, and
\[
\sum_{\substack{u\ne\varepsilon\\ [u_{1:|u|-1}]_\sim=k,\,u_{|u|}=a}}
\mu(u)\,\mathbf 1\{\rho(k,a)\ne F(u)_{|u|}\}
=
\mu(k,a)\bigl(1-q_{k,a}(\rho(k,a))\bigr)
\ge
\mu(k,a)\bigl(1-\max_oq_{k,a}(o)\bigr),
\]
with equality exactly when $\rho(k,a)$ is a mode of $q_{k,a}$.  Summing over
$(k,a)$ and choosing the mode at every pair proves the displayed formula;
groups with $\mu(k,a)=0$ contribute $0$ under any choice of $\rho$ and may be
assigned arbitrarily.  Minimizing further over $\sim$ of index at most $M$
gives $\ComRD(M)$.

For the zero threshold: $\ComRD(M)=0$ for a given $\sim$ if and only if every
visited pair $(k,a)$ has $q_{k,a}$ a point mass, i.e.\ $F(u)_{|u|}$ is
$\mu$-almost surely determined by $([u_{1:|u|-1}]_\sim,u_{|u|})$, which is
precisely the determination condition of Definition~\ref{def:pair-determination-index};
hence $\ComRD(M)=0$ exactly for $M\ge\kappa_{\mathrm{pair}}(F,\mu)$.  The
inequality $\kappa_{\mathrm{pair}}(F,\mu)\le\kappa_{\det}(F)$ holds because
the full Myhill--Nerode congruence $\sim_F$ determines the one-step output
through the pair: $[u_{1:|u|-1}]_{\sim_F}$ determines the residual
$F_{u_{1:|u|-1}}$, whose value on the one-letter continuation $u_{|u|}$ is
$F(u)_{|u|}$, so the residual automaton of Theorem~\ref{thm:commitment-spec}
is one-step exact on all of $\mathcal I^*$.  When $\mu$ has full support the
reverse inequality holds as well: a pair-determining right congruence then
forces $u\sim v$ to imply $F_u=F_v$, because each continuation output
$F(uw)_{|u|+j}$ is a function of $([uw_{1:j-1}]_\sim,w_j)$ and right
invariance identifies the two arguments along $u$ and $v$.  The two indices
then coincide, and strictness of the inequality is a support phenomenon, as
in Remark~\ref{rem:pair-vs-class}.
\end{proof}

\begin{remark}[Why the Pair, Not the Class]
\label{rem:pair-vs-class}
Determination by the pair $([u_{1:|u|-1}]_\sim,u_{|u|})$ is strictly weaker
than determination by the class $[u]_\sim$ alone: the pair refines the
class, since $[u]_\sim=[u_{1:|u|-1}]_\sim\cdot u_{|u|}$, but a single class
is reached by many distinct pairs, so a class-level representative is a
stronger requirement than a pair-level output rule.  For the identity
transduction $F(u)=u$ over a binary alphabet under a full-support $\mu$,
every residual coincides, so $\kappa_{\det}(F)=1$: the one-state rule
$\rho(*,x)=x$ is one-step exact, $\ComRD(1)=0$, and
$\kappa_{\mathrm{pair}}(F,\mu)=1$, whereas determining $u_{|u|}$ from the
class $[u]_\sim$ alone requires two classes.  The gap between
$\kappa_{\mathrm{pair}}$ and $\kappa_{\det}$, by contrast, is a support
phenomenon: for the transduction $F(u)_t=u_t\oplus u_{t-1}$ (with $u_0=0$)
the two indices coincide at $2$ under a full-support law, but if the input
process is supported on strictly alternating words, then $u_{t-1}=1-u_t$ on
the support, the one-step output is constant there, and
$\kappa_{\mathrm{pair}}(F,\mu)=1$.
\end{remark}

\begin{remark}[Relation to the Three Commitment Notions]
\label{rem:com-rd-scope}
Three distinct commitment quantities are now in play, kept apart by
notation. $\Comex(M)$ (Theorem~\ref{thm:commitment-spec}) is the exact
worst-case Boolean threshold, testing equality of full residual functions on
every history.  $\ComRD(M)$ (Definition~\ref{def:com-rd-gap}) is a
distributional, single-step quantitative loss averaged over a fixed history
law $\mu$, whose zero set can be strictly larger than that of $\Comex$
because it only demands agreement on the very next symbol and only on the
$\mu$-visited part of history space.  $\ComGame(M)$
(Section~\ref{subsec:quant-simul}) is a simultaneous-move strategic gap
against an adversarial input sequence chosen online, and is not an averaging
or a threshold quantity at all.  No general inequality between $\ComRD$ and
$\ComGame$ is asserted; they are posed on different protocols.
\end{remark}

%----------------------------------------------------------------------
\subsection{General Safety Properties}
\label{subsec:safety}
%----------------------------------------------------------------------

For a safety property
\[
\Lambda\subseteq(\mathcal I\times\mathcal O)^*,
\]
there may be many safe output choices.  A deterministic specification $F$ is
\textbf{admissible} for $\Lambda$ if every finite input-output prefix generated
by $F$ lies in $\Lambda$.  We write this as
\[
F\subseteq\Lambda.
\]
A BSG enforces $\Lambda$ if its induced specification is admissible for every
input sequence.

The complexity is the minimal Myhill--Nerode index of a deterministic
admissible strategy.

\begin{definition}[Controller Complexity of a Safety Property]
\label{def:controller-complexity}
\[
\kappa_{\mathrm{ctrl}}(\Lambda)
=
\min_{F\subseteq\Lambda}
\kappa_{\det}(F),
\]
where the minimum is over deterministic specifications whose graphs lie inside
$\Lambda$.  If no finite-state admissible specification exists, set
$\kappa_{\mathrm{ctrl}}(\Lambda)=\infty$.
\end{definition}

\begin{theorem}[Exact Safety Enforcement]
\label{thm:safety-enforce}
A deterministic BSG with $M$ states can enforce $\Lambda$ iff
\[
M\ge\kappa_{\mathrm{ctrl}}(\Lambda).
\]
If $\Lambda$ admits a unique admissible residual specification $F_\Lambda$,
then
\[
\kappa_{\mathrm{ctrl}}(\Lambda)
=
\kappa_{\det}(F_\Lambda).
\]
\end{theorem}

\begin{proof}
If a BSG enforces $\Lambda$, it implements some admissible deterministic
specification $F\subseteq\Lambda$.  Its number of states is at least
$\kappa_{\det}(F)$, hence at least $\kappa_{\mathrm{ctrl}}(\Lambda)$.
Conversely, if $M\ge\kappa_{\mathrm{ctrl}}(\Lambda)$, choose an admissible
specification $F\subseteq\Lambda$ attaining the minimum and implement it by its
residual automaton.

If $\Lambda$ admits a unique admissible residual specification $F_\Lambda$,
then the minimum is attained uniquely up to residual equivalence, so
\[
\kappa_{\mathrm{ctrl}}(\Lambda)
=
\kappa_{\det}(F_\Lambda).
\]
\end{proof}

The distinction between $\kappa_{\det}(F)$ and
$\kappa_{\mathrm{ctrl}}(\Lambda)$ is that a safety property may admit many
strategies, and the minimal controller may be smaller than the complexity of an
arbitrary admissible specification.

\begin{theorem}[Finite-Monitor Winning Region for Regular Safety]
\label{thm:finite-monitor}
Suppose $\Lambda\subseteq(\mathcal I\times\mathcal O)^*$ is a regular safety
property, recognized by a deterministic finite automaton
$\mathcal M=(\mathcal Q,q_0,\delta_{\mathcal M},\mathrm{Acc})$ over the joint
alphabet $\mathcal I\times\mathcal O$, in the sense that $w\in\Lambda$ iff
every prefix of $w$ has its $\delta_{\mathcal M}$-run staying inside
$\mathrm{Acc}$.  Define the \textbf{safe winning region}
$W\subseteq\mathcal Q$ by the greatest-fixpoint recursion
\[
W_0=\mathrm{Acc},
\qquad
W_{n+1}
=
\Bigl\{
q\in W_n:\forall a\in\mathcal I\ \exists b\in\mathcal O\ \
\delta_{\mathcal M}(q,(a,b))\in W_n
\Bigr\},
\qquad
W=\bigcap_{n\ge0}W_n.
\]
Since $\mathcal Q$ is finite the recursion stabilizes after at most
$|\mathcal Q|$ steps.  Then:
\begin{enumerate}[label=(\roman*)]
\item A finite-state admissible strategy enforcing $\Lambda$ from $q_0$
exists if and only if $q_0\in W$.
\item When $q_0\in W$, the residual automaton on state space $W$ with output
rule $\beta(q,a)\in\{b\in\mathcal O:\delta_{\mathcal M}(q,(a,b))\in W\}$ (any
fixed selection) enforces $\Lambda$ and has at most $|W|\le|\mathcal Q|$
states; consequently
\[
\kappa_{\mathrm{ctrl}}(\Lambda)\le|W|\le|\mathcal Q|.
\]
\item $W$ is computable in time polynomial in
$|\mathcal Q|\cdot|\mathcal I|\cdot|\mathcal O|$ by the stated fixpoint
iteration.
\end{enumerate}
\end{theorem}

\begin{proof}
This is the classical attractor / safety-game construction, specialized to a
Mealy controller choosing $b$ in response to an adversarial $a$.
Monotonicity of the recursion ($W_{n+1}\subseteq W_n$) and finiteness of
$\mathcal Q$ give stabilization in at most $|\mathcal Q|$ steps, since each
strict decrease removes at least one state.  If $q_0\in W$, then from every
$q\in W$ and every $a\in\mathcal I$ there is $b\in\mathcal O$ with
$\delta_{\mathcal M}(q,(a,b))\in W$ by construction, so the selection $\beta$
keeps every play inside $W\subseteq\mathrm{Acc}$ forever, which is exactly
admissibility for $\Lambda$; this proves the ``if'' direction of~(i)
and~(ii).  Conversely, suppose a finite-state admissible strategy enforces
$\Lambda$ from $q_0$.  An induction on $n$ shows every state reachable by
that strategy lies in $W_n$ for every $n$: a violation of membership in
$W_{n+1}$ at a reachable state would exhibit an adversary input $a$ after
which every output $b$ leaves $W_n$, and admissibility of the strategy would
then force the next reached state to already violate $W_n$, contradicting
the inductive hypothesis at the previous step.  Hence every reachable state,
in particular $q_0$, lies in $\bigcap_nW_n=W$, proving the ``only if''
direction of~(i).  Claim~(iii) is immediate from the fixpoint iteration,
which performs at most $|\mathcal Q|$ rounds each costing
$O(|\mathcal Q|\cdot|\mathcal I|\cdot|\mathcal O|)$ to test the two
quantifiers over $\mathcal Q\times\mathcal I\times\mathcal O$.
\end{proof}

\begin{remark}[Relation to $\kappa_{\mathrm{ctrl}}$]
\label{rem:finite-monitor-scope}
Theorem~\ref{thm:finite-monitor} gives a constructive upper bound
$\kappa_{\mathrm{ctrl}}(\Lambda)\le|W|$ computable directly from a
recognizing automaton for $\Lambda$, complementing the abstract
characterization of Definition~\ref{def:controller-complexity} and
Theorem~\ref{thm:safety-enforce}.  The bound $|W|$ need not equal
$\kappa_{\mathrm{ctrl}}(\Lambda)$: further residual minimization of the
$W$-automaton, exactly as in Theorem~\ref{thm:commitment-spec}, can strictly
reduce the state count while preserving admissibility.
\end{remark}

%----------------------------------------------------------------------
\subsection{Quantitative Commitment: Alternating Moves}
\label{subsec:quant-alt}
%----------------------------------------------------------------------

Let a deterministic reward automaton be
\[
\mathcal R=(\mathcal Q,q_0,\delta_{\mathcal R},r),
\]
with bounded reward
\[
r:\mathcal Q\times\mathcal I\times\mathcal O\to\mathbb R
\]
and discount $\gamma\in(0,1)$.  In the alternating-move interpretation, the
adversary chooses an input $a\in\mathcal I$, the agent observes $a$, and then
chooses an output $b\in\mathcal O$.

Two value functions occur below and must be kept apart, since they belong to
different protocols.

The \textbf{alternating-move} value is the state function
\[
\Valt(q)
=
\min_{a\in\mathcal I}
\max_{b\in\mathcal O}
\left[
r(q,a,b)+\gamma\,\Valt(\delta_{\mathcal R}(q,a,b))
\right],
\]
in which the adversary moves first and minimizes and the agent, having observed
the current input, maximizes.  Write $\Valt(\infty)=\Valt(q_0)$ for its value
at the initial reward state; the agent here is unrestricted in memory, since
$\Valt$ is computed on the reward automaton itself.

The simultaneous protocol cannot be realized by a Mealy BSG.  A Mealy output
map has the type $\lambda:\mathcal S\times\mathcal I\to\mathcal O$, so it reads
the current input $a_t$; a policy that must commit to $b_t$ before seeing
$a_t$ is a different object and requires a different machine model.

\begin{definition}[Pre-Input-Output Controller]
\label{def:pio-controller}
A \textbf{pre-input-output controller} of budget $M$ is a tuple
\[
\pi=(\mathcal M,m_0,\beta,\eta),
\qquad
|\mathcal M|\le M,
\]
with output map $\beta:\mathcal M\to\mathcal O$ and update map
$\eta:\mathcal M\times\mathcal I\times\mathcal O\to\mathcal M$, acting by
\[
b_t=\beta(m_t),
\qquad
m_{t+1}=\eta(m_t,a_t,b_t).
\]
The output at time $t$ depends on the memory state alone, so it is committed
before $a_t$ is revealed; the memory may then absorb $a_t$.  Write $\Pi_M$ for
the set of such controllers.  A Mealy BSG is \emph{not} a pre-input-output
controller, since its output map is a function of $(\mathcal S,\mathcal I)$.
\end{definition}

The \textbf{simultaneous-move} value at budget $M$ is
\[
\Vsim(M)
=
\sup_{\pi\in\Pi_M}
\ \inf_{a_{0:\infty}\in\mathcal I^{\mathbb N}}
\ \sum_{t\ge0}\gamma^t\,r(q_t,a_t,b_t),
\]
the trajectory $(q_t)$ being generated on the reward automaton by the realized
pairs $(a_t,b_t)$.

The $M$-state commitment gap is the difference between the two protocols,
\[
\ComGame(M)=\Valt(\infty)-\Vsim(M).
\]
This is the \textbf{simultaneous-move strategic gap}, kept notationally
distinct from the exact Boolean gap $\Comex(M)$ of
Theorem~\ref{thm:commitment-spec} and from the distributional
rate-distortion gap $\ComRD(M)$ of Section~\ref{subsec:commitment-rd}: it
compares two action protocols on the same reward automaton, not two feasible
classes of the same protocol.  It is an \textbf{observation-and-memory gap}: a positive value combines the
finite-memory restriction with the loss of access to the current input, and
the two contributions are not separated by this definition.  Even at
$M=\infty$ the gap need not vanish, since the alternating protocol grants the
agent information that no pre-input-output controller possesses.

%----------------------------------------------------------------------
\subsection{Quantitative Commitment: Simultaneous Moves}
\label{subsec:quant-simul}
%----------------------------------------------------------------------

In simultaneous-move games, the adversary chooses $a_t$ and the agent chooses
$b_t$ without observing each other's current choice.  The agent may use finite
memory and past observations, but not the current adversary action.

Define the strategic spread
\[
\alpha(q)
=
\Valt(q)
-
\max_{b\in\mathcal O}
\min_{a\in\mathcal I}
\left[
r(q,a,b)+\gamma\,\Valt(\delta_{\mathcal R}(q,a,b))
\right].
\]
The first term is the alternating-move value, in which the agent observes the
current input before responding.  The second term is the best one-step
simultaneous-move guarantee using a pure output $b$.  By the elementary minimax
inequality,
\[
\max_b\min_a \le \min_a\max_b,
\]
so $\alpha(q)\ge0$.

\begin{theorem}[Strategic Spread Lower Bound, Precise Version]
\label{thm:strategic-spread}
Let $\Valt$ be the alternating-move value function of
Section~\ref{subsec:quant-simul}.
For a pre-input-output controller $\pi\in\Pi_M$ of
Definition~\ref{def:pio-controller}, let $b_t=\beta(m_t)$ be its output at
time $t$, committed before $a_t$ is revealed.  Define the myopic worst-case adversary by
\[
a_t
\in
\argmin_{a\in\mathcal I}
\left[
r(q_t,a,b_t)+\gamma \Valt(\delta_{\mathcal R}(q_t,a,b_t))
\right].
\]
Let $q_t^\pi$ be the reward-automaton trajectory generated by $\pi$ and this
adversary.  Define
\[
\alpha_\gamma(\mathcal R,M)
=
\inf_{\pi\in\Pi_M}
\sum_{t=0}^{\infty}
\gamma^t
\alpha(q_t^\pi),
\]
where the infimum is over pre-input-output controllers of budget $M$
(Definition~\ref{def:pio-controller}), these being the policies admissible in
the simultaneous protocol.

Then the finite-memory simultaneous-move commitment gap satisfies
\[
\ComGame(M)
\ge
\alpha_\gamma(\mathcal R,M).
\]
\end{theorem}

\begin{proof}
For each $t$,
\[
\Valt(q_t)
\ge
\min_a
\left[
r(q_t,a,b_t)+\gamma \Valt(q_{t+1})
\right],
\]
because $\Valt(q_t)=\min_a\max_b[\cdots]$ and the maximum over $b$ is at least the
value at the particular output $b_t$.

By the choice of $a_t$,
\[
r(q_t,a_t,b_t)+\gamma \Valt(q_{t+1})
=
\min_a
\left[
r(q_t,a,b_t)+\gamma \Valt(\delta_{\mathcal R}(q_t,a,b_t))
\right].
\]
Therefore,
\[
\Valt(q_t)-r(q_t,a_t,b_t)-\gamma \Valt(q_{t+1})
\ge
\alpha(q_t).
\]
Multiply by $\gamma^t$ and sum from $t=0$ to $T-1$:
\[
\Valt(q_0)
-
\sum_{t=0}^{T-1}\gamma^t r(q_t,a_t,b_t)
-
\gamma^T \Valt(q_T)
\ge
\sum_{t=0}^{T-1}\gamma^t\alpha(q_t).
\]
Since rewards are bounded and $\gamma<1$,
\[
\gamma^T \Valt(q_T)\to0.
\]
Thus,
\[
\Valt(q_0)
-
\sum_{t=0}^{\infty}\gamma^t r(q_t,a_t,b_t)
\ge
\sum_{t=0}^{\infty}\gamma^t\alpha(q_t).
\]

For each deterministic finite-memory policy $\pi$, the adversary constructed
above is admissible in the simultaneous-move value because $\pi$ is
deterministic and therefore its current output is predictable from the
history.  Hence
\[
\inf_{\text{adv}} \operatorname{Payoff}(\pi,\text{adv})
\le
\operatorname{Payoff}(\pi,\text{myopic adversary}).
\]
For every $\pi$,
\[
\operatorname{Payoff}(\pi,\text{myopic adversary})
\le
\Valt(q_0)-\sum_{t=0}^{\infty}\gamma^t\alpha(q_t^\pi).
\]
Therefore
\[
\max_{\pi}
\inf_{\text{adv}} \operatorname{Payoff}(\pi,\text{adv})
\le
\Valt(q_0)
-
\inf_{\pi}
\sum_{t=0}^{\infty}\gamma^t\alpha(q_t^\pi).
\]
The left-hand side is $\Vsim(M)$.  Hence
\[
\Vsim(M)
\le
\Valt(q_0)-\alpha_\gamma(\mathcal R,M),
\]
and therefore
\[
\ComGame(M)
=
\Valt(q_0)-\Vsim(M)
\ge
\alpha_\gamma(\mathcal R,M).
\]
\end{proof}

\begin{remark}[Memory Tracking and Adversary Specification]
\label{rem:strategic-proof-tracking}
The proof tracks the machine memory state $m_t$ through the augmented state
$(q_t,m_t)$.  The lower bound applies to every finite-memory deterministic
policy, and therefore to every finite-state BSG.  The adversary is explicitly
specified as the myopic worst-case response to the deterministic current
output $b_t$.
\end{remark}

\begin{corollary}[Stateless Games]
\label{cor:stateless}
For a stateless game,
\[
\ComGame(M)
=
\frac{\alpha}{1-\gamma}
\]
for all $M\ge1$, the budget being irrelevant here because a pre-input-output
controller in a stateless game emits an output depending only on its memory
state, against which the adversary best-responds; additional memory therefore
cannot improve the guarantee.  Here $\alpha$ is the \emph{one-stage} strategic
spread
\[
\alpha
=
\min_{a\in\mathcal I}\max_{b\in\mathcal O}r(a,b)
\ -\
\max_{b\in\mathcal O}\min_{a\in\mathcal I}r(a,b).
\]
Both terms are one-stage quantities in the immediate reward $r$.  Neither is
the discounted alternating-move value $\Valt$, which in the stateless case
satisfies $\Valt=\bigl(\min_a\max_br(a,b)\bigr)/(1-\gamma)$.
\end{corollary}

\begin{proof}
Write $m_1=\min_a\max_br(a,b)$ and $m_2=\max_b\min_ar(a,b)$ for the two
one-stage values, so that $\alpha=m_1-m_2\ge0$ by the minimax inequality.

In a stateless game the reward-automaton state is constant, so the Bellman
equation of Section~\ref{subsec:quant-simul} reads $\Valt=m_1+\gamma\Valt$,
whence
\[
\Valt=\frac{m_1}{1-\gamma}.
\]
The strategic spread at the unique state then evaluates to
\[
\alpha(q)
=
\Valt-\max_b\min_a\bigl[r(a,b)+\gamma\Valt\bigr]
=
\Valt-\gamma\Valt-\max_b\min_ar(a,b)
=
(1-\gamma)\Valt-m_2
=
m_1-m_2,
\]
which is $\alpha$ and is independent of $t$.  The lower bound of
Theorem~\ref{thm:strategic-spread} gives
\[
\ComGame(M)
\ge
\sum_{t=0}^\infty
\gamma^t\alpha
=
\frac{\alpha}{1-\gamma}.
\]
For the reverse inequality, choose a constant output
\[
b^\star
\in
\arg\max_{b}
\min_a r(a,b).
\]
A single-state pre-input-output controller emitting $b^\star$ at every time
guarantees discounted payoff
\[
\frac{1}{1-\gamma}
\max_b\min_a r(a,b).
\]
Since the alternating-move value is $\Valt=m_1/(1-\gamma)$, the gap is at most
\[
\frac{m_1-\max_b\min_a r(a,b)}{1-\gamma}
=
\frac{m_1-m_2}{1-\gamma}
=
\frac{\alpha}{1-\gamma}.
\]
Thus equality holds.  Finite memory does not improve the guarantee, because
the deterministic policy's current output is predictable from the history, and
the adversary can choose the myopic worst-case response in each period.
\end{proof}

%----------------------------------------------------------------------
\subsection{Commitment as the Booleanized $\alpha=0$ Vertex}
\label{subsec:commitment-boolean-vertex}
%----------------------------------------------------------------------

The exact zero-error commitment cost is not raw R\'enyi order-$0$ divergence.
Raw $D_0$ tests one-sided support inclusion, not equality, and is not
Boolean-valued.  The construction uses symmetrization followed by Boolean
valuation.

For probability vectors $p,q$, recall the classical order-$0$ R\'enyi
divergence
\[
D_0(p\|q)
=
-\log q(\supp p),
\]
with the usual extended-value conventions.  Define the symmetrized support
mismatch
\[
S_0(p,q)
=
D_0(p\|q)+D_0(q\|p),
\]
and a Boolean valuation
\[
v(t)=
\begin{cases}
0,&t=0,\\
1,&t>0,
\end{cases}
\]
or, in the extended convention, $v(t)=\infty$ for $t>0$.  The Booleanized
support cost is
\[
B_0(p,q)
=
v(S_0(p,q)).
\]

For probability vectors,
\[
D_0(p\|q)=0
\iff
q(\supp p)=1
\iff
\supp q\subseteq\supp p
\]
up to null sets.  Therefore
\[
S_0(p,q)=0
\iff
\supp p=\supp q,
\]
and hence
\[
B_0(p,q)=0
\iff
\supp p=\supp q.
\]

Commitment residuals are usually Boolean or symbolic functions rather than
probability vectors.  Let a residual row be a Boolean function
\[
r_u:\Sigma^*\to\{0,1\}.
\]
Its support is
\[
\supp r_u
=
\{w\in\Sigma^*:r_u(w)=1\}.
\]
Equality of Boolean functions is equality of supports:
\[
r_u=r_v
\iff
\supp r_u=\supp r_v.
\]
Thus the Booleanized support cost recovers the commitment zero-error test:
\[
B_0(r_u,r_v)=0
\iff
r_u=r_v.
\]

For general deterministic residuals with non-Boolean output alphabet, one may
encode the residual as the Boolean support of its graph:
\[
\supp F_u
=
\{(w,o):F_u(w)=o\}.
\]
Equality of residuals is again equality of supports.  The primitive commitment
cost remains equality; the Booleanized support construction is an encoding,
not a literal probabilistic divergence on arbitrary residual objects.

Consequently, if $\kappaMN$ is the Myhill--Nerode index, the commitment gap
satisfies
\[
\Com(M)=0
\iff
M\ge\kappaMN,
\]
and under the normalized $0/1$ convention,
\[
\Com(M)=1
\quad
\text{if }M<\kappaMN.
\]

\begin{remark}[Scope of the Boolean Vertex]
\label{rem:commitment-alpha0-scope}
Raw one-sided $D_0$ tests support inclusion, not equality, and is not
Boolean-valued.  The construction uses symmetrization and Boolean valuation.
The relevant object is Boolean residual equality / Myhill--Nerode index, not
ordinary real or complex rank.
\end{remark}

\begin{remark}[Not Ordinary Schatten Rank]
\label{rem:commitment-not-schatten}
The commitment vertex is a Boolean-rank / Myhill--Nerode threshold.  Ordinary
Schatten $p\to0$ limits count ordinary singular-value rank, not deterministic
Boolean realization rank.  Boolean Hankel matrices need not define compact
Hilbert-space operators, and no ordinary Schatten-norm limit is used to derive
the commitment threshold.
\end{remark}

\begin{remark}[The $p\to0$ Limit and Why It Is Not the Boolean Vertex]
\label{rem:p-to-zero}
It is natural to ask whether commitment is recovered from the Schatten family
by letting $p\downarrow0$.  It is not, and the obstruction is worth isolating,
because it explains why the Boolean vertex is a separate construction rather
than a limiting case of Definition~\ref{def:response}.

For a finite-rank operator $A$ the \emph{unpowered} tail sum does converge to a
rank count:
\[
\lim_{p\downarrow0}\ \sum_{i>r}\sigma_i(A)^p
=
\#\{i>r:\sigma_i(A)>0\},
\]
since $\sigma^p\to1$ for every $\sigma>0$.  Composing with the threshold
$t\mapsto\mathbf 1_{\{t>0\}}$ then yields a Boolean valuation.  That
composition is an additional operation, not part of the Schatten family.

The tail functional actually used in the template,
$\Phi^{(p)}_r(A)=\bigl(\sum_{i>r}\sigma_i(A)^p\bigr)^{1/p}$, does \emph{not}
converge to a $\{0,\infty\}$-valued map.  If the tail has $k\ge2$ nonzero
singular values then $\sum_{i>r}\sigma_i^p\to k>1$, so
$\Phi^{(p)}_r(A)=\bigl(\sum_{i>r}\sigma_i^p\bigr)^{1/p}\to+\infty$ irrespective
of the magnitudes: for $\sigma=(3,\tfrac12,\tfrac14)$ and $r=1$ one has
$\Phi^{(1)}_1=0.75$ but $\Phi^{(0.01)}_1\approx4.5\times10^{29}$.  If the tail
has exactly one nonzero entry the limit is that entry, and if the tail is zero
the limit is $0$.  The limit is therefore not two-valued, and the divergence is
driven by the tail \emph{count} rather than by any norm.

Accordingly, commitment replaces ordinary rank by the deterministic Boolean
rank $\rankBdet N_\Lambda$ and the real-valued norm by a Boolean realizability
valuation.  It is the Boolean-rank vertex of the rank-control principle, not a
corollary of Mirsky's theorem, and it does not require $N_\Lambda$ to define a
compact Hilbert-space operator.
\end{remark}

\begin{remark}[Commitment Threshold Behavior]
\label{rem:commitment-degeneracy}
Exact Boolean commitment does not produce a smooth spectral-tail budget curve.
The gap is a threshold governed by $\kappaMN$.  Smooth bias--variance behavior
appears only in quantitative commitment variants, such as the discounted
reward models above, or in approximate-commitment losses explicitly defined
elsewhere.
\end{remark}

%======================================================================
\section{Grounding}
\label{sec:grounding}
%======================================================================

The grounding regime concerns approximation of stochastic channels or residual
kernels.  The spectral converse is a theorem about linear finite-rank
realizations, not about deterministic Mealy machines.  Exact symbolic
deterministic grounding is a different problem.  The regime also contains
intrinsic lower bounds for deterministic approximation: a superpolynomial
determinism gap and a stochasticity floor.

%----------------------------------------------------------------------
\subsection{Grounding Instantiation}
\label{subsec:grounding-instantiation}
%----------------------------------------------------------------------

In the grounding instantiation, residual objects are bounded operators,
stochastic kernels, or elements of a Banach or $C^*$-algebraic residual space.
When a Banach representation $B$ is chosen, the cost is typically an
operator-norm or Banach-space distance:
\[
\mathcal E(\nu,\widehat\nu)
=
\norm{\nu-\widehat\nu}_{\mathrm{op}}.
\]
The scalar impulse response used by the Hankel theory is obtained through a
bounded linear functional
\[
\phi:B\to\mathbb C,
\qquad
h(u)=\phi(\nu_u).
\]

The spectral converse concerns linear finite-rank realizations of the induced
Hankel operator.  It does not assert that a deterministic Mealy machine with
the same state count realizes or approximates the channel.  Deterministic
symbolic grounding and linear finite-rank grounding are distinct approximation
problems.

\begin{remark}[Residual Norm versus Hankel Norm]
\label{rem:grounding-residual-hankel}
As in Remark~\ref{rem:hilbert-residual-hankel}, stated here for the grounding
instantiation: the norm used for residual-operator comparisons is the norm of
the chosen Banach representation $B$, whereas the Hankel gap uses the $\ell^2$
operator norm of the Hankel operator $H_\nu$.  No domination between these
norms is assumed unless explicitly stated.
\end{remark}

%----------------------------------------------------------------------
\subsection{Linear Finite-Rank Spectral Converse}
\label{subsec:grounding-spectral}
%----------------------------------------------------------------------

\begin{remark}[Algebraic Hankel Rank Versus Bounded-Operator Rank]
\label{rem:algebraic-vs-operator-rank}
A finite-dimensional weighted automaton factors its formal Hankel matrix as
\[
h(uv)=\alpha^\top T_uT_v\eta,
\]
which bounds the \emph{algebraic} matrix rank by the state dimension.  This
does not by itself imply that the infinite matrix defines a bounded, let alone
compact, operator on $\ell^2$: the factor rows and columns need not be square
summable.  The spectral statements below are therefore conditional on
compactness of $H_\nu$, and algebraic-rank hypotheses do not substitute for
it.
\end{remark}

\begin{theorem}[Spectral Bound for Linear Grounding]
\label{thm:spectral-grounding}
Let $\nu$ be a stochastic channel whose impulse response defines a compact
Hankel operator $H_\nu$ on $\ell^2((\mathcal I\times\mathcal O)^*)$.
Define the unrestricted linear rank-$M$ gap by
\[
\Dunres(M)
=
\inf_{\operatorname{rank}B\le M}
\norm{H_\nu-B}_{\mathrm{op}}.
\]
Then
\[
\Dunres(M)
=
\sigma_{M+1}(H_\nu).
\]
Define the Hankel-structured finite-rank realization gap by
\[
\DHankstr(M)
=
\inf_{\substack{
\operatorname{rank}B\le M\\
B\ \text{Hankel}
}}
\norm{H_\nu-B}_{\mathrm{op}}.
\]
Then
\[
\DHankstr(M)
\ge
\sigma_{M+1}(H_\nu).
\]
Equality holds if an optimal rank-$M$ truncation can be chosen Hankel.
Theorem~\ref{thm:aak-equality} below identifies a checkable hypothesis under
which equality is automatic.
\end{theorem}

\begin{proof}
The unrestricted equality is Meta-Theorem~\ref{meta:spectral}, applied to the
Hilbert-module task theory induced by $\nu$.  By definition,
\[
\Dunres(M)
=
\inf_{\operatorname{rank}B\le M}
\norm{H_\nu-B}_{\mathrm{op}}.
\]
Eckart--Young--Mirsky gives
\[
\inf_{\operatorname{rank}B\le M}
\norm{H_\nu-B}_{\mathrm{op}}
=
\sigma_{M+1}(H_\nu).
\]
The Hankel-structured feasible set is a subset of the unrestricted feasible
set, so restricting the infimum cannot decrease the value.  Hence
\[
\DHankstr(M)
\ge
\Dunres(M)
=
\sigma_{M+1}(H_\nu).
\]
\end{proof}

\begin{corollary}[Spectral Converse under Operator-Norm Domination]
\label{cor:grounding-domination}
Let $\mathcal L_\infty$ be a grounding cost and let $\GrdLin(M)$ denote the
linear finite-rank residual-cost gap.  If
\[
\mathcal L_\infty(\nu,\widehat\nu)
\ \ge\
\|H_\nu-H_{\widehat\nu}\|_{\mathrm{op}}
\]
for every finite-rank residual approximant $\widehat\nu$, then
\[
\GrdLin(M)\ \ge\ \sigma_{M+1}(H_\nu).
\]
\end{corollary}

\begin{proof}
Every budget-$M$ approximant induces a residual operator $H_{\widehat\nu}$ of
rank at most $M$.  By hypothesis the grounding cost dominates the operator-norm
distance, so the infimum of the cost over such approximants is at least the
infimum of $\|H_\nu-H_{\widehat\nu}\|_{\mathrm{op}}$ over rank-$M$ operators,
which is $\sigma_{M+1}(H_\nu)$ by Eckart--Young--Mirsky
\cite{eckart1936,mirsky1960}, as in Theorem~\ref{thm:spectral-grounding}.
\end{proof}

\begin{remark}[Domination Is the Only Input]
\label{rem:grounding-domination-scope}
Corollary~\ref{cor:grounding-domination} is the grounding instance of the
domination pattern used throughout: a converse in a spectral quantity follows
from a single inequality between the task cost and an operator norm, with no
further structure on $\mathcal L_\infty$.  It supplies a lower bound only; the
reverse inequality requires the achievability hypotheses discussed in
Section~\ref{sec:oracle}.
\end{remark}

\begin{corollary}[Finite-Rank and Compact Limits]
\label{cor:grounding-finite-rank}
Under the compactness hypothesis of Theorem~\ref{thm:spectral-grounding},
\[
\Dunres(M)\longrightarrow0
\qquad\text{as }M\to\infty .
\]
If moreover $\operatorname{rank}(H_\nu)=r<\infty$, then $\Dunres(M)=0$ for all
$M\ge r$.  The corresponding statements for $\DHankstr(M)$ require the
additional Hankel-approximation hypotheses of
Theorems~\ref{thm:aak-equality} and~\ref{thm:aak-multiletter}.
\end{corollary}

\begin{proof}
Compactness of $H_\nu$ gives $\sigma_{M+1}(H_\nu)\to0$, and
Theorem~\ref{thm:spectral-grounding} identifies $\Dunres(M)$ with
$\sigma_{M+1}(H_\nu)$.  If $\operatorname{rank}(H_\nu)=r$ then
$\sigma_{M+1}(H_\nu)=0$ for every $M\ge r$.
\end{proof}

\begin{proposition}[Finite-Section Certificates for Linear Grounding]
\label{prop:grounding-finite-section}
Let $H$ be a compact Hankel operator on $\ell^2(\Sigma^*)$, and let $P_N$
denote the orthogonal projection onto the span of the words of length at
most $N$.  Put $H_N=P_NHP_N$.  Then, for every $M\ge0$,
\[
\sigma_{M+1}(H_N)
\le
\sigma_{M+1}(H)
=
\Dunres(M)
\le
\sigma_{M+1}(H_N)+\norm{H-H_N}_{\mathrm{op}}.
\]
Consequently $\sigma_{M+1}(H_N)\to\sigma_{M+1}(H)$, with the displayed norm
error, and the same lower certificate applies to the structured gap:
\[
\DHankstr(M)\ge\sigma_{M+1}(H_N).
\]
If the finite matrix entries of $H_N$ are computable (respectively
rational), these are computable (respectively algebraic) finite
certificates.  Any independent bound on $\norm{H-H_N}_{\mathrm{op}}$ turns
the certificate into a two-sided error interval.
\end{proposition}

\begin{proof}
For every rank-$\le M$ operator $B$,
\[
\norm{H-B}_{\mathrm{op}}
\ge
\norm{P_N(H-B)P_N}_{\mathrm{op}}
\ge
\inf_{\rank C\le M}\norm{H_N-C}_{\mathrm{op}}
=
\sigma_{M+1}(H_N),
\]
where the last equality is finite-dimensional Eckart--Young--Mirsky.  Taking
the infimum over $B$ proves the first inequality.  The reverse estimate
follows from the Lipschitz bound
$|\sigma_j(A)-\sigma_j(B)|\le\norm{A-B}_{\mathrm{op}}$.  Since $P_N\to I$
strongly and $H$ is compact, $\norm{H-P_NHP_N}_{\mathrm{op}}\to0$.  Finally,
the structured feasible set is a subset of the unrestricted one, so
restricting the infimum in $\DHankstr(M)$ cannot decrease it below
$\sigma_{M+1}(H_N)$ by the same argument applied to $\DHankstr$ in place of
$\Dunres$.
\end{proof}

\begin{proposition}[Exact Zero Threshold for the Structured Linear Problem]
\label{prop:grounding-structured-zero}
Let $H$ be a bounded Hankel operator.  Then
\[
\DHankstr(M)=0
\quad\Longleftrightarrow\quad
\operatorname{rank}(H)\le M.
\]
Thus, whenever $H$ has finite rank $r$, both the unrestricted and the
Hankel-structured linear gaps vanish exactly for $M\ge r$.
\end{proposition}

\begin{proof}
If $\operatorname{rank}(H)\le M$, then $H$ itself is an admissible structured
approximant, so $\DHankstr(M)=0$.  Conversely, suppose $\DHankstr(M)=0$;
choose rank-$\le M$ Hankel operators $B_n$ with
$\norm{H-B_n}_{\mathrm{op}}\to0$.  The set of bounded operators of rank at
most $M$ is norm closed: if $\operatorname{rank}(H)>M$, choose $M+1$ vectors
whose images under $H$ are linearly independent; their Gram determinant is
positive and remains positive under sufficiently small operator-norm
perturbations, so no operator of rank at most $M$ can lie arbitrarily close
to $H$ in operator norm.  Hence $\operatorname{rank}(H)\le M$.
\end{proof}

\begin{theorem}[Hankel-Structured Equality via Adamjan--Arov--Krein]
\label{thm:aak-equality}
We treat the scalar case first.  Suppose $|\Sigma|=1$, so that
$\ell^2(\Sigma^*)\cong\ell^2(\mathbb N)$ carries a single unilateral shift $S$
of multiplicity one, and an operator $B$ is Hankel exactly when $S^{*}B=BS$.
Let $S_+$ be the unilateral shift on the Hardy space $H^2$ of the disc.  (For
$|\Sigma|>1$ see Theorem~\ref{thm:aak-multiletter} and
Remark~\ref{rem:aak-multialphabet-scope}: the free monoid $\Sigma^*$ does
\emph{not} carry a natural shift of multiplicity one, and the scalar theory
below does not apply directly.)  Assume there is a unitary
$U:\ell^2(\Sigma^*)\to H^2$ such that
\begin{enumerate}[label=(\alph*)]
\item $U$ \textbf{intertwines the shifts}: $US=S_+U$; consequently $B$ is
Hankel in the sense above if and only if $UBU^{*}$ satisfies
$S_+^{*}(UBU^{*})=(UBU^{*})S_+$, i.e.\ if and only if $UBU^{*}$ is a
classical Hardy-space Hankel operator, so $U$ carries the feasible set
\[
\{B:\operatorname{rank}B\le M,\ B\ \text{Hankel}\}
\]
\emph{bijectively} onto the classical rank-$\le M$ Hankel feasible set;
\item $UH_\nu U^{*}=H_\varphi$ for some symbol
$\varphi\in L^\infty(\mathbb T)$ with $H_\varphi$ compact, equivalently
$\varphi\in H^\infty+C(\mathbb T)$.
\end{enumerate}
Then for every $M\ge0$
\[
\DHankstr(M)
=
\Dunres(M)
=
\sigma_{M+1}(H_\nu).
\]
Moreover an optimal approximant may be taken of the form $H_\psi$ with
$\mathbb P_-\psi$ rational of degree at most $M$, and the optimal error is
attained.  Uniqueness of the optimal rank-$\le M$ Hankel approximant is
\emph{not} asserted in general; it holds under the classical simplicity
hypothesis $\sigma_{M+1}(H_\varphi)>\sigma_{M+2}(H_\varphi)$, equivalently
that the Schmidt pair at level $M+1$ is unique up to a unimodular scalar
\cite[Ch.~4]{peller2003}.  Absent that hypothesis, existence and the value
$\sigma_{M+1}$ are the safe AAK claims made here.  The two descriptions of the feasible
set agree by Kronecker's theorem: a Hankel operator has finite rank exactly
when the antianalytic part of its symbol is rational, and then
\[
\rank H_\psi=\deg\mathbb P_-\psi ,
\]
where for a rational $r=p/q$ in lowest terms $\deg r=\max(\deg p,\deg q)$,
equivalently the sum of the multiplicities of the poles of $r$ including a
possible pole at infinity \cite[Sec.~4]{peller2003}.

\emph{Indexing convention.}  Singular values are indexed here from
$\sigma_1=\norm{H_\nu}_{\mathrm{op}}$, so the rank-at-most-$M$ error is
$\sigma_{M+1}$.  The operator-theory literature indexes from
$s_0(T)=\norm{T}$, defining $s_m(T)=\inf\{\norm{T-R}:\rank R\le m\}$; the two
conventions are related by $\sigma_{m+1}=s_m$, and the statement above is the
case $m=M$.  Rank budget and symbol degree carry the \emph{same} index $M$,
because Kronecker's theorem is an equality rather than an inequality; there is
no off-by-one between them.
\end{theorem}

\begin{proof}
Under the stated hypotheses the Adamjan--Arov--Krein theorem
\cite{aak1971} states exactly that the distance from $H_\varphi$ to the set of
Hankel operators of rank at most $M$ equals $\sigma_{M+1}(H_\varphi)$, and that
the infimum is attained by a Hankel operator $\Gamma_M$ of rank at most $M$,
equivalently one whose symbol has rational antianalytic part of degree at most
$M$.  In the source convention this reads $\norm{\Gamma-\Gamma_m}=s_m(\Gamma)$
with $s_m=\sigma_{m+1}$.  We do not reproduce the proof, which is classical;
see also \cite{peller2003} for a modern treatment.

Transporting by $U$ preserves singular values and the operator norm, as for
any unitary.  Preservation of \emph{Hankel structure}, by contrast, is not
automatic for a general unitary and is supplied here by hypothesis~(a): the
intertwining relation $US=S_+U$ gives, for any bounded $B$,
\[
S^{*}B=BS
\quad\Longleftrightarrow\quad
S_+^{*}(UBU^{*})=(UBU^{*})S_+,
\]
so $U$ maps the Hankel feasible set onto the classical Hankel feasible set and
$U^{*}$ maps it back.  Since $U$ also preserves rank, the two constrained
infima defining $\DHankstr(M)$ and its Hardy-space counterpart are over
corresponding feasible sets and have equal values.  Hence the identity
descends to $H_\nu$.
Combining with the unrestricted equality of
Theorem~\ref{thm:spectral-grounding} gives the chain
\[
\sigma_{M+1}(H_\nu)
=
\Dunres(M)
\le
\DHankstr(M)
=
\sigma_{M+1}(H_\nu),
\]
so all three coincide.
\end{proof}

\begin{remark}[Two Rank Conventions, and Why They Differ]
\label{rem:rank-conventions}
Two spectral-tail conventions appear in this paper.  They are both correct and
must not be conflated; the difference is structural, not notational.

\emph{(G) Operator-approximation form.}  In the grounding regime the budget $M$
constrains an operator directly, $\rank B\le M$, and Eckart--Young--Mirsky or
Adamjan--Arov--Krein give the error $\sigma_{M+1}$.  Budget $M$, rank $M$,
tail starting at index $M+1$.

\emph{(R) Quotient form.}  In the retention regime the budget $M$ constrains a
\emph{partition} into $M$ blocks, and the induced between-block covariance
$B_\phi$ has $\rank B_\phi\le M-1$: the $M$ centroids satisfy the single affine
relation $\sum_k\pi(C_k)(\bar p_{C_k}-\bar p)=0$, so centring removes one
degree of freedom.  Ky Fan then bounds $\tr(B_\phi)$ by
$\sum_{i=1}^{M-1}\lambda_i$ and the tail is $\sum_{i\ge M}\lambda_i$.  Budget
$M$, effective rank $M-1$, tail starting at index $M$.

The one-step offset between the two is therefore the codimension of the
centring constraint, and it is a property of the feasible set rather than of
the indexing.  Writing the retention tail as $\sum_{i>M}\lambda_i$ would
silently import convention~(G) into a regime where it is false, and would not
recover the converse.  Conversely, a grounding tail written $\sum_{i>M}
\sigma_i$ is the same as $\sigma_{M+1}$ onward and is correct.  These are the
``effective rank budgets'' $M$ and $M-1$ recorded in
Theorem~\ref{thm:schatten-nogo}, and they are one of the reasons no
unconditional cross-regime inequality holds.
\end{remark}

\begin{theorem}[Multi-Letter Case, Conditional Form]
\label{thm:aak-multiletter}
Let $|\Sigma|>1$.  The free monoid $\Sigma^*$ carries $|\Sigma|$ prefix
shifts $S_a$, $a\in\Sigma$, jointly an isometry of multiplicity $|\Sigma|$,
rather than a single unilateral shift of multiplicity one; consequently the
classical scalar Adamjan--Arov--Krein theorem does not apply to $H_\nu$ as
stated.  Suppose nevertheless that there exist a Hilbert space $\mathcal K$
carrying a multi-shift or free-semigroup Hankel structure, a unitary
$U:\ell^2(\Sigma^*)\to\mathcal K$ intertwining the respective shift systems and
carrying
\[
\{B:\operatorname{rank}B\le M,\ B\ \text{Hankel on }\ell^2(\Sigma^*)\}
\]
bijectively onto the corresponding rank-$\le M$ Hankel class on $\mathcal K$,
and an Adamjan--Arov--Krein/Nehari-type finite-rank approximation theorem valid
on $\mathcal K$.  Then
\[
\DHankstr(M)=\Dunres(M)=\sigma_{M+1}(H_\nu).
\]
Absent such a structure-preserving embedding, only the Eckart--Young--Mirsky
inequality
\[
\DHankstr(M)\ \ge\ \sigma_{M+1}(H_\nu)
\]
is asserted.
\end{theorem}

\begin{proof}
Identical to the scalar case: the hypothesised $U$ preserves rank, operator
norm, and by assumption the Hankel feasible sets, so the two constrained
infima agree, and the assumed approximation theorem on $\mathcal K$ supplies
the value $\sigma_{M+1}$.  Nothing beyond transport of structure is used.
\end{proof}

\begin{remark}[Multi-Alphabet Scope]
\label{rem:aak-multialphabet-scope}
The hypothesis of Theorem~\ref{thm:aak-multiletter} is strong and is
\emph{not} known to be checkable in general.  Since this manuscript routinely
uses the joint alphabet $\Sigma=\mathcal I\times\mathcal O$, for which
$|\Sigma|>1$ in every nondegenerate case, the scalar AAK theorem should not be
presented as automatically applicable.  Identifying the correct free-semigroup
or vector-valued Hardy model, and determining whether finite-rank AAK
approximation survives the equivalence, is an operator-theoretic
problem; it is recorded as
Open Problem~\ref{open:hankel-multiletter}.
\end{remark}

\begin{openproblem}[Intrinsic Multi-Alphabet Hankel Equality]
\label{open:hankel-multiletter}
For $|\Sigma|>1$, characterize intrinsically -- in terms of the channel $\nu$
-- when the Hankel operator on $\ell^2(\Sigma^*)$ admits a structure-preserving
embedding into a vector-valued Hardy or full Fock space Hankel class carrying
an Adamjan--Arov--Krein/Nehari finite-rank approximation theorem, and hence
when $\DHankstr(M)=\sigma_{M+1}(H_\nu)$.  \emph{Established baseline:}
Theorem~\ref{thm:aak-equality} settles the scalar case $|\Sigma|=1$
unconditionally, and Corollary~\ref{cor:hankel-strict} converts a strict
gap into an obstruction, but only in that scalar model.
\emph{Falsifiable success criterion:} an intrinsic (i.e.\ stated purely in
terms of $\nu$, with no auxiliary embedding datum) necessary and sufficient
condition for $\DHankstr(M)=\sigma_{M+1}(H_\nu)$ to hold for all $M$
simultaneously, or a proof that no such $\nu$-intrinsic condition exists
(e.g.\ by exhibiting two channels with identical $\nu$-intrinsic invariants
of every fixed finite type but differing equality status).
\end{openproblem}

\begin{corollary}[When the Structured Gap Is Strict]
\label{cor:hankel-strict}
If $\DHankstr(M)>\sigma_{M+1}(H_\nu)$ for some $M$, then no
\emph{shift-intertwining} Hardy-space symbol as in
Theorem~\ref{thm:aak-equality} exists; that is, either no unitary intertwines
the two shifts, or no compact Hardy symbol is available.  Consequently a strict
structured gap is an obstruction to Hardy-space representability, not merely a
failure of the approximation argument.
\end{corollary}

\begin{proof}
Contrapositive of Theorem~\ref{thm:aak-equality}.
\end{proof}

\begin{remark}[Scope: Scalar Shift Model Only]
\label{rem:hankel-strict-scope}
Corollary~\ref{cor:hankel-strict} is stated for, and its proof uses, the
scalar shift model of Theorem~\ref{thm:aak-equality}, i.e.\ $|\Sigma|=1$.  In
the multi-letter setting $|\Sigma|>1$, Theorem~\ref{thm:aak-multiletter}
requires an explicit multi-shift/free-semigroup embedding hypothesis that is
not established in general; without it the contrapositive has no content,
since the hypothesis of Theorem~\ref{thm:aak-equality} does not apply and a
strict structured gap is then not automatically an obstruction to any
representability notion beyond the one stated in
Theorem~\ref{thm:aak-multiletter}.  Corollary~\ref{cor:hankel-strict} should
not be read as an intrinsic representability obstruction outside the scalar
model absent the extra structure assumed there.
\end{remark}

\begin{remark}[Unrestricted versus Hankel-Restricted Feasible Sets]
\label{rem:grounding-unrestricted-restricted}
The equality
\[
\Dunres(M)
=
\sigma_{M+1}(H_\nu)
\]
is for the unrestricted rank-$M$ operator-norm approximation problem.  The
Hankel-structured finite-rank realization gap satisfies
\[
\DHankstr(M)
\ge
\sigma_{M+1}(H_\nu),
\]
because restricting the feasible set cannot decrease the infimum.  Equality is
not guaranteed without additional structure.
\end{remark}

\begin{remark}[Infinite Support Does Not Imply Unboundedness]
\label{rem:infinite-support-grounding}
A deterministic Mealy machine may induce an impulse response with infinite
support.  Infinite support alone does not imply that the associated Hankel
operator is unbounded.  Boundedness requires decay or summability conditions.
For example,
\[
h(u)=\gamma^{|u|}
\]
with $0<\gamma<1/|\Sigma|$ defines a bounded Hankel-type operator by the
Schur-test domination of Proposition~\ref{prop:stable-decay}.  Bounded symbolic
impulse responses that fail to decay may or may not define bounded Hankel
operators; boundedness is checked case by case.
\end{remark}

\begin{remark}[AAK Scope]
\label{rem:grounding-aak}
The default spectral converse is Eckart--Young--Mirsky.  Adamjan--Arov--Krein
applies only if a Hardy-space embedding is separately established.  The
grounding theory does not assume such an embedding by default.
\end{remark}

\begin{openproblem}[Exact Symbolic Grounding Gap]
\label{open:symbolic-grounding}
Determine the exact symbolic grounding gap
\[
\Delta_{\mathrm{grd}}^{\mathrm{symb}}(M)
=
\Delta_{\mathrm{grd}}(M;\gamma)
\]
of Definition~\ref{def:symbolic-grounding-gap}, for deterministic Mealy
machines.  The spectral tail
$\sigma_{M+1}(H_\nu)$ is an exact theorem for the unrestricted linear
finite-rank relaxation, not for exact symbolic deterministic approximation.
Relating the linear relaxation gap to the symbolic deterministic gap requires
additional structure, such as Boolean realization theorems, stochasticity
floors, or superpolynomial determinism gaps.  \emph{Established baseline:}
Theorem~\ref{thm:observable-floor} and
Proposition~\ref{prop:symbolic-finite-horizon} give the $M\to\infty$ limit and
certified finite-horizon intervals; see the discussion in
Section~\ref{subsec:open-problems}.  \emph{Falsifiable success criterion:}
either an exact closed form for $\Delta_{\mathrm{grd}}^{\mathrm{symb}}(M)$ at
finite $M$, or a proof that no such closed form exists in a stated class of
formulas (e.g.\ rational functions of the transition/emission data and
$\gamma$).
\end{openproblem}

%----------------------------------------------------------------------
\subsection{Max-Divergence and Operator-Norm Distance}
\label{subsec:grounding-maxdiv}
%----------------------------------------------------------------------

The exponent correspondence organizes grounding by the $\alpha=\infty$
supremum vertex.  The $\alpha=\infty$ R\'enyi vertex is max-divergence.  The
exact grounding cost is operator-norm distance.  These are distinct
functionals, but they are locally equivalent under bounded-overlap hypotheses.

\begin{definition}[Symmetrized Max-Divergence]
\label{def:dmax}
For finite-dimensional positive definite states $\rho,\sigma$, define
\[
\Dmax(\rho\|\sigma)
=
\log
\left\|
\sigma^{-1/2}\rho\sigma^{-1/2}
\right\|_\infty,
\]
with $+\infty$ if $\supp\rho\not\subseteq\supp\sigma$.
Define the symmetrized version
\[
\Dmaxsym(\rho\|\sigma)
=
\max\{
\Dmax(\rho\|\sigma),
\Dmax(\sigma\|\rho)
\}.
\]
\end{definition}

\begin{proposition}[Local Equivalence with Operator-Norm Distance]
\label{prop:grounding-local-equivalence}
Let $\rho,\sigma$ be finite-dimensional states.
\begin{enumerate}[label=(\alph*)]
\item Suppose $\rho,\sigma\ge mI$ for some $m>0$, and
\[
\norm{\rho-\sigma}_\infty\le\varepsilon.
\]
Then
\[
\Dmaxsym(\rho\|\sigma)
\le
\log\left(1+\frac{\varepsilon}{m}\right).
\]
\item Conversely, if
\[
\Dmaxsym(\rho\|\sigma)\le\delta,
\]
then
\[
e^{-\delta}\sigma
\le
\rho
\le
e^\delta\sigma,
\]
and therefore
\[
\norm{\rho-\sigma}_\infty
\le
(e^\delta-1)
\norm{\sigma}_\infty.
\]
For density operators, $\norm{\sigma}_\infty\le1$, so
\[
\norm{\rho-\sigma}_\infty
\le
e^\delta-1.
\]
\end{enumerate}
\end{proposition}

\begin{proof}
For (a), if $\norm{\rho-\sigma}_\infty\le\varepsilon$, then
\[
\rho\le\sigma+\varepsilon I.
\]
Since $\sigma\ge mI$, we have $I\le\sigma/m$, hence
\[
\rho
\le
\sigma+\frac{\varepsilon}{m}\sigma
=
\left(1+\frac{\varepsilon}{m}\right)\sigma.
\]
Therefore
\[
\Dmax(\rho\|\sigma)
\le
\log\left(1+\frac{\varepsilon}{m}\right).
\]
Exchanging $\rho$ and $\sigma$ gives the symmetrized bound.

For (b), $\Dmaxsym(\rho\|\sigma)\le\delta$ means both
\[
\rho\le e^\delta\sigma
\]
and
\[
\sigma\le e^\delta\rho.
\]
The second inequality is equivalent to
\[
e^{-\delta}\sigma\le\rho.
\]
Thus
\[
-(1-e^{-\delta})\sigma
\le
\rho-\sigma
\le
(e^\delta-1)\sigma.
\]
Since $1-e^{-\delta}\le e^\delta-1$, we obtain
\[
-(e^\delta-1)\sigma
\le
\rho-\sigma
\le
(e^\delta-1)\sigma.
\]
Taking operator norms gives
\[
\norm{\rho-\sigma}_\infty
\le
(e^\delta-1)
\norm{\sigma}_\infty.
\]
For density operators, $\norm{\sigma}_\infty\le1$.
\end{proof}

\begin{theorem}[Grounding Vertex]
\label{thm:grounding-vertex}
The grounding cost is the $C^*$-norm / Schatten-$\infty$ distance
\[
\text{grounding cost}
=
\norm{a-b}_\infty
\]
between residual operators or kernels.  The $\alpha=\infty$ R\'enyi vertex is
the max-divergence $\Dmax$, or its symmetrization $\Dmaxsym$.  These are
distinct functionals.

On commuting diagonal states,
\[
\Dmax(p\|q)
=
\log\max_i\frac{p_i}{q_i},
\qquad
\norm{\operatorname{diag}(p)-\operatorname{diag}(q)}_\infty
=
\max_i|p_i-q_i|.
\]
They are related by Proposition~\ref{prop:grounding-local-equivalence} under
bounded-overlap hypotheses, but they are not identical.
\end{theorem}

\begin{remark}[Scope of the Grounding-Vertex Identification]
\label{rem:grounding-vertex-scope}
Theorem~\ref{thm:grounding-vertex} identifies the grounding cost with the
$C^*$-norm / Schatten-$\infty$ distance only in the special operator-norm
grounding model of Section~\ref{subsec:grounding-instantiation}; the
preceding general instantiation allows an arbitrary Banach residual norm, for
which no such identification is asserted absent the domination hypothesis of
Corollary~\ref{cor:grounding-domination}.  Moreover,
Proposition~\ref{prop:grounding-local-equivalence} compares
\emph{finite-dimensional density operators} locally in $\Dmaxsym$ versus
$\norm{\cdot}_\infty$; it does \emph{not} compare $\Dmax$ to the
infinite-dimensional Hankel-operator distance
$\norm{H_\nu-H_{\widehat\nu}}_{\mathrm{op}}$, and no such comparison is
established in this manuscript.
\end{remark}

\begin{remark}[Supremum Vertex as Organization]
\label{rem:grounding-supremum-organization}
The max-divergence vertex is the divergence-face analogue of worst-case or
supremum aggregation.  It organizes the grounding regime conceptually, but the
exact spectral converse remains the Schatten-$\infty$ / operator-norm tail for
the unrestricted linear finite-rank Hankel relaxation:
\[
\Dunres(M)
=
\sigma_{M+1}(H_\nu).
\]
The Hankel-structured finite-rank realization gap satisfies the lower bound
\[
\DHankstr(M)
\ge
\sigma_{M+1}(H_\nu),
\]
with equality only under additional structural hypotheses.
\end{remark}

%----------------------------------------------------------------------
\subsection{Determinism Gap for Formal Hankel Rank}
\label{subsec:exp-determinism}
%----------------------------------------------------------------------

The linear finite-rank spectral converse can be small even when deterministic
symbolic approximation is hard.  The following theorem gives a
superpolynomial separation between stochastic linear rank and deterministic
state complexity; see Remark~\ref{rem:exp-gap-scope} for why ``exponential''
would overstate the proved rate.

\begin{theorem}[Determinism Gap for Formal Hankel Rank]
\label{thm:exp-gap}
Let $\mathcal S(k)$ denote any lower bound on the number of states of a
deterministic finite automaton recognizing a language that some $k$-state
probabilistic finite automaton recognizes with a fixed isolated cutpoint.  Then
there exists a family of causal channels $\{\nu_k\}$ whose formal Hankel
matrices have \emph{algebraic} rank at most $k$, such that any deterministic
Mealy machine approximating $\nu_k$ with worst-case per-step total-variation
error $<1/4$ requires at least $\mathcal S(k)/|\mathcal O|$ states.

The best unconditional value of $\mathcal S$ available in the literature is
that of Ambainis \cite{ambainis1996},
\[
\mathcal S(k)
=
\Omega\!\left(2^{\,k\log\log k/\log k}\right),
\]
which is superpolynomial but \emph{sub}-exponential, since the exponent is
$o(k)$.  A genuinely exponential $\mathcal S(k)=2^{\Omega(k)}$ is not
established by the classical references and is not assumed here.
\end{theorem}

\begin{proof}
In the isolated-cutpoint framework of \cite{rabin1963,freivalds1981}, for every
$k$ there exists a regular language
$L_k$ recognized by a $k$-state probabilistic finite automaton with bounded
error at least $1/4$:
\[
x\in L_k
\implies
p_{\mathrm{acc}}(x)>3/4,
\]
and
\[
x\notin L_k
\implies
p_{\mathrm{acc}}(x)<1/4,
\]
while the minimal deterministic finite automaton for $L_k$ has at least
$\mathcal S(k)$ states.

Construct a causal binary-output channel $\nu_k$ that, after reading prefix
$x$, emits a Bernoulli random variable with parameter
$p_{\mathrm{acc}}(x)$.  This channel is realizable by a $k$-state stochastic
automaton.  Its scalar impulse response has a $k$-dimensional linear
representation
\[
h_k(uv)=\alpha^\top T_uT_v\eta,
\]
so the formal Hankel matrix $\mathsf H_k(u,v)=h_k(uv)$ has \emph{algebraic}
rank at most $k$.  This factorization does \emph{not} assert that
$\mathsf H_k$ defines a bounded, let alone compact, operator on unweighted
$\ell^2$; see Remark~\ref{rem:algebraic-vs-operator-rank}.

A deterministic Mealy machine outputs a bit $b(x)\in\{0,1\}$ at each prefix
$x$.  Its per-step total-variation error for a Bernoulli output is
\[
|p_{\mathrm{acc}}(x)-b(x)|.
\]
If this error is $<1/4$ for all prefixes $x$, then:
\begin{itemize}
\item if $x\in L_k$, then $p_{\mathrm{acc}}(x)>3/4$, so $b(x)=1$;
\item if $x\notin L_k$, then $p_{\mathrm{acc}}(x)<1/4$, so $b(x)=0$.
\end{itemize}
Thus the deterministic machine decides $L_k$ on all prefixes.

To conclude a state lower bound we pass from the Mealy machine to an acceptor.
A Mealy machine assigns outputs to transitions rather than to states, so the
conversion is not immediate.  Given a Mealy machine $(\mathcal S,s_0,\tau,
\lambda)$ whose final emitted symbol decides membership, form the acceptor with
state set $\mathcal S\times\mathcal O$, initial state $(s_0,0)$, transitions
$\bigl((s,y),x\bigr)\mapsto\bigl(\tau(s,x),\lambda(s,x)\bigr)$, and accepting
set $\mathcal S\times\{1\}$.  This records the last emitted symbol in the
state, recognizes exactly the language decided by the Mealy machine, and
multiplies the state count by $|\mathcal O|$, a constant.

Hence an acceptor for $L_k$ with $|\mathcal O|\cdot|\mathcal S|$ states exists,
so $|\mathcal O|\cdot|\mathcal S|\ge\mathcal S(k)$ and
$|\mathcal S|\ge\mathcal S(k)/|\mathcal O|$, which for fixed $\mathcal O$ is
$\Omega(\mathcal S(k))$.
\end{proof}

\begin{remark}[Scope of the Superpolynomial Gap]
\label{rem:exp-gap-scope}
The theorem concerns deterministic Mealy approximation of a stochastic channel
under worst-case per-step total-variation error.  It does not contradict the
linear finite-rank spectral converse: the channel has low \emph{algebraic}
Hankel rank as a formal linear object, but deterministic symbolic realization
may require exponentially many states.

Two scope restrictions are essential.  First, the separation is between
\emph{formal linear realization dimension} and deterministic state complexity;
it is not an instance of the compact-operator spectral theorem, because
boundedness of $\mathsf H_k$ on unweighted $\ell^2$ is not established.
Second, the separation is quantitative rather than exponential.  The classical
references \cite{rabin1963,freivalds1981} supply the isolated-cutpoint
framework but no rate; the best proved rate is
$\mathcal S(k)=\Omega(2^{k\log\log k/\log k})$ \cite{ambainis1996}, whose
exponent is $o(k)$.  The gap is therefore superpolynomial and sub-exponential,
and describing it as ``exponential'' would overstate what is known.  Whether a
truly exponential family exists is open.
\end{remark}

%----------------------------------------------------------------------
\subsection{Stochasticity Floor}
\label{subsec:stochasticity-floor}
%----------------------------------------------------------------------

Even when state complexity is not the obstruction, deterministic machines
cannot eliminate intrinsic stochasticity.

For a fixed input sequence $x$, let
\[
h_t=(x_{1:t},Y_{1:t-1})
\]
be the history available before emitting $Y_t$, and define the local
stochastic spread
\[
s(h_t)
=
1-\max_{b\in\mathcal O}p_{h_t}(b),
\]
where
\[
p_{h_t}(b)
=
\mathbb P(Y_t=b\mid h_t).
\]
The stochasticity coefficient is
\[
\sigma_\gamma(\nu)
=
\sup_{x\in\mathcal I^{\mathbb N}}
\sum_{t=1}^\infty
\gamma^{t-1}
\mathbb E[s(h_t)\mid X=x].
\]

Define the discounted Kantorovich--Rubinstein / Hamming distance between the
channel output law $\nu\mid x$ and a deterministic output sequence
$y=A(x)$ by
\[
d_{\mathrm{KR},\gamma}
\bigl(
\nu\mid x,y
\bigr)
=
\mathbb E_{Y\sim\nu\mid x}
\sum_{t=1}^\infty
\gamma^{t-1}
\mathbf 1\{Y_t\neq y_t\}.
\]
This is the Kantorovich--Rubinstein distance for the discounted Hamming ground
metric.

\begin{definition}[Discounted Symbolic Grounding Gap]
\label{def:symbolic-grounding-gap}
The \textbf{discounted symbolic grounding gap} over deterministic Mealy
machines with at most $M$ states is
\[
\Delta_{\mathrm{grd}}(M;\gamma)
=
\inf_{\substack{
\mathcal A\ \text{deterministic Mealy}\\
|\mathcal S_{\mathcal A}|\le M
}}
\ \sup_{x\in\mathcal I^{\mathbb N}}
d_{\mathrm{KR},\gamma}
\bigl(
\nu\mid x,\mathcal A(x)
\bigr).
\]
This is the symbolic counterpart of the linear relaxation gaps
$\Dunres(M)$ and $\DHankstr(M)$: the feasible set consists of deterministic
machines rather than finite-rank operators, and the cost is a discounted
Hamming distance rather than an operator norm.  Determining it exactly is
Open Problem~\ref{open:symbolic-grounding}, where it is written
$\Delta_{\mathrm{grd}}^{\mathrm{symb}}(M)$ when the discount is suppressed.
\end{definition}

\begin{theorem}[Stochasticity Lower Bound]
\label{thm:stochasticity}
Assume $0<\gamma<1$.  For every deterministic Mealy machine $\mathcal A$,
\[
\sup_{x\in\mathcal I^{\mathbb N}}
d_{\mathrm{KR},\gamma}
\bigl(
\nu\mid x,\mathcal A(x)
\bigr)
\ge
\sigma_\gamma(\nu).
\]
Consequently, with $\Delta_{\mathrm{grd}}(M;\gamma)$ as in
Definition~\ref{def:symbolic-grounding-gap},
\[
\Delta_{\mathrm{grd}}(M;\gamma)
\ge
\sigma_\gamma(\nu)
\]
for every $M$.  Moreover $\sigma_\gamma(\nu)>0$ precisely when there exist an
input sequence $x$ and a time $t$ such that, with positive probability under
$\nu\mid x$, the conditional output law at time $t$ is not a point mass.
Nondegeneracy on unreachable or null histories is insufficient.
\end{theorem}

\begin{proof}
Let $b_t^{\mathcal A}$ be the deterministic output emitted by $\mathcal A$ at
time $t$ under input sequence $x$.  Conditional on the history $h_t$,
\[
\mathbb E[
\mathbf 1_{\{Y_t\neq b_t^{\mathcal A}\}}
\mid h_t
]
=
1-p_{h_t}(b_t^{\mathcal A})
\ge
1-\max_{b\in\mathcal O}p_{h_t}(b)
=
s(h_t).
\]
Therefore
\[
\mathbb E[
\mathbf 1_{\{Y_t\neq b_t^{\mathcal A}\}}
\mid X=x
]
\ge
\mathbb E[
s(h_t)\mid X=x
].
\]
Multiplying by $\gamma^{t-1}$, summing over $t$, and taking the supremum over
$x$ gives
\[
\sup_x
d_{\mathrm{KR},\gamma}
\bigl(
\nu\mid x,\mathcal A(x)
\bigr)
\ge
\sigma_\gamma(\nu).
\]
Taking the infimum over $M$-state deterministic machines gives the stated
lower bound on $\Delta_{\mathrm{grd}}(M;\gamma)$.

For the stated equivalence, suppose first that there exist $x$ and $t$ such
that, with positive probability under $\nu\mid x$, the conditional law
$p_{h_t}(\cdot)$ is not a point mass, i.e.\ $s(h_t)>0$ on an event of
positive $\nu\mid x$-probability.  Then $\mathbb E[s(h_t)\mid X=x]>0$, and
since every summand of $\sigma_\gamma(\nu)$ is non-negative, the $t$-th term
alone gives $\sigma_\gamma(\nu)\ge\gamma^{t-1}\mathbb E[s(h_t)\mid X=x]>0$.
Conversely, if for every $x$ and every $t$ the conditional law
$p_{h_t}(\cdot)$ is a point mass whenever $h_t$ occurs with positive
$\nu\mid x$-probability, then $s(h_t)=0$ on every event of positive
probability, so $\mathbb E[s(h_t)\mid X=x]=0$ for every $x,t$ and
$\sigma_\gamma(\nu)=0$.  This proves the claimed criterion directly, in terms
of positive probability on reachable histories, without invoking any further
undefined non-degeneracy notion.
\end{proof}

\begin{theorem}[Observable Deterministic Floor and Its Finite-State Limit]
\label{thm:observable-floor}
Assume $0<\gamma<1$ and a finite input alphabet.  For an input word
$u\in\mathcal I^t$, let
\[
q_u(b)=\mathbb P(Y_t=b\mid X_{1:t}=u),
\]
where all earlier outputs have been integrated out, and define
\[
F_\gamma(\nu)=
\sup_{x\in\mathcal I^{\mathbb N}}
\sum_{t\ge1}\gamma^{t-1}
\bigl(1-\max_{b\in\mathcal O}q_{x_{1:t}}(b)\bigr).
\]
Then every deterministic causal input--output map $A$ satisfies
\[
\sup_x d_{\mathrm{KR},\gamma}(\nu\mid x,A(x))\ge F_\gamma(\nu).
\]
Equality is attained among unrestricted deterministic causal maps by
choosing, at each input history $u$, a symbol in $\arg\max_bq_u(b)$.
Moreover,
\[
\lim_{M\to\infty}\Delta_{\mathrm{grd}}(M;\gamma)=F_\gamma(\nu).
\]
In particular, $\sigma_\gamma(\nu)\le F_\gamma(\nu)$, and the inequality can
be strict, because $\sigma_\gamma(\nu)$ conditions its local spread on past
random outputs $Y_{<t}$ unavailable to a Mealy machine acting on the input
alone, whereas $q_u$ has those outputs integrated out.
\end{theorem}

\begin{proof}
For a deterministic map $A$, its time-$t$ symbol is a function $a(u)$ of the
observable input history $u=x_{1:t}$.  Hence
\[
\mathbb P(Y_t\ne a(u)\mid X_{1:t}=u)=1-q_u(a(u))
\ge1-\max_bq_u(b).
\]
Summation over $t$ and the supremum over $x$ prove the lower bound.
Pointwise modal choices attain equality because the choice at one input
history changes no other summand.  To obtain finite-state approximants,
store the input prefix through depth $N$, use those modal choices through
time $N$, and thereafter enter one sink state with arbitrary fixed output.
This is a finite Mealy machine with at most
$1+|\mathcal I|+\cdots+|\mathcal I|^{N}$ states, and its excess discounted
loss over $F_\gamma(\nu)$ is at most $\gamma^N/(1-\gamma)$, since every
post-depth-$N$ summand is bounded by $\gamma^{t-1}\cdot1$ and
$\sum_{t>N}\gamma^{t-1}=\gamma^N/(1-\gamma)$.  Letting $N\to\infty$ proves
the limit, using also the lower bound
$\Delta_{\mathrm{grd}}(M;\gamma)\ge F_\gamma(\nu)$ for every $M$ from the
first part.  Finally, conditional Jensen applied to the map
$y\mapsto\max_bp_{h_t}(b\mid Y_{<t}=y)$ gives
\[
\mathbb E\bigl[\max_b\mathbb P(Y_t=b\mid X_{1:t},Y_{<t})\ \big|\ X_{1:t}\bigr]
\ge
\max_b\mathbb P(Y_t=b\mid X_{1:t})=\max_bq_{X_{1:t}}(b),
\]
so summing $\gamma^{t-1}\bigl(1-\mathbb E[\max_bp_{h_t}(b)\mid X_{1:t}]\bigr)
\le\gamma^{t-1}(1-\max_bq_{X_{1:t}}(b))$ and comparing to the definitions of
$\sigma_\gamma$ and $F_\gamma$ gives $\sigma_\gamma(\nu)\le F_\gamma(\nu)$.
\end{proof}

\begin{remark}[Scope: Input-Only Deterministic Class]
\label{rem:observable-floor-scope}
Theorem~\ref{thm:observable-floor} is stated for the manuscript's
input-only deterministic Mealy class of
Definition~\ref{def:symbolic-grounding-gap}.  If the intended deterministic
approximant instead observes feedback of past outputs, both the definition
of $q_u$ and the formula for $F_\gamma$ must be changed consistently; no such
extension is claimed here.
\end{remark}

\begin{proposition}[Finite-Horizon Symbolic Certificates]
\label{prop:symbolic-finite-horizon}
Assume $0<\gamma<1$.  Define
\[
\Delta^{(T)}_{\mathrm{grd}}(M;\gamma)
=
\min_{\substack{A\ \mathrm{deterministic\ Mealy}\\|A|\le M}}
\max_{x\in\mathcal I^T}
\sum_{t=1}^T\gamma^{t-1}
\mathbb P(Y_t\ne A(x)_t\mid X_{1:t}=x_{1:t}).
\]
Then
\[
\Delta^{(T)}_{\mathrm{grd}}(M;\gamma)
\le
\Delta_{\mathrm{grd}}(M;\gamma)
\le
\Delta^{(T)}_{\mathrm{grd}}(M;\gamma)+\frac{\gamma^T}{1-\gamma}.
\]
If the target is a finite-state stochastic channel with rational transition
and emission probabilities and $\gamma$ is rational, then
$\Delta^{(T)}_{\mathrm{grd}}(M;\gamma)$ is a decidable rational number: one
may enumerate the finitely many labeled $M$-state Mealy machines and the
finitely many input words of length $T$, evaluating each displayed loss by
finite forward recursion.
\end{proposition}

\begin{proof}
For each fixed machine and input word, the omitted tail
$\sum_{t>T}\gamma^{t-1}\mathbb P(Y_t\ne A(x)_t\mid X_{1:t}=x_{1:t})$ lies in
$[0,\gamma^T/(1-\gamma)]$, since each summand is a probability in $[0,1]$
scaled by $\gamma^{t-1}$.  Taking first the maximum over input words of
length $T$ and then the minimum over $M$-state machines gives both
inequalities: the machine and input word optimal for the truncated problem,
extended arbitrarily beyond depth $T$, gives the upper bound on
$\Delta_{\mathrm{grd}}(M;\gamma)$, while restricting any machine's supremum
over infinite input sequences to length-$T$ prefixes and truncating the sum
gives the lower bound on $\Delta^{(T)}_{\mathrm{grd}}(M;\gamma)$.  Under the
stated rational finite-state assumptions, every forward recursion computing
$\mathbb P(Y_t\ne A(x)_t\mid X_{1:t}=x_{1:t})$ from the rational transition
and emission probabilities is a finite composition of rational operations,
and every discounted sum over $t\le T$ of rational $\gamma$-powers and
rational probabilities is again rational, while the candidate machine set
(finitely many labeled $M$-state Mealy machines) and input set
($|\mathcal I|^T$ words) are finite, so the minimum and maximum are attained
and computable.
\end{proof}

\begin{corollary}[Exact Memoryless Stochasticity Floor]
\label{cor:memoryless-floor}
Let $\nu$ be a memoryless output channel emitting i.i.d.\ outputs with law
$p$ on $\mathcal O$, independent of history and input.  Then for every
$M\ge1$,
\[
\Delta_{\mathrm{grd}}(M;\gamma)
=
\frac{1-\max_{b\in\mathcal O}p_b}{1-\gamma}.
\]
In particular, for a Bernoulli$(p)$ channel,
\[
\Delta_{\mathrm{grd}}(M;\gamma)
=
\frac{\min\{p,1-p\}}{1-\gamma}.
\]
\end{corollary}

\begin{proof}
For a memoryless channel, the local stochastic spread is constant:
\[
s(h_t)=1-\max_{b\in\mathcal O}p_b.
\]
Therefore
\[
\sigma_\gamma(\nu)
=
\sum_{t=1}^\infty
\gamma^{t-1}
\left(1-\max_{b}p_b\right)
=
\frac{1-\max_b p_b}{1-\gamma}.
\]
The lower bound follows from Theorem~\ref{thm:stochasticity}.

For the upper bound, choose
\[
b^\star\in\arg\max_{b\in\mathcal O}p_b
\]
and use the one-state deterministic machine that emits $b^\star$ at every
time.  Its discounted Hamming error is exactly
\[
\sum_{t=1}^\infty
\gamma^{t-1}
(1-p_{b^\star})
=
\frac{1-\max_b p_b}{1-\gamma}.
\]
Thus the lower bound is attained.
\end{proof}


\begin{remark}[Deterministic Machines Cannot Remove Intrinsic Randomness]
\label{rem:stochasticity-meaning}
The stochasticity floor is independent of the state budget $M$.  No
deterministic Mealy machine, however large, can match a genuinely stochastic
channel in discounted Hamming/Kantorovich--Rubinstein distance better than the
channel's intrinsic stochastic spread.  This obstruction persists even with
unbounded deterministic memory.
\end{remark}

\begin{remark}[Relation to the Linear Spectral Converse]
\label{rem:stochasticity-vs-linear}
The stochasticity floor concerns deterministic symbolic grounding.  The linear
spectral converse concerns finite-rank linear realizations.  A channel may have
small linear Hankel gap while having a positive stochasticity floor for
deterministic machines.  The two results measure different approximation
regimes.
\end{remark}

%----------------------------------------------------------------------
\subsection{State Compression Above the Floor}
\label{subsec:grounding-tracking}
%----------------------------------------------------------------------

Theorem~\ref{thm:stochasticity} bounds the symbolic grounding gap below by a
quantity independent of the budget $M$.  It therefore says nothing about how
the gap varies with $M$.  The following decomposition separates the
budget-independent floor from a budget-dependent term, and exhibits that term
as a mode-prediction analogue of the retention gap.

Throughout this subsection the channel is a stationary finite-state channel
with state set $\mathcal S$, stationary law $\pi$, and predictive laws
$P_s=\mathcal L(Y_t\mid S_t=s)$ on a finite output alphabet, and $\phi$ ranges
over maps $\Splus\to\mathcal K$ with $|\mathcal K|\le M$.  Write
$\pi(C_k)=\sum_{s\in C_k}\pi_s$ for the mass of the block $C_k=\phi^{-1}(k)$.

\begin{definition}[Block Symbol and Tracking Deficit]
\label{def:tracking-deficit}
For a block $C_k$ of positive mass, the \textbf{block symbol} is
\[
\hat b_k
\in
\argmax_{b\in\mathcal O}\ \sum_{s\in C_k}\pi_s P_s(b),
\]
ties broken by any fixed order, and the \textbf{tracking deficit} of $\phi$ is
\[
D(\phi)
=
\sum_{s\in\Splus}\pi_s
\Bigl[
\max_{b}P_s(b)-P_s\bigl(\hat b_{\phi(s)}\bigr)
\Bigr].
\]
The one-step stationary floor is
$\sigma_1=\sum_{s\in\Splus}\pi_s\bigl(1-\max_bP_s(b)\bigr)$.
\end{definition}

\begin{proposition}[Floor Plus Tracking Decomposition]
\label{prop:grounding-tracking}
The following is first an \emph{oracle} state-compression identity: the
emitted symbol may use the current latent state $S_t$ through its block
label $\phi(S_t)$, which is not in general available to a deterministic
Mealy machine acting on the input alone.  It becomes a realizable
deterministic Mealy construction only under the implementability condition
of Proposition~\ref{prop:block-mode-implementability} below.  Let $\phi$ be
any state compression and let $\mathcal A_\phi$ denote the (possibly
non-implementable) rule that, given oracle access to $S_t$, occupies block
$\phi(S_t)$ and emits $\hat b_{\phi(S_t)}$.  Then the stationary one-step
error of $\mathcal A_\phi$ decomposes as
\[
\sum_{s\in\Splus}\pi_s\Bigl(1-P_s\bigl(\hat b_{\phi(s)}\bigr)\Bigr)
=
\sigma_1+D(\phi),
\]
and
\begin{enumerate}[label=(\roman*)]
\item $D(\phi)\ge0$, with equality if and only if $\hat b_{\phi(s)}$ attains
$\max_bP_s(b)$ for every $s\in\Splus$;
\item $D$ vanishes on the singleton compression, so
$\min_{|\mathcal K|\le M}D=0$ whenever $M\ge|\Splus|$;
\item $D$ is monotone under refinement: if $\phi$ refines $\psi$, then
\[
D(\phi)\le D(\psi).
\]
Consequently $\min_{|\mathcal K|\le M}D(\phi)$ is nonincreasing in $M$.
\end{enumerate}
\end{proposition}

\begin{proof}
The identity is immediate from
$1-P_s(\hat b_{\phi(s)})=\bigl(1-\max_bP_s(b)\bigr)+\bigl(\max_bP_s(b)-P_s(\hat b_{\phi(s)})\bigr)$,
weighted by $\pi_s$ and summed.  For~(i), each bracket is a maximum minus a
coordinate of the same vector, hence nonnegative, and the sum of nonnegative
terms vanishes exactly when each does.  For~(ii), the singleton compression has
$\hat b_{\{s\}}=\argmax_bP_s(b)$, so every bracket is zero.

For~(iii), write $w_s(b)=\pi_sP_s(b)$.  The statewise first term
$\sum_s\pi_s(1-\max_bP_s(b))$ in the decomposition is constant across
partitions, whereas the blockwise accuracy is
\[
\sum_{C\in\mathcal P}\max_b\sum_{s\in C}w_s(b),
\]
so that $D(\phi)=\sigma_1-\sum_{C\in\mathcal P_\phi}\max_b\sum_{s\in C}w_s(b)$
for the partition $\mathcal P_\phi$ induced by $\phi$, and likewise for
$\psi$.  If a block $C\in\mathcal P_\psi$ is split by the refinement $\phi$
into $C_1,\dots,C_r\in\mathcal P_\phi$, then, since a maximum of a sum is at
most the sum of the maxima taken termwise over a partition of the summation
range,
\[
\sum_{j=1}^r\max_b\sum_{s\in C_j}w_s(b)
\ge
\max_b\sum_{j=1}^r\sum_{s\in C_j}w_s(b)
=
\max_b\sum_{s\in C}w_s(b).
\]
Applying this to every block of $\mathcal P_\psi$ and summing gives
$\sum_{C'\in\mathcal P_\phi}\max_b\sum_{s\in C'}w_s(b)
\ge\sum_{C\in\mathcal P_\psi}\max_b\sum_{s\in C}w_s(b)$, and subtracting from
$\sigma_1$ reverses the inequality to $D(\phi)\le D(\psi)$.  Since every
budget-$M$ feasible partition is refined by every budget-$M'$ feasible
partition for $M'\ge M$ that further splits it (nesting of feasible sets, as
in the previous version of this claim), monotonicity of $D$ under
refinement upgrades the earlier feasible-set nesting argument to the sharper
conclusion that $\min_{|\mathcal K|\le M}D(\phi)$ is nonincreasing in $M$.
\end{proof}

\begin{proposition}[Implementability of a Block-Mode Policy]
\label{prop:block-mode-implementability}
Suppose the channel is input-driven and synchronized: there is a map
$s:\mathcal I^*\to\mathcal S$ with $S_t=s(X_{1:t})$.  If $\phi$ is a right
congruence for this update, namely there exists
$\bar\tau:\mathcal K\times\mathcal I\to\mathcal K$ such that
\[
\phi(s(ux))=\bar\tau(\phi(s(u)),x)
\qquad(u\in\mathcal I^*,\ x\in\mathcal I),
\]
then the block symbols $\hat b_k$ define a deterministic Mealy machine with
at most $|\mathcal K|$ states, and its stationary one-step error is
$\sigma_1+D(\phi)$ as in Proposition~\ref{prop:grounding-tracking}.  Without
this condition the same formula is only the error of a genie that reveals
$\phi(S_t)$ to the predictor.
\end{proposition}

\begin{proof}
The displayed identity makes the block label $\phi(S_t)=\phi(s(X_{1:t}))$ a
deterministic finite-state update from the input history alone, via
$\bar\tau$, so the machine with state set $\mathcal K$, transition
$\bar\tau$, and output rule $k\mapsto\hat b_k$ is a well-typed Mealy machine
of the form of Section~\ref{subsec:bsg-intro}, and outputting $\hat b_k$ at
state $k$ realizes exactly the oracle rule of
Proposition~\ref{prop:grounding-tracking}.  Conditioning on $S_t=s$ gives
error $1-P_s(\hat b_{\phi(s)})$; averaging over $\pi$ and using
$1-P_s(\hat b_{\phi(s)})=(1-\max_bP_s(b))+
(\max_bP_s(b)-P_s(\hat b_{\phi(s)}))$ proves the formula.  If the label
$\phi(S_t)$ is not a function of the input history alone -- for instance if
$S_t$ depends on the channel's own latent randomness beyond what
$X_{1:t}$ determines -- no input-only Mealy machine can in general implement
$\hat b_{\phi(S_t)}$, and the decomposition remains only a genie bound.
\end{proof}

\begin{remark}[What the Deficit Is, and What It Is Not]
\label{rem:tracking-deficit-scope}
Two points fix the scope of Proposition~\ref{prop:grounding-tracking}.

First, the deficit must be measured against the \emph{block symbol} evaluated
under the true state law, $P_s(\hat b_{\phi(s)})$, and not against the mode of
the block mixture, $\max_bP_{\phi(s)}(b)$.  The latter substitution yields a
quantity that is not a deficit at all: for a two-state block with
$P_{s_1}=(\tfrac12,\tfrac12)$, $P_{s_2}=(\tfrac9{10},\tfrac1{10})$ and equal
weights, the mixture is $(\tfrac7{10},\tfrac3{10})$ and
$\max_bP_{s_1}(b)-\max_bP_{\phi(s_1)}(b)=-\tfrac15<0$.  With the block symbol
the corresponding bracket is $\tfrac12-\tfrac12=0$.

Second, a positive deficit is not forced by $M<|\Splus|$.  Distinct predictive
laws may share a mode, and then merging them costs nothing in this currency:
for $\mathcal O=\{1,2,3\}$, three states of equal weight with
$P_{s_1}=(\tfrac35,\tfrac15,\tfrac15)$,
$P_{s_2}=(\tfrac7{10},\tfrac15,\tfrac1{10})$ and
$P_{s_3}=(\tfrac45,\tfrac1{10},\tfrac1{10})$, the one-block compression has
$\hat b=1$ and $D=0$ although $M=1<3=|\Splus|$.  The vanishing condition is
mode preservation, which is strictly weaker than separating the predictive
states; contrast Theorem~\ref{thm:retention-zero}, where the KL currency does
force one block per distinct predictive law.
\end{remark}

\begin{remark}[Relation to Retention]
\label{rem:tracking-vs-retention}
Proposition~\ref{prop:grounding-tracking} gives the symbolic grounding problem
above the floor the same shape as retention: a compression $\phi$ of the
stationary state set, a representative per block, and a $\pi$-weighted cost of
merging.  The currencies differ.  Retention charges
$D_{\mathrm{KL}}(P_s\|\bar P_k)$ against the mixture centroid, which by
Theorem~\ref{thm:retention-zero} vanishes only when each block carries a single
predictive law.  Grounding charges $\max_bP_s(b)-P_s(\hat b_k)$ against a
single symbol, which vanishes under the weaker mode-preservation condition of
Remark~\ref{rem:tracking-deficit-scope}.  This comparison of zero sets is
itself an oracle-state comparison unless the implementability condition of
Proposition~\ref{prop:block-mode-implementability} holds for a shared family
of partitions on both sides; under that condition the zero-set inclusion is
valid (mode preservation is weaker than equality of predictive laws), because
both problems then range over the same implementable feasible class.  No
stronger claim should be read into it: it does not by itself show that the
grounding compression problem is ``never harder'' than retention in any
global sense, only that its zero set is no smaller.  No inequality between
the two \emph{gap values} follows, the costs being incomparable; this is the
Independence Theorem~\ref{thm:independence} in the present instance.
\end{remark}

%======================================================================
\section{The Conditional Schatten Converse Template}
\label{sec:schatten}
%======================================================================

The analytic converses for grounding and quadratic retention have a common
operator-approximation mechanism.  This section isolates that mechanism in a
conditional template.  The Schatten tail form follows once a regime supplies a
rank-controlled response operator and a Schatten domination inequality.
Contextwise $\ell^p$ aggregation does not by itself imply Schatten-$p$
domination.

\begin{proposition}[Comparison of Schatten Tails of a Common Operator]
\label{prop:common-operator}
Let $A$ have finite rank $\kappa$ and let $0\le r<\kappa$.  For
$1\le p\le q\le\infty$, set $n=\kappa-r$.  Then
\[
\Phi^{(q)}_r(A)
\le
\Phi^{(p)}_r(A)
\le
n^{1/p-1/q}\,
\Phi^{(q)}_r(A),
\]
with the convention $1/\infty=0$.  In particular
\[
\Phi^{(\infty)}_r(A)
\le
\Phi^{(2)}_r(A)
\le
\sqrt n\,\Phi^{(\infty)}_r(A),
\qquad
\Phi^{(\infty)}_r(A)
\le
\Phi^{(1)}_r(A)
\le
n\,\Phi^{(\infty)}_r(A).
\]
\end{proposition}

\begin{proof}
Apply the monotonicity and equivalence inequalities for $\ell^p$ norms on
$\mathbb R^n$, namely $\norm{x}_q\le\norm{x}_p\le n^{1/p-1/q}\norm{x}_q$ for
$1\le p\le q\le\infty$ \cite[Ch.~2]{cover}, to the finite tail vector
$(\sigma_{r+1}(A),\dots,\sigma_\kappa(A))$, whose length is $n=\kappa-r$.  The
displayed special cases are $q=\infty$ with $p=2$ and $p=1$.
\end{proof}

\begin{remark}[Scope of the Comparison]
\label{rem:common-operator-scope}
Proposition~\ref{prop:common-operator} compares different Schatten exponents
for \emph{one} operator and one rank budget.  It does not compare the
different regimes, whose response operators, effective rank budgets, and
constants differ; that absence is the content of
Theorem~\ref{thm:schatten-nogo}.
\end{remark}

\begin{remark}[Relation to Aggregation and Domination]
\label{rem:response-what-supplies}
The definition separates three levels:
\begin{enumerate}[label=(\arabic*)]
\item contextwise aggregation, for example an $\ell^p$-type norm over
contexts;
\item operator approximation, for example a Schatten-$p$ norm over singular
values;
\item finite-state rank control, for example rank $\le M$, rank $\le M-1$,
or Boolean rank $\le M$.
\end{enumerate}
Mirsky's theorem governs the passage from operator approximation and rank
control to a spectral tail.  The implication
\[
\text{contextwise }\ell^p\text{ aggregation}
\quad\Longrightarrow\quad
\text{Schatten-}p\text{ domination}
\]
is not automatic.  The exponent $p$ in the tail is the exponent of the
operator norm appearing in the domination bridge, not automatically the
exponent of contextwise aggregation.

A sufficient way to prove domination is to exhibit a linear measurement or
embedding map from response-operator errors to context errors with a
frame-type lower bound
\[
\norm{\text{context error}}_{\ell^p}
\ge
\frac1{c_p}
\norm{A_\delta-A_{\widehat\delta}}_{\Sp p}.
\]
Without such a bridge, the template does not yield a regime-specific converse.
\end{remark}

\begin{remark}[Nonempty Instantiation]
\label{rem:frobenius-nonempty}
The schema is nonempty.  Let the exact object be a Hilbert--Schmidt operator
$A\in\mathcal S_2$, let budget-$M$ approximants be arbitrary operators $B$
with $\rank(B)\le M$, and define the task cost by
\[
\mathcal L(A,B)
=
\norm{A-B}_{\Sp 2}.
\]
Then
\[
A_\delta=A,
\qquad
A_{\widehat\delta}=B,
\qquad
r(M)=M,
\qquad
c_2=1,
\]
and
\[
\Delta(M)
=
\Phi^{(2)}_M(A)
=
\left(
\sum_{i>M}
\sigma_i(A)^2
\right)^{1/2}.
\]
This is an operator-approximation instantiation.  It does not assert the
existence of a finite-state $\mathcal L^2$ regime satisfying Schatten-$2$
domination.
\end{remark}

\begin{remark}[Conditional Status of the Template]
\label{rem:schatten-conditional-status}
The template is conditional.  It yields a Schatten-tail lower bound only after
a regime supplies:
\begin{enumerate}[label=(\arabic*)]
\item a response operator $A_\delta$;
\item an effective rank budget $r_{\mathbb T}(M)$;
\item a domination inequality relating task loss to a Schatten norm.
\end{enumerate}
Without these data, the template does not apply.
\end{remark}

For a compact operator $A$ on a Hilbert space, write its singular values in
nonincreasing order:
\[
\sigma_1(A)\ge\sigma_2(A)\ge\cdots\ge0.
\]
For $1\le p<\infty$, if $A$ belongs to the Schatten ideal $\mathcal S_p$, the
Schatten-$p$ norm is
\[
\norm{A}_{\Sp{p}}
=
\left(
\sum_i \sigma_i(A)^p
\right)^{1/p},
\]
and for a compact operator outside $\mathcal S_p$ we use the extended value
$\norm{A}_{\Sp{p}}=+\infty$, so that all displayed inequalities are read in
$[0,\infty]$.  Finally
\[
\norm{A}_{\Sp{\infty}}
=
\sigma_1(A).
\]
For an integer rank budget $r\ge0$, define the Schatten-$p$ tail
\[
\Phi^{(p)}_r(A)
=
\begin{cases}
\left(
\sum_{i>r}
\sigma_i(A)^p
\right)^{1/p},
& 1\le p<\infty,\\[2ex]
\sigma_{r+1}(A),
& p=\infty,
\end{cases}
\]
with the convention $\Phi^{(p)}_r(A)=0$ if $r\ge\rank(A)$ in the finite-rank
case.

\begin{lemma}[Linear Measurement Bridges]
\label{lem:frame-bridge}
Suppose the residual difference $z=\delta-\widehat\delta$ enters the response
operators through a \emph{linear} map $T$, so that
$A_\delta-A_{\widehat\delta}=T(z)$, and suppose the task cost is a contextwise
weighted $\ell^p$ aggregation, $\mathcal L_{\mathbb T}=\norm{z}_{\ell^p(w)}$.
If $T$ satisfies the frame-type bound
\[
\norm{T(z)}_{\Sp p}\ \le\ C_p\,\norm{z}_{\ell^p(w)}
\qquad\text{for all }z,
\]
then the domination inequality of Definition~\ref{def:response} holds with
modulus $c_p=C_p$.  Conversely, any modulus $c_p$ valid for all $z$ yields such
a bound with $C_p=c_p$, so for linear $T$ the two conditions are the same.
\end{lemma}

\begin{proof}
Immediate: $\mathcal L_{\mathbb T}=\norm{z}_{\ell^p(w)}\ge
C_p^{-1}\norm{T(z)}_{\Sp p}=c_p^{-1}\norm{A_\delta-A_{\widehat\delta}}_{\Sp p}$,
and the converse is the same inequality with the roles of the two sides
exchanged.
\end{proof}

\begin{proposition}[No Dimension-Free Hankel Bridge]
\label{prop:no-dimension-free-bridge}
Let $T$ be the Hankel embedding sending a sequence $z=(z_0,\dots,z_{2n-2})$ to
the $n\times n$ matrix $H(z)_{ij}=z_{i+j}$, and let the aggregation be
unweighted $\ell^p$.  Then the optimal constants satisfy
\[
C_2(n)=\sqrt n,
\qquad
C_1(n)\ \ge\ n,
\qquad
C_\infty(n)\ \ge\ n .
\]
In particular no dimension-free frame bridge exists for the Hankel embedding at
$p\in\{1,2,\infty\}$, and any modulus $c_p$ obtained this way degrades with the
budget.
\end{proposition}

\begin{proof}
For $p=2$, $\norm{H(z)}_{\Sp2}^2=\sum_k m_k|z_k|^2$ with multiplicity
$m_k=\#\{(i,j):i+j=k\}=\min(k+1,n,2n-1-k)$, so the ratio is maximized by the
atom $z=e_{n-1}$ on the main anti-diagonal, where $m_{n-1}=n$; hence
$C_2(n)=\sqrt n$ exactly.

For $p=1$ take the same atom.  Then $H(e_{n-1})$ is the anti-identity, whose
$n$ singular values all equal $1$, so $\norm{H(e_{n-1})}_{\Sp1}=n$ while
$\norm{e_{n-1}}_{\ell^1}=1$, giving $C_1(n)\ge n$.

For $p=\infty$ take $z\equiv1$.  Then $H(z)$ is the all-ones matrix, of rank
one with operator norm $n$, while $\norm{z}_{\ell^\infty}=1$, giving
$C_\infty(n)\ge n$.
\end{proof}

\begin{remark}[Consequence for the Template]
\label{rem:bridge-consequence}
Lemma~\ref{lem:frame-bridge} identifies exactly what a bridge must supply, and
Proposition~\ref{prop:no-dimension-free-bridge} shows that the most natural
candidate does not supply a constant one.  Two readings are available.

If one insists on a dimension-free (constant) modulus, the linear-measurement
route fails for Hankel embeddings and a different mechanism is needed; this
is the content of the open problem on domination bridges.  Definition~\ref{def:response}
and Theorem~\ref{thm:unified} are stated for the general budget-dependent
modulus $c_p(M)$ precisely to accommodate this case: the template still
applies with $c_p(M)$ in place of a constant $c_p$, but the converse it
yields is correspondingly weaker, and the weakening must be carried
explicitly into every downstream tail bound rather than absorbed into a
constant.

Two scoping points prevent this from being read too broadly.

First, the diagonal embedding $z\mapsto\operatorname{diag}(z)$ is a Schatten
isometry for every $p$, since the singular values of $\operatorname{diag}(z)$
are the $|z_i|$; the retention instance therefore admits a dimension-free
bridge with $C_p=1$.  The obstruction is specific to the overlap structure of
Hankel operators, in which one sequence entry occupies up to $n$ matrix
positions.

Second, the grounding instance of Corollary~\ref{cor:grd-schatten} is
\emph{not} exposed to the obstruction.  There the cost is itself the
operator-norm distance $\norm{H_\nu-H_{\widehat\nu}}_{\mathrm{op}}$ and the
response operator is the same object, so domination holds with equality and
$c_\infty=1$, with no frame constant to pay.
Proposition~\ref{prop:no-dimension-free-bridge} bears on a different
situation: a regime whose cost aggregates residual \emph{sequence entries} in
$\ell^p$, which is then to be compared with a Schatten norm of their Hankel
embedding.  That is precisely the implication ``contextwise $\ell^p$ implies
Schatten-$p$ domination'' which is not automatic, and the proposition
quantifies its failure.
\end{remark}

\begin{definition}[Schatten-$p$ Response Representation]
\label{def:response}
Fix $p\in[1,\infty]$.  A task theory $\mathbb T$ has a
\textbf{Schatten-$p$ response representation} with response operator
$A_\delta$, effective rank budget
\[
r_{\mathbb T}:\mathbb N\to\mathbb N\cup\{0\},
\]
and modulus
\[
c_p:\mathbb N\to(0,\infty)
\]
if every budget-$M$ approximant $\widehat\delta$ is assigned a response
operator $A_{\widehat\delta}$ such that:
\begin{enumerate}[label=(\alph*)]
\item $\rank(A_{\widehat\delta})\le r_{\mathbb T}(M)$;
\item the domination inequality holds:
\[
\mathcal L_{\mathbb T}(\delta,\widehat\delta)
\ge
\frac1{c_p(M)}
\norm{A_\delta-A_{\widehat\delta}}_{\Sp{p}}.
\]
\end{enumerate}
A \textbf{constant-modulus} representation is the special case
$c_p(M)\equiv c_p$ for all $M$; this is the case used by
Corollary~\ref{cor:grd-schatten}, where $c_\infty\equiv1$, and it is the only
case in which the bracketed notation $c_p$ (without an argument) is used
below.
\end{definition}

\begin{theorem}[Eckart--Young--Mirsky]
\label{thm:mirsky}
The following is classical and is quoted without proof; see
\cite{eckart1936,mirsky1960}, and \cite{peller2003} for the operator setting.
Let $A$ be a compact operator on a Hilbert space, and let
$\vertiii{\cdot}$ be a unitarily invariant norm.  For every integer $r\ge0$,
\[
\min_{\rank B\le r}
\vertiii{A-B}
=
\vertiii{
\operatorname{diag}
\bigl(
\sigma_{r+1}(A),
\sigma_{r+2}(A),
\dots
\bigr)
}.
\]
In particular, for every $1\le p\le\infty$,
\[
\min_{\rank B\le r}
\norm{A-B}_{\Sp{p}}
=
\Phi^{(p)}_r(A).
\]
\end{theorem}

\begin{theorem}[Conditional Unified Spectral Converse]
\label{thm:unified}
Suppose a task theory $\mathbb T$ has a Schatten-$p$ response representation
in the sense of Definition~\ref{def:response}, with response operator
$A_\delta$, effective rank budget $r_{\mathbb T}(M)$, and modulus
$c_p(\cdot)$.  Then
\[
\Delta_{\mathbb T}(M)
\ge
\frac1{c_p(M)}
\Phi^{(p)}_{r_{\mathbb T}(M)}(A_\delta).
\]
Equivalently, for $1\le p<\infty$,
\[
\Delta_{\mathbb T}(M)
\ge
\frac1{c_p(M)}
\left(
\sum_{i>r_{\mathbb T}(M)}
\sigma_i(A_\delta)^p
\right)^{1/p},
\]
and for $p=\infty$,
\[
\Delta_{\mathbb T}(M)
\ge
\frac1{c_\infty(M)}
\sigma_{r_{\mathbb T}(M)+1}(A_\delta).
\]
When the modulus is constant, $c_p(M)\equiv c_p$, these recover the familiar
budget-independent bounds, and we then write $c_p$ for $c_p(M)$ as in
Corollary~\ref{cor:grd-schatten}.
\end{theorem}

\begin{proof}
For any budget-$M$ approximant $\widehat\delta$,
\[
\rank(A_{\widehat\delta})
\le
r_{\mathbb T}(M).
\]
Hence, by Theorem~\ref{thm:mirsky},
\[
\norm{A_\delta-A_{\widehat\delta}}_{\Sp{p}}
\ge
\min_{\rank B\le r_{\mathbb T}(M)}
\norm{A_\delta-B}_{\Sp{p}}
=
\Phi^{(p)}_{r_{\mathbb T}(M)}(A_\delta).
\]
By domination at budget $M$,
\[
\mathcal L_{\mathbb T}(\delta,\widehat\delta)
\ge
\frac1{c_p(M)}
\Phi^{(p)}_{r_{\mathbb T}(M)}(A_\delta).
\]
Taking the infimum over budget-$M$ approximants gives the result.
\end{proof}

\begin{remark}[Numerical Verification Scope]
\label{rem:schatten-numerical-scope}
Numerical checks are consistency checks, not proofs of the hypotheses.
\begin{itemize}[leftmargin=1.6em]
\item Mirsky/Eckart--Young identities at $p\in\{1,2,\infty\}$ can be checked
numerically on random finite matrices.  Such checks verify the
operator-approximation theorem, not the domination hypotheses.
\item The $\ell^p$ tail comparisons of Proposition~\ref{prop:common-operator}
can be checked on random finite spectra.  Such checks verify only
elementary norm monotonicity.
\item The retention check uses the tail
\[
\sum_{i\ge M}
\lambda_i(\Sigma_\pi),
\]
not $\sum_{i>M}\lambda_i(\Sigma_\pi)$.  A check using the latter does
not recover the retention converse.
\item Numerical experiments cannot prove Hankel domination for grounding,
cannot prove that a given KL or finite-state cost satisfies Schatten
domination, cannot prove Myhill--Nerode rank identities, and cannot
prove the Independence Theorem.
\end{itemize}
Such checks verify operator identities and indexing; they do not verify full
applicability of the template.
\end{remark}

\begin{remark}[Equality versus Lower Bound]
\label{rem:unified-equality}
If the task cost is exactly the scaled Schatten-$p$ response error and the
feasible response operators include the truncated singular-value approximants,
then Theorem~\ref{thm:unified} becomes an equality.  In general, domination
gives only a lower bound.
\end{remark}

\subsection{Grounding as the Schatten-$\infty$ Instance}
\label{subsec:schatten-grounding}

\begin{corollary}[Grounding as $p=\infty$]
\label{cor:grd-schatten}
Let $H_\nu$ be the response Hankel operator of the target channel.  For the
unrestricted linear finite-rank Hankel relaxation gap,
\[
\Dunres(M)
=
\inf_{\rank B\le M}
\norm{H_\nu-B}_{\Sp{\infty}}.
\]
Theorem~\ref{thm:unified} with
\[
A_\delta=H_\nu,
\qquad
r_{\mathrm{grd}}(M)=M,
\qquad
c_\infty=1
\]
gives
\[
\Dunres(M)
=
\Phi^{(\infty)}_M(H_\nu)
=
\sigma_{M+1}(H_\nu).
\]

For the Hankel-structured finite-rank realization gap,
\[
\DHankstr(M)
=
\inf_{\substack{
\rank B\le M\\
B\ \text{Hankel}
}}
\norm{H_\nu-B}_{\Sp{\infty}},
\]
the feasible set is restricted, and therefore
\[
\DHankstr(M)
\ge
\sigma_{M+1}(H_\nu).
\]
Equality requires additional structure, for example an optimal Hankel
truncation or a valid Hardy-space Adamjan--Arov--Krein embedding.
\end{corollary}

\begin{proof}
The unrestricted gap is by definition the rank-$M$ operator-norm approximation
problem for $H_\nu$, so Eckart--Young--Mirsky
\cite{eckart1936,mirsky1960} gives equality with $\sigma_{M+1}(H_\nu)$;
equivalently, this is Theorem~\ref{thm:unified} applied with $A_\delta=H_\nu$,
$r_{\mathrm{grd}}(M)=M$ and $c_\infty=1$, for which the domination inequality
holds with equality.

The Hankel-structured gap restricts the feasible set from all rank-$M$
operators to rank-$M$ \emph{Hankel} operators.  An infimum over a subset is at
least the infimum over the ambient set, so
$\DHankstr(M)\ge\Dunres(M)=\sigma_{M+1}(H_\nu)$.  Nothing forces equality: the
optimal rank-$M$ truncation supplied by Eckart--Young--Mirsky need not be
Hankel.  Adamjan--Arov--Krein is not invoked here, and is available only when
an independent Hardy-space embedding is supplied
(Theorem~\ref{thm:aak-equality}).

For the residual-cost gap $\GrdLin(M)$, apply Theorem~\ref{thm:unified} with
the assumed domination inequality, as in
Corollary~\ref{cor:grounding-domination}.
\end{proof}

\begin{remark}[Scope of the Grounding Instance]
\label{rem:grd-schatten-scope}
This is a converse for the \emph{unrestricted} linear finite-rank grounding
relaxation.  The Hankel-structured finite-rank realization gap satisfies the
same singular-value lower bound, since restricting the feasible set to Hankel
operators cannot decrease the infimum, but equality there requires the
additional structural hypotheses of Theorems~\ref{thm:aak-equality}
and~\ref{thm:aak-multiletter}.  Transfer to an exact symbolic deterministic gap
requires a separate comparison theorem.
\end{remark}

\subsection{Retention as the Schatten-$1$ Instance with Effective Rank $M-1$}
\label{subsec:schatten-retention}

\begin{corollary}[Retention as $p=1$ with Rank Budget $M-1$]
\label{cor:ret-schatten}
In the Gaussian quadratic restricted retention regime, let
$\Sigma_\pi\succeq0$ be the stationary Fisher covariance.  For $M\ge1$,
\[
\RetQuad(M)
\ge
\frac12
\Phi^{(1)}_{M-1}(\Sigma_\pi)
=
\frac12
\sum_{i\ge M}
\lambda_i(\Sigma_\pi).
\]
If the retention gap is exactly this quadratic surrogate gap, this recovers
the retention converse.  If the retention gap is an expected KL gap in a
small-radius regime, the same formula is the second-order spectral lower bound
under the regularity assumptions of Theorem~\ref{thm:local-full-kl}.
\end{corollary}

\begin{proof}
Work in Fisher coordinates, with weighted predictive points $\{y_s\}$ and
weights $\pi_s$ summing to one.  Write
\[
\bar y=\sum_s\pi_sy_s,
\qquad
\Sigma_\pi=\sum_s\pi_s(y_s-\bar y)\otimes(y_s-\bar y).
\]
A lumpable quotient with at most $M$ blocks induces a partition into
$K\le M$ nonempty cells $C_1,\dots,C_K$ with weights
$\Pi_k=\sum_{s\in C_k}\pi_s$ and centroids $m_k$.  Define the between-cluster
covariance
\[
B=\sum_{k=1}^K\Pi_k(m_k-\bar y)\otimes(m_k-\bar y).
\]
Since $\sum_{k}\Pi_k(m_k-\bar y)=0$, the centred centroid vectors satisfy one
nontrivial linear relation, so
\[
\rank(B)\le K-1\le M-1 .
\]

For the quadratic distortion with representatives $c_k$, completing the square
inside each cell gives
\[
E
=
\frac12\sum_{k=1}^K\sum_{s\in C_k}\pi_s\norm{y_s-c_k}^2
\ \ge\
\frac12\sum_{k=1}^K\sum_{s\in C_k}\pi_s\norm{y_s-m_k}^2
=
\frac12\tr(\Sigma_\pi-B),
\]
the last step being the ANOVA identity $\Sigma_\pi=W+B$ with $W\succeq0$ the
within-cluster covariance.

Because $\Sigma_\pi-B=W\succeq0$, its trace coincides with its nuclear norm,
\[
\tr(\Sigma_\pi-B)=\norm{\Sigma_\pi-B}_{\Sp{1}},
\]
which is precisely what places this instance at the Schatten exponent $p=1$.
Since $\rank(B)\le M-1$, Theorem~\ref{thm:mirsky} gives
\[
\norm{\Sigma_\pi-B}_{\Sp{1}}
\ \ge\
\min_{\rank P\le M-1}\norm{\Sigma_\pi-P}_{\Sp{1}}
=
\sum_{i>M-1}\lambda_i(\Sigma_\pi)
=
\sum_{i\ge M}\lambda_i(\Sigma_\pi).
\]
Hence $E\ge\frac12\sum_{i\ge M}\lambda_i(\Sigma_\pi)$, and taking the infimum
over feasible quotients gives the stated converse, with effective rank $M-1$
and $\Phi^{(1)}_{M-1}(\Sigma_\pi)=\sum_{i\ge M}\lambda_i(\Sigma_\pi)$.
\end{proof}

\begin{remark}[Why the Exponent Is $1$]
\label{rem:retention-why-p1}
The retention cost is not literally an $\ell^1$ norm of residuals.  In the
quadratic regime it is a mean of squared errors.  The Schatten-$1$ tail appears
because the residual covariance is positive semidefinite, so its trace equals
its nuclear norm.  The exponent $p=1$ is justified by the covariance-residual
structure, and does not follow from the aggregation being a mean.
\end{remark}

\subsection{A Conditional Schatten-$2$ / Frobenius Regime}
\label{subsec:schatten-frobenius}

\begin{corollary}[Conditional Frobenius Converse]
\label{cor:l2}
If a task theory has a Schatten-$2$ response representation with response
operator $A_\delta$, effective rank budget $r_{\mathbb T}(M)$, and modulus
$c_2$, then
\[
\Delta_{\mathbb T}(M)
\ge
\frac1{c_2}
\Phi^{(2)}_{r_{\mathbb T}(M)}(A_\delta)
=
\frac1{c_2}
\left(
\sum_{i>r_{\mathbb T}(M)}
\sigma_i(A_\delta)^2
\right)^{1/2}.
\]
\end{corollary}

\begin{proof}
This is Theorem~\ref{thm:unified} with $p=2$.
\end{proof}

\subsection{Commitment as the Boolean-Rank Vertex}
\label{subsec:schatten-commitment}

Commitment is not a Schatten-norm instance.  It is the Boolean realizability
vertex of the rank-control principle.

\begin{definition}[Deterministic Boolean Hankel Rank]
\label{def:boolean-hankel-rank}
Let $\Lambda\subseteq\Sigma^*$ be a language.  Its Boolean Hankel matrix is
\[
N_\Lambda(u,v)
=
\mathbf 1_{\{uv\in\Lambda\}}.
\]
Define
\[
\rankBdet N_\Lambda
\]
to be the minimum number of states of a deterministic finite automaton
recognizing $\Lambda$.  By the Myhill--Nerode theorem,
\[
\rankBdet N_\Lambda
=
\kappaMN(\Lambda),
\]
the number of Myhill--Nerode equivalence classes.  If $\Lambda$ is not
regular, set $\kappaMN(\Lambda)=\infty$.
\end{definition}

\begin{remark}[Scope of the Boolean Realization Rank]
\label{rem:boolean-rank-not-ordinary}
The quantity $\rankBdet N_\Lambda$ is the \emph{deterministic} Boolean
realization rank, the invariant governing commitment.  It is distinct from the
real or complex rank of $N_\Lambda$, and from Boolean matrix ranks admitting
nondeterministic factorizations, all three of which may differ from it.
\end{remark}

\begin{proposition}[Commitment Exact Threshold]
\label{prop:com-threshold}
Use the $0/1$ commitment loss:
\[
\mathcal E_{\mathrm{com}}(f,g)
=
\begin{cases}
0,& f=g,\\
1,& f\neq g.
\end{cases}
\]
Then
\[
\Com(M)
=
\begin{cases}
0,& M\ge \kappaMN(\Lambda),\\
1,& M<\kappaMN(\Lambda),
\end{cases}
\]
equivalently $\Com(M)=v_0\bigl(\kappaMN(\Lambda)-M\bigr)$ with
\[
v_0(t)=
\begin{cases}
0,& t\le0,\\
1,& t>0.
\end{cases}
\]
Under the extended-real Boolean valuation, in which the nonzero value is
$\infty$ rather than $1$, replace $1$ by $\infty$ throughout; the threshold
$\kappaMN(\Lambda)$ is unchanged.
\end{proposition}

\begin{proof}
By the Myhill--Nerode theorem, a deterministic finite automaton recognizes
$\Lambda$ with $M$ states iff $M\ge\kappaMN(\Lambda)$.  Therefore the
zero-error commitment gap is zero exactly in that case, and positive otherwise.
With $0/1$ loss, the positive value is $1$.
\end{proof}

\subsection{Scope Boundary: No Unconditional Cross-Regime Ordering}
\label{subsec:schatten-nogo}

\begin{theorem}[No Universal Cross-Regime Inequality]
\label{thm:schatten-nogo}
The comparison inequalities of Proposition~\ref{prop:common-operator} require
three things simultaneously: a common response operator, a common effective
rank budget, and comparable domination moduli.  The three regimes supply none
of them.  Their response operators are
\[
H_\nu,
\qquad
\Sigma_\pi,
\qquad
N_\Lambda,
\]
their effective rank budgets are
\[
M,
\qquad
M-1,
\qquad
M,
\]
and their domination moduli $c_p$ are unrelated.

Therefore no unconditional inequality of the form
\[
\Delta_{\mathrm{grd}}(M)\le \Delta_{\mathrm{ret}}(M),
\qquad
\Delta_{\mathrm{ret}}(M)\le \Delta_{\mathrm{com}}(M),
\]
or any fixed monotone relation among the three gaps holds for all tasks.  The
unified Mirsky template unifies the functional \emph{form} of the converses
under shared response representations; it does not unify the operators, the
effective rank budgets, or the numerical gaps.
\end{theorem}

\begin{proof}
This is a direct consequence of the Independence Theorem of
Section~\ref{sec:joint}, which constructs product tasks whose grounding,
retention, and commitment obstructions vary independently.  That construction
is self-contained and does not invoke the present theorem, so the forward
reference is one of presentation order only.
\end{proof}

%======================================================================
\section{The Two-Axis Oracle Inequality}
\label{sec:oracle}
%======================================================================

The offline approximation theory and the online estimation theory are two
distinct axes.  The offline axis measures the deficit of the best budget-$M$
predictor in a regime-specific hypothesis class.  The online axis measures
regret against that budget-$M$ class.  This section composes the two axes,
within a fixed regime, into an adaptive structural-risk principle.

The bridge is type-correct: the online hypothesis class matches the offline
approximation class.  Retention uses stochastic finite-state predictors;
grounding spectral relaxations use linear finite-rank predictors; commitment
uses deterministic Mealy machines.  Spectral tails enter online rates only
through explicit spectral-rate assumptions.

%----------------------------------------------------------------------
\subsection{Type-Correct Axes on One Clock}
\label{subsec:oracle-setup}
%----------------------------------------------------------------------

Fix a regime $\mathsf r$.  Let
\[
\Hcal_1^{\mathsf r}
\subseteq
\Hcal_2^{\mathsf r}
\subseteq
\cdots
\]
be the nested budget classes in that regime, and let
\[
\Hcal_\infty^{\mathsf r}
=
\bigcup_{M\ge1}\Hcal_M^{\mathsf r},
\]
or its appropriate closure when stochastic or linear predictors are involved.

The regime-specific typing is as follows:
\begin{itemize}[leftmargin=1.8em]
\item \textbf{Commitment.}
$\Hcal_M^{\mathrm{com}}$ is the class of deterministic Mealy machines with at
most $M$ states.  The loss is the $0/1$ mismatch loss.

\item \textbf{Retention.}
$\Hcal_M^{\mathrm{ret}}$ is a class of stochastic finite-state predictors,
for example stationary controlled causal machines whose admissible quotients
are lumpable quotients of the stationary support $\Splus$ with at most $M$
blocks.  The loss is log-loss or a bounded proper surrogate.

\item \textbf{Grounding spectral relaxation.}
$\Hcal_M^{\mathrm{grd,lin}}$ is the class of linear finite-rank realizations
of rank at most $M$, matching the linear finite-rank grounding converse.  The
exact symbolic deterministic Mealy grounding problem is a different object and
is not substituted here.
\end{itemize}

Let
\[
L_T(h)=\sum_{t=1}^T \ell_t(h)
\]
denote cumulative loss.  The estimation-axis envelopes invoked below are of the
standard concentration type \cite{boucheron2013}.

\begin{definition}[Finite-State Benchmark]
\label{def:finite-state-benchmark}
For regime $\mathsf r$, define
\[
L_T^{\star,\mathsf r}
=
\inf_{h\in \Hcal_\infty^{\mathsf r}}
\mathbb E\,L_T(h).
\]
This is the best loss achievable by the unrestricted finite-state model class
in that regime.  It equals the Bayes-optimal true-process loss only if the
Bayes predictor lies in the closure of $\Hcal_\infty^{\mathsf r}$.
\end{definition}

\begin{definition}[State Rate]
\label{def:state-rate}
The \emph{state rate} of a finite-state factor with at most $M$ classes is
\[
R(M)=\log M ,
\]
the description length of a state index, and the associated
state-rate distortion function of a task theory $\mathbb T$ is
$D_{\mathbb T}(R)=\Delta_{\mathbb T}(\lfloor e^{R}\rfloor)$.

Throughout, curves are parameterized by the budget $M$ rather than by $R$,
which is a change of variable and affects no theorem.

The object $D_{\mathbb T}$ is a lossy-compression trade-off in the informal
sense: a source $\delta_{\mathbb T}$, a description budget $R=\log M$, and a
fidelity criterion $\mathcal E_{\mathbb T}$.  Three properties separate it
from a Shannon rate--distortion function, and the third is the substantive one.
The quantity $\log M$ is a one-shot description length for the memory, not a
per-symbol transmission rate.  The distortion axis is a regime-specific cost,
not a single-letter expected distortion.  And $D_{\mathbb T}$ need not be
convex, so the operational apparatus that convexity underwrites --- in
particular the time-sharing argument that makes the Shannon curve the lower
convex envelope of achievable pairs --- is unavailable here
(Proposition~\ref{prop:rd-nonconvex}).  The trade-off is therefore genuine but
is not an instance of Shannon's theory, and the results below are stated for
$\Delta_{\mathbb T}(M)$ directly.
\end{definition}

\begin{proposition}[The Finite-State Curve Need Not Be Convex]
\label{prop:rd-nonconvex}
There is a retention instance, with every partition lumpable, for which
$R\mapsto D_{\mathbb T}(R)$ is not convex on $\{\log M:1\le M\le|\Splus|\}$.
\end{proposition}

\begin{proof}
Take $|\Splus|=5$, $\mathcal O=\{1,2,3\}$, the reset transition
$\tau(s_i,\ell)=s_\ell$ of Theorem~\ref{thm:retention-reset-np}, so that every
partition is lumpable, stationary weights proportional to
$(0.0344,0.3506,0.1906,0.2176,0.2068)$, and predictive laws the row
normalizations of
\[
(0.4805,0.4113,0.1082),\ (0.2746,0.4018,0.3236),\ (0.2960,0.5426,0.1614),
\]
\[
(0.2334,0.5019,0.2648),\ (0.6498,0.0548,0.2954).
\]
Minimizing $\RetKL$ over all partitions with at most $M$ blocks gives
\[
D(0)=0.0948616,\quad
D(\log2)=0.0148089,\quad
D(\log3)=0.0049099,
\]
\[
D(\log4)=0.0021747,\quad
D(\log5)=0 ,
\]
whose successive chord slopes are
$-0.115492$, $-0.024414$, $-0.009508$, $-0.009746$.  Convexity requires the
slopes to be nondecreasing; the last decreases, so $D$ is not convex.  The
computation is exact rational-plus-logarithm arithmetic, carried out at $60$
significant digits, and the final slope decrement is $-2.38\times10^{-4}$,
far above any rounding effect.
\end{proof}

\begin{remark}[Why Convexity Fails]
\label{rem:rd-nonconvex-mechanism}
The Shannon curve is convex because rates may be time-shared: given codes at
rates $R_1$ and $R_2$ one randomizes between them and attains every point of
the chord.  The finite-state budget admits no such operation.  The constraint
$|\mathcal K|\le M$ is a hard cardinality bound on a \emph{deterministic}
lumpable quotient, and a randomization between quotients of sizes $M_1$ and
$M_2$ is not itself a quotient of any intermediate size.  Convexity therefore
has no mechanism to hold, and by
Proposition~\ref{prop:rd-nonconvex} it indeed fails; over randomly generated
retention instances with $|\Splus|\le5$ it fails in roughly one case in eight.

The failure is not an artefact of measuring the budget logarithmically.
Convexity is a property of the pair (rate axis, distortion axis), so one might
hope that some other parameterization of the budget restores it.  It does not.
Take $|\Splus|=4$, $\mathcal O=\{1,2,3\}$, stationary weights proportional to
$(17,18,22,21)$, and predictive laws the row normalizations of
\[
(20,25,2),\quad(2,34,30),\quad(37,1,27),\quad(20,9,1).
\]
The optimal costs are
$D=0.2887482,\ 0.1601480,\ 0.0145216,\ 0$ at $M=1,2,3,4$.  Against the raw
budget $R=M$ the chord slopes are $-0.128600$, $-0.145626$, $-0.014522$, and
against $R=\log M$ they are $-0.185531$, $-0.359159$, $-0.050478$; both
sequences decrease at the first step, so the curve is non-convex under either
parameterization.  Since the obstruction is the absence of a convex
combination of feasible quotients, and that is a property of the feasible set
rather than of the axis, no rescaling of the rate can supply a convexity
guarantee.
\end{remark}

\begin{remark}[Necessity of the Feasibility Restriction]
\label{rem:unifilar-feasibility}
The restriction to feasible triples in Definition~\ref{def:unifilar-lumpable}
is not dispensable.  If
$(s,x,y)$ has emission probability zero the update $\tau(s,x,y)$ is
unconstrained by the process, so requiring the intertwining there would reject
quotients that are lumpable in every realizable sense.  Over randomly generated unifilar machines with $|\mathcal S|\le4$,
$|\mathcal I|\le2$ and $|\mathcal O|\le3$ together with random partitions,
approximately fifteen per cent of the quotients satisfying
Definition~\ref{def:unifilar-lumpable} fail the corresponding condition stated
without the feasibility restriction, so the two conditions are genuinely
distinct.
\end{remark}

\begin{proposition}[Right Congruences on Joint Histories]
\label{prop:unifilar-lumpability}
Let a unifilar machine have finite synchronization depth $L$ in joint
histories, so that $S_t$ is a function $\sigma$ of the last $L$ pairs
$(X_u,Y_u)$, and let $F$ be the pruned history system of feasible joint
histories of length at least $L$
(Definition~\ref{def:history-system}).  Then:
\begin{enumerate}[label=(\roman*)]
\item if $\phi$ is unifilar-lumpable with block-uniform support, the relation
$u\sim_\phi v$ defined by $\phi(\sigma(u))=\phi(\sigma(v))$ is a
support-relative right congruence on $F$ in the sense of
Definition~\ref{def:support-right-cong}.  Without a support hypothesis it
still satisfies clause~(ii) on commonly feasible events, but clause~(i) may
fail, and then the quotient carries no well-defined feasible-event set
(Remark~\ref{rem:support-clause-i-needed});
\item conversely, if $Z$-predictive equivalence $\sim_Z$ on $F$ is itself a
support-relative right congruence, then every finite-index support-relative
right congruence $\sim$ on $F$ that is \emph{coarser} than $\sim_Z$ induces a
well-defined unifilar-lumpable quotient of $\Splus$.  The hypothesis on
$\sim_Z$ holds whenever $Z$ is the full controlled future.
\end{enumerate}
\end{proposition}

\begin{proof}
For~(i), let $u\sim_\phi v$, and write $k$ for the common value
$\phi(\sigma(u))=\phi(\sigma(v))$.  Block-uniform support gives
$\supp P_{\sigma(u)}^x=\supp P_{\sigma(v)}^x$ for every input $x$, so an event
$(x,y)$ is feasible from $u$ exactly when it is feasible from $v$; this is
clause~(i) of Definition~\ref{def:support-right-cong}.  For such an event,
unifilar lumpability applied to the pair $\sigma(u),\sigma(v)$, which have
equal $\phi$-image, gives
$\phi(\tau(\sigma(u),x,y))=\tau_{\mathcal K}(k,x,y)=\phi(\tau(\sigma(v),x,y))$,
and by unifilarity $\sigma(u\cdot(x,y))=\tau(\sigma(u),x,y)$; hence
$u\cdot(x,y)\sim_\phi v\cdot(x,y)$, which is clause~(ii).  Iterating gives
closure under feasible words.  Without block-uniform support the displayed
computation still applies to events feasible from both histories, which is the
weaker assertion recorded in the statement.

For~(ii), the only issue is that $\sim$ descend to states, that is, that
$\sigma(u)=\sigma(v)$ imply $u\sim v$; otherwise the induced map on
$\mathcal S$ is not well defined.  This is where the hypothesis enters.  If
$\sigma(u)=\sigma(v)$ then the two histories reach the same state and
therefore carry the same $Z$-predictive law, so they are $Z$-predictive
equivalent; since $\sim$ is coarser than $Z$-predictive equivalence,
$u\sim v$.  Define $\bar\phi(\sigma(u))$ to be the $\sim$-class of $u$, which
is now unambiguous, and let $\mathcal K$ be the set of classes, finite by
hypothesis.  Right-congruence of $\sim$ under a feasible event $(x,y)$ is
exactly the intertwining
$\bar\phi(\tau(s,x,y))=\tau_{\mathcal K}(\bar\phi(s),x,y)$ required by
Definition~\ref{def:unifilar-lumpable}.

The calculation is the one in Proposition~\ref{prop:lumpability} with events
$(x,y)$ in place of letters, and support-relative feasibility handled as in
the feasible-input clause of that proof.

It remains to check the final assertion, that $\sim_Z$ is a support-relative
right congruence when $Z$ is the full controlled future.  Clause~(ii) is the
usual argument: conditioning an equality of conditional future laws on a
commonly realized next event preserves the equality.  Clause~(i) is not an
extra hypothesis but a consequence.  If $u\sim_Zv$ then the two conditional
laws of the controlled future agree; taking the one-step marginal under the
input program that begins with $x$ shows $P_{\sigma(u)}^x=P_{\sigma(v)}^x$, so
the two emission laws have the same support and $u,v$ have the same feasible
events.  Thus equality of full-future laws already forces equality of feasible
sets, and no separate common-feasibility assumption is needed.
\end{proof}

\begin{remark}[Which Hypothesis Does the Work in (ii)]
\label{rem:unifilar-converse-hypothesis}
Coarseness, not finite index, is what makes the converse true.  Dropping it
breaks descent: take a single-state machine with two feasible outputs, so
every joint history reaches the same state, and let $u\sim v$ when $|u|$ and
$|v|$ have the same parity.  That relation is a right congruence of index two,
yet $\sigma(u)=\sigma(v)$ for all $u,v$ while $u\not\sim v$ whenever the
lengths differ in parity, so no induced map on $\mathcal S$ exists.  Here
$Z$-predictive equivalence is the all-relation and the parity congruence is
strictly finer, so the hypothesis of Proposition~\ref{prop:unifilar-lumpability}(ii)
correctly excludes it.
\end{remark}

\begin{proposition}[Input-Driven Specialization]
\label{prop:input-driven-specialization}
Let the machine be input-driven, $\tau(s,x,y)=\tau_0(s,x)$, and
input-synchronized in the sense of
Definition~\ref{def:synchronized-realization}.
\begin{enumerate}[label=(\roman*)]
\item Lumpability in the sense of Definition~\ref{def:lumpable-quotient}
implies unifilar lumpability.  The converse holds if $\phi$ has connected
support, in particular if it has block-uniform support, and in particular if
every $P_s^x$ has full support; it fails in general
(Remark~\ref{rem:unifilar-support-not-automatic}).
\item The relation $u\sim_\phi v\iff\phi(\sigma(u))=\phi(\sigma(v))$ on
feasible joint histories depends only on the input projections of $u$ and $v$,
and its projection to input histories is exactly the right congruence of
Proposition~\ref{prop:lumpability}.
\end{enumerate}
\end{proposition}

\begin{proof}
\emph{(i)} Given a quotient map $\tau_{\mathcal K}^0$ witnessing
Definition~\ref{def:lumpable-quotient}, set
$\tau_{\mathcal K}(k,x,y)=\tau_{\mathcal K}^0(k,x)$; the left-hand side of the
unifilar condition is independent of $y$, so the identity holds at every
feasible triple.

For the converse, fix a block $C$ and an input $x$, and suppose $s,s'\in C$
share a feasible output $y\in\supp P_s^x\cap\supp P_{s'}^x$.  Unifilar
lumpability at the two triples $(s,x,y)$ and $(s',x,y)$ gives
\[
\phi(\tau_0(s,x))
=
\tau_{\mathcal K}(k,x,y)
=
\phi(\tau_0(s',x)),
\]
so $\phi\circ\tau_0(\cdot,x)$ is constant along every edge of the
support-overlap graph on $C$.  If that graph is connected, the function is
constant on $C$, which is Definition~\ref{def:lumpable-quotient}.
Block-uniform support makes the graph complete, hence connected, and full
supports make it complete for the same reason.  Connectedness cannot be
dropped: in the witness of
Remark~\ref{rem:unifilar-support-not-automatic} the graph on the block
$\{s_0,s_2\}$ has no edges, and the implication fails.

\emph{(ii)} By Definition~\ref{def:synchronized-realization},
$\sigma(u)=\sigma_L(\operatorname{suffix}_L(u_{\mathcal I}))$ depends only on
the input projection $u_{\mathcal I}$, hence so does $\phi(\sigma(u))$.  For an
input extension $x$,
\[
\phi(\sigma(u_{\mathcal I}x))
=
\phi(\tau_0(\sigma(u_{\mathcal I}),x))
=
\tau_{\mathcal K}^0\bigl(\phi(\sigma(u_{\mathcal I})),x\bigr),
\]
which is the right-congruence computation of
Proposition~\ref{prop:lumpability}.  Passing to joint extensions appends
output letters that change neither the state nor the input projection.
\end{proof}

\begin{remark}[Relation to $\epsilon$-Machines]
\label{rem:epsilon-machine-relation}
An $\epsilon$-machine in the sense of computational mechanics
\cite{shalizi2001} is a unifilar predictive representation whose states are
the causal states: classes of pasts carrying identical conditional laws of the
future, minimal among such representations.  A unifilar controlled machine
(Definition~\ref{def:unifilar-machine}) in which $Z$ is the full controlled
future, and whose states are the $\sim_Z$-classes of feasible joint histories,
is a controlled $\epsilon$-machine in that sense.  The input-driven machines of
Definition~\ref{def:controlled-markov} form the subclass with
output-independent transitions; in the synchronized case their causal states
factor through input histories
(Proposition~\ref{prop:input-driven-specialization}(ii)).  References
elsewhere in this article to lumped $\epsilon$-machines are to be read in this
sense.
\end{remark}

\begin{definition}[Finite-Horizon Approximation Deficit]
\label{def:approx-deficit}
The budget-$M$ approximation deficit is
\[
\Aapp^{\mathsf r}(M)
=
\inf_{h\in\Hcal_M^{\mathsf r}}
\mathbb E\,L_T(h)
-
L_T^{\star,\mathsf r}.
\]
The per-symbol deficit is
\[
\Delta_T^{\mathsf r}(M)
=
\frac1T\Aapp^{\mathsf r}(M).
\]
\emph{Quantification over processes.}  As written, $\Aapp^{\mathsf r}(M)$
depends on the process $P$ generating the expectation, and the minimax
statements below require an envelope that does not.  Two readings make this
precise, and the choice must be fixed once and used consistently.  Either fix a
class $\mathcal P$ of processes and take the worst case,
\[
a_M(\mathcal P)
=
\sup_{P\in\mathcal P}
\Bigl[
\inf_{h\in\Hcal_M^{\mathsf r}}\mathbb E_P L_T(h)-L_T^{\star,\mathsf r}(P)
\Bigr],
\]
or fix for each $M$ an explicit hard subclass $\mathcal P_M$ on which the
deficit is bounded below.  Throughout the oracle section the first reading is
in force, with $\mathcal P$ the process class of the regime $\mathsf r$;
displays such as
$\sup_P\Reg_T\ge\sup_M\min\{\Aapp^{\mathsf r}(M),\Est_M^{\mathsf r}(T)\}$ are
to be read with $\Aapp^{\mathsf r}(M)=a_M(\mathcal P)$, so that both sides
quantify over the same class and the inequality is between deterministic
envelopes.  Where a single fixed $P$ is intended, as in the fixed-budget
decomposition, the process is named explicitly.

This definition is finite-horizon and does not presuppose stationarity.  Under
stationary ergodic conditions one has
\[
\Aapp^{\mathsf r}(M)
=
T\Delta^{\mathsf r}(M)+o(T),
\]
where $\Delta^{\mathsf r}(M)$ is the stationary per-symbol gap of the
corresponding regime; the identification of $\Delta_T^{\mathsf r}(M)$ with a
stationary gap requires that hypothesis and is not automatic.
\end{definition}

\begin{definition}[Online Estimation Rate]
\label{def:estimation-rate}
An estimation rate for budget $M$ is any function $\Est_M^{\mathsf r}(T)$ for
which there exists a base learner $\mathcal A_M$ satisfying
\[
\mathbb E\,L_T(\mathcal A_M)
-
\inf_{h\in\Hcal_M^{\mathsf r}}
\mathbb E\,L_T(h)
\le
\Est_M^{\mathsf r}(T).
\]
\end{definition}

%----------------------------------------------------------------------
\subsection{Exact Fixed-Budget Decomposition}
\label{subsec:oracle-fixed}
%----------------------------------------------------------------------

\begin{proposition}[Fixed-Budget Bias--Variance Decomposition]
\label{prop:fixed}
For any learner $\mathcal A_M$ whose predictions lie in or are generated from
$\Hcal_M^{\mathsf r}$, if
\[
\mathbb E\,L_T(\mathcal A_M)
-
\inf_{h\in\Hcal_M^{\mathsf r}}
\mathbb E\,L_T(h)
\le
\Est_M^{\mathsf r}(T),
\]
then its total excess loss against the unrestricted finite-state benchmark
satisfies
\[
\Reg_T(\mathcal A_M)
\eqdef
\mathbb E\,L_T(\mathcal A_M)-L_T^{\star,\mathsf r}
\le
\Aapp^{\mathsf r}(M)+\Est_M^{\mathsf r}(T).
\]
\end{proposition}

\begin{proof}
The in-class infimum need not be attained, so fix $\varepsilon>0$ and choose an
$\varepsilon$-approximate in-class optimum
$h_M^\varepsilon\in\Hcal_M^{\mathsf r}$ with
\[
\mathbb E\,L_T(h_M^\varepsilon)
\le
\inf_{h\in\Hcal_M^{\mathsf r}}\mathbb E\,L_T(h)+\varepsilon .
\]
Inserting it and splitting,
\[
\Reg_T(\mathcal A_M)
=
\bigl(
\mathbb E\,L_T(\mathcal A_M)-\mathbb E\,L_T(h_M^\varepsilon)
\bigr)
+
\bigl(
\mathbb E\,L_T(h_M^\varepsilon)-L_T^{\star,\mathsf r}
\bigr).
\]
The first bracket is at most $\Est_M^{\mathsf r}(T)+\varepsilon$ by the
estimation hypothesis, and the second is at most
$\Aapp^{\mathsf r}(M)+\varepsilon$ by the choice of $h_M^\varepsilon$.  Hence
\[
\Reg_T(\mathcal A_M)
\le
\Aapp^{\mathsf r}(M)+\Est_M^{\mathsf r}(T)+2\varepsilon ,
\]
and letting $\varepsilon\downarrow0$ gives the claim.  When the infimum is
attained one may take $\varepsilon=0$ directly.
\end{proof}

%----------------------------------------------------------------------
\subsection{Adaptive Oracle Inequalities}
\label{subsec:oracle-adaptive}
%----------------------------------------------------------------------

\begin{theorem}[Adaptive Agnostic Oracle Inequality]
\label{thm:oracle-agnostic}
Assume losses are bounded in $[0,1]$.  For each $M\ge1$, let
$\mathcal A_M$ be a base learner satisfying
\[
\mathbb E\,L_T(\mathcal A_M)
-
\inf_{h\in\Hcal_M^{\mathsf r}}
\mathbb E\,L_T(h)
\le
\Est_M^{\mathsf r}(T).
\]
There exists an aggregated learner $\mathcal A^{\mathrm{agg}}$ such that for
every $M\ge1$,
\[
\Reg_T(\mathcal A^{\mathrm{agg}})
\le
\Aapp^{\mathsf r}(M)
+
\Est_M^{\mathsf r}(T)
+
\PsiAgg_M(T),
\]
where
\[
\PsiAgg_M(T)
\le
C\left(
\sqrt{T\ln(eM)}
+
\ln(eM)
\right)
\]
for a universal constant $C$, uniformly in $T$.  This is sharper than the
$\ln(eM)$ form obtainable from a learning-rate grid: the
second-order term is $\ln(eM)$, with no dependence on $T$.  Consequently, if
\[
M_T^\star
\in
\argmin_{M\ge1}
\left[
\Aapp^{\mathsf r}(M)+\Est_M^{\mathsf r}(T)
\right],
\]
then
\[
\Reg_T(\mathcal A^{\mathrm{agg}})
\le
\min_{M\ge1}
\left[
\Aapp^{\mathsf r}(M)+\Est_M^{\mathsf r}(T)
\right]
+
\PsiAgg_{M_T^\star}(T).
\]
The learner does not know $M_T^\star$, $\Aapp^{\mathsf r}(\cdot)$, or the
spectral decay.
\end{theorem}

\begin{proof}
Treat the base learners $\{\mathcal A_M\}_{M\ge1}$ as a countable expert set
with prior
\[
\pi_M=\frac{6}{\pi^2M^2},
\qquad
\ln(1/\pi_M)=2\ln M+O(1).
\]
Aggregate them with a parameter-free prior-weighted expert algorithm enjoying
the standard second-order guarantee
\[
\mathbb E\,L_T(\mathcal A^{\mathrm{agg}})
-
\mathbb E\,L_T(\mathcal A_M)
\le
C\sqrt{T\ln(1/\pi_M)}
+
C\ln(1/\pi_M)
\]
simultaneously for every $M$, with no tuning to $M$ or $T$
\cite{cesabianchi2006,derooij2014}.  Substituting
$\ln(1/\pi_M)=2\ln M+O(1)$ gives an aggregation penalty
\[
\PsiAgg_M(T)
\le
C\left(
\sqrt{T\ln(eM)}
+
\ln(eM)
\right).
\]
Adding the fixed-budget decomposition of Proposition~\ref{prop:fixed} and
taking the infimum over $M$ gives the result.

A single exponential-weights instance carries one learning rate, and a grid of
rates $\eta_j$ does not by itself yield a per-expert tuned bound; a
parameter-free or adaptive-rate aggregator is therefore required.
\end{proof}

\begin{theorem}[Realizable $0/1$ Model Selection]
\label{thm:oracle-realizable}
Assume $0/1$ loss.  Suppose each base learner $\mathcal A_M$ makes expected
mistakes at most
\[
\Est_M^{\mathrm{real}}(T)
\]
more than the best predictor in $\Hcal_M^{\mathsf r}$.  There exists a
weighted-majority aggregation over the base learners, with prior
$w_M=6/(\pi^2M^2)$ on the budget index, such that for every $M$,
\[
\Reg_T(\mathcal A^{\mathrm{wm}})
\le
C\,\Aapp^{\mathsf r}(M)
+
C\,\Est_M^{\mathrm{real}}(T)
+
C\ln(eM),
\]
for a universal constant $C$.
\end{theorem}

\begin{proof}
The weighted-majority mistake bound is standard
\cite{littlestone1988,cesabianchi2006} and is not reproved here.  If expert $M$ makes $m_M$
mistakes and has prior weight $w_M$, weighted majority with a constant
multiplicative update makes at most
\[
C_1 m_M+C_2\ln(1/w_M)
\]
mistakes in expectation, up to finite-alphabet constants.  Since
\[
m_M
\le
\inf_{h\in\Hcal_M}\mathbb E L_T(h)
+
\Est_M^{\mathrm{real}}(T),
\]
subtracting $L_T^{\star,\mathsf r}$ gives the claim.
\end{proof}

\begin{remark}[Realizable and Agnostic Laws Are Separate]
\label{rem:realizable-not-agnostic}
The realizable theorem does not use the bounded-loss $\sqrt T$ aggregation
overhead.  Realizable and agnostic budget laws are therefore stated
separately.
\end{remark}

\begin{theorem}[Mixable Log-Loss Oracle Inequality]
\label{thm:oracle-logloss}
Let $\Hcal_M^{\mathsf r}$ consist of stochastic predictors emitting
probability distributions over a finite output alphabet.  Assume either
positivity or clipped log-loss so that the loss is mixable.  If each base
learner $\mathcal A_M$ satisfies
\[
\mathbb E\,L_T(\mathcal A_M)
-
\inf_{h\in\Hcal_M^{\mathsf r}}
\mathbb E\,L_T(h)
\le
\Est_M^{\mathsf r}(T),
\]
then there exists an aggregating learner $\mathcal A^{\mathrm{mix}}$ such that
for every $M$,
\[
\Reg_T(\mathcal A^{\mathrm{mix}})
\le
\Aapp^{\mathsf r}(M)
+
\Est_M^{\mathsf r}(T)
+
C\ln(eM).
\]
\end{theorem}

\begin{proof}
For mixable losses, Vovk's aggregating algorithm \cite{vovk1990} with prior weights $w_M$ gives
\[
L_T(\mathcal A^{\mathrm{mix}})
\le
\min_M
\left[
L_T(\mathcal A_M)+\ln(1/w_M)
\right]
\]
up to mixability constants.  Taking expectations and applying
Proposition~\ref{prop:fixed} yields the result.
\end{proof}

%----------------------------------------------------------------------
\subsection{Spectral Input: Converses versus Achievability}
\label{subsec:oracle-spectral}
%----------------------------------------------------------------------

The offline schema supplies spectral converses.  These are lower bounds.  They
do not by themselves give online regret upper rates.

\begin{definition}[Spectral-Rate Assumption]
\label{def:spectral-rate}
We say that regime $\mathsf r$ satisfies a spectral-rate assumption with tail
$S^{\mathsf r}(M)$ if there exist constants $c,C>0$ such that
\[
c\,S^{\mathsf r}(M)
\le
\Delta_T^{\mathsf r}(M)
\le
C\,S^{\mathsf r}(M)+o(1)
\]
uniformly in the relevant asymptotic range.  Equivalently,
\[
\Aapp^{\mathsf r}(M)
\le
T\,C\,S^{\mathsf r}(M)+o(T).
\]
The lower inequality is the converse.  The upper inequality is an
achievability bound.
\end{definition}

For grounding, the unrestricted relaxation satisfies
\[
S^{\mathrm{grd}}(M)
=
\sigma_{M+1}(H_\nu),
\qquad
\Dunres(M)
=
\sigma_{M+1}(H_\nu),
\]
while the Hankel-structured gap satisfies
\[
\DHankstr(M)
\ge
\sigma_{M+1}(H_\nu).
\]
For quadratic retention,
\[
S^{\mathrm{ret}}(M)
=
\sum_{i\ge M}\lambda_i(\Sigma_\pi)
=
\Phi^{(1)}_{M-1}(\Sigma_\pi).
\]

%----------------------------------------------------------------------
\subsection{Budget Laws under Spectral-Rate Assumptions}
\label{subsec:oracle-laws}
%----------------------------------------------------------------------

Define the oracle objective
\[
\Fbudget_T^{\mathsf r}(M)
=
\Aapp^{\mathsf r}(M)+\Est_M^{\mathsf r}(T).
\]

For the remainder of this subsection write
\[
a^{\mathsf r}(M)
=
\frac1T\Aapp^{\mathsf r}(M)
\]
for the per-round approximation deficit.  Every displayed upper rate below
assumes a corresponding \emph{achievability} bound on $a^{\mathsf r}(M)$, not
merely a spectral converse.

\begin{corollary}[Agnostic Geometric Decay]
\label{cor:geo-agnostic}
Suppose the agnostic estimation rate is
\[
\Est_M^{\mathsf r}(T)
=
O\!\left(\sqrt{T\,M\log(eM)}\right),
\]
and suppose
\[
a^{\mathsf r}(M)\le C\rho^M,
\qquad
0<\rho<1.
\]
Then
\[
M_T^\star
=
\Theta(\log T),
\]
and the adaptive aggregated learner satisfies
\[
\Reg_T
=
\widetilde O(\sqrt T).
\]
If, in addition, a minimax lower bound of the same order is available for the
specified protocol and process class, then
\[
\Reg_T
=
\widetilde\Theta(\sqrt T).
\]
\end{corollary}

\begin{proof}
Balance
\[
T\rho^M
\asymp
\sqrt{T\,M\log M}.
\]
Squaring gives
\[
T^2\rho^{2M}
\asymp
T\,M\log M,
\]
so
\[
\rho^{2M}
\asymp
\frac{M\log M}{T}.
\]
Thus
\[
M
=
\frac{1}{2\ln(1/\rho)}
\ln T
+
O(\ln\ln T).
\]
Substitution gives both terms of order $\sqrt T$ times polylogarithmic factors.
The adaptive overhead from Theorem~\ref{thm:oracle-agnostic} is absorbed into
the $\widetilde O$ notation.
\end{proof}

\begin{corollary}[Agnostic Polynomial Decay]
\label{cor:poly-agnostic}
Suppose
\[
a^{\mathsf r}(M)\le C M^{-\alpha},
\qquad
\alpha>0.
\]
Then
\[
M_T^\star
=
\widetilde\Theta\!\left(T^{1/(2\alpha+1)}\right),
\]
and
\[
\Reg_T
=
\widetilde O\!\left(T^{(\alpha+1)/(2\alpha+1)}\right).
\]
If, in addition, a minimax lower bound of the same order is available for the
specified protocol and process class, then
\[
\Reg_T
=
\widetilde\Theta\!\left(T^{(\alpha+1)/(2\alpha+1)}\right).
\]
\end{corollary}

\begin{proof}
Balance
\[
T M^{-\alpha}
\asymp
\sqrt{T\,M\log M}.
\]
Squaring gives
\[
T^2M^{-2\alpha}
\asymp
T\,M\log M,
\]
so
\[
M^{2\alpha+1}
\asymp
\frac{T}{\log M}.
\]
Hence
\[
M_T^\star
=
\widetilde\Theta\!\left(T^{1/(2\alpha+1)}\right).
\]
The balanced value is
\[
T M^{-\alpha}
=
\widetilde\Theta\!\left(
T^{1-\alpha/(2\alpha+1)}
\right)
=
\widetilde\Theta\!\left(
T^{(\alpha+1)/(2\alpha+1)}
\right).
\]
\end{proof}

\begin{corollary}[Realizable Geometric Decay]
\label{cor:geo-realizable}
If the realizable or mixable estimation rate is
\[
\Est_M^{\mathsf r}(T)
=
O(M\log(eM)),
\]
and
\[
a^{\mathsf r}(M)\le C\rho^M,
\qquad
0<\rho<1,
\]
then
\[
M_T^\star
=
\Theta(\log T),
\]
and
\[
\Reg_T
=
\widetilde O(\log T).
\]
If, in addition, a minimax lower bound of the same order is available for the
specified protocol and process class, then
\[
\Reg_T
=
\widetilde\Theta(\log T).
\]
\end{corollary}

\begin{proof}
Balance
\[
T\rho^M
\asymp
M\log M.
\]
Thus
\[
M
=
\frac{1}{\ln(1/\rho)}
\ln T
+
O(\ln\ln T),
\]
and the common objective value is
\[
M\log M
=
\Theta(\log T\log\log T).
\]
The logarithmic model-selection overhead is lower order.
\end{proof}

\begin{corollary}[Realizable Polynomial Decay]
\label{cor:poly-realizable}
If
\[
a^{\mathsf r}(M)\le C M^{-\alpha},
\qquad
\alpha>0,
\]
and the realizable or mixable estimation rate is
\[
\Est_M^{\mathsf r}(T)
=
O(M\log(eM)),
\]
then
\[
M_T^\star
=
\widetilde\Theta\!\left(T^{1/(\alpha+1)}\right),
\]
and
\[
\Reg_T
=
\widetilde O\!\left(T^{1/(\alpha+1)}\right).
\]
If, in addition, a minimax lower bound of the same order is available for the
specified protocol and process class, then
\[
\Reg_T
=
\widetilde\Theta\!\left(T^{1/(\alpha+1)}\right).
\]
\end{corollary}

\begin{proof}
Balance
\[
T M^{-\alpha}
\asymp
M\log M.
\]
Thus
\[
M^{\alpha+1}
\asymp
\frac{T}{\log M},
\]
so
\[
M_T^\star
=
\widetilde\Theta\!\left(T^{1/(\alpha+1)}\right).
\]
The balanced objective is
\[
M\log M
=
\widetilde\Theta\!\left(T^{1/(\alpha+1)}\right).
\]
\end{proof}

\begin{remark}[Different Exponents Are Structural]
\label{rem:exponents-structural}
The agnostic exponent
\[
\frac{\alpha+1}{2\alpha+1}
\]
and the realizable exponent
\[
\frac{1}{\alpha+1}
\]
are not interchangeable.  They reflect the different costs of estimation:
\[
\sqrt{T M\log M}
\quad\text{versus}\quad
M\log M.
\]
\end{remark}

%----------------------------------------------------------------------
\subsection{Continuous Bias--Variance Sandwich and Lower-Bound Status}
\label{subsec:oracle-sandwich}
%----------------------------------------------------------------------

The oracle inequality proved above is an upper bound.  A separate question is
whether the resulting bias--variance envelope is minimax-optimal.  This
subsection establishes an envelope comparison valid under regularity
hypotheses, and distinguishes it from the stronger minimax statement of
Theorem~\ref{thm:oracle-minimax-lower}.

\begin{proposition}[Necessity of the Aggregation Penalty]
\label{prop:aggregation-necessary}
Consider online aggregation over $M$ base predictors with losses in $[0,1]$,
in which the aggregator observes each base predictor's loss each round and
competes with the best of them.  There is a universal $c>0$ such that for all
$M\ge2$ and all $T\ge\log_2M$, every aggregation rule, deterministic or
randomized, suffers
\[
\sup\ \Bigl[\,\mathbb E\,L_T(\text{aggregator})-\min_{m\le M}L_T(m)\Bigr]
\ \ge\
c\sqrt{T\ln M}
\]
on some loss sequence.  Consequently the $\sqrt{T\ln(eM)}$ term of
$\PsiAgg_M(T)$ cannot be removed at the level of generality at which the
oracle inequality is stated.
\end{proposition}

\begin{proof}
This is the classical lower bound for prediction with expert advice; we
indicate the standard argument and cite it rather than reproduce it in full.
Draw each base predictor's loss on each round as an independent fair coin in
$\{0,1\}$.  Every aggregation rule then has expected cumulative loss exactly
$T/2$, while the expected minimum over $M$ independent symmetric random walks
of length $T$ is $T/2-\Theta(\sqrt{T\ln M})$ for $M\le2^{T}$, by the standard
maximal-inequality estimate for the maximum of $M$ sub-Gaussian variables.
The gap is therefore $\Omega(\sqrt{T\ln M})$ in expectation over the random
sequence, hence for some fixed sequence.  See
\cite[Thm.~3.7]{cesabianchi2006} and \cite[Ch.~2]{boucheron2013} for the
precise constants.
\end{proof}

\begin{remark}[What the Lower Bound Does and Does Not Settle]
\label{rem:aggregation-scope}
Proposition~\ref{prop:aggregation-necessary} closes the question of necessity
at the generality of \emph{arbitrary} base predictors treated as experts: no
adaptive rule competing with $M$ unrelated predictors avoids
$\Omega(\sqrt{T\ln M})$.  It does not settle whether the penalty is necessary
for the \emph{nested} families arising here, where the base predictors are the
budget-$M'$ learners for $M'\le M$ and are therefore highly dependent: the best
predictor in a nested family is monotone in the budget, and an aggregator might
exploit that structure.  Whether $\PsiAgg_M(T)$ can be improved for nested
finite-state classes remains open.
\end{remark}

\begin{assumption}[Regular Bias--Variance Envelope]
\label{ass:regular-bv-envelope}
Let $M$ be relaxed to a continuous budget variable $x\in[1,\infty)$.  Define
an approximation envelope $a_T(x)$ and an estimation envelope $b_T(x)$ as
follows.

The approximation envelope is a continuous interpolation of the spectral
approximation tail, with the index convention
\[
a_T(x)
=
T\,\Delta(x),
\qquad
\Delta(x)
\approx
\frac12
\sum_{i\ge \lceil x\rceil}
\lambda_i,
\]
or the corresponding regime-specific spectral tail.  In the quadratic
retention regime, the discrete tail is
\[
\Delta(M)
=
\frac12
\sum_{i\ge M}
\lambda_i(\Sigma_\pi),
\]
with effective rank $M-1$.

The estimation envelope $b_T(x)$ is a continuous interpolation of the
type-correct online estimation cost, including all logarithmic factors and,
where relevant, the adaptive aggregation penalty
\[
\PsiAgg_M(T)
=
O\!\left(
\sqrt{T\ln(eM)}
+
\ln(eM)
\right).
\]
Assume:
\begin{enumerate}[label=(\roman*)]
\item $a_T(x)>0$ and $a_T$ is continuous and nonincreasing;
\item $b_T(x)>0$ and $b_T$ is continuous and nondecreasing;
\item there exists a crossing point $x_0$ such that
\[
a_T(x_0)=b_T(x_0).
\]
\end{enumerate}
\end{assumption}

\begin{proposition}[Continuous Bias--Variance Sandwich]
\label{prop:continuous-bv-sandwich}
Under Assumption~\ref{ass:regular-bv-envelope}, define
\[
B_T
=
\sup_{x\ge1}
\min\{a_T(x),b_T(x)\},
\]
and
\[
A_T
=
\inf_{x\ge1}
\bigl[a_T(x)+b_T(x)\bigr].
\]
Then
\[
B_T
\le
A_T
\le
2B_T.
\]
\end{proposition}

\begin{proof}
Let $x_0$ be a crossing point with
\[
a_T(x_0)=b_T(x_0)=B_T.
\]
Because $a_T$ is nonincreasing and $b_T$ is nondecreasing, for $x\le x_0$ we
have $a_T(x)\ge B_T$, while for $x\ge x_0$ we have $b_T(x)\ge B_T$.  Hence
for every $x$,
\[
a_T(x)+b_T(x)
\ge
B_T.
\]
Taking the infimum over $x$ gives $A_T\ge B_T$.

Conversely, evaluating the sum at the crossing point gives
\[
A_T
\le
a_T(x_0)+b_T(x_0)
=
2B_T.
\]
Therefore
\[
B_T
\le
A_T
\le
2B_T.
\]
\end{proof}

\begin{remark}[The Sandwich Is Not by Itself a Minimax Lower Bound]
\label{rem:sandwich-not-minimax}
Proposition~\ref{prop:continuous-bv-sandwich} is an envelope comparison.  It
shows that, under regularity, the sum envelope $A_T$ and the sup--min envelope
$B_T$ are within a factor of $2$.  It does not by itself prove
\[
\inf_{\text{learner}}
\sup_{P}
\Reg_T
\ge
B_T.
\]
Promoting the sandwich to a minimax lower bound requires explicit
approximation floors, explicit estimation floors, and a reduction valid for
arbitrary adaptive learners.  Such floors are supplied, as hypotheses, by
Assumption~\ref{ass:oracle-floors}; they are not established unconditionally
for the regimes considered here.
\end{remark}

\begin{remark}[Conditions Governing the Envelope Comparison]
\label{rem:envelope-conditions}
Four conditions govern the use of the bias--variance sandwich, and each is
required for the statements of this section to be read correctly.
\begin{enumerate}[label=(\arabic*)]
\item The continuous sandwich is valid under
Assumption~\ref{ass:regular-bv-envelope}, that is, for envelopes that are
positive, continuous, monotone in opposite directions, and crossing.
\item The retention spectral tail uses the index convention
\[
\sum_{i\ge M},
\]
not $\sum_{i>M}$: the effective rank is $M-1$, because centring removes one
degree of freedom (Theorem~\ref{thm:spectral-retention}).
\item The estimation envelope must retain its logarithmic factors and, where
model selection is adaptive, the aggregation penalty $\PsiAgg_M(T)$.
\item Rounding a continuous optimal budget to an integer one is not universally
harmless; it requires regularity near the optimum, for example bounded
adjacent spectral ratios or the absence of large spectral gaps.
\end{enumerate}
The envelope comparison is an inequality between deterministic envelopes.  It
becomes a minimax lower bound under the explicit floors of
Assumption~\ref{ass:oracle-floors}, which yield
Theorem~\ref{thm:oracle-minimax-lower}.  The minimax optimality of the adaptive
oracle envelope is therefore conditional on those floors.
\end{remark}

\begin{assumption}[Explicit Approximation and Estimation Floors]
\label{ass:oracle-floors}
Fix a regime $\mathsf r$ and horizon $T$.  Assume:
\begin{enumerate}[label=(\roman*)]
\item \textbf{Packing floor.}
For each $M$ there are finitely many processes $P_M^{1},\dots,P_M^{m}$ in the
class, $m=m(M)$, and a quantity $\Delta_M>0$, such that, writing
\[
R_i(h)
=
\mathbb E_{P_M^{i}}\ell_T(h)-L_T^{\star,\mathsf r}(P_M^{i}),
\]
every predictor $h$, adaptive or not, satisfies
\[
\frac1m\sum_{i=1}^{m}R_i(h)\ \ge\ \Delta_M .
\]
That is, no single predictor is simultaneously near-optimal across the family,
the \emph{average} shortfall being at least $\Delta_M$.

\item \textbf{Envelope calibration.}
The packing parameter dominates the envelope minimum,
\[
\Delta_M\ \ge\ \min\bigl\{\Aapp^{\mathsf r}(M),\Est_M^{\mathsf r}(T)\bigr\} .
\]

The two clauses must be kept apart, and conflating them is a genuine hazard.
Defining $\Delta_M$ to \emph{equal} $\min\{\Aapp^{\mathsf r}(M),
\Est_M^{\mathsf r}(T)\}$ would make the assumption vacuous in every realizable
regime.  There the model class contains the truth, so
$\Aapp^{\mathsf r}(M)=0$ for all $M$ at or above the true state count, the
minimum is $0$, and the floor asserts only $\frac1m\sum_iR_i(h)\ge0$, which is
automatic.  The packing floor and the envelope calibration are therefore
separate demands: the first is a statement about the family, proved by
exhibiting mutually distinguishable processes; the second is a comparison
between that family's information content and the envelope whose optimality is
at issue.  Only the conjunction yields a matching lower bound, and in a
realizable regime the calibration clause is where the content resides, since
the relevant comparison is then against $\Est_M^{\mathsf r}(T)$ alone.

The averaging is over a family whose size may grow with $M$, and this is
essential rather than cosmetic.  A two-point version of this condition, with
$R_0(h)+R_1(h)\ge\Delta_M$, is unsatisfiable once $\Delta_M$ exceeds
$2\log2$.  Indeed under log loss $R_i(h)=\KL(P^i\|Q_h)$ for the transcript law
$Q_h$ induced by $h$, and
\[
\min_{Q}\bigl[\KL(P^0\|Q)+\KL(P^1\|Q)\bigr]
=
2\,\mathrm{JS}(P^0,P^1)
\ \le\
2\log2,
\]
the minimum being attained at the mixture $\tfrac12(P^0+P^1)$.  Since
$\Est_M^{\mathsf r}(T)$ diverges in $T$ in every regime considered here, a
two-point floor would fail for all large $T$.  With $m$ processes the same
computation gives
\[
\min_Q\frac1m\sum_{i=1}^m\KL(P^i\|Q)
=
I(V;Y),
\qquad V\sim\mathrm{Unif}[m],
\]
which is at most $\log m$ and equals $\log m$ when the transcript laws are
mutually singular.  Carrying $\Delta_M=\Theta(M\log M)$ therefore requires
$m=\exp(\Theta(M\log M))$, of the order of $|\mathcal H_M|$.
\item \textbf{Monotonicity.}
$\Aapp^{\mathsf r}(\cdot)$ is nonincreasing and
$\Est^{\mathsf r}_{(\cdot)}(T)$ is nondecreasing.
\end{enumerate}
\end{assumption}

\begin{remark}[What a Single Family Can and Cannot Deliver]
\label{rem:floors-two-axes}
Clauses~(i) and~(ii) constrain a single packing.  The consequence for regimes
in which \emph{both} envelopes are active is the following.

A packing internal to $\mathcal H_M$ certifies estimation difficulty: its
members are all realizable at budget $M$, so it lower-bounds the price of not
knowing which member is in force, and calibrates against
$\Est_M^{\mathsf r}(T)$.  A packing certifying approximation difficulty must
instead place its members \emph{outside} the reach of $\mathcal H_M$, so that
$\inf_{h\in\Hcal_M^{\mathsf r}}R_P(h)\ge a_M>0$.  These are different families
with different type signatures, and exhibiting each separately does not by
itself certify that one adversary forces both burdens at once: a learner facing
the approximation-hard family may pay $a_M$ while facing no estimation burden,
and conversely.

It does not follow, however, that a bound on the \emph{sum} is out of reach
without a direct-sum construction.  The sum-form envelope $\inf_M[a_M+b_M]$ is
the quantity of interest, and it is already reachable.  The discrete
sandwich of Lemma~\ref{lem:discrete-bv-sandwich} already converts between the
two forms: with $a(M)=\Aapp^{\mathsf r}(M)$, $b(M)=\Est_M^{\mathsf r}(T)$,
$B=\sup_M\min\{a(M),b(M)\}$ and $A=\inf_M[a(M)+b(M)]$, it gives
$B\le A\le(1+\kappa)B$ under its nondegeneracy and crossing conditions.
Consequently a floor calibrated against the \emph{minimum} automatically yields
a bound against the \emph{sum},
\[
\sup_P\Reg_T\ \ge\ c\,B\ \ge\ \frac{c}{1+\kappa}\,\inf_{M}\bigl[
\Aapp^{\mathsf r}(M)+\Est_M^{\mathsf r}(T)\bigr],
\]
which is the second display of Theorem~\ref{thm:oracle-minimax-lower}.  The
loss is the jump ratio $\kappa$ at the crossing index, and nothing more.

The role of a direct-sum construction is therefore narrower than it may appear.
It is not needed to reach the sum-form envelope at all; it is needed only to
remove the factor $(1+\kappa)$.  That factor is itself harmless:
by Lemma~\ref{lem:kappa-bounded} the jump ratio of every estimation envelope
occurring here is bounded by an absolute constant --- at most $1+\log2$ in the
agnostic regime and $2(1+\log2)$ in the realizable regime --- independently of
the horizon $T$ and of where the crossing falls.  The sum-form minimax lower
bound therefore holds with an absolute constant, and \emph{no} direct-sum
construction is required for it.

What a both-axes construction would still buy is a bound in which the two
components are certified by a single adversary, which is a statement about the
\emph{source} of the difficulty rather than about the size of the constant.
Accordingly Assumption~\ref{ass:oracle-floors} is stated for the single
quantity $\Delta_M$ calibrated against the minimum, which is what
Lemma~\ref{lem:discrete-bv-sandwich} consumes, and nothing is lost in passing
to the sum.
\end{remark}

\begin{lemma}[The Jump Ratio Is an Absolute Constant]
\label{lem:kappa-bounded}
Let the estimation envelope have the regularly varying form
\[
b(M)=\gamma\,M^{\alpha}\bigl(\log eM\bigr)^{\beta},
\qquad
\alpha,\beta\ge0,\ (\alpha,\beta)\ne(0,0),\ \gamma>0 .
\]
Then for every $M\ge2$,
\[
\frac{b(M)}{b(M-1)}
\ \le\
\frac{b(2)}{b(1)}
=
2^{\alpha}\bigl(1+\log2\bigr)^{\beta},
\]
and consequently the jump ratio of Lemma~\ref{lem:discrete-bv-sandwich}
satisfies
\[
\kappa\ \le\ 2^{\alpha}(1+\log2)^{\beta}
\]
whatever the crossing index.  The bound depends only on the exponents: it is
independent of the horizon $T$, of the scale $\gamma$, and of the
approximation envelope.

If $b$ is a sum of finitely many such terms, $b=\sum_jb_j$ with exponents
$(\alpha_j,\beta_j)$, then
$\kappa\le\max_j2^{\alpha_j}(1+\log2)^{\beta_j}$.
\end{lemma}

\begin{proof}
Write $b(M)/b(M-1)=u(M)\,v(M)$ with
\[
u(M)=\Bigl(\tfrac{M}{M-1}\Bigr)^{\alpha},
\qquad
v(M)=\Bigl(\tfrac{\log eM}{\log e(M-1)}\Bigr)^{\beta}.
\]

\emph{Both factors are nonincreasing on $M\ge2$.}  For $u$ this is immediate,
since $M\mapsto M/(M-1)=1+1/(M-1)$ is decreasing and $t\mapsto t^\alpha$ is
increasing for $\alpha\ge0$.  For $v$ it suffices to treat $\beta=1$ and put
$\ell(x)=\log ex=1+\log x$, which is positive, increasing and concave on
$x\ge1$.  For a positive concave $\ell$ the ratio $\ell(x)/\ell(x-1)$ is
nonincreasing: writing $x\mapsto\ell(x)/\ell(x-1)$ and differentiating, the
sign of the derivative is that of
$\ell'(x)\ell(x-1)-\ell(x)\ell'(x-1)$, and concavity gives
$\ell'(x)\le\ell'(x-1)$ while monotonicity gives $\ell(x-1)\le\ell(x)$, so with
$\ell'>0$ the expression is at most
$\ell'(x-1)\bigl[\ell(x-1)-\ell(x)\bigr]\le0$.  Raising to the power $\beta\ge0$
preserves the monotonicity.

\emph{Conclusion.}  The product of two nonincreasing positive functions is
nonincreasing, so $b(M)/b(M-1)$ attains its supremum over $M\ge2$ at $M=2$,
where its value is $2^{\alpha}(1+\log2)^{\beta}$ because $\log e\cdot1=1$ and
$\log e\cdot2=1+\log2$.  Since $\kappa$ is by definition $b(M^\ast)/b(M^\ast-1)$
for the crossing index $M^\ast\ge2$, or $1$ when $M^\ast=1$, the stated bound
holds in all cases.

\emph{Sums.}  For $b=\sum_jb_j$ the mediant inequality gives
\[
\frac{\sum_jb_j(M)}{\sum_jb_j(M-1)}
\ \le\
\max_j\frac{b_j(M)}{b_j(M-1)}
\ \le\
\max_j2^{\alpha_j}(1+\log2)^{\beta_j},
\]
since for positive $x_j,y_j$ one has
$(\sum_jx_j)/(\sum_jy_j)\le\max_j(x_j/y_j)$.
\end{proof}

\begin{remark}[The Jump Ratio for the Envelopes of This Paper]
\label{rem:kappa-values}
Applying Lemma~\ref{lem:kappa-bounded} to the estimation envelopes used here:

\begin{center}
\begin{tabular}{lcc}
\toprule
estimation envelope & $(\alpha,\beta)$ & $\kappa\le$ \\
\midrule
$\sqrt{T\log eM}+\log eM$ & $(0,\tfrac12),(0,1)$ & $1+\log2\approx1.6931$ \\
$\sqrt{TM\log eM}$        & $(\tfrac12,\tfrac12)$ & $\sqrt{2(1+\log2)}\approx1.8402$ \\
$M\log eM$                & $(1,1)$               & $2(1+\log2)\approx3.3863$ \\
\bottomrule
\end{tabular}
\end{center}

\noindent
Each is an absolute constant.  Consequently the second display of
Theorem~\ref{thm:oracle-minimax-lower} holds with an absolute constant in front
of the sum-form envelope: for instance $c/(1+\kappa)\ge c/2.70$ in the agnostic
finite-class regime and $c/4.39$ in the realizable regime.

Two points of scope.  First, the normalization $\log eM$ rather than $\log M$
is what makes $b$ positive at $M=1$, as condition~(a) of
Lemma~\ref{lem:discrete-bv-sandwich} requires; the bare $M\log M$ vanishes
there.  If one insists on $M\log M$ the crossing index is at least $3$ and the
supremum of consecutive ratios over $M\ge3$ is $3\log3/(2\log2)\approx2.3774$,
again absolute.  Second, the lemma constrains only the estimation envelope:
no regularity is imposed on the approximation envelope beyond the monotonicity
already assumed.
\end{remark}

\begin{lemma}[Discrete Bias--Variance Sandwich]
\label{lem:discrete-bv-sandwich}
Let $a,b:\mathbb N_{\ge1}\to(0,\infty)$ with $a$ nonincreasing and $b$
nondecreasing.  Assume
\begin{enumerate}[label=(\alph*)]
\item $b(1)\le a(1)$;
\item $b(M)\ge a(M)$ for some $M$.
\end{enumerate}
Let $M^\ast=\min\{M:b(M)\ge a(M)\}$, which exists by~(b), and set
\[
\kappa
=
\begin{cases}
b(M^\ast)/b(M^\ast-1), & M^\ast\ge2,\\
1, & M^\ast=1,
\end{cases}
\]
the \emph{jump ratio} of the estimation envelope at the crossing index; note
$\kappa\ge1$ since $b$ is nondecreasing and positive.  Put
\[
B=\sup_{M\ge1}\min\{a(M),b(M)\},
\qquad
A=\inf_{M\ge1}\bigl[a(M)+b(M)\bigr].
\]
Then
\[
B\ \le\ A\ \le\ (1+\kappa)\,B .
\]
No continuity and no exact crossing point are required.
\end{lemma}

\begin{proof}
\emph{Left inequality.}  Fix $M$ and let $N$ be arbitrary.  If $N\ge M$ then
$b(N)\ge b(M)\ge\min\{a(M),b(M)\}$ because $b$ is nondecreasing; if $N<M$ then
$a(N)\ge a(M)\ge\min\{a(M),b(M)\}$ because $a$ is nonincreasing.  Since both
$a$ and $b$ are positive, in either case
\[
a(N)+b(N)\ \ge\ \min\{a(M),b(M)\}.
\]
Taking the infimum over $N$ gives $A\ge\min\{a(M),b(M)\}$, and then the
supremum over $M$ gives $A\ge B$.

\emph{Right inequality.}  Suppose first $M^\ast\ge2$.  By minimality of
$M^\ast$ we have $b(M^\ast-1)<a(M^\ast-1)$, so
\[
B\ \ge\ \min\{a(M^\ast-1),b(M^\ast-1)\}=b(M^\ast-1),
\]
while $b(M^\ast)\ge a(M^\ast)$ gives
\[
B\ \ge\ \min\{a(M^\ast),b(M^\ast)\}=a(M^\ast).
\]
Hence, using $b(M^\ast)=\kappa\,b(M^\ast-1)$,
\[
A
\ \le\
a(M^\ast)+b(M^\ast)
\ =\
a(M^\ast)+\kappa\,b(M^\ast-1)
\ \le\
B+\kappa B
\ =\
(1+\kappa)B .
\]
If instead $M^\ast=1$, then $b(1)\ge a(1)$, which with~(a) forces
$b(1)=a(1)$; thus $B\ge\min\{a(1),b(1)\}=a(1)$ and
$A\le a(1)+b(1)=2a(1)\le 2B=(1+\kappa)B$ since $\kappa=1$.
\end{proof}

\begin{remark}[Role of the Jump Ratio]
\label{rem:sandwich-jump}
The constant cannot be improved to a universal one.  For $a=(10,1)$ and
$b=(1,9)$ on $\{1,2\}$ one has $B=1$ and $A=10$, so $A/B=10=1+\kappa$ with
$\kappa=9$.  What the continuous statement
(Proposition~\ref{prop:continuous-bv-sandwich}) buys by assuming continuity and
an exact crossing is precisely $\kappa=1$, giving the factor $2$.  In the
discrete setting the estimation envelope typically satisfies
$b(M)/b(M-1)=1+O(1/M)$ -- this holds for
$b(M)=\sqrt{T\ln(eM)}+\ln(eM)$ -- so $\kappa\to1$ and the factor tends to $2$.
\end{remark}

\begin{lemma}[The Packing Floor Is an Information Condition]
\label{lem:packing-criterion}
Under log loss, let $P^1,\dots,P^m$ be processes in the class and let $V$ be
uniform on $[m]$.  Then for every predictor $h$, with $Q_h$ the transcript law
it induces,
\[
\frac1m\sum_{i=1}^{m}R_i(h)
\ \ge\
\min_Q\frac1m\sum_{i=1}^{m}\KL(P^i\|Q)
=
I(V;Y_{1:T}),
\]
the minimum being attained at the mixture $\bar P=\frac1m\sum_iP^i$.
Consequently Assumption~\ref{ass:oracle-floors}(i) holds with a given
$\Delta_M$ if and only if some finite subfamily of the class satisfies
\[
I(V;Y_{1:T})\ \ge\ \Delta_M .
\]
\end{lemma}

\begin{proof}
Under log loss $R_i(h)=\KL(P^i\|Q_h)$, so the average is
$\frac1m\sum_i\KL(P^i\|Q_h)$.  For any $Q$,
\[
\frac1m\sum_i\KL(P^i\|Q)
=
\frac1m\sum_i\KL(P^i\|\bar P)+\KL(\bar P\|Q)
=
I(V;Y_{1:T})+\KL(\bar P\|Q),
\]
by the compensation identity, and $\KL(\bar P\|Q)\ge0$ with equality iff
$Q=\bar P$.  The middle expression is the mutual information because $V$ is
uniform and $Y_{1:T}\mid V=i$ has law $P^i$.
\end{proof}

\begin{remark}[What Discharging the Floor Requires in Each Regime]
\label{rem:packing-per-regime}
Lemma~\ref{lem:packing-criterion} reduces the floor to exhibiting a subfamily
carrying $\Delta_M$ nats about its index by horizon $T$.  The regimes differ
sharply in how hard that is.

In the \emph{deterministic} realizable regimes the transcripts of distinct
members are eventually distinct, so the induced laws are mutually singular and
$I(V;Y)=\log m$ exactly, as soon as the horizon exceeds the separating length.
Proposition~\ref{prop:floors-instance} carries this out.

In the \emph{stochastic} regimes distinctness is statistical rather than
logical, and $I(V;Y)$ grows only as fast as the family can be estimated.  For
the natural retention packing --- $M$ hidden Bernoulli biases spaced
$\Theta(1/M)$ apart, indexed by $b\in[M]^M$, so $\log m=M\log M$ --- each
coordinate needs $\Omega(M^2)$ samples to resolve, hence
$T=\Omega(M^3)$ before $I(V;Y)$ approaches $\log m$.  At $T=\Theta(M\log M)$
the mutual information is a small fraction of $\log m$.

Consequently the floor in a stochastic regime should be read with an explicit
horizon hypothesis, and $\Delta_M$ taken as the information actually available
at that horizon rather than as $\log m$.  Constructing packings that are
efficient in $T$ for the retention, grounding, and commitment regimes is open
(Section~\ref{subsec:open-problems}).
\end{remark}

\begin{proposition}[The Packing Floor Is Satisfied in the Realizable Stream Regime]
\label{prop:floors-instance}
Let $\mathsf r$ be the realizable persistent-stream regime with
$\Est_M^{\mathsf r}(T)=\Theta(M\log M)$, and let $\Gact_M$ be the gated family
of Definition~\ref{def:gated-active-family} on the forcing stream $w$ of
Theorem~\ref{thm:stream-lower-bound}, of length $O(M\log M)$.  Take
$\{P_M^{i}\}$ to be the $m=M^{M}$ laws induced by the members of $\Gact_M$.
Then the transcript laws are mutually singular, so
\[
\min_Q\frac1m\sum_{i=1}^{m}\KL(P_M^{i}\|Q)
=
I(V;Y)
=
\log m
=
M\log M,
\]
and the packing floor of Assumption~\ref{ass:oracle-floors}(i) holds with
$\Delta_M=M\log M$.

Moreover the envelope calibration of
Assumption~\ref{ass:oracle-floors}(ii) holds in this regime.  The regime is
realizable, so $\Aapp^{\mathsf r}(M)=0$ and
$\min\{\Aapp^{\mathsf r}(M),\Est_M^{\mathsf r}(T)\}=0\le\Delta_M$; the
substantive comparison is with the estimation envelope, and
$\Delta_M=M\log M=\Theta(\Est_M^{\mathsf r}(T))$ matches it in order.
\end{proposition}

\begin{proof}
Two distinct maps $g\ne g'$ differ at some $v\in Q$ and some coordinate $j$.
The stream $w$ contains the block $w_v\,\mathtt c\,\mathtt d^{L}$, whose
readout rounds emit the coordinates of $g(v)$ in order, and by the gating
property of Definition~\ref{def:gated-active-family} those rounds report $g$ at
the argument $v$ and at no other.  Hence the transcripts of $g$ and $g'$ differ
in the $j$-th readout of that block, and the $m$ transcripts are pairwise
distinct.  Being deterministic machines on a fixed input stream, each induces a
point mass, so the $m$ laws are mutually singular and
$I(V;Y)=H(V)=\log m$ with $V$ uniform.  Finally $|\Gact_M|=M^{M}$ gives
$\log m=M\log M$.

The reduction to the floor is Lemma~\ref{lem:packing-criterion}: every
predictor satisfies $\frac1m\sum_iR_i(h)\ge I(V;Y)=M\log M$.  The horizon
suffices because the forcing stream has length $O(M\log M)$ and separates every
pair of members within it.
\end{proof}

\begin{remark}[Status of the Floor Across Regimes]
\label{rem:floors-status}
Proposition~\ref{prop:floors-instance} discharges
Assumption~\ref{ass:oracle-floors}(i) and~(ii) for the realizable
persistent-stream regime, so Theorem~\ref{thm:oracle-minimax-lower} is
unconditional there.  Four qualifications delimit what this does and does not
settle.

Zeroth, and prior to the other three, the loss types differ and must not be
silently identified.  Lemma~\ref{lem:packing-criterion} and
Proposition~\ref{prop:floors-instance} are statements about logarithmic loss,
with $R_i(h)=\KL(P^i\|Q_h)$ and the packing certified by mutual singularity of
transcript laws; Theorem~\ref{thm:stream-lower-bound}, by contrast, is a $0/1$
mistake theorem.  The mutual-information computation therefore does not by
itself yield a $0/1$ lower bound, and the two should be read as separate
results about the same family: an unconditional $\Theta(M\log M)$ mistake bound
in the realizable stream protocol, and a log-loss packing statement with
$\Delta_M=M\log M$.  Neither is derived from the other here.

First, the regime is realizable, so the approximation envelope vanishes and the
floor is calibrated against the estimation envelope alone.  The proposition
therefore certifies the estimation axis; it says nothing about the
approximation axis; by Remark~\ref{rem:floors-two-axes} the sum-form envelope
is nevertheless reached through the discrete sandwich, at the cost of the jump
ratio $(1+\kappa)$, which Lemma~\ref{lem:kappa-bounded} bounds by an absolute
constant.

Second, the corresponding constructions for the retention, grounding, and
commitment regimes are not carried out here and remain open
(Section~\ref{subsec:open-problems}).  Those regimes are stochastic, and by
Remark~\ref{rem:packing-per-regime} a packing there must additionally be
efficient in the horizon.

Third, what the proposition establishes in general terms is that the hypothesis
is satisfiable and has at least one verified instance, which the two-point form
of the floor did not: that form is unsatisfiable for $\Delta_M>2\log2$, as
recorded in Assumption~\ref{ass:oracle-floors}(i).
\end{remark}

\begin{theorem}[Minimax Lower Bound for the Oracle Envelope]
\label{thm:oracle-minimax-lower}
Under Assumption~\ref{ass:oracle-floors} the bounds below hold with $c=1$;
they are written with a general constant $c>0$ since only its positivity is
used.  For every learner, adaptive over $M$ or not,
\[
\sup_{P}\Reg_T
\;\ge\;
c\,
\sup_{M\ge1}
\min\bigl\{
\Aapp^{\mathsf r}(M),\,
\Est_M^{\mathsf r}(T)
\bigr\}.
\]
If in addition the discrete nondegeneracy and crossing conditions~(a),~(b) of
Lemma~\ref{lem:discrete-bv-sandwich} hold for
$a(M)=\Aapp^{\mathsf r}(M)$ and $b(M)=\Est_M^{\mathsf r}(T)$, with jump ratio
$\kappa=\kappa(T)$ at the crossing index, then
\[
\sup_{P}\Reg_T
\;\ge\;
\frac{c}{1+\kappa}
\inf_{M\ge1}
\left[
\Aapp^{\mathsf r}(M)+\Est_M^{\mathsf r}(T)
\right],
\]
which matches the adaptive upper bound of
Theorem~\ref{thm:oracle-agnostic} up to the aggregation penalty
$\PsiAgg_M(T)$ and the factor $1+\kappa(T)$.

The matching is up to \emph{universal} constants precisely when
$\sup_T\kappa(T)<\infty$; otherwise the constant degrades with $\kappa(T)$ and
the two-sided statement is correspondingly weaker.  For the estimation
envelopes arising here the condition holds with room to spare: for
$\Est_M^{\mathsf r}(T)=\sqrt{T\ln(eM)}+\ln(eM)$ one has
$\kappa(T)=1+O(1/M^\ast)$, so $\kappa(T)\to1$ and the factor tends to $2$,
recovering the constant of the continuous sandwich
(Proposition~\ref{prop:continuous-bv-sandwich}) without its continuity and
exact-crossing hypotheses.
\end{theorem}

\begin{proof}
Fix $M$ and let $\mathcal A$ be any learner, deterministic or randomized,
adaptive or not.  A learner is a measurable map from transcripts to
predictions, so it is in particular an admissible $h$ in the packing floor.
Applying that floor to $h=\mathcal A$ gives
\[
\frac1m\sum_{i=1}^{m}R_i(\mathcal A)\ \ge\ \Delta_M ,
\]
and a maximum dominates an average, so
\[
\max_{1\le i\le m}R_i(\mathcal A)\ \ge\ \Delta_M .
\]
Since the adversary may select the process,
\[
\sup_P\Reg_T(\mathcal A)\ \ge\ \Delta_M ,
\]
and as $M$ was arbitrary,
\[
\sup_P\Reg_T
\ \ge\
\sup_{M\ge1}
\min\bigl\{\Aapp^{\mathsf r}(M),\Est_M^{\mathsf r}(T)\bigr\},
\]
which is the first display with $c=1$.  The averaged form is in fact stronger
than the two-point form it replaces: the latter yielded only
$\tfrac12\Delta_M$ through $\max\{u,v\}\ge\tfrac12(u+v)$, whereas here the
maximum dominates the average outright.

For the second display, put $a(M)=\Aapp^{\mathsf r}(M)$ and
$b(M)=\Est_M^{\mathsf r}(T)$.  These are positive, $a$ nonincreasing and $b$
nondecreasing, by the monotonicity clause of
Assumption~\ref{ass:oracle-floors}.  Monotonicity alone is not enough:
Lemma~\ref{lem:discrete-bv-sandwich} additionally requires the nondegeneracy
condition $b(1)\le a(1)$ and the crossing condition $b(M)\ge a(M)$ for some
$M$, and these are hypotheses of the present clause, not consequences of
monotonicity.  Under them, and with $\kappa$ the jump ratio at the crossing
index,
\[
\sup_M\min\{a,b\}=B\ \ge\ \frac{1}{1+\kappa}A
=\frac{1}{1+\kappa}\inf_M\bigl[a(M)+b(M)\bigr].
\]
Combining with the first display, and recalling $c=1$, yields
\[
\sup_P\Reg_T
\ \ge\
\frac{1}{1+\kappa}
\inf_{M\ge1}
\left[\Aapp^{\mathsf r}(M)+\Est_M^{\mathsf r}(T)\right].
\]
The continuous route through
Assumption~\ref{ass:regular-bv-envelope} and
Proposition~\ref{prop:continuous-bv-sandwich} gives the same conclusion with
$1+\kappa$ replaced by $2$, but at the price of assuming continuity of the
interpolated envelopes and an exact crossing point; the discrete route assumes
neither.
\end{proof}

\begin{remark}[Sum-Separation versus Two-Point Testing]
\label{rem:oracle-no-lecam}
Two distinct hypotheses yield a minimax lower bound of this type, and they
require different arguments.  Under the packing floor of
Assumption~\ref{ass:oracle-floors}, which constrains \emph{every} predictor,
the bound follows directly and no indistinguishability argument enters; the
floor is an averaged Fano-type condition, and
Proposition~\ref{prop:floors-instance} exhibits a family satisfying it.  Under
the weaker hypothesis that only the \emph{optimal} predictors of the two
processes are separated, a Le Cam or Hellinger two-point argument is required,
together with a demonstration that any learner with small regret on both
processes would distinguish the transcript laws; see
\cite[Ch.~2]{tsybakov2009}.  Combining a total-variation affinity condition
with a packing condition is redundant, since the latter alone suffices.
The formulation adopted here is the packing form.
\end{remark}

\begin{remark}[What Remains Genuinely Open]
\label{rem:oracle-minimax-residual}
Theorem~\ref{thm:oracle-minimax-lower} closes the two-sided characterization
\emph{up to the aggregation penalty} $\PsiAgg_M(T)$ and universal constants,
conditional on the explicit floors of Assumption~\ref{ass:oracle-floors}.  Two
finer questions remain.  First, whether the penalty
$\PsiAgg_M(T)=O(\sqrt{T\log(eM)})$ is itself necessary, i.e.\ whether adaptive
aggregation over an unknown budget is strictly harder than the oracle problem;
the second-order $\log(eM)$ term in the parameter-free bound suggests a
separation, though a small one.  Second, whether the floors of
Assumption~\ref{ass:oracle-floors} hold unconditionally for the specific
finite-state regimes of this manuscript, which is a statement about spectral
decay of the associated operators rather than about the oracle mechanism.
\end{remark}

%----------------------------------------------------------------------
\subsection{Scope and Conditions}
\label{subsec:oracle-scope}
%----------------------------------------------------------------------

\begin{itemize}[leftmargin=1.6em]
\item \textbf{Type correctness.}
The bridge is stated within a fixed regime and a fixed hypothesis-class type.

\item \textbf{Benchmark.}
The benchmark $L_T^{\star,\mathsf r}$ is the best finite-state predictor in
the regime-specific union.  It is called the true-process or Bayes benchmark
only under an explicit richness assumption.

\item \textbf{Converse versus achievability.}
Spectral converses are lower bounds on $\Aapp^{\mathsf r}(M)$.  They become
regret upper rates only under the spectral-rate assumption of
Definition~\ref{def:spectral-rate}.  The oracle upper bounds of
Theorems~\ref{thm:oracle-agnostic}, \ref{thm:oracle-realizable} and
\ref{thm:oracle-logloss} are stated in terms of the true approximation deficit
$\Aapp^{\mathsf r}(M)$; expressing the resulting regret as a spectral rate
requires an \emph{upper} bound on that deficit by a spectral tail, which a
converse does not supply.

\item \textbf{Active synchronization.}
The active realizable estimation rate is
\[
\Est_M^{\mathrm{active}}(T)
=
\Theta(M\log M)
\]
unconditionally on the full class, by Corollary~\ref{cor:active-theta}.  The
oracle inequality remains valid with this rate.

\item \textbf{No cross-regime unification.}
The bridge composes the offline and online axes within one regime.  It does
not unify the three regimes and does not produce a universal ordering of
grounding, retention, and commitment gaps.
\end{itemize}

%======================================================================
\section{The Linear Second-Order Safety Surrogate}
\label{sec:pos-linear}
%======================================================================

This section develops a linear second-order safety surrogate.  The surrogate
is exact for a linear block-local problem, but it is not an exact formula for
the discrete right-congruence Price of Safety.  The distinction is
type-theoretic: pinching and majorization live in a linear operator surrogate,
whereas the discrete Price of Safety is a constrained optimization over right
congruences.

%----------------------------------------------------------------------
\subsection{Type-Correct Objects}
\label{subsec:pos-objects}
%----------------------------------------------------------------------

Let $\mathcal S$ be a finite state set with stationary weights
$\pi_s>0$, $\sum_s\pi_s=1$.  Let $y_s\in\mathbb R^d$ be the predictive or
Fisher feature of state $s$, and define the centered features
\[
x_s=y_s-\bar y,
\qquad
\bar y=\sum_s\pi_s y_s.
\]
Let
\[
D_\pi=\operatorname{diag}(\pi_s)
\]
and let $X$ be the matrix with rows $x_s^\top$.  The feature covariance is
\[
\Sigma_\pi
=
X^\top D_\pi X.
\]
The state-indexed centered Gram matrix is
\[
G
=
D_\pi^{1/2}XX^\top D_\pi^{1/2},
\]
with entries
\[
G_{st}
=
\sqrt{\pi_s\pi_t}\,
\langle x_s,x_t\rangle.
\]
The matrix $G$ is positive semidefinite and satisfies
\[
G\sqrt{\pi}=0,
\qquad
\sqrt{\pi}:=(\sqrt{\pi_s})_{s\in\mathcal S}.
\]

Let
\[
\mathcal B=\{B_1,\dots,B_r\}
\]
be a safety partition of the state set.  Let $P_b$ be the coordinate projection
onto $B_b$.  The block-diagonal algebra is
\[
\mathcal A
=
\left\{
X:\;X=\sum_b P_bXP_b
\right\}.
\]
The associated pinching map is
\[
\EA(X)
=
\sum_b P_bXP_b.
\]
This map is unital, positive, trace-preserving, idempotent, and, in the
complexified setting, completely positive.

For an arbitrary matrix $X$, let
\[
s_1(X)\ge s_2(X)\ge\cdots\ge0
\]
denote its singular values, and define the Ky Fan $r$-norm by
\[
\KyF{X}{r}
=
\sum_{i=1}^r
s_i(X),
\qquad
\KyF{X}{0}=0.
\]
If $X\succeq0$ this coincides with the eigenvalue form
\[
\KyF{X}{r}
=
\sum_{i=1}^r
\lambda_i^\downarrow(X),
\]
so the two conventions agree on $G$ and $\EA G$.  The singular-value form is
the one used for the indefinite difference $G-\EA G$ in
Corollary~\ref{cor:pos-coupling}, where it is a genuine norm and in
particular satisfies the triangle inequality.

\begin{remark}[A More Partition-Faithful Linear Safe Relaxation]
\label{rem:partition-faithful}
If one wants a linear relaxation closer to safe partitions, require the safe
linear projection to contain the safety-block label subspace.  Let
\[
g_b=P_b\sqrt{\pi}
\]
be the weighted indicator of safety block $B_b$, and let
\[
\mathcal S_{\mathcal A}
=
\operatorname{span}\{g_b:b=1,\dots,r\}.
\]
A partition-faithful safe linear relaxation is
\[
\mathrm{Safe}_{\mathrm{lin}}^{\mathrm{part}}(M)
=
\frac12
\max_{\substack{
P=P^\top=P^2,\;P\in\mathcal A\\
\rank P\le M,\;
\mathcal S_{\mathcal A}\subseteq\operatorname{range}P
}}
\tr(PG),
\]
when the feasible set is nonempty, in particular when $M$ is at least the
number of safety blocks.  If the feasible set is empty, the constrained
maximum is infeasible.

This constrained PCA problem respects the fact that a safe partition must at
least distinguish safety blocks.  It does not reduce to the simple pinching
formula $\KyF{\EA G}{M-1}$.  The pinching law of this section is the
block-local surrogate, not this partition-faithful relaxation.
\end{remark}

\begin{remark}[Compatibility with Causal States]
\label{rem:pos-causal-compatibility}
If the retention variables are causal states, the safety blocks should be
unions of causal states, i.e.\ the causal-state equivalence should refine the
safety equivalence.  In the stationary formulation, the safety blocks should
be compatible with lumpable quotients of $\Splus$.

If instead safety splits causal states, the construction below may still be
applied to a history-indexed Gram matrix, but the resulting quantity is a
history-level safety surrogate rather than the causal-state Price of Safety,
and the two should not be identified.
\end{remark}

%----------------------------------------------------------------------
\subsection{The Linear Second-Order Block-Local Surrogate}
\label{subsec:pos-surrogate}
%----------------------------------------------------------------------

\begin{definition}[Free Linear Retained Information]
\label{def:free-linear}
For an $M$-state budget, set $r=M-1$.  The free linear second-order retained
information is
\[
\mathrm{Free}_{\mathrm{lin}}(M)
=
\frac12
\KyF{G}{M-1}.
\]
\end{definition}

\begin{definition}[Block-Local Safe Linear Retained Information]
\label{def:safe-linear-local}
A block-local linear summary of dimension $r$ is an orthogonal projection $P$
of rank at most $r$ lying in the block algebra $\mathcal A$.  Define
\[
\mathrm{Safe}_{\mathrm{lin}}^{\mathrm{loc}}(M)
=
\frac12
\max_{\substack{
P=P^\top=P^2\\
\rank P\le M-1\\
P\in\mathcal A
}}
\tr(PG).
\]
\end{definition}

\begin{lemma}[Safe Block-Local Maximum as a Ky Fan Norm]
\label{lem:safe-local-kyfan}
For $r=M-1$,
\[
\max_{\substack{
P=P^\top=P^2\\
\rank P\le r\\
P\in\mathcal A
}}
\tr(PG)
=
\KyF{\EA G}{r}.
\]
Consequently,
\[
\mathrm{Safe}_{\mathrm{lin}}^{\mathrm{loc}}(M)
=
\frac12
\KyF{\EA G}{M-1}.
\]
\end{lemma}

\begin{proof}
If $P\in\mathcal A$, then $P=\sum_b P_bPP_b$, and therefore
\[
\tr(PG)
=
\tr(P\EA G).
\]
Since $\EA G$ is block diagonal, its spectral projections can be chosen block
diagonal, i.e.\ inside $\mathcal A$.  Therefore the maximum over block-diagonal
projections of rank at most $r$ equals the unrestricted Ky Fan maximum for
$\EA G$, namely $\KyF{\EA G}{r}$.
\end{proof}

\begin{remark}[What the Surrogate Compares]
\label{rem:poslin-measures}
The surrogate of Definition~\ref{def:poslin} compares unconstrained linear
summaries with \emph{block-local} linear summaries.  This is not the safe
right-congruence problem.  A safe right congruence is a discrete partition
whose classes lie inside safety blocks; such a partition may distinguish whole
safety blocks from one another, and the corresponding centered contrast vectors
are then not block-local.  The block-local constraint is therefore strictly
more restrictive than safety, and the surrogate is a different optimization
object, related to the discrete problem only through the relaxation gaps of
Proposition~\ref{prop:pos-relaxation-identity}.
\end{remark}

\begin{definition}[Linear Block-Local Price of Safety Surrogate]
\label{def:poslin}
The linear surrogate is
\[
\PoSlin(M)
=
\mathrm{Free}_{\mathrm{lin}}(M)
-
\mathrm{Safe}_{\mathrm{lin}}^{\mathrm{loc}}(M)
=
\frac12
\KyF{G}{M-1}
-
\frac12
\KyF{\EA G}{M-1}.
\]
\end{definition}

%----------------------------------------------------------------------
\subsection{Pinching Law, Majorization, and Coupling}
\label{subsec:pos-pinching}
%----------------------------------------------------------------------

\begin{theorem}[Pinching Majorization and Surrogate Nonnegativity]
\label{thm:pinching}
For positive semidefinite $G$,
\[
\lambda(\EA G)\prec\lambda(G).
\]
Hence for every $r\ge0$,
\[
\KyF{\EA G}{r}
\le
\KyF{G}{r},
\]
and therefore
\[
\PoSlin(M)\ge0.
\]
\end{theorem}

\begin{proof}
The pinching map is an average over block sign or block phase unitaries.  In
the complex case,
\[
\EA(X)
=
\int_{[0,2\pi]^q}
U_\theta XU_\theta^{*}\,d\theta,
\qquad
U_\theta=\sum_{b=1}^q e^{i\theta_b}P_b,
\]
with $d\theta$ normalized Haar measure on the torus; in the real case one
averages instead over the $2^q$ block sign flips
$U_\varepsilon=\sum_b\varepsilon_bP_b$, $\varepsilon_b\in\{\pm1\}$.  In either
case the off-block entries average to zero and the diagonal blocks are
preserved, which is the asserted identity for $\EA$.  Ky Fan norms are
unitarily invariant and convex, so
\[
\KyF{\EA G}{r}
\le
\int
\KyF{U_\theta G U_\theta^*}{r}
\,d\theta
=
\KyF{G}{r}.
\]
This is majorization of the pinched spectrum by the original spectrum.
\end{proof}

\begin{corollary}[Coupling Upper Bound]
\label{cor:pos-coupling}
For every $M\ge1$,
\[
\PoSlin(M)
\le
\frac12
\KyF{G-\EA G}{M-1}.
\]
\end{corollary}

\begin{proof}
By the Ky Fan triangle inequality,
\[
\KyF{G}{r}
\le
\KyF{\EA G}{r}
+
\KyF{G-\EA G}{r}.
\]
Rearranging and multiplying by $1/2$ gives the claim.
\end{proof}

\begin{remark}[What the Coupling Bound Measures]
\label{rem:pos-coupling-design}
The quantity $\tfrac12\KyF{G-\EA G}{M-1}$ is the Ky Fan size of the
cross-block coupling removed by safety pinching, and it bounds the linear
surrogate from above.  The surrogate may be strictly smaller: the bound is
tight only when the top $(M-1)$-dimensional eigenspaces of $G$ and $\EA G$ are
aligned, and misalignment strictly decreases $\PoSlin(M)$ relative to the
coupling term.  The quantity is therefore a design diagnostic --- it identifies
which cross-block interactions safety destroys --- rather than an evaluation of
the surrogate.
\end{remark}

\begin{proposition}[Relaxation-Gap Decomposition of the Price of Safety]
\label{prop:pos-relaxation-identity}
Let $\mathrm{Free}_{\mathrm{quad}}(M)$ and $\mathrm{Safe}_{\mathrm{quad}}(M)$ denote the free and safe optima of the
discrete right-congruence problem at budget $M$, safety being understood in the
sense of Definition~\ref{def:safe-right-cong} and $M\ge r$ so that both are
finite (Proposition~\ref{prop:pos-quad-consistent}), and let $\mathrm{Free}_{\mathrm{lin}}(M)$ and
$\mathrm{Safe}_{\mathrm{lin}}^{\mathrm{loc}}(M)$ denote the corresponding linear block-local surrogates, so that
\[
\PoSquad(M)=\mathrm{Free}_{\mathrm{quad}}(M)-\mathrm{Safe}_{\mathrm{quad}}(M),
\qquad
\PoSlin(M)=\mathrm{Free}_{\mathrm{lin}}(M)-\mathrm{Safe}_{\mathrm{lin}}^{\mathrm{loc}}(M).
\]
Define the \emph{signed discrepancies}
\[
\rho_{\mathrm{free}}(M)=\mathrm{Free}_{\mathrm{lin}}(M)-\mathrm{Free}_{\mathrm{quad}}(M),
\qquad
\rho_{\mathrm{safe}}(M)=\mathrm{Safe}_{\mathrm{lin}}^{\mathrm{loc}}(M)-\mathrm{Safe}_{\mathrm{quad}}(M).
\]
Neither is a relaxation gap: the block-local linear safe value is neither an
upper nor a lower bound on the discrete safe value, so $\rho_{\mathrm{safe}}$
may take either sign.  In the two-state singleton example of
Remark~\ref{rem:safe-discrete-not-captured} one has $\rho_{\mathrm{safe}}(2)=-\tfrac14$.
Then, identically,
\[
\PoSquad(M)
=
\PoSlin(M)
+
\rho_{\mathrm{safe}}(M)
-
\rho_{\mathrm{free}}(M).
\]
Consequently:
\begin{enumerate}[label=(\roman*)]
\item the surrogate is exact, $\PoSquad(M)=\PoSlin(M)$, precisely when
$\rho_{\mathrm{free}}(M)=\rho_{\mathrm{safe}}(M)$;
\item in general
$\bigl|\PoSquad(M)-\PoSlin(M)\bigr|=\bigl|\rho_{\mathrm{safe}}(M)-\rho_{\mathrm{free}}(M)\bigr|$,
so a two-sided bound on the two discrepancies transfers to a bound on the
surrogate error.  No one-sided bound follows from either discrepancy alone,
since both may be negative.
\end{enumerate}
This is an algebraic identity together with its immediate consequences; it
supplies no unconditional comparison between $\PoSquad$ and $\PoSlin$.
\end{proposition}

\begin{proof}
The identity is immediate from the definitions:
\[
\PoSlin+\rho_{\mathrm{safe}}-\rho_{\mathrm{free}}
=
(\mathrm{Free}_{\mathrm{lin}}-\mathrm{Safe}_{\mathrm{lin}}^{\mathrm{loc}})+(\mathrm{Safe}_{\mathrm{lin}}^{\mathrm{loc}}-\mathrm{Safe}_{\mathrm{quad}})-(\mathrm{Free}_{\mathrm{lin}}-\mathrm{Free}_{\mathrm{quad}})
=
\mathrm{Free}_{\mathrm{quad}}-\mathrm{Safe}_{\mathrm{quad}}
=
\PoSquad .
\]
Claims~(i) and~(ii) are immediate from the identity, since
$\PoSquad-\PoSlin=\rho_{\mathrm{safe}}-\rho_{\mathrm{free}}$.
\end{proof}

\begin{remark}[What This Reduces the Open Problem To]
\label{rem:pos-reduction}
Proposition~\ref{prop:pos-relaxation-identity} does not evaluate $\PoSquad$.
It reduces the comparison between the discrete Price of Safety and its linear
surrogate to a single signed quantity,
$\rho_{\mathrm{safe}}-\rho_{\mathrm{free}}$.  Because neither discrepancy has a
determined sign, this is a two-sided reduction only, and no unconditional
inequality between $\PoSquad$ and $\PoSlin$ follows.  What is unconditional in
this section is $\PoSlin(M)\ge0$ by pinching majorization
(Theorem~\ref{thm:pinching}) and the Ky Fan coupling bound
(Corollary~\ref{cor:pos-coupling}).  Whether
$\rho_{\mathrm{safe}}-\rho_{\mathrm{free}}$ admits a bound in terms of spectral
data of $G$ is open.
\end{remark}

%----------------------------------------------------------------------
\subsection{Equality Conditions and Free-Lunch Criteria}
\label{subsec:pos-equality}
%----------------------------------------------------------------------

\begin{lemma}[Equality in Pinching Majorization]
\label{lem:equality-pinching}
Let $X$ be Hermitian.  If
\[
\lambda(\EA X)=\lambda(X)
\]
as multisets, then
\[
X=\EA X.
\]
\end{lemma}

\begin{proof}
Equality of spectra implies equality of traces of squares:
\[
\tr(X^2)=\tr((\EA X)^2).
\]
Write $X_{bc}=P_bXP_c$.  Then
\[
\tr(X^2)
=
\sum_{b,c}\tr(X_{bc}X_{cb})
=
\sum_b
\norm{X_{bb}}_F^2
+
\sum_{b\neq c}
\norm{X_{bc}}_F^2,
\]
while
\[
\tr((\EA X)^2)
=
\sum_b
\norm{X_{bb}}_F^2.
\]
Equality forces $\norm{X_{bc}}_F=0$ for all $b\neq c$, hence
\[
X=\sum_bX_{bb}=\EA X.
\]
\end{proof}

\begin{theorem}[Free Lunch for All Budgets]
\label{thm:pos-allM}
The following are equivalent:
\begin{enumerate}[label=(\roman*)]
\item $\PoSlin(M)=0$ for every $M\ge1$;
\item $\KyF{G}{r}=\KyF{\EA G}{r}$ for every $r\ge0$;
\item $\lambda(G)=\lambda(\EA G)$;
\item $G=\EA G$, i.e.\ $G$ is block diagonal with respect to the safety
partition.
\end{enumerate}
\end{theorem}

\begin{proof}
The equivalence of (i) and (ii) is the definition of $\PoSlin$.  Equality of
all Ky Fan norms is equivalent to equality of spectra, giving
(ii) $\Leftrightarrow$ (iii).  Lemma~\ref{lem:equality-pinching} gives
(iii) $\Rightarrow$ (iv), and (iv) trivially implies (i).
\end{proof}

\begin{theorem}[Fixed-Budget Free Lunch]
\label{thm:pos-fixedM}
Fix $M\ge1$ and set $r=M-1$.  Then
\[
\PoSlin(M)=0
\]
if and only if
\[
\KyF{G}{r}=\KyF{\EA G}{r}.
\]
Equivalently, there exists a block-local projection $P\in\mathcal A$ of rank
at most $r$ such that
\[
\tr(PG)=\KyF{G}{r}.
\]
Thus a free optimal linear summary can be chosen block-local.  If the top
eigenspaces have multiplicities, this statement is existential: not every free
optimal summary need be block-local.
\end{theorem}

\begin{proof}
If $\PoSlin(M)=0$, then the Ky Fan norms are equal.  Let $P$ be a top spectral
projection of $\EA G$ of rank $r$.  Since $\EA G$ is block diagonal, $P$ may
be chosen in $\mathcal A$.  Then
\[
\tr(PG)
=
\tr(P\EA G)
=
\KyF{\EA G}{r}
=
\KyF{G}{r},
\]
so $P$ is a free optimizer.  The converse is immediate: if such a block-local
free optimizer exists, then
\[
\KyF{\EA G}{r}
\ge
\tr(P\EA G)
=
\tr(PG)
=
\KyF{G}{r},
\]
and the reverse inequality is Theorem~\ref{thm:pinching}.
\end{proof}

\begin{remark}[Block Diagonality Is Sufficient but Not Necessary at Fixed $M$]
\label{rem:pos-fixed-not-blockdiag}
For a fixed budget $M$, safety can be free even if $G$ has cross-block
coupling, provided that coupling lies outside the top $M-1$ spectral directions
relevant to the free optimum.  Full block diagonality is the necessary and
sufficient condition for safety to be free simultaneously for all budgets.
\end{remark}

%----------------------------------------------------------------------
\subsection{Relation to the Discrete Price of Safety}
\label{subsec:pos-discrete}
%----------------------------------------------------------------------

The discrete Price of Safety is an optimization over right congruences.  In
the quadratic surrogate, a partition
$\mathcal C=\{C_1,\dots,C_K\}$ of states has cluster weights
\[
\Pi_k=\sum_{s\in C_k}\pi_s
\]
and centroids
\[
c_k=\Pi_k^{-1}\sum_{s\in C_k}\pi_s y_s.
\]
Its second-order retained information is
\[
I_{\mathrm{quad}}(\mathcal C;Z)
=
\frac12
\sum_{k=1}^K
\Pi_k
\norm{c_k-\bar y}^2.
\]
Define the quadratic discrete free and safe values by
\[
\mathrm{Free}_{\mathrm{quad}}(M)
=
\max_{\substack{
\text{free right congruences}\\
\operatorname{index}\le M
}}
I_{\mathrm{quad}}(\mathcal C;Z),
\]
and
\[
\mathrm{Safe}_{\mathrm{quad}}(M)
=
\max_{\substack{
\text{safe right congruences}\\
\operatorname{index}\le M
}}
I_{\mathrm{quad}}(\mathcal C;Z).
\]
The quadratic Price of Safety is
\[
\PoSquad(M)
=
\mathrm{Free}_{\mathrm{quad}}(M)
-
\mathrm{Safe}_{\mathrm{quad}}(M).
\]

The safe feasible set can be empty, so the second maximum requires a
convention.  The notion of safety used here is the one described in
Remark~\ref{rem:poslin-measures}, stated for reference.

\begin{definition}[Safe Right Congruence]
\label{def:safe-right-cong}
Let $\mathcal B=\{B_1,\dots,B_r\}$ be a safety partition of $\Splus$, with $r$
the number of blocks of positive stationary mass.  A right congruence $\sim$ on
$\Splus$ is \textbf{safe} if every $\sim$-class is contained in some $B_b$;
equivalently, if $\sim$ refines the safety equivalence
$s\sim_{\mathcal B}s'\iff s,s'$ lie in a common $B_b$.  When no safe right
congruence of index at most $M$ exists we set
$\mathrm{Safe}_{\mathrm{quad}}(M)=-\infty$, so that $\PoSquad(M)=+\infty$.
\end{definition}

\begin{proposition}[Feasibility and Range of the Discrete Price of Safety]
\label{prop:pos-quad-consistent}
With Definition~\ref{def:safe-right-cong}:
\begin{enumerate}[label=(\roman*)]
\item a safe right congruence of index at most $M$ exists if and only if
$M\ge r$;
\item $\PoSquad(M)=+\infty$ exactly when $M<r$; in particular, for $r\ge2$ the
one-class congruence is not safe and $\PoSquad(1)=+\infty$;
\item for $M\ge r$ one has
$0\le\PoSquad(M)\le\mathrm{Free}_{\mathrm{quad}}(M)$, and $\PoSquad$ is finite.
\end{enumerate}
\end{proposition}

\begin{proof}
\emph{(i)} If $M\ge r$, the safety partition $\mathcal B$ itself is safe, each
class being a block, and has index $r\le M$.  Conversely, a safe congruence has
every class inside a single $B_b$, so each of the $r$ positive-mass blocks
contains at least one class, whence the index is at least $r$.

\emph{(ii)} Immediate from~(i) and the convention of
Definition~\ref{def:safe-right-cong}.  For $r\ge2$ the unique congruence of
index $1$ has the single class $\Splus$, which is contained in some $B_b$ only
if $r=1$.

\emph{(iii)} For $M\ge r$ the safe feasible set is nonempty and is contained in
the free feasible set, so
$\mathrm{Safe}_{\mathrm{quad}}(M)\le\mathrm{Free}_{\mathrm{quad}}(M)$ and
$\PoSquad(M)\ge0$.  Both values are finite, being maxima of
$I_{\mathrm{quad}}$ over finitely many partitions, and
$I_{\mathrm{quad}}\ge0$ gives the upper bound.
\end{proof}

\begin{remark}[The Convention Is Not Inherited from the Surrogate]
\label{rem:pos-convention-independent}
Definition~\ref{def:safe-right-cong} is a choice about the discrete problem and
is not forced by the linear surrogate, which imposes the strictly stronger
block-local constraint $P\in\mathcal A$.  The two differ already on four states
with safety blocks $B_1=\{s_1,s_2\}$, $B_2=\{s_3,s_4\}$: the congruence
$\{\{s_1,s_2\},\{s_3,s_4\}\}$ is safe, every class lying inside a block, yet its
centred block contrasts separate $B_1$ from $B_2$ and are therefore not
block-diagonal.  This is the discrepancy recorded in
Remark~\ref{rem:poslin-measures}, and it is why the surrogate and the discrete
problem are related only through the signed discrepancies of
Proposition~\ref{prop:pos-relaxation-identity} rather than by inheritance of
feasible sets.  The partition-faithful linear relaxation
$\mathrm{Safe}_{\mathrm{lin}}^{\mathrm{part}}$ adopts the matching feasibility
convention, being declared infeasible on the same range $M<r$.
\end{remark}

\begin{proposition}[Free Discrete Objective Is Bounded by the Free Linear Surrogate]
\label{prop:free-discrete-bound}
For every $M$,
\[
\mathrm{Free}_{\mathrm{quad}}(M)
\le
\frac12
\KyF{G}{M-1}.
\]
\end{proposition}

\begin{proof}
Let $\mathcal C=\{C_1,\dots,C_K\}$ be a partition with $K\le M$, with cell
weights $\Pi_k$ and centroids $c_k$, and define the between-cluster covariance
\[
B_{\mathcal C}
=
\sum_{k=1}^K\Pi_k(c_k-\bar y)(c_k-\bar y)^\top ,
\]
so that
\[
I_{\mathrm{quad}}(\mathcal C;Z)=\tfrac12\tr(B_{\mathcal C}).
\]
The ANOVA decomposition gives $\Sigma_\pi=W_{\mathcal C}+B_{\mathcal C}$ with
$W_{\mathcal C}\succeq0$ the within-cluster covariance, whence
\[
0\preceq B_{\mathcal C}\preceq\Sigma_\pi .
\]
Because $\sum_k\Pi_k(c_k-\bar y)=0$, the centred centroids satisfy one linear
relation, so $\rank(B_{\mathcal C})\le K-1\le M-1$.  Ky Fan's maximum
principle \cite{kyfan1951} therefore gives
\[
\tr(B_{\mathcal C})\ \le\ \KyF{\Sigma_\pi}{M-1}.
\]
The nonzero spectra of $\Sigma_\pi$ and $G$ coincide, since $G$ is the Gram
matrix of the same weighted centred vectors, so
$\KyF{\Sigma_\pi}{M-1}=\KyF{G}{M-1}$ and
\[
I_{\mathrm{quad}}(\mathcal C;Z)\ \le\ \tfrac12\KyF{G}{M-1}.
\]
Taking the maximum over free right congruences of index at most $M$ gives the
claim.
\end{proof}

\begin{remark}[Safe Discrete Objective Is Not Captured by the Block-Local Surrogate]
\label{rem:safe-discrete-not-captured}
The block-local safe linear value
\[
\frac12
\KyF{\EA G}{M-1}
\]
is not generally equal to, an upper bound on, or a lower bound on
$\mathrm{Safe}_{\mathrm{quad}}(M)$.  Safe discrete partitions may distinguish
whole safety blocks, and the corresponding contrast vectors need not be
block-local.  Conversely, block-local linear summaries need not correspond to
any hard partition or right congruence.

For example, take two states with weights $1/2,1/2$, scalar features
$y_1=1$, $y_2=-1$, and singleton safety blocks.  Then
\[
G=
\begin{pmatrix}
1/2 & -1/2\\
-1/2 & 1/2
\end{pmatrix},
\qquad
\EA G=
\begin{pmatrix}
1/2 & 0\\
0 & 1/2
\end{pmatrix}.
\]
For $M=2$, the free linear value is $1/2$, while the block-local safe linear
value is $1/4$, so
\[
\PoSlin(2)=1/4.
\]
But the discrete quadratic free and safe partitions both separate the two
states, so
\[
\mathrm{Free}_{\mathrm{quad}}(2)
=
\mathrm{Safe}_{\mathrm{quad}}(2)
=
1/2,
\]
and
\[
\PoSquad(2)=0.
\]
Thus the block-local surrogate is not an exact formula for the discrete Price
of Safety.
\end{remark}

\begin{remark}[Relaxation-Gap Identity]
\label{rem:relaxation-gap}
Suppose one defines some linear safe relaxation
$\mathrm{Safe}_{\mathrm{lin}}(M)$ and knows relaxation gaps
\[
\rho_{\mathrm{free}}(M)
=
\frac12
\KyF{G}{M-1}
-
\mathrm{Free}_{\mathrm{quad}}(M),
\]
and
\[
\rho_{\mathrm{safe}}(M)
=
\mathrm{Safe}_{\mathrm{lin}}(M)
-
\mathrm{Safe}_{\mathrm{quad}}(M),
\]
with the latter well-defined as a genuine relaxation gap.  Then
\[
\PoSquad(M)
=
\left[
\frac12
\KyF{G}{M-1}
-
\mathrm{Safe}_{\mathrm{lin}}(M)
\right]
+
\rho_{\mathrm{safe}}(M)
-
\rho_{\mathrm{free}}(M).
\]
Thus the discrete Price of Safety equals the linear surrogate plus the
difference of safe and free relaxation gaps.  Without bounds on those gaps, the
surrogate does not control the discrete Price of Safety.
\end{remark}

%----------------------------------------------------------------------
\subsection{Scope and Conditions}
\label{subsec:pos-scope}
%----------------------------------------------------------------------

\begin{itemize}[leftmargin=1.6em]
\item \textbf{Exact domain of the theorem.}
The identities and inequalities above are exact for the linear block-local
second-order surrogate.  They are not exact statements about the discrete
right-congruence Price of Safety.

\item \textbf{Full KL / information-bottleneck objective.}
For the full KL objective, the surrogate is a second-order local model.

\item \textbf{Right-congruence constraint.}
The linear surrogate drops the right-congruence/lumpability constraint.

\item \textbf{State-indexed covariance is canonical.}
The state-indexed Gram matrix $G$ is the object for safety pinching.  The
feature covariance $\Sigma_\pi$ has the same nonzero spectrum and may be used
for the free term, but pinching $\Sigma_\pi$ by feature-coordinate blocks is
not invariant under orthogonal reparameterization and is not canonically
induced by a safety partition of states.
\end{itemize}

%======================================================================
\section{The Aggregation Exponent Vertex Correspondence}
\label{sec:exponent}
%======================================================================

This section states the relation between R\'enyi order, Schatten exponent, and
aggregation exponent.  The claim is a vertex correspondence, not a
representation theorem identifying the manuscript's costs with raw R\'enyi
vertices.

%----------------------------------------------------------------------
\subsection{Layers and Scope}
\label{subsec:exponent-root}
%----------------------------------------------------------------------

There are three distinct exponent-like layers:
\begin{enumerate}[label=(\arabic*)]
\item \textbf{R\'enyi order $\alpha$}, which indexes a family of
divergences;
\item \textbf{Schatten exponent $p$}, which indexes operator-norm tails used
in spectral converses;
\item \textbf{Aggregation exponent or vertex label $e$}, which informally
distinguishes counting/support, averaging/entropy, and
worst-case/supremum behavior.
\end{enumerate}

The sandwiched R\'enyi identity shows that a R\'enyi divergence of order
$\alpha$ can be written as a logarithmic transform of a Schatten-$\alpha$ norm
(a quasi-norm when $0<\alpha<1$)
of an auxiliary $\alpha$-dependent operator.  That is a genuine algebraic fact.
But it does not by itself prove that the Boolean commitment cost is $D_0$, nor
that the operator-norm grounding cost is $D_\infty$, nor that the converse
exponent $p$ is derived from the R\'enyi order $\alpha$.

The correspondence is as follows:
\begin{itemize}[leftmargin=1.6em]
\item retention is exactly the $\alpha=1$ R\'enyi/KL vertex;
\item commitment is recovered exactly after symmetrizing the $\alpha=0$
support test and applying a Boolean valuation;
\item grounding is organized by the $\alpha=\infty$ worst-case/supremum
vertex, but the exact grounding cost is the $C^*$- or operator-norm
distance, not the max-divergence.
\end{itemize}

%----------------------------------------------------------------------
\subsection{R\'enyi Divergences and the Schatten Representation}
\label{subsec:renyi-schatten}
%----------------------------------------------------------------------

Work first on finite alphabets or finite-dimensional Hilbert spaces.  Domain
conditions are stated explicitly.

\begin{definition}[Classical R\'enyi Divergence]
\label{def:classical-renyi}
For probability vectors $p,q$ on a finite alphabet and
$\alpha\in(0,1)\cup(1,\infty)$,
\[
\Ren{\alpha}(p\Vert q)
=
\frac{1}{\alpha-1}
\log
\sum_i p_i^\alpha q_i^{1-\alpha}.
\]
The vertices are defined by limits:
\[
\Ren{1}(p\Vert q)=\KL(p\Vert q),
\]
\[
\Ren{0}(p\Vert q)
=
-\log q(\supp p),
\]
and
\[
\Ren{\infty}(p\Vert q)
=
\log\max_i\frac{p_i}{q_i},
\]
with the usual extended-value conventions.
\end{definition}

\begin{definition}[Sandwiched R\'enyi Divergence]
\label{def:sandwiched-renyi}
Let $\rho,\sigma\succeq0$ with $\Tr\rho=\Tr\sigma=1$.  For
$\alpha\in(0,1)\cup(1,\infty)$, define
\[
\sRen{\alpha}(\rho\Vert\sigma)
=
\frac{1}{\alpha-1}
\log
\Tr\!\left[
\left(
\sigma^{\frac{1-\alpha}{2\alpha}}
\rho
\sigma^{\frac{1-\alpha}{2\alpha}}
\right)^\alpha
\right],
\]
where for $\alpha>1$ one assumes $\supp\rho\subseteq\supp\sigma$ for
finiteness, and where generalized inverses or support projections are used
when $\sigma$ is not full rank.
\end{definition}

\begin{remark}[Commuting States]
\label{rem:renyi-commuting}
If $\rho$ and $\sigma$ commute they are simultaneously diagonalizable, the
conjugation by $\sigma^{(1-\alpha)/2\alpha}$ acts diagonally, and
$\sRen{\alpha}$ reduces to the classical R\'enyi divergence of the two
eigenvalue distributions.  The noncommutative content of
Definition~\ref{def:sandwiched-renyi} is therefore confined to the
noncommuting case, and every classical statement in this section is the
commuting specialization.
\end{remark}

\begin{proposition}[Schatten-Norm Representation]
\label{prop:schatten-rep}
For $\alpha\neq1$, let
\[
K_\alpha
=
\sigma^{\frac{1-\alpha}{2\alpha}}
\rho
\sigma^{\frac{1-\alpha}{2\alpha}}
\succeq0.
\]
Then
\[
\sRen{\alpha}(\rho\Vert\sigma)
=
\frac{\alpha}{\alpha-1}
\log
\norm{K_\alpha}_{\Sp{\alpha}},
\]
where
\[
\norm{K_\alpha}_{\Sp{\alpha}}
=
\left(\Tr |K_\alpha|^\alpha\right)^{1/\alpha}
\]
is the Schatten-$\alpha$ norm for $\alpha\ge1$, and the Schatten-$\alpha$
\emph{quasi-norm} for $0<\alpha<1$, where the triangle inequality fails and
only a quasi-triangle inequality holds.  The identity above is an algebraic
identity valid in both ranges; no norm property is used.
\end{proposition}

\begin{proof}
Since $K_\alpha\succeq0$,
\[
\Tr K_\alpha^\alpha
=
\norm{K_\alpha}_{\Sp{\alpha}}^\alpha.
\]
Taking logarithms and multiplying by $1/(\alpha-1)$ gives
\[
\frac{1}{\alpha-1}\log\Tr K_\alpha^\alpha
=
\frac{\alpha}{\alpha-1}
\log
\norm{K_\alpha}_{\Sp{\alpha}}.
\]
\end{proof}

\begin{remark}[Scope of the Schatten-Norm Representation]
\label{rem:schatten-rep-scope}
The sandwiched R\'enyi divergence is not a norm.  It is not homogeneous and not
a distance.  The operator $K_\alpha$ depends on $\alpha$, $\rho$, and
$\sigma$, and is not the response-operator difference used in the Schatten
converse template.  The proposition shows that the R\'enyi order and a Schatten
exponent appear in the same algebraic expression.  It does not by itself prove
that the R\'enyi order explains the converse exponent.
\end{remark}

\begin{proposition}[Standard Vertex Limits]
\label{prop:renyi-limits}
The following limits hold under the usual support and positivity conventions.
\begin{enumerate}[label=(\roman*)]
\item As $\alpha\to1$,
\[
\sRen{\alpha}(\rho\Vert\sigma)
\to
\Tr\rho(\log\rho-\log\sigma),
\]
the Umegaki relative entropy.  On commuting states this is
$\KL(p\Vert q)$.

\item As $\alpha\to0$ on the classical face,
\[
\Ren{0}(p\Vert q)
=
-\log q(\supp p).
\]
This tests one-sided support inclusion:
\[
\Ren{0}(p\Vert q)=0
\iff
q(\supp p)=1
\iff
\supp q\subseteq\supp p
\]
up to null sets.

\item As $\alpha\to\infty$,
\[
\sRen{\alpha}(\rho\Vert\sigma)
\to
\Dmax(\rho\Vert\sigma)
=
\log
\left\|
\sigma^{-1/2}\rho\sigma^{-1/2}
\right\|_\infty,
\]
the max-relative divergence.  On commuting states this is
\[
\Dmax(p\Vert q)
=
\log\max_i\frac{p_i}{q_i}.
\]
\end{enumerate}
\end{proposition}

\begin{proof}
For $\alpha\to1$, write
$D_\alpha(p\Vert q)=\frac{1}{\alpha-1}\log Z(\alpha)$ with
$Z(\alpha)=\sum_ip_i^\alpha q_i^{1-\alpha}$.  Since $Z(1)=1$, both numerator
and denominator vanish at $\alpha=1$, and
\[
Z'(\alpha)=\sum_ip_i^\alpha q_i^{1-\alpha}\log\frac{p_i}{q_i},
\qquad
Z'(1)=\sum_ip_i\log\frac{p_i}{q_i}=\KL(p\Vert q).
\]
L'H\^opital's rule therefore gives
$D_\alpha\to Z'(1)/Z(1)=\KL(p\Vert q)$.  The corresponding statement for the
sandwiched quantum family is Umegaki's relative entropy limit, for which we
refer to the literature rather than reproduce the argument.

For $\alpha\to\infty$, the Schatten quasi-norms of a fixed positive operator
converge to the operator norm,
$\norm{K_\alpha}_{\Sp{\alpha}}\to\norm{\sigma^{-1/2}\rho\sigma^{-1/2}}_\infty$,
by monotone convergence of $\ell^\alpha$ norms of the singular-value sequence,
and $\alpha/(\alpha-1)\to1$; combining the two gives $\Dmax$.

For $\alpha\to0$ in the classical case, $p_i^\alpha\to\mathbf 1\{p_i>0\}$, so
$Z(\alpha)\to q(\supp p)$ and
$\frac{1}{\alpha-1}\log Z(\alpha)\to-\log q(\supp p)=D_0(p\Vert q)$.

The operator-face limit as $\alpha\to0$ requires separate treatment because
the relevant supports need not be nested; it is not used in what follows.
\end{proof}

%----------------------------------------------------------------------
\subsection{Regime Vertices}
\label{subsec:exponent-regimes}
%----------------------------------------------------------------------

\subsubsection{Retention: Exact $\alpha=1$ Vertex}
\label{subsubsec:exponent-retention}

\begin{theorem}[Retention Is the $\alpha=1$ Vertex]
\label{thm:retention-alpha1}
The retention cost is KL divergence between predictive distributions.  Hence
on the classical measure face it is exactly
\[
\text{retention cost}
=
\Ren{1}(p\Vert q)
=
\KL(p\Vert q).
\]
On the noncommutative face, under the usual support conditions, it is the
Umegaki relative entropy
\[
\Tr\rho(\log\rho-\log\sigma).
\]
\end{theorem}

\begin{remark}[Cost versus Converse]
\label{rem:retention-cost-converse}
The retention cost is KL / relative entropy.  The retention spectral converse
in the Schatten template uses a trace-norm / Schatten-$1$ tail:
\[
\RetQuad(M)
\ge
\frac12
\sum_{i\ge M}
\lambda_i(\Sigma_\pi)
=
\frac12
\Phi^{(1)}_{M-1}(\Sigma_\pi),
\]
with the effective rank $M-1$.  The trace norm appears in the converse, not as
the retention cost itself.
\end{remark}

\subsubsection{Commitment: Booleanized Symmetrized $\alpha=0$ Vertex}
\label{subsubsec:exponent-commitment}

The raw order-$0$ R\'enyi divergence is not the commitment cost.  It is
one-sided and not Boolean-valued.  The construction is a symmetrized support
test followed by a Boolean valuation.

\begin{definition}[Symmetrized Support Mismatch]
\label{def:sym-support}
For probability vectors $p,q$, define
\[
S_0(p,q)
=
\Ren{0}(p\Vert q)+\Ren{0}(q\Vert p),
\]
with extended value $+\infty$ when either support overlap has zero mass.
Define the Boolean valuation
\[
v(t)
=
\begin{cases}
0,& t=0,\\
1,& t>0,
\end{cases}
\]
or, in extended convention, $v(t)=\infty$ for $t>0$.  The Booleanized support
cost is
\[
B_0(p,q)
=
v(S_0(p,q)).
\]
\end{definition}

\begin{proposition}[Boolean Support Equality]
\label{prop:B0}
For probability vectors $p,q$,
\[
B_0(p,q)=0
\iff
\supp p=\supp q
\]
up to null sets.
\end{proposition}

\begin{proof}
Since
\[
\Ren{0}(p\Vert q)=0
\iff
q(\supp p)=1
\iff
\supp q\subseteq\supp p,
\]
we have
\[
S_0(p,q)=0
\iff
\supp q\subseteq\supp p
\text{ and }
\supp p\subseteq\supp q
\iff
\supp p=\supp q.
\]
Applying $v$ preserves the zero set.
\end{proof}

\begin{theorem}[Commitment as Booleanized $\alpha=0$ Vertex]
\label{thm:commitment-alpha0}
Let the commitment residual rows be Boolean functions
\[
r_x:\Sigma^*\to\{0,1\}.
\]
Define their supports by
\[
\supp r_x=\{w\in\Sigma^*:r_x(w)=1\}.
\]
The Booleanized support cost
\[
B_0(r_x,r_y)
=
v(\mathbf 1_{\{\supp r_x\neq\supp r_y\}})
\]
satisfies
\[
B_0(r_x,r_y)=0
\iff
r_x=r_y.
\]
Equivalently, when Boolean rows are represented as normalized indicator
vectors on a finite continuation set, this is the Booleanization of the
symmetrized order-$0$ R\'enyi support test.

Consequently, if $\kappaMN$ is the Myhill--Nerode index, the commitment gap
satisfies
\[
\Com(M)=0
\iff
M\ge\kappaMN,
\]
and $\Com(M)=1$ otherwise under the $0/1$ convention, or $\infty$ under the
extended Boolean convention.
\end{theorem}

\begin{proof}
For Boolean rows, the support is exactly the set of continuations accepted by
the residual language.  Equality of Boolean functions is equality of their
supports.  Proposition~\ref{prop:B0} gives the support-equality test in the
probabilistic encoding, and the direct Boolean definition gives the same zero
set without normalization.  The Myhill--Nerode theorem identifies the number of
distinct residuals with the minimal deterministic automaton size.  Hence a
budget-$M$ machine realizes the language exactly iff $M\ge\kappaMN$.
\end{proof}

\begin{remark}[Scope of the Boolean Vertex]
\label{rem:exponent-alpha0-scope}
The scope restrictions of Remark~\ref{rem:commitment-alpha0-scope} apply
verbatim at this vertex: raw one-sided $D_0$ supports inclusion rather than
equality and is not Boolean-valued, so the construction uses symmetrization and
Boolean valuation, and the operative object is Boolean residual equality, that
is the Myhill--Nerode index, not ordinary real or complex rank.
\end{remark}

\subsubsection{Grounding: Supremum Vertex and Operator-Norm Distance}
\label{subsubsec:exponent-grounding}

The $\alpha=\infty$ R\'enyi vertex is max-divergence.  The grounding cost is
an operator-norm / $C^*$-norm distance.  These are not the same functional.

\begin{definition}[Symmetrized Max-Divergence]
\label{def:dmax-exponent}
This is Definition~\ref{def:dmax}, repeated here so that the exponent layer is
self-contained; the hypotheses are the same, namely finite-dimensional positive
definite states.  For such $\rho,\sigma$, define
\[
\Dmax(\rho\Vert\sigma)
=
\log
\left\|
\sigma^{-1/2}\rho\sigma^{-1/2}
\right\|_\infty,
\]
with $+\infty$ if $\supp\rho\not\subseteq\supp\sigma$.  Define the
symmetrized version
\[
\Dmaxsym(\rho\Vert\sigma)
=
\max\{
\Dmax(\rho\Vert\sigma),
\Dmax(\sigma\Vert\rho)
\}.
\]
\end{definition}

The local equivalence with operator-norm distance was proved in
Proposition~\ref{prop:grounding-local-equivalence}.

\begin{theorem}[Grounding Vertex]
\label{thm:grounding-alpha-infty}
This is Theorem~\ref{thm:grounding-vertex}, restated at the exponent vertex.
The grounding cost is the $C^*$-norm / Schatten-$\infty$ distance
\[
\text{grounding cost}
=
\norm{a-b}_\infty
\]
between residual operators or kernels, while the $\alpha=\infty$ R\'enyi vertex
is the max-divergence $\Dmax$, or its symmetrization $\Dmaxsym$.  These are
distinct functionals:
\[
\Dmax(p\Vert q)
=
\log\max_i\frac{p_i}{q_i},
\qquad
\norm{\operatorname{diag}(p)-\operatorname{diag}(q)}_\infty
=
\max_i|p_i-q_i|.
\]
They are related by Proposition~\ref{prop:grounding-local-equivalence} under
bounded-overlap hypotheses, but they are not identical.
\end{theorem}

\begin{remark}[Two Possible Grounding Interpretations]
\label{rem:grounding-interpretations}
One may either:
\begin{enumerate}[label=(\arabic*)]
\item keep the grounding regime as the operator-norm distance regime, and
regard $\Dmax$ as the divergence-face analogue of its supremum
aggregation; or
\item define a separate max-divergence grounding regime, distinct from the
$C^*$-norm grounding regime.
\end{enumerate}
The manuscript uses interpretation (1).  The exact spectral converse is the
unrestricted linear finite-rank Hankel relaxation gap
\[
\Dunres(M)=\sigma_{M+1}(H_\nu),
\]
while the Hankel-structured finite-rank realization gap satisfies
\[
\DHankstr(M)\ge\sigma_{M+1}(H_\nu),
\]
with equality only under additional structural hypotheses.
\end{remark}

%----------------------------------------------------------------------
\subsection{Data Processing and Monotonicity}
\label{subsec:exponent-dpi}
%----------------------------------------------------------------------

\begin{proposition}[Data Processing, Scoped]
\label{prop:DPI-scoped}
The following data-processing statements are classical and are quoted without
proof; see \cite{cover} for the classical face and \cite{tsybakov2009} for the
divergence-theoretic background.
\begin{enumerate}[label=(\alph*)]
\item \textbf{Classical measure face.}
For finite alphabets and stochastic maps $\Phi$, classical R\'enyi divergence
satisfies
\[
\Ren{\alpha}(\Phi p\Vert \Phi q)
\le
\Ren{\alpha}(p\Vert q)
\]
in the standard ranges, in particular for $\alpha=1$ and $\alpha=\infty$, and
for $\alpha=0$ under the usual support conventions.

\item \textbf{Sandwiched operator face.}
For completely positive trace-preserving maps $\Phi$, the sandwiched R\'enyi
divergence satisfies
\[
\sRen{\alpha}(\Phi\rho\Vert\Phi\sigma)
\le
\sRen{\alpha}(\rho\Vert\sigma)
\]
for $\alpha\in[1/2,\infty]$, under the usual support/domain conditions.
\end{enumerate}
\end{proposition}

\begin{remark}[Monotonicity Meta-Theorem Is Not Exactly DPI]
\label{rem:monotonicity-not-dpi}
The monotonicity meta-theorem says that the gap is nonincreasing in the state
budget $M$.  The primary reason is nested feasible sets:
\[
\{\text{machines of index }\le M\}
\subseteq
\{\text{machines of index }\le M+1\}.
\]
Therefore the optimized deficit cannot increase.  In the retention regime there
is an additional information-processing interpretation: if a coarser machine
label is a deterministic function of a finer one, then
\[
I(K_{\mathrm{fine}};Z)
\ge
I(K_{\mathrm{coarse}};Z).
\]
Thus retention monotonicity can be viewed as DPI for mutual information / KL.
But the full meta-theorem across all regimes is not literally the DPI of the
R\'enyi family.
\end{remark}

%----------------------------------------------------------------------
\subsection{Relation to the Schatten Converse Exponent}
\label{subsec:exponent-schatten-relation}
%----------------------------------------------------------------------

The Schatten template assigns analytic converses to Schatten exponents:
\[
\text{grounding: }p=\infty,
\qquad
\text{retention: }p=1,
\qquad
\text{commitment: }p=0,
\]
with the following qualifications:
\begin{itemize}[leftmargin=1.6em]
\item grounding is the unrestricted linear finite-rank relaxation, for which
\[
\Dunres(M)=\sigma_{M+1}(H_\nu),
\]
while the Hankel-structured finite-rank realization gap satisfies
\[
\DHankstr(M)\ge\sigma_{M+1}(H_\nu);
\]
\item retention uses the effective rank $M-1$;
\item commitment is a Boolean-rank / Myhill--Nerode threshold, not an
ordinary Schatten-rank limit.
\end{itemize}

The present section assigns cost vertices as follows:
\[
\text{grounding: }\alpha=\infty,
\qquad
\text{retention: }\alpha=1,
\qquad
\text{commitment: }\alpha=0,
\]
again with qualifications: commitment is Booleanized symmetrized $D_0$, and
grounding's exact cost is operator-norm distance rather than max-divergence.

\begin{theorem}[Formal Vertex Consistency, Not Derivation]
\label{thm:vertex-consistency}
The labels $\alpha$ and $p$ agree at the three vertices:
\[
\alpha=p\in\{0,1,\infty\}.
\]
This is a formal consistency between the cost-face vertex correspondence and
the Schatten converse correspondence.  It is not a derivation of the Schatten
converse theorem from the R\'enyi identity.
\end{theorem}

\begin{proof}
The R\'enyi cost-face analysis gives the vertex labels
$\alpha\in\{0,1,\infty\}$.  The Schatten template gives the converse-face
labels $p\in\{0,1,\infty\}$.  The labels match vertex by vertex.  However, the
sandwiched R\'enyi identity uses an auxiliary $\alpha$-dependent operator
$K_\alpha$, whereas the Mirsky converse uses a regime-specific response
operator difference $A_\delta-A_{\widehat\delta}$.  No implication from one to
the other is established without additional domination or achievability
hypotheses.
\end{proof}

\begin{remark}[What the Matching Labels Do Explain]
\label{rem:vertex-two-ingredients}
Although the correspondence is not a derivation, it is not a coincidence
either, and Definition~\ref{def:response} already isolates the reason.  A
Schatten-$p$ response representation consists of two independent data: a
\emph{rank} condition, clause~(a), asserting
$\rank(A_{\widehat\delta})\le r_{\mathbb T}(M)$, and a \emph{norm} condition,
clause~(b), asserting that the task cost dominates
$\norm{A_\delta-A_{\widehat\delta}}_{\Sp p}$.  Theorem~\ref{thm:unified}
shows that once both are supplied the value is forced, by
Eckart--Young--Mirsky, to be the Schatten-$p$ tail
$\Phi^{(p)}_{r_{\mathbb T}(M)}(A_\delta)$.

The three converses therefore have parallel form because they share the
optimization ingredient~(a) and differ only in the norm selected by~(b):
grounding's cost \emph{is} an operator norm, giving $p=\infty$; retention's
cost dominates a trace-norm tail of the covariance residual, giving $p=1$; and
commitment's threshold is Boolean rank, the $p\to0$ vertex.  The equality
$\alpha=p$ records that the R\'enyi order selecting the divergence and the
Schatten exponent selecting the tail agree at the three vertices.  What remains
unavailable, and what clause~(e) of
Meta-Theorem~\ref{meta:typed-rate-distortion} and
Theorem~\ref{thm:schatten-nogo} record, is any transfer between regimes: the
operators $H_\nu$, $\Sigma_\pi$, $N_\Lambda$, the rank budgets $M$, $M-1$, $M$
and the moduli are unrelated, so the shared form yields no shared value.
\end{remark}

%----------------------------------------------------------------------
\subsection{Scope and Conditions}
\label{subsec:exponent-scope}
%----------------------------------------------------------------------

\begin{itemize}[leftmargin=1.6em]
\item \textbf{No cross-regime gap inequality.}
The vertex correspondence unifies functional organization at isolated orders.
It yields no inequality relating the numerical grounding, retention, and
commitment gaps.  The Independence Theorem remains intact.

\item \textbf{Scope of the R\'enyi family.}
The operational regimes identified here occupy the vertices
$\alpha\in\{0,1,\infty\}$.  Intermediate orders
$\alpha\in(0,\infty)\setminus\{1\}$ serve as interpolation inequalities and
diagnostics; promoting one to a finite-state approximation regime requires a
task theory and a domination bridge supplied for that order.

\item \textbf{Domain regularity.}
For $\alpha>1$, sandwiched R\'enyi divergence requires support conditions,
typically $\supp\rho\subseteq\supp\sigma$, or generalized inverses.  The
operator-face $\alpha\to0$ limit is subtle and is not used to define
commitment.

\item \textbf{Inherited conditionality.}
The alignment with the Schatten converse exponent is conditional on the
Schatten template.  If the Schatten correspondence is a conditional template
rather than an unconditional theorem, then the present vertex consistency is
likewise conditional.

\item \textbf{Conditional status.}
This section is a vertex correspondence.  Its validity depends on the
hypotheses of the preceding sections.
\end{itemize}

%======================================================================
\section{Joint, Comparative, Interaction, and Independence Results}
\label{sec:joint}
%======================================================================

The regime-specific gaps can be composed, compared, and constrained jointly.
This section develops the joint zero-error memory theorem, the comparative
state-complexity sandwich, the independence theorem, and the interaction
between exact commitment and approximate retention.

%----------------------------------------------------------------------
\subsection{Joint Zero-Error Memory}
\label{subsec:joint-zero-error}
%----------------------------------------------------------------------

Suppose task specifications $T_1,\dots,T_m$ are posed on the same history
space and induce Nerode right congruences
\[
\sim_1,\dots,\sim_m.
\]
A joint zero-error approximant must be exact for every component task.

Define the meet of the congruences by intersection:
\[
\bigwedge_{i=1}^m \sim_i
\eqdef
\bigcap_{i=1}^m \sim_i.
\]
This is the finest right congruence refining every $\sim_i$.

\begin{theorem}[Exact Joint Zero-Error Memory]
\label{thm:joint}
Let task specifications $T_1,\dots,T_m$ induce right congruences
$\sim_1,\dots,\sim_m$.  Then
\[
\Delta_{T_1,\dots,T_m}(M)=0
\iff
M\ge
\operatorname{index}
\left(
\bigwedge_{i=1}^m\sim_i
\right).
\]
Moreover,
\[
\operatorname{index}
\left(
\bigwedge_i\sim_i
\right)
\ge
\max_i\operatorname{index}(\sim_i),
\]
and in the worst case equals the product
\[
\prod_i\operatorname{index}(\sim_i).
\]
\end{theorem}

\begin{proof}
A joint zero-error approximant must assign the same state only to histories
that are equivalent for every component task.  Therefore its congruence must
refine each $\sim_i$, equivalently it must refine the meet
$\bigwedge_i\sim_i$.  Hence zero error is possible only if
\[
M\ge
\operatorname{index}
\left(
\bigwedge_i\sim_i
\right).
\]
Conversely, quotienting by the meet gives a deterministic finite-state machine
that refines every component Nerode congruence and therefore realizes every
component task exactly.

The inequality
\[
\operatorname{index}
\left(
\bigwedge_i\sim_i
\right)
\ge
\max_i\operatorname{index}(\sim_i)
\]
is immediate because the meet refines each $\sim_i$.  The map
\[
[u]_{\bigwedge_i\sim_i}
\mapsto
\bigl(
[u]_{\sim_1},\dots,[u]_{\sim_m}
\bigr)
\]
is injective, so
\[
\operatorname{index}
\left(
\bigwedge_i\sim_i
\right)
\le
\prod_i\operatorname{index}(\sim_i).
\]
The upper bound is tight when the component congruences are independent, for
example independent modular counters with relatively prime moduli.
\end{proof}

%----------------------------------------------------------------------
\subsection{Comparative State Complexity}
\label{subsec:comparative}
%----------------------------------------------------------------------

For each regime or component task $i$, define the approximate state complexity
\[
M_i(\varepsilon_i)
=
\min
\left\{
\operatorname{index}(\sim):
\Delta_{\mathbb T_i}(\sim)\le\varepsilon_i
\right\}.
\]
For a vector of distortion budgets
\[
\vec\varepsilon=(\varepsilon_1,\dots,\varepsilon_m),
\]
define the joint state complexity
\[
M_{\mathrm{joint}}(\vec\varepsilon)
=
\min
\left\{
\operatorname{index}(\sim):
\Delta_{\mathbb T_i}(\sim)\le\varepsilon_i
\text{ for all }i
\right\}.
\]

The following product bound concerns only component tasks whose admissible
approximants are deterministic right-congruence quotients of one common
history system and whose representatives can be combined componentwise.  It
does \emph{not} apply to the unrestricted linear finite-rank Hankel
relaxation of Section~\ref{sec:grounding}, whose resource is operator rank
rather than quotient index, and for which no representative map of the
required componentwise-combinable form is available; see
Remark~\ref{rem:comparative-grounding-excluded} below.  It applies directly
to commitment and retention, and to symbolic grounding once the
implementability condition of
Proposition~\ref{prop:block-mode-implementability} is imposed so that
symbolic-grounding approximants are themselves right-congruence quotients.

\begin{theorem}[Comparative State Complexity]
\label{thm:comparative}
Assume every component task $i=1,\dots,m$ has admissible approximants given
by right-congruence quotients of a common history system, that each
$M_i(\varepsilon_i)$ is attained by some congruence $\sim_i$ and
representative map $\Psi_i$, and that the product representative remains
admissible with componentwise losses unchanged, namely
\[
\Psi([u]_\sim)
=
\bigl(
\Psi_1([u]_{\sim_1}),\dots,\Psi_m([u]_{\sim_m})
\bigr)
\]
is itself an admissible approximant of loss
$\Delta_{\mathbb T_i}(\sim_i)$ in each coordinate $i$.  Then for every
$\vec\varepsilon=(\varepsilon_1,\dots,\varepsilon_m)$,
\[
\max_i M_i(\varepsilon_i)
\le
M_{\mathrm{joint}}(\vec\varepsilon)
\le
\prod_i M_i(\varepsilon_i).
\]
Without attainment the upper bound holds with each $\varepsilon_i$ replaced
by $\varepsilon_i+\eta_i$, for arbitrary $\eta_i>0$.
\end{theorem}

\begin{remark}[Why the Unrestricted Linear Relaxation Is Excluded]
\label{rem:comparative-grounding-excluded}
The proof of Theorem~\ref{thm:comparative} builds the joint approximant as
the quotient by the meet $\bigwedge_i\sim_i$ of the component
right-congruences, and controls the index of that meet by the product of
component indices, as in Theorem~\ref{thm:joint}.  The unrestricted linear
finite-rank grounding gap $\Dunres(M)$ of
Theorem~\ref{thm:spectral-grounding} is a statement about compact-operator
rank in $\ell^2(\Sigma^*)$, not about a right congruence on input histories;
``meeting'' two rank-$M$ conditions has no analogue of
Theorem~\ref{thm:joint}'s exact index formula, and no product construction
combining two finite-rank Hankel approximants into one of controlled rank is
supplied here.  Consequently Theorem~\ref{thm:comparative} and
Theorem~\ref{thm:independence} say nothing about joint budgets that mix an
unrestricted linear grounding coordinate with a commitment or retention
coordinate, absent a bridge of the kind discussed in
Section~\ref{sec:oracle}.
\end{remark}

\begin{proof}
The lower bound is immediate because a joint approximant is, in particular, an
individual approximant for each regime.

For the upper bound, choose congruences $\sim_i$ of index
$M_i(\varepsilon_i)$ attaining (or, absent attainment, $\eta_i$-approximating)
the individual optima.  Let
\[
u\sim v
\iff
\bigwedge_i u\sim_i v.
\]
Then
\[
\operatorname{index}(\sim)
\le
\prod_i\operatorname{index}(\sim_i)
=
\prod_i M_i(\varepsilon_i).
\]
The quotient by $\sim$ maps into the product of the component quotients.  If
$\Psi_i$ is the optimal representative map for component $i$, define the
product representative
\[
\Psi([u]_\sim)
=
\bigl(
\Psi_1([u]_{\sim_1}),
\dots,
\Psi_m([u]_{\sim_m})
\bigr).
\]
This realizes all componentwise approximations simultaneously.
\end{proof}

\begin{remark}[Tightness]
\label{rem:comparative-tightness}
The lower and upper bounds can both be tight.  The lower bound is tight when a
single congruence simultaneously optimizes all components.  The upper bound is
tight, up to constants, when the component congruences are independent and the
joint state must encode the tuple of component states.
\end{remark}

\begin{remark}[Algorithmic Consequence]
\label{rem:comparative-algorithmic}
The product construction gives a universal joint approximant, but it may incur
exponential blowup.  Optimizing the joint congruence directly can be
computationally hard, as shown by the interaction complexity theorem below.
\end{remark}

%----------------------------------------------------------------------
\subsection{Independence of Regime Obstructions}
\label{subsec:independence}
%----------------------------------------------------------------------

The comparative sandwich does not impose a functional relation among the
individual regime gaps.  The following independence theorem shows that the
regime obstructions can be varied independently.

\begin{theorem}[Independence of Regime Obstructions]
\label{thm:independence}
There exist product sequential tasks with designated retention, commitment,
and grounding coordinates such that each regime's gap or obstruction depends
only on its designated coordinate.  In particular:
\begin{enumerate}[label=(\roman*),leftmargin=2.0em]
\item Retention can be zero while commitment is infinite.
\item Commitment can be zero while retention is positive.
\item Retention and commitment can be zero while deterministic symbolic
grounding has a positive stochasticity floor.
\item Commitment complexity can be made any finite integer $N$ or infinite.
\item Retention gap can be made zero or positive, and can be varied by
varying predictive separation.
\item Deterministic grounding obstruction can be made zero or positive by
varying intrinsic stochasticity.
\end{enumerate}
Consequently, no nontrivial universal inequality or functional dependence
holds among the grounding, retention, and commitment obstructions across all
tasks.
\end{theorem}

\begin{proof}
We construct product tasks whose output stream is a tuple of independent
coordinates.  Each regime is measured on its designated coordinate.

\medskip
\noindent\textbf{Retention zero, commitment infinite.}
Let the retention coordinate be a memoryless fair coin:
\[
Y_t^{\mathrm{ret}}\sim\mathrm{Bernoulli}(1/2),
\]
independent of history.  Its predictive distribution is the same after every
history, so it has one causal state and
\[
\RetKL(1)=0.
\]
Let the commitment coordinate be the square-counting language after a delimiter
$\#$:
\[
L_{\mathrm{sq}}
=
\{
\#1^n : n \text{ is a perfect square}
\}.
\]
For histories $\#1^n$ and $\#1^m$ with $n<m$, choose
\[
t=m^2-n.
\]
Then $n+t=m^2$ is square, while
\[
m+t=m^2+(m-n)
\]
lies strictly between $m^2$ and $(m+1)^2$, because
\[
0<m-n<m<2m+1.
\]
Thus $m+t$ is not square.  Hence the residuals after $\#1^n$ and $\#1^m$ are
distinct for every $n\neq m$, so
\[
\kappaMN(L_{\mathrm{sq}})=\infty.
\]
The product task has zero one-state retention gap but infinite exact
commitment complexity.

\medskip
\noindent\textbf{Commitment zero, retention positive.}
Let the commitment coordinate be the constant safe output, with
\[
\kappaMN=1.
\]
Let the retention coordinate be the running-parity process with fair i.i.d.\
inputs and two causal states, as in Example~\ref{ex:parity}.  Then
\[
\RetKL(2)=0,
\qquad
\RetKL(1)=\log 2.
\]
The product task has exact commitment with one state but positive one-state
retention gap.

\medskip
\noindent\textbf{Retention and commitment zero, deterministic grounding
obstruction positive.}
Let the output be a memoryless fair coin independent of input, and impose no
nontrivial commitment constraint.  The retention coordinate has one causal
state, so retention is zero.  The commitment coordinate is trivial, so
commitment is zero.  However, a deterministic Mealy machine emits a
deterministic output bit at each time.  For a memoryless fair coin, the local
stochastic spread is
\[
s(h_t)=1-\max\{1/2,1/2\}=1/2
\]
at every time.  Therefore the stochasticity floor of
Theorem~\ref{thm:stochasticity} is
\[
\sigma_\gamma
=
\sum_{t=1}^\infty
\gamma^{t-1}\frac12
=
\frac{1}{2(1-\gamma)}
>0.
\]
Thus deterministic symbolic grounding has a positive obstruction even though
retention and commitment vanish.  The associated linear Hankel object has rank
one, illustrating the separation between linear finite-rank grounding and
deterministic symbolic grounding.

\medskip
\noindent\textbf{Finite and tunable complexity.}
Replacing $L_{\mathrm{sq}}$ by the modular language
\[
L_N=\{\#1^n:n\equiv0\pmod N\}
\]
gives commitment complexity exactly $N$.  Varying the predictive separation in
a two-state retention process varies the retention gap continuously from $0$
to its maximum.  Varying the Bernoulli bias $p$ in a memoryless grounding
coordinate varies the stochasticity floor as
\[
\sigma_\gamma(p)
=
\frac{\min\{p,1-p\}}{1-\gamma}.
\]
Thus the coordinates can be tuned independently.
\end{proof}

\begin{remark}[Consequence for Cross-Regime Comparisons]
\label{rem:independence-use}
The independence theorem is the structural reason why the unified Schatten
template and the exponent vertex correspondence cannot produce unconditional
cross-regime inequalities.  They unify functional forms under shared response
representations; they do not unify the numerical gaps.  The product
construction measures each regime on its designated coordinate; it does not
identify the machine types.
\end{remark}

%----------------------------------------------------------------------
\subsection{Interaction of Exact Commitment and Approximate Retention}
\label{subsec:interaction}
%----------------------------------------------------------------------

Let commitment be exact and retention be approximate.  Let
$\sim_{\mathrm{com}}$ be the commitment congruence, with index
\[
\kappa=\operatorname{index}(\sim_{\mathrm{com}}).
\]
Exact commitment requires the machine congruence to refine
$\sim_{\mathrm{com}}$:
\[
\sim\subseteq\sim_{\mathrm{com}}.
\]

\begin{theorem}[Exact Commitment with Approximate Retention]
\label{thm:interaction}
For every $\varepsilon\ge0$,
\[
M_{\mathrm{joint}}(0,\varepsilon)
=
\min
\left\{
\operatorname{index}(\sim):
\sim\subseteq\sim_{\mathrm{com}},
\;
\Delta_{\mathrm{ret}}(\sim)\le\varepsilon
\right\}.
\]
\end{theorem}

\begin{proof}
Exact commitment requires $\sim\subseteq\sim_{\mathrm{com}}$.  The same
congruence $\sim$ governs retention.  The result follows by minimizing the
retention cost over safe congruences.
\end{proof}

\begin{lemma}[Constrained Information-Bottleneck Form]
\label{lem:constrained-ib}
In the synchronized finite-history subclass, suppose
\[
\sim_\delta\subseteq\sim\subseteq\sim_{\mathrm{com}},
\]
so that $\sim$ is coarser than causal equivalence and refines the commitment
congruence.  Then
\[
\Delta_{\mathrm{ret}}^{\mathrm{KL}}(\sim)
=
I(S;Z\mid K_\sim)
=
I(S;Z)-I(K_\sim;Z).
\]
Consequently,
\[
\Delta_{\mathrm{ret}}^{\mathrm{safe}}(M)
=
I(S;Z)
-
\max_{\substack{
\sim\subseteq\sim_{\mathrm{com}}\\
\operatorname{index}(\sim)\le M
}}
I(K_\sim;Z).
\]
In the stationary formulation of
Definition~\ref{def:controlled-markov}, the same identity holds with the
feasible set written as safe lumpable quotients
\[
\phi:\Splus\to\mathcal K
\]
whose induced quotient refines the commitment congruence:
\[
\Delta_{\mathrm{ret}}^{\mathrm{safe}}(M)
=
I(S;Z)
-
\max_{\substack{
\phi\ \text{safe lumpable}\\
|\mathcal K|\le M
}}
I(K_\phi;Z).
\]
\end{lemma}

\begin{proof}
This is Meta-Theorem~\ref{meta:ib}, with the feasible set restricted to safe
congruences $\sim\subseteq\sim_{\mathrm{com}}$ in the synchronized subclass,
or equivalently to safe lumpable quotients of the stationary support
$\Splus$ in the stationary formulation.
\end{proof}

\begin{remark}[Why the Lower Constraint Is Only Sufficient]
\label{rem:no-lower-constraint}
Lemma~\ref{lem:constrained-ib} assumes
\[
\sim_\delta\subseteq\sim\subseteq\sim_{\mathrm{com}},
\]
which guarantees that the machine state is a function of the causal state, so
that $K_\sim$ is a quotient of $S$.  This is a \emph{sufficient} condition for
the information-bottleneck form, not a constraint that may be imposed in
general.

Requiring $\sim_\delta\subseteq\sim$ in the definition of the joint problem
would be incorrect.  A joint machine may need to \emph{split} a causal state
in order to carry commitment information, since exact commitment is a
constraint on histories rather than on predictive laws.  In that case the
machine state is not a function of $S$, and the feasible set of
Theorem~\ref{thm:interaction} is strictly larger than the set of quotients of
$S$.

The correct general formulation therefore optimizes over finite-state history
factors $K$ that satisfy exact commitment.  If $K$ is such a factor and $S$ is
predictively sufficient with $Z\perp K\mid S$, then
\[
\mathbb E\,D_{\mathrm{KL}}(P_{Z\mid S}\|P_{Z\mid K})
=
I(S;Z\mid K)
=
I(S;Z)-I(K;Z),
\]
so the safe retention cost is
\[
\Delta_{\mathrm{ret}}^{\mathrm{safe}}(M)
=
I(S;Z)
-
\sup_{K}
I(K;Z),
\]
the supremum being over at-most-$M$-state history factors whose state
congruence refines the commitment congruence.  When commitment is constant on
each causal state, this reduces to the quotient formulation above.
\end{remark}

\begin{corollary}[Equality Condition]
\label{cor:interaction-equality}
Let
\[
M_{\mathrm{ret}}(\varepsilon)
=
\min
\left\{
\operatorname{index}(\sim):
\Delta_{\mathrm{ret}}(\sim)\le\varepsilon
\right\}.
\]
Then
\[
M_{\mathrm{joint}}(0,\varepsilon)
\ge
\max(\kappa,M_{\mathrm{ret}}(\varepsilon)).
\]
Equality holds iff there exists a safe congruence
$\sim\subseteq\sim_{\mathrm{com}}$ with
\[
\operatorname{index}(\sim)
=
\max(\kappa,M_{\mathrm{ret}}(\varepsilon))
\]
and
\[
\Delta_{\mathrm{ret}}(\sim)\le\varepsilon.
\]
Equivalently:
\begin{enumerate}[label=(\arabic*)]
\item If $M_{\mathrm{ret}}(\varepsilon)\le\kappa$, equality holds iff
\[
\Delta_{\mathrm{ret}}(\sim_{\mathrm{com}})\le\varepsilon.
\]
\item If $M_{\mathrm{ret}}(\varepsilon)>\kappa$, equality holds iff there
exists an optimal retention congruence $\sim^*$ with
\[
\operatorname{index}(\sim^*)=M_{\mathrm{ret}}(\varepsilon),
\]
\[
\Delta_{\mathrm{ret}}(\sim^*)\le\varepsilon,
\]
and
\[
\sim^*\subseteq\sim_{\mathrm{com}}.
\]
\end{enumerate}
\end{corollary}

\begin{proof}
The lower bound follows because any safe congruence has index at least $\kappa$
and any feasible retention congruence has index at least
$M_{\mathrm{ret}}(\varepsilon)$.

If $M_{\mathrm{ret}}(\varepsilon)\le\kappa$, then a safe congruence of index
$\kappa$ must be exactly $\sim_{\mathrm{com}}$, because any refinement of
$\sim_{\mathrm{com}}$ has index at least $\kappa$, and equality of finite index
forces equality of partitions.  Hence equality holds exactly when the
commitment congruence itself is retention-feasible.

If $M_{\mathrm{ret}}(\varepsilon)>\kappa$, then equality requires a retention
optimal congruence of index $M_{\mathrm{ret}}(\varepsilon)$ that is also safe.
\end{proof}

\begin{corollary}[Price of Safety]
\label{cor:price-safety}
Let $K$ range over at-most-$M$-state deterministic history factors whose state
congruence refines the exact commitment congruence, and assume $S$ is
predictively sufficient, so that $Z\perp K\mid S$.  Then
\[
\Delta_{\mathrm{ret}}^{\mathrm{safe}}(M)
=
I(S;Z)
-
\sup_{K}
I(K;Z).
\]
If commitment is constant on every predictive state, the feasible factors may
be represented as safe lumpable quotients $K_\sim$ of $S$, and the supremum
reduces to the quotient form
\[
\Delta_{\mathrm{ret}}^{\mathrm{safe}}(M)
=
I(S;Z)
-
\max_{\substack{
\sim\subseteq\sim_{\mathrm{com}}\\
\operatorname{index}(\sim)\le M
}}
I(K_\sim;Z).
\]
See Remark~\ref{rem:no-lower-constraint}.
The price of safety is
\[
\operatorname{PoS}(M)
=
\Delta_{\mathrm{ret}}^{\mathrm{safe}}(M)
-
\Delta_{\mathrm{ret}}^{\mathrm{free}}(M)
\ge0.
\]
\end{corollary}

\begin{proof}
This follows from Lemma~\ref{lem:constrained-ib}.  The safe feasible set is a
subset of the free feasible set, so the constrained maximum of $I(K_\sim;Z)$
is no larger than the free maximum.  Therefore the safe retention gap is no
smaller than the free retention gap.
\end{proof}

\begin{remark}[No Universal Strict Overhead]
\label{rem:no-universal-overhead}
Overhead beyond
\[
\max(\kappa,M_{\mathrm{ret}}(\varepsilon))
\]
is not structurally forced.  It occurs exactly when no safe congruence of index
\[
\max(\kappa,M_{\mathrm{ret}}(\varepsilon))
\]
satisfies the retention distortion constraint.  In particular, if
\[
M_{\mathrm{ret}}(\varepsilon)\le\kappa,
\]
equality requires
\[
\Delta_{\mathrm{ret}}(\sim_{\mathrm{com}})\le\varepsilon.
\]
If that fails, then
\[
M_{\mathrm{joint}}(0,\varepsilon)>\kappa
\]
even though
\[
M_{\mathrm{ret}}(\varepsilon)<\kappa.
\]
\end{remark}

\begin{theorem}[Complexity of Interaction]
\label{thm:interaction-complexity}
For rational data --- rational predictive means and a rational tolerance
$\varepsilon$ --- the decision problem
\[
M_{\mathrm{joint}}(0,\varepsilon)\le M
\]
is NP-hard, even for trivial safety constraints $\kappa=1$, in the Gaussian
quadratic restricted retention regime.  The qualifier is inherited from
Lemma~\ref{lem:kmeans-gap} and
Theorems~\ref{thm:retention-reset-np}--\ref{thm:retention-definite-np}, on
which the reduction below depends, and without it the threshold
$\varepsilon$ constructed in the proof need not be representable in polynomial
bit complexity.
\end{theorem}

\begin{proof}
For $\kappa=1$, the commitment constraint is vacuous: every congruence refines
the one-class commitment congruence, so $M_{\mathrm{joint}}(0,\varepsilon)\le k$
holds exactly when some congruence of index at most $k$ has retention gap at
most $\varepsilon$.

The joint problem is stated with a non-strict threshold and the retention
problem with a strict one, so the two are not interchangeable at a common
$\theta$; a threshold strictly separating the two cases must be chosen.  The
gap-amplification lemma supplies one.  By Lemma~\ref{lem:kmeans-gap} the
transformed instances satisfy the promise
\[
\mathrm{OPT}_a\le T
\qquad\text{or}\qquad
\mathrm{OPT}_a\ge T+2n ,
\]
and the reductions of Theorems~\ref{thm:retention-reset-np}
and~\ref{thm:retention-definite-np} carry this separation to the retention
objective, giving rational values $\theta_{\mathrm{yes}}<\theta_{\mathrm{no}}$
with $\RetQuad(k)\le\theta_{\mathrm{yes}}$ in the YES case and
$\RetQuad(k)\ge\theta_{\mathrm{no}}$ in the NO case.  Set $M=k$ and
\[
\varepsilon=\tfrac12\bigl(\theta_{\mathrm{yes}}+\theta_{\mathrm{no}}\bigr),
\]
a rational of polynomial bit complexity.  Then
$M_{\mathrm{joint}}(0,\varepsilon)\le k$ holds in the YES case and fails in the
NO case, so the reduction is correct, and it is computable in polynomial time.
NP-hardness of the retention problem in the Gaussian quadratic restricted
regime is Theorems~\ref{thm:retention-reset-np}
and~\ref{thm:retention-definite-np}.
\end{proof}

\begin{remark}[Scope of the Interaction Complexity]
\label{rem:interaction-complexity-scope}
The theorem is a theorem of the Gaussian quadratic restricted retention regime.
It does not establish NP-hardness for unrestricted full-KL retention.  It also
does not assert that the exact Price of Safety is given by the linear pinching
surrogate of Section~\ref{sec:pos-linear}.
\end{remark}

%======================================================================
\section{The Temporal Axis}
\label{sec:temporal}
%======================================================================

The offline approximation gaps and the joint/interaction results are static.
This section adds the temporal axis: online learning of finite-state
predictors.  The theory distinguishes passive/adversarial-input models from
active/query models.  It also separates realizable mistake bounds from
agnostic regret bounds.

The temporal results supply the estimation rates used in the oracle bridge of
Section~\ref{sec:oracle}.

%----------------------------------------------------------------------
\subsection{Passive and Active Models}
\label{subsec:temporal-models}
%----------------------------------------------------------------------

Let $\mathcal H_M$ denote the class of input-output functions realized by
deterministic Mealy machines with at most $M$ states over fixed finite
alphabets $\mathcal I$ and $\mathcal O$, with $|\mathcal I|\ge2$ and
$|\mathcal O|\ge2$.

The temporal theory distinguishes two information structures.

\begin{itemize}[leftmargin=1.8em]
\item \textbf{Passive / adversarial-input model.}
The learner does not choose inputs.  At each round $t$, an adversary chooses
$x_t\in\mathcal I$, the learner predicts an output $\widehat y_t$, and then
the true output $y_t$ is revealed.  The input sequence may be adversarial.

\item \textbf{Active / query model.}
The learner chooses inputs $x_t$ and may apply synchronizing or
distinguishing experiments.  The learner still predicts the output before
observing it, and mistakes are counted in the same way.
\end{itemize}

The two models have different synchronization properties.  In the passive
model, the learner cannot force the system to reveal its state.  In the active
model, the learner may be able to synchronize or identify the initial state by
choosing suitable input words.

%----------------------------------------------------------------------
\subsection{Littlestone Dimension of Finite-State Machines}
\label{subsec:littlestone}
%----------------------------------------------------------------------

\begin{lemma}[Littlestone Dimension]
\label{lem:littlestone}
For fixed alphabets with $|\mathcal I|\ge2$ and $|\mathcal O|\ge2$,
\[
\Ldim(\mathcal H_M)
=
\Theta(M\log M),
\]
with constants depending only on $|\mathcal I|$ and $|\mathcal O|$, not on $M$
or the horizon $T$.

Here $\mathcal H_M$ carries the \emph{reset-word} protocol: each instance is a
fresh word evaluated from the designated initial state.  Write
$\mathcal H_M^{\mathrm{reset}}$ when the distinction matters, and
$\mathcal H_M^{\mathrm{stream}}$ for the same machine family evaluated along a
single persistent stream, as in Corollary~\ref{cor:stream-ldim}.  These are
different online classes and their sequential dimensions are established by
different arguments.
\end{lemma}

\begin{remark}[From Word Classifiers to Transducers]
\label{rem:dfa-to-mealy}
The VC estimate of \cite{ishigami1993} concerns deterministic \emph{acceptors}:
$n$-state automata classifying independently presented words.  Its transfer to
Mealy transductions uses the following reduction.  Given an acceptor
$(\mathcal S,s_0,\delta,F)$ with $|\mathcal S|=n$, adjoin a delimiter
$\#\notin\mathcal I$ and define a Mealy machine on the same state set with
\[
\lambda(s,x)=0\ \ (x\in\mathcal I),
\qquad
\lambda(s,\#)=\mathbf 1_{\{s\in F\}},
\qquad
\tau(s,\#)=s_0,
\]
and $\tau=\delta$ on $\mathcal I$.  Reading $u\#$ from $s_0$ emits the
classification of $u$ on the final symbol and returns to $s_0$, so a shattered
set of words for the acceptor yields a shattered set of instances $u\#$ for the
transducer with the same cardinality and $n$ states.

If the delimiter is not admitted into the alphabet, encode $\mathcal I\cup\{\#\}$
into fixed-length binary blocks; this multiplies the state count by a constant
depending only on $|\mathcal I|$, which is absorbed into the constants of the
lemma.
\end{remark}

\begin{proof}[Proof of Lemma~\ref{lem:littlestone}]
\emph{Upper bound.}  A member of $\mathcal H_M$ is specified by a designated
initial state, a transition table $\tau:Q\times\mathcal I\to Q$, and an output
table $\lambda:Q\times\mathcal I\to\mathcal O$, so
\[
|\mathcal H_M|
\ \le\
M\cdot M^{M|\mathcal I|}\cdot|\mathcal O|^{M|\mathcal I|},
\qquad
\log_2|\mathcal H_M|
\ \le\
\bigl(M|\mathcal I|+1\bigr)\log_2M+M|\mathcal I|\log_2|\mathcal O| ,
\]
which is $O(M\log M)$ for fixed alphabets.  The inequality is not an equality,
since distinct labelled tables may realize the same transduction; only the
upper bound is needed.  The halving algorithm predicts by
majority over the version space and at least halves it on every mistake, so no
shattered mistake tree has depth exceeding $\log_2|\mathcal H_M|$.  Hence
$\Ldim(\mathcal H_M)=O(M\log M)$.

\emph{Lower bound.}  Every shattered set yields a shattered tree, presenting
its points in a fixed order, so $\Ldim\ge\operatorname{VCdim}$.  The automata
VC estimate of \cite{ishigami1993} gives
$\operatorname{VCdim}=\Omega(n\log n)$ for $n$-state acceptors over fixed
alphabets, and Remark~\ref{rem:dfa-to-mealy} converts an $n$-state acceptor
into an $O(n)$-state Mealy machine preserving shattering, so
$\operatorname{VCdim}(\mathcal H_M^{\mathrm{reset}})=\Omega(M\log M)$.

That estimate concerns classification of \emph{separately presented} words and
does not by itself bound the sequential complexity of a single persistent
transducer stream, which is a different protocol.  The lower bound appropriate
to the persistent-stream protocol is the explicit forcing construction of
Theorem~\ref{thm:stream-lower-bound}, which is proved without reference to this
lemma.  The present lemma is therefore stated for the ordinary Littlestone
dimension in the reset-word protocol; persistent-stream statements rest on the
explicit construction rather than on this lemma.

For $|\mathcal O|>2$ the halving step is the multiclass version
\cite{daniely2014,cesabianchi2006}, which changes only the alphabet-dependent
constants.
\end{proof}

\begin{remark}[Protocol Scope of Lemma~\ref{lem:littlestone}]
\label{rem:ldim-protocol-scope}
The upper bound is a counting bound and is protocol independent.  The lower
bound is inherited from the reset-word VC construction
\cite{ishigami1993}, in which each instance is a fresh word evaluated from the
designated initial state.  Consequently $\Ldim(\mathcal H_M)=\Theta(M\log M)$
should be read in the reset-word protocol.  Transfer to a single persistent
stream requires a mistake tree realizable \emph{inside one input stream};
Theorem~\ref{thm:stream-lower-bound} constructs such a tree, so the transfer is
valid and the persistent-stream bound is also $\Theta(M\log M)$.  See
Remark~\ref{rem:stream-lb-mechanism}.
\end{remark}

The constants depend on the alphabets but not on $M$ or the horizon $T$.

%----------------------------------------------------------------------
\subsection{Passive Realizable Setting}
\label{subsec:passive-realizable}
%----------------------------------------------------------------------

In the passive realizable setting, the target input-output function
$f^\star$ belongs to $\mathcal H_M$, and the learner makes predictions against
an adversarial input sequence.

\begin{theorem}[Passive Realizable Mistakes]
\label{thm:passive-realizable}
In the passive realizable setting, with known or unknown initial state, the
minimax mistake count on a single persistent stream satisfies the
unconditional upper bound
\[
\inf_{\mathcal A}
\sup_{f^\star\in\mathcal H_M}
\sup_{x\in\mathcal I^{\mathbb N}}
\sum_{t=1}^T
\mathbf 1\{y_t\neq f^\star(x)_t\}
=
O(M\log M).
\]
Moreover the matching lower bound holds in the \emph{same} persistent-stream
protocol: by Theorem~\ref{thm:stream-lower-bound} there are fixed alphabets
with $|\mathcal I|=4$, $|\mathcal O|=2$ and, for infinitely many $M$, an
adversary strategy on one never-reset stream forcing
$\Omega(M\log M)$ mistakes.  Hence
\[
\inf_{\mathcal A}
\sup_{f^\star\in\mathcal H_M}
\sup_{x\in\mathcal I^{\mathbb N}}
\sum_{t=1}^T
\mathbf 1\{y_t\neq f^\star(x)_t\}
=
\Theta(M\log M)
\]
for $T$ at least the length of the forcing stream.  The same conclusion holds
a fortiori in the reset-word protocol, where a reset is a special case of the
input letter $\mathtt r$ below.
\end{theorem}

\begin{proof}
For known initial state, the optimal mistake bound is the Littlestone dimension
of $\mathcal H_M$ by Littlestone's theorem \cite{littlestone1988}.  By
Lemma~\ref{lem:littlestone}, this is $O(M\log M)$.  The matching lower bound
is Theorem~\ref{thm:stream-lower-bound}, proved below; its proof is
independent of the present theorem, so there is no circularity.

For unknown initial state, the effective hypothesis class is enlarged by at
most a factor $M$, because each $M$-state machine has at most $M$ possible
initial states.  The logarithm of the enlarged class size increases by at most
$\log M$, which is lower order relative to $M\log M$.  More formally, the
Littlestone dimension of a finite union of $M$ classes of Littlestone dimension
$O(M\log M)$ remains $O(M\log M)$, while the lower bound remains
$\Omega(M\log M)$ because the known-initial-state subclass is contained in the
unknown-initial-state class.  Hence the upper bound $O(M\log M)$ holds in
general, and it is tight in the reset-word protocol.
\end{proof}

\begin{theorem}[Persistent-Stream Mistake Lower Bound]
\label{thm:stream-lower-bound}
Fix $\mathcal I=\{\mathtt r,\mathtt e,\mathtt d,\mathtt c\}$, so
$|\mathcal I|=4$, and $\mathcal O=\{0,1\}$.  For every $L\ge1$ and $M=2^L$
there is an adversary strategy on a \emph{single persistent input stream},
never reset by any external mechanism, and a target
$f^\star\in\mathcal H_M$ consistent with all revealed outputs, such that:
\begin{enumerate}[label=(\roman*)]
\item every \emph{deterministic} learner makes at least
\[
M\log_2M=\Omega(M\log M)
\]
mistakes;
\item for every \emph{randomized} learner there is a fixed target in the
family with
\[
\mathbb E[\text{mistakes}]
\ \ge\
\tfrac12 M\log_2M
=
\Omega(M\log M),
\]
the expectation being over the learner's internal randomness alone.  This
follows from the bound $\mathbb E_{g}\mathbb E[\text{mistakes}]\ge\tfrac12
M\log_2M$ under the uniform target distribution on $Q^Q$, since an average is
attained by some target.  Under that same uniform distribution, for every
$\eta\in(0,1)$
\[
\Pr\left[
\text{mistakes}\ \ge\ \bigl(\tfrac12-\eta\bigr)M\log_2M
\right]
\ \ge\
1-\exp\!\left(-\tfrac12\eta^2M\log_2M\right).
\]
\end{enumerate}
The forcing stream has length $O(M\log M)$.  The construction is stated for
$|\mathcal I|=4$; Theorem~\ref{thm:stream-lb-binary} extends it to binary
input by block encoding, so the bound holds for all
$|\mathcal I|\ge2$, $|\mathcal O|\ge2$.
\end{theorem}

\begin{proof}
\textbf{The machine class.}
Let $Q=\{0,1\}^L$, so $|Q|=M=2^L$.  We exhibit a subfamily of $\mathcal H_M$
indexed by an arbitrary function $g:Q\to Q$.  Write a state as a bit string
$v=v_1\cdots v_L$.  Transitions and Mealy outputs are
\[
\begin{array}{llll}
\mathtt r:&\tau(v,\mathtt r)=0^L, &\lambda(v,\mathtt r)=0
&\text{(return to a fixed state)},\\[2pt]
\mathtt e:&\tau(v,\mathtt e)=v\oplus 0^{L-1}1, &\lambda(v,\mathtt e)=0
&\text{(flip the last bit)},\\[2pt]
\mathtt d:&\tau(v,\mathtt d)=v_2\cdots v_Lv_1, &\lambda(v,\mathtt d)=v_1
&\text{(cyclic shift, emit the bit shifted out)},\\[2pt]
\mathtt c:&\tau(v,\mathtt c)=g(v), &\lambda(v,\mathtt c)=0
&\text{(the unknown map)}.
\end{array}
\]
Only $\mathtt c$ depends on $g$; the letters $\mathtt r,\mathtt e,\mathtt d$
are the same in every member of the family.  Each such machine has exactly $M$
states, so the family is a subset of $\mathcal H_M$, and it has $M^M$ members.

\textbf{Two elementary facts.}
First, the word $\mathtt d^L$ acts as the identity on states, since $L$ cyclic
shifts restore $v$, and its output sequence is exactly $v_1v_2\cdots v_L$.
Thus $\mathtt d^L$ \emph{reads out} the current state without changing it.
Second, from any state the word
\[
w_v
=
\mathtt r\,
\prod_{i=1}^{L}\bigl(\mathtt d\,\mathtt e^{\,v_i}\bigr)
\]
drives the machine to $v$: starting from $0^L$, each round shifts left and
optionally sets the new last bit, so after $L$ rounds the state is
$v_1\cdots v_L$.  All outputs along $w_v$ are determined by
$\mathtt r,\mathtt e,\mathtt d$ alone, hence known to the learner; they are not
counted as forced mistakes.  Note $|w_v|\le 2L+1$.

\textbf{The adversary.}
Enumerate $Q=\{v^{(1)},\dots,v^{(M)}\}$.  The adversary maintains a partial
function $g$, initially empty, and processes the states in order.  For each
$v=v^{(k)}$ it plays the block
\[
w_v\;\mathtt c\;\mathtt d^L .
\]
The prefix $w_v$ drives the machine to $v$; the letter $\mathtt c$ moves it to
the still-unspecified state $g(v)$; the suffix $\mathtt d^L$ reads out the $L$
bits of $g(v)$ one per round, returning to $g(v)$.

During those $L$ readout rounds the learner must predict a bit before seeing
it.  The adversary answers with the \emph{opposite} bit each time, and records
those $L$ answers as the bits of $g(v)$.  Every answer is legal: at the moment
of answering, the corresponding bit of $g(v)$ has not been constrained by any
earlier round, because $g(v)$ is queried only in this block and $g$ is an
arbitrary function.  Hence the learner is forced into $L$ mistakes per block.

After the block the machine sits at $g(v)$; the next block begins with
$\mathtt r$, which returns it to $0^L$.

The precise content of the ``no reset'' claim should be stated carefully, since
the letter $\mathtt r$ is a synchronizing word for the family.  What is
excluded is an \emph{external} restart: there is no mechanism outside the
interaction that reinitializes the machine, no episode boundary, and no point
at which the learner's state or the mistake count is cleared.  The stream is
one uninterrupted sequence of rounds, every round is scored, and the learner's
knowledge carries across block boundaries.  What is \emph{not} claimed is that
the input alphabet lacks a synchronizing letter: $\mathtt r$ is an ordinary
input symbol that happens to send every state to $0^L$, and the learner may
play it or observe it like any other.  The lower bound is therefore a statement
about the persistent-stream protocol with a synchronizing input available, not
about families in which no synchronizing word exists.  This is the stronger
reading for the purpose at hand, since the adversary's power is not derived
from resets: the $M\log_2M$ forced mistakes occur in the readout rounds
$\mathtt d^{L}$, and the $\mathtt r$ letters cost nothing and reveal nothing.

\textbf{Consistency.}
After all $M$ blocks the adversary has defined $g$ on all of $Q$, and every
$g(v)\in Q$.  The resulting machine $f^\star$ lies in the family, hence in
$\mathcal H_M$, and by construction reproduces every output revealed during the
stream: the $\mathtt r,\mathtt e,\mathtt d$ outputs are forced by the fixed
part of the transition table, and the $\mathtt d^L$ readouts are exactly the
bits of $g(v)$ as recorded.  The interaction is therefore realizable, and the
protocol is realizable rather than agnostic.

\textbf{Count.}
Each of the $M$ blocks forces $L$ mistakes, giving
\[
M\cdot L
=
M\log_2 M
\]
mistakes in total.  The stream length is
$M(|w_v|+1+L)\le M(3L+2)=O(M\log M)$.

For randomized learners we use Yao's minimax principle rather than an adaptive
adversary, which avoids any independence assumption about the learner's
predictions.  Draw the target map $g:Q\to Q$ uniformly at random from the
$M^M$ possibilities and fix in advance the deterministic input stream
consisting of the blocks $w_v\,\mathtt c\,\mathtt d^L$ in a fixed enumeration
of $Q$.  This stream does not depend on the learner.

Enumerate the readout rounds $i=1,\dots,N$ with $N=M\log_2M$, and let
$\mathcal F_{i-1}$ be the $\sigma$-algebra generated by the transcript before
round $i$ together with the learner's internal randomness.  The bit revealed
at round $i$ is a specific bit of $g(v)$ for the block currently being
processed, and each such bit is queried exactly once along the stream.  Under
the uniform law on $Q^Q$ the bits of $g$ are i.i.d.\ fair bits, so the bit
revealed at round $i$ is a fair bit \emph{independent of} $\mathcal F_{i-1}$.
The learner's prediction at round $i$ is $\mathcal F_{i-1}$-measurable, hence
independent of that bit, so writing $X_i$ for the mistake indicator,
\[
\mathbb E\bigl[X_i\mid\mathcal F_{i-1}\bigr]=\tfrac12
\qquad\text{almost surely.}
\]
Note that this conditional statement is all that is required; the $X_i$
themselves need not be independent.  Summing,
\[
\mathbb E\Bigl[\sum_{i=1}^NX_i\Bigr]=\frac N2 ,
\]
and since the average over $g$ is at most the worst case over $g$, some fixed
$g\in Q^Q$ forces expected mistakes at least $N/2$ for the given learner.

For concentration, set $Y_i=X_i-\mathbb E[X_i\mid\mathcal F_{i-1}]$.  Then
$(Y_i)$ is a martingale difference sequence with $|Y_i|\le1$, so
Azuma--Hoeffding \cite[Ch.~2]{boucheron2013} gives, for $\eta\in(0,1)$,
\[
\Pr\Bigl[\sum_{i=1}^NX_i\le\bigl(\tfrac12-\eta\bigr)N\Bigr]
=
\Pr\Bigl[\sum_{i=1}^NY_i\le-\eta N\Bigr]
\le
\exp\!\left(-\frac{\eta^2N}{2}\right).
\]
This is the stated high-probability bound, with no independence assumption on
the rounds.

No deterministic per-run bound is available for randomized learners: a learner
that happens to guess every bit correctly makes no mistakes on that run, an
event of probability $2^{-N}>0$.  Claim~(i) is therefore restricted to
deterministic learners, and~(ii) is the correct randomized statement.

Finally, for $M$ not a power of two, apply the construction with
$L=\lfloor\log_2M\rfloor$ and leave the remaining fewer than $M/2$ states
unreachable; the bound changes only by a constant factor.
\end{proof}

\begin{remark}[What the Construction Supplies]
\label{rem:stream-lb-mechanism}
The obstruction identified in Remark~\ref{rem:ldim-protocol-scope} is that a
reset-word shattering argument does not embed into one persistent stream
without a mechanism for changing state, since the learner cannot be
repositioned between instances.  The
construction above removes that obstruction by making navigation part of the
input alphabet: $\mathtt r$ and $\mathtt e,\mathtt d$ furnish an explicit
\emph{transport system} that reaches any state from any state using only
outputs the learner can already predict, so the transport rounds are free.  The
$\Omega(M\log M)$ information is then extracted by the readout word
$\mathtt d^L$, which is output-revealing and state-preserving.  In effect the
persistent stream simulates the reset protocol at no mistake cost, which is
exactly what the transfer required.
\end{remark}

\begin{remark}[Alphabet Size]
\label{rem:stream-lb-alphabet}
The construction uses $|\mathcal I|=4$ and $|\mathcal O|=2$, and is unchanged
for any larger alphabets, so the bound holds whenever $|\mathcal I|\ge4$ and
$|\mathcal O|\ge2$.  The binary-input case is handled separately in
Theorem~\ref{thm:stream-lb-binary} by block encoding, at the cost of a
constant factor.
\end{remark}

\begin{theorem}[Binary-Input Persistent-Stream Lower Bound]
\label{thm:stream-lb-binary}
Let $\mathcal I=\{0,1\}$ and $\mathcal O=\{0,1\}$.  There is a constant
$c>0$ such that for infinitely many $M'$ there is an adversary on a single
persistent, never-reset binary input stream forcing every deterministic
learner to make at least
\[
c\,M'\log_2M'
=
\Omega(M'\log M')
\]
mistakes, and every randomized learner to make that many in expectation up to
a factor $\tfrac12$.  One may take any $c<\tfrac13$ for $M'$ large.
\end{theorem}

\begin{proof}
Encode the four letters of Theorem~\ref{thm:stream-lower-bound} as two-bit
blocks
\[
\mathtt r\mapsto00,\qquad
\mathtt e\mapsto01,\qquad
\mathtt d\mapsto10,\qquad
\mathtt c\mapsto11 .
\]
Let $A$ be an $M$-state machine of the family $\mathcal G_M$, $M=2^L$.  Define
a binary-input Mealy machine $A'$ with state set
\[
Q'=Q\times\{-,0,1\},
\]
where the second coordinate buffers the first bit of the current block.  On
reading the first bit $b$ of a block, $A'$ moves from $(q,-)$ to $(q,b)$ and
emits $0$; on reading the second bit $b'$, it applies the letter encoded by
$(b,b')$, moves to $(\tau(q,a),-)$, and emits $\lambda(q,a)$.  Then
$|Q'|=3M$, and $A'$ reproduces exactly the output of $A$ on the second bit of
each block, emitting the uninformative symbol $0$ on every first bit.

The adversary runs the strategy of Theorem~\ref{thm:stream-lower-bound}
through this encoding.  First bits are perfectly predictable, so they force no
mistakes and cost nothing; each readout of $A$ still occurs on exactly one
round of the binary stream, namely the second bit of a $\mathtt d$ block, and
is still a forced mistake by the same lazy-adversary argument.  Hence the
number of forced mistakes is unchanged at $M\log_2M$.

Writing $M'=3M$ for the state count of the binary machine,
\[
\frac{M\log_2M}{M'\log_2M'}
=
\frac{M\log_2M}{3M(\log_2M+\log_23)}
=
\frac13\cdot\frac{L}{L+\log_23}
\xrightarrow[L\to\infty]{}
\frac13 ,
\]
and the ratio is increasing in $L$, exceeding $0.23$ already at $L=3$.  Thus
the forced mistake count is at least $c\,M'\log_2M'$ for any $c<\tfrac13$ once
$M'$ is large, which is $\Omega(M'\log M')$.  The randomized statement follows
exactly as in Theorem~\ref{thm:stream-lower-bound}(ii).
\end{proof}

\begin{lemma}[Subsequence-to-All-$M$ Extension]
\label{lem:subsequence-allM}
Let $C:\mathbb N\to[0,\infty)$ be nondecreasing.  Suppose there are
$M_1<M_2<\cdots$ with $M_{j+1}\le\alpha M_j$ for a constant $\alpha>1$, and
$C(M_j)\ge\beta M_j\log M_j$ for all $j$.  Then there is $\beta'>0$ with
\[
C(M)\ \ge\ \beta'M\log M
\]
for all sufficiently large $M$.
\end{lemma}

\begin{proof}
Given large $M$, pick $j$ with $M_j\le M<M_{j+1}$; then $M_j\ge M/\alpha$.
Monotonicity gives
\[
C(M)\ \ge\ C(M_j)\ \ge\ \beta M_j\log M_j
\ \ge\ \frac{\beta}{\alpha}M\bigl(\log M-\log\alpha\bigr).
\]
For $M\ge\alpha^2$ we have $\log M-\log\alpha\ge\tfrac12\log M$, so
$C(M)\ge\frac{\beta}{2\alpha}M\log M$.
\end{proof}

\begin{corollary}[Persistent-Stream Bound for All $M$]
\label{cor:stream-all-M}
The minimax realizable mistake complexity on a persistent stream is
nondecreasing in $M$, since $\mathcal H_M\subseteq\mathcal H_{M+1}$.  The
constructions of Theorems~\ref{thm:stream-lower-bound}
and~\ref{thm:stream-lb-binary} give lower bounds along the subsequences
$M=2^L$ and $M'=3\cdot2^L$ respectively, each with ratio $\alpha=2$.  By
Lemma~\ref{lem:subsequence-allM}, therefore,
\[
\operatorname{Mistakes}^{\mathrm{stream}}_{\mathrm{realizable}}(M)
=
\Omega(M\log M)
\]
for \emph{all} sufficiently large $M$, and hence, with the counting upper
bound, $\Theta(M\log M)$ for all sufficiently large $M$ and every fixed
alphabet pair with $|\mathcal I|\ge2$, $|\mathcal O|\ge2$.
\end{corollary}

\begin{remark}[Binary Input Closes the Alphabet Gap]
\label{rem:stream-lb-binary}
Theorems~\ref{thm:stream-lower-bound} and~\ref{thm:stream-lb-binary} together
give $\Omega(M\log M)$ on a persistent stream for \emph{every} fixed alphabet
pair with $|\mathcal I|\ge2$ and $|\mathcal O|\ge2$, matching the finite-class
upper bound.  The block encoding costs only the constant factor $3$ in state
count and nothing in forced mistakes, because the padding rounds are
perfectly predictable and therefore free.
\end{remark}

\begin{remark}[Scope of the Streaming Lower Bound]
\label{rem:streaming-lower-scope}
The upper bound above is a finite-class bound and is unconditional.  The
lower bound is inherited from $\Ldim(\mathcal H_M)$, and standard
VC-dimension lower bounds for automata are stated for classification of
independently presented words with resets.  A single persistent transducer
stream is a different protocol: the learner cannot reset the machine, and
successive rounds are dependent.  Transferring the lower bound therefore
requires a shattering or mistake-tree construction realized \emph{within one
input stream}.  Absent such a construction,
Theorem~\ref{thm:passive-realizable} yields an unconditional $O(M\log M)$
upper bound together with an $\Omega(M\log M)$ lower bound for the reset-word
protocol only.
\end{remark}

\begin{remark}[No Passive Synchronization Assumption]
\label{rem:passive-no-sync}
The unknown-initial-state proof does not assume that the learner can force
synchronization.  In the passive model, inputs are adversarial, and the learner
cannot in general apply a synchronizing word.  The bound follows from the size
and Littlestone dimension of the enlarged hypothesis class, not from active
state identification.
\end{remark}

\begin{remark}[Identification versus State Tracking]
\label{rem:identification-vs-tracking}
What is identified or predicted is the correct input-output function, not
necessarily the future internal state trajectory.  In deterministic realizable
prediction, correct output prediction over all continuations may require
identifying the residual equivalence class of the current history, but the
learner is not required to track the machine's internal state label.
\end{remark}

%----------------------------------------------------------------------
\subsection{Active Realizable Setting}
\label{subsec:active-realizable}
%----------------------------------------------------------------------

In the active realizable setting, the learner chooses inputs.  This can reduce
or eliminate the unknown-initial-state obstruction if the machine admits an
output-aware synchronizing or distinguishing experiment.

\begin{definition}[Active Identification Objectives]
\label{def:residual-knowledge}
Three objectives must be separated, because they carry different costs and are
achieved against different prior knowledge.  In each case the learner is said
to \emph{attain} the objective at the first time $\tau$ at which the stated
quantity is determined by the transcript up to $\tau$, uniformly over all
targets in the class consistent with that transcript.

\textbf{(SI) State identification in a known skeleton.}  The transition and
output tables are given; only the initial state is unknown.  The learner
attains~(SI) when it can name the current state \emph{up to observational
equivalence}, that is, when all states consistent with the transcript have the
same continuation function.  The qualification is needed because a nonminimal
skeleton may contain distinct states with identical behaviour, which no
input-output experiment can separate; for a minimal skeleton it is the same as
naming the state.

\textbf{(RI) Residual identification.}  The machine is unknown.  The learner
attains~(RI) when it can identify the current Myhill--Nerode residual class,
equivalently the full continuation function from the current history.

\textbf{(MI) Model identification from a known initial state.}  The initial
state is given; the tables are unknown.  The learner attains~(MI) when it can
name the residual class of every reachable history, not merely of the current
one.

Each objective is \emph{prediction-closing}: once attained, the continuation
function from the current history is determined, so under the realizable
assumption every subsequent output is predictable with no further mistakes.
For~(SI) this holds because the tables are given and the state is now known;
for~(RI) and~(MI) it is immediate from the definition.  Consequently no
objective admits a nontrivial post-attainment continuation term, and the total
cost of a learner meeting the objective is exactly the cost of attaining it.

Throughout,
\[
\MistRI(M)
=
\inf_{\mathcal A\in\mathcal R}\
\sup_{A\in\mathcal H_M,\ q\in Q_A}\
\operatorname{tot}(\mathcal A,A,q),
\]
where $\operatorname{tot}$ counts \emph{all} prediction mistakes over the whole
interaction and $\mathcal R$ is the set of active learners that attain
objective~(RI) at some finite time on every instance.  Both $\MistRI$ and
$\Esync$ are minimax over the same learner set; they differ in what they count.
$\Esync(M)$ of Definition~\ref{def:sync-mistake-complexity} counts only the
mistakes made \emph{before} attainment, so
\[
\MistRI(M)\ \ge\ \Esync(M)
\]
holds by definition, a learner being free a priori to err after attaining~(RI).
The reverse inequality is not definitional: it is
Theorem~\ref{thm:active-certified}(I), and its proof turns on the
prediction-closing property of Definition~\ref{def:residual-knowledge}, which
forbids post-attainment mistakes.  Were~(RI) replaced by an objective that is
not prediction-closing, the two quantities would separate.

The requirement is not vacuous: as Remark~\ref{rem:active-objective-needed}
observes, a learner that repeatedly plays an input with known constant output
makes no mistakes and learns nothing, so without an identification requirement
the unconstrained quantity is $0$.  All active lower bounds in this subsection
are therefore identification complexities.

The objectives differ in what must be paid \emph{before} attainment, and are
not ordered by inclusion: (SI) presupposes the tables and asks only for the
state, whereas~(MI) presupposes the state and asks for the tables.  (RI)
presupposes neither, and is the objective against which
Definition~\ref{def:sync-mistake-complexity} and all lower bounds of this
subsection are stated.
\end{definition}

\begin{remark}[Which Objective Each Bound Uses]
\label{rem:objective-bookkeeping}
Because all three objectives of Definition~\ref{def:residual-knowledge} are
prediction-closing, a mistake bound in this subsection is always a bound on
the cost of \emph{attainment}, never a sum of an attainment cost and a
subsequent learning cost.  In particular no $O(M\log M)$ continuation term is
appended after attainment under any of the three objectives; where such a term
appears it is part of the attainment cost itself.

The complexity $\Esync(M)$ of Definition~\ref{def:sync-mistake-complexity} is
defined against~(RI), and the upper bound it yields is $O(\Esync(M))$
(Theorem~\ref{thm:active-certified}(I)).  The length-form bound of
Proposition~\ref{prop:active-length-upper} is also stated against~(RI); the
$M\log M$ term there arises because the universal experiment realising
$\Lsyncu(M)$ certifies the residual class only after a further model-search
phase, and is therefore part of the attainment cost, not a continuation term.
The lower-bound constructions of
Theorems~\ref{thm:active-explicit-directsum} and~\ref{thm:active-two-phase}
are likewise stated against~(RI); their additive content lies in the two
\emph{independent unknown maps}, not in a synchronization-plus-Littlestone
split.
\end{remark}

\begin{definition}[Adaptive Output-Aware Synchronization Length]
\label{def:output-aware-sync}
For a deterministic Mealy machine $A$ with unknown initial state, an adaptive
output-aware synchronization experiment is a finite decision tree.  Each
internal node is labeled by an input $x\in\mathcal I$; each outgoing edge is
labeled by an observed output $y\in\mathcal O$.  A leaf is reached after
observing the corresponding input-output transcript.

The experiment identifies the residual class if every leaf corresponds to a
single Myhill--Nerode residual class of $A$.  Let
\[
L_{\mathrm{sync}}^{\mathrm{adapt}}(A)
\]
be the minimum worst-case depth of such an experiment \emph{for the machine
$A$}.  Define
\[
\Lsync(M)
=
\sup_{A\in\mathcal H_M}
L_{\mathrm{sync}}^{\mathrm{adapt}}(A),
\]
with value $\infty$ if some $M$-state machine admits no finite adaptive
synchronization experiment.

This quantity is \emph{machine-specific}: the experiment realizing
$L_{\mathrm{sync}}^{\mathrm{adapt}}(A)$ may depend on $A$.  A learner that does
not know $A$ cannot in general run it.  For such a learner define instead the
\textbf{universal adaptive synchronization depth}
\[
\Lsyncu(M)
=
\min_{\mathcal T}\ \sup_{A\in\mathcal H_M}\ \operatorname{depth}_A(\mathcal T),
\]
the minimum over \emph{single} decision trees $\mathcal T$, implementable
without knowledge of $A$, of the worst-case depth at which the transcript
produced by $\mathcal T$ determines the \emph{current state} of $A$ within each
machine still consistent with it.

The two quantities certify different things, and the distinction is what makes
the length-form bound of Proposition~\ref{prop:active-length-upper} come out as
it does.  $L_{\mathrm{sync}}^{\mathrm{adapt}}(A)$ is defined for a known
machine $A$ and certifies its residual class, which for known tables is the
same as certifying the state.  A universal tree faces an unknown machine, so
the strongest thing it can certify from an output transcript is the state
within each consistent candidate; the tables, and hence the residual class,
remain to be resolved.  This is objective~(SI) of
Definition~\ref{def:residual-knowledge} relativized to the consistent set, not
objective~(RI), and it is why a subsequent model-search phase appears in the
proof.

Always
\[
\Lsync(M)\ \le\ \Lsyncu(M),
\]
and the inequality may be strict, since the left side optimizes the experiment
per machine while the right side must fix one experiment in advance.

\emph{Identification is up to observational equivalence.}  The phrase
``determines the current state'' must be read modulo Myhill--Nerode
equivalence.  The class $\mathcal H_M$ contains
non-minimal machines carrying distinct states $s\ne s'$ with
$\lambda(s,w)=\lambda(s',w)$ for every input word $w$.  No experiment whatsoever
can separate such a pair, since separation would require an observable
difference and there is none.  Under a literal state-identification reading one
would therefore have $\Lsyncu(M)=\infty$ for trivial reasons, and the quantity
would carry no content.

Accordingly, a transcript is said to determine the current state when it
determines its class under the observational equivalence
\[
s\sim s'
\iff
\lambda(s,w)=\lambda(s',w)\ \text{for all }w\in\mathcal I^*,
\]
equivalently when it determines the state of the minimal machine equivalent to
$A$.  On minimal machines the two readings agree, since there $\sim$ is the
identity; the same convention applies to $\Lsync$, to objective~(SI) of
Definition~\ref{def:residual-knowledge}, and to $\EsyncSI$.  With this reading
$\Lsyncu(M)$ is finite whenever some universal experiment exists, and
Lemma~\ref{lem:moore-separation} supplies the separating words that make it so.
\end{definition}

\begin{definition}[Synchronization Mistake Complexity]
\label{def:sync-mistake-complexity}
For an active learner $\mathcal A$, a machine $A\in\mathcal H_M$ and an initial
state $q$, let $\mistk(\mathcal A,A,q)$ denote the number of prediction
mistakes $\mathcal A$ incurs before it attains objective~(RI) of
Definition~\ref{def:residual-knowledge}, with the value $+\infty$ if it never
does.  Define the \emph{synchronization mistake complexity}
\[
\Esync(M)
=
\inf_{\mathcal A}\
\sup_{A\in\mathcal H_M,\ q\in Q_A}\
\mistk(\mathcal A,A,q),
\]
the infimum being over all active learners, deterministic or randomized, with
$\mistk$ replaced by its expectation in the randomized case.

The order of the operators is the substance of the definition.  The infimum
comes first, so $\Esync(M)$ is the cost incurred by the \emph{best} learner
against the worst instance: it is simultaneously an upper bound achievable by
some learner and a lower bound valid for every learner.  This two-sided
reading is what licenses its use on both sides of the bounds below.  Reversing
the operators would define a different and much larger quantity.

Despite the subscript, $\Esync(M)$ is the cost of attaining
objective~(RI), which on an unknown machine includes identifying the
continuation function and not merely synchronizing the state.  It reduces to a
state-synchronization cost exactly when the tables are known, as in
objective~(SI).

For a subclass $\mathcal C_M\subseteq\mathcal H_M$ write
\[
\Esync(\mathcal C_M)
=
\inf_{\mathcal A}\
\sup_{A\in\mathcal C_M,\ q\in Q_A}\
\mistk(\mathcal A,A,q),
\]
so that $\Esync(M)=\Esync(\mathcal H_M)$.

Two distinct complexities are in play and must not be conflated.  Alongside the
residual-identification quantity above, define the
\emph{state-identification complexity}
\[
\EsyncSI(\mathcal C_M)
=
\inf_{\mathcal A}\
\sup_{A\in\mathcal C_M,\ q\in Q_A}\
\mistk^{\mathrm{SI}}(\mathcal A,A,q),
\]
where $\mistk^{\mathrm{SI}}$ counts mistakes before objective~(SI) is attained,
the tables being given.  Since attaining~(RI) on an unknown machine requires
resolving the tables as well as the state,
\[
\EsyncSI(\mathcal C_M)\ \le\ \Esync(\mathcal C_M),
\]
and the two may differ by the cost of model search.  The additive statements
below decompose $\Esync$ into a state-identification part and a continuation
part, so it is $\EsyncSI$ that pairs additively with the continuation cost.
The distinction matters because on
the full class the complexity is pinned down outright
(Corollary~\ref{cor:active-theta}), so the additive form
$M\log M+\Esync(M)$ collapses to $\Theta(\Esync(M))$ there; it is on
subclasses, where the counting bound $\log_2|\mathcal C_M|$ need not be
$O(M\log M)$, that the two terms can genuinely separate.  Statements below
carrying a direct-sum or operational-equivalence hypothesis are to be read with
$\Esync$ evaluated on the relevant subclass.
\end{definition}

\begin{proposition}[State Identification Costs Only Logarithmically Many Mistakes]
\label{prop:esyncsi-log}
Let the transition and output tables of a \emph{minimal} $M$-state machine be
given and only the initial state unknown, as in objective~(SI).  Then
\[
\EsyncSI(M)\ \le\ \bigl\lfloor\log_2M\bigr\rfloor ,
\]
for \emph{every} output alphabet with $|\mathcal O|\ge2$.  In particular the
bound carries no dependence on $|\mathcal O|$.
\end{proposition}

\begin{proof}
Maintain the version space $V\subseteq Q$ of states consistent with the
transcript, initially $V=Q$, so $|V|\le M$.  While $|V|>1$, minimality
guarantees that some input word separates two survivors; let $x$ be the first
letter of such a word on which the survivors do not all agree in output, which
exists after advancing $V$ along finitely many letters with agreed output.
Predict the output symbol receiving a plurality of votes among $\{\lambda(s,x):
s\in V\}$ and play $x$.

Partition $V$ by the value of $\lambda(\cdot,x)$ into output classes of sizes
$c_1\ge c_2\ge\dots\ge c_r$, where $r\le|\mathcal O|$ and $\sum_ic_i=|V|$, and
let the learner predict the plurality symbol, the one realizing $c_1$.

If the prediction is wrong, the observed symbol realizes some class $c_j$ with
$j\ge2$, and the survivors are exactly that one class, of size
$c_j\le c_2$.  The key point is that the survivors form a \emph{single} class
rather than the union of all non-predicted classes.  Since $c_1\ge c_2$ and
$c_1+c_2\le|V|$,
\[
2c_2\ \le\ c_1+c_2\ \le\ |V|,
\qquad\text{so}\qquad
c_2\ \le\ \tfrac12|V| .
\]
Every mistake therefore at least halves the version space, whatever the size of
the output alphabet.  If the prediction is right, $V$ can only shrink.
Advancing every survivor along $x$ preserves $|V|$ or decreases it, since
$\tau(\cdot,x)$ is a function.

Starting from $|V|\le M$ and halving at each mistake, with $|V|\ge1$
throughout, the number of mistakes is at most $\lfloor\log_2M\rfloor$.  When
$|V|=1$ the current state is determined, which is objective~(SI); since
$\EsyncSI$ is an infimum over learners, the bound follows.
\end{proof}

\begin{remark}[Why the Alphabet Drops Out]
\label{rem:halving-alphabet-free}
A coarser count replaces ``the survivors are one class of size $c_2$'' by
``the survivors are everything not predicted'', giving the weaker factor
$1-1/|\mathcal O|$ and the bound
$\log_{|\mathcal O|/(|\mathcal O|-1)}M$, which degrades as the alphabet grows:
at $M=1024$ it reads $10.0$, $17.1$, $24.1$ and $51.9$ for
$|\mathcal O|=2,3,4,8$.  That is an artefact of the counting, not of the
problem.  Because a deterministic machine sends each surviving state to exactly
one output symbol, an erring learner is left with a single class, and the
plurality rule forces that class to be no larger than half the version space.

The same counting applies wherever the version space consists of candidates
each of which predicts a single deterministic symbol.  In particular it applies
to the machine--state pairs of Theorem~\ref{thm:active-halving}, since a
deterministic machine paired with a state emits one symbol under a given input;
the alphabet-dependent factor is unnecessary there as well.
Enumerating all integer class profiles with $|V|\le25$ and $r\le5$ confirms
$c_2\le|V|/2$ without exception, with equality exactly when $c_1=c_2=|V|/2$ and
$r=2$; larger alphabets only fragment the complement further and cannot help
the adversary.
\end{remark}

\begin{theorem}[State Identification Costs Logarithmically Many Mistakes, Exactly]
\label{thm:esyncsi-theta}
For every output alphabet with $|\mathcal O|\ge2$ and every input alphabet
with $|\mathcal I|\ge1$,
\[
\EsyncSI(M)\ =\ \bigl\lfloor\log_2M\bigr\rfloor ,
\]
with \emph{no} dependence on $|\mathcal O|$ or $|\mathcal I|$ beyond the
requirement $|\mathcal O|\ge2$.  In particular $\EsyncSI(2^L)=L$.

The lower bound is witnessed by a minimal $M$-state machine on which every
deterministic learner is forced to make $\lfloor\log_2M\rfloor$ mistakes before
identifying the current state; every randomized learner makes at least
$\tfrac12\lfloor\log_2M\rfloor$ mistakes in expectation against a fixed initial
state.
\end{theorem}

\begin{proof}
The upper bound is Proposition~\ref{prop:esyncsi-log}, which gives
$\lfloor\log_2M\rfloor$ for every $|\mathcal O|\ge2$.  For the lower bound fix
$L\ge1$, put $M=2^L$, and let
\[
Q=\{0,1\}^L,
\qquad
\mathcal O=\{0,1\}.
\]
Use a single input letter $\mathtt d$ acting as the left cyclic shift,
\[
\tau(v_1v_2\cdots v_L,\mathtt d)=v_2\cdots v_Lv_1,
\qquad
\lambda(v_1v_2\cdots v_L,\mathtt d)=v_1 ,
\]
adjoining a second letter acting as the identity with constant output if
$|\mathcal I|\ge2$ is required; that letter reveals nothing and can be ignored.
The tables are known to the learner and only the initial state is unknown, as
objective~(SI) requires.  The construction uses only two of the available
output symbols, so it is available verbatim for every $|\mathcal O|\ge2$; a
larger alphabet neither helps the learner nor the adversary here.

\emph{Minimality.}  Two distinct states $v\ne w$ differ in some coordinate $j$,
and the word $\mathtt d^{\,j-1}$ drives them to states whose first coordinates
are $v_j\ne w_j$, so the next output separates them.  Hence all $M$ states are
pairwise distinguishable and the machine is minimal.

\emph{The transcript is exactly the initial state.}  Applying $\mathtt d$
repeatedly from $v$ emits $v_1,v_2,\dots,v_L$ in order and returns to $v$.
Thus after $L$ rounds the state is determined, and before round $t\le L$ the
transcript constrains only $v_1,\dots,v_{t-1}$: the remaining $2^{L-t+1}$
states consistent with the transcript are exactly those with the observed
prefix, and every completion is a legal initial state.

\emph{Deterministic lower bound.}  Run the adversary argument on the version
space.  At round $t\le L$ the learner must commit to a prediction
$\widehat y_t\in\{0,1\}$ for the bit $v_t$.  Both completions $v_t=0$ and
$v_t=1$ remain consistent, so the adversary sets $v_t=1-\widehat y_t$, forcing
a mistake, and records the bit.  After $L$ rounds the adversary has specified
an initial state $v^\star\in Q$ consistent with every revealed output, and the
learner has made $L=\log_2M$ mistakes.  Until round $L$ the version space has
size at least $2$, so the learner has not attained~(SI); the mistakes are
therefore all incurred before identification.

\emph{Randomized lower bound.}  Let the initial state be uniform on $Q$.
Conditioned on the transcript through round $t-1$, the bit $v_t$ is a fair coin
independent of the learner's internal randomness, so any prediction errs with
probability exactly $\tfrac12$ and the expected number of mistakes over the
first $L$ rounds is $L/2$.  Averaging over the prior, some fixed
$v^\star\in Q$ forces at least $\tfrac12\log_2M$ expected mistakes against the
learner's own randomness.

For general $M$ no padding is required, because $\mathcal H_M$ consists of the
machines with \emph{at most} $M$ states.  Put $L=\lfloor\log_2M\rfloor$, so
$2^L\le M$.  The cyclic-shift machine above has exactly $2^L$ states and is
minimal, hence belongs to $\mathcal H_M$, and against it every deterministic
learner is forced into $L$ mistakes.  Since $\EsyncSI(M)$ is a supremum over
$\mathcal H_M$, this gives $\EsyncSI(M)\ge L=\lfloor\log_2M\rfloor$, matching
the upper bound.
\end{proof}

\begin{remark}[Why No Padding Is Needed]
\label{rem:esyncsi-intermediate}
Equality at intermediate budgets is a consequence of the ``at most $M$
states'' convention for $\mathcal H_M$, not of an interpolating construction.
The witness is built at $2^L$ states, and for any $M$ with $2^L\le M$ that same
machine is already a member of $\mathcal H_M$; no padding to exactly $M$ states
is required, and none is available, since adjoining duplicate states destroys
minimality while adjoining fresh reachable ones destroys the fair-bit property.

Were $\mathcal H_M$ instead read as machines with \emph{exactly} $M$ states,
the argument would give only $\EsyncSI(M)=\Theta(\log M)$, with the constant
lost to monotonicity.  The two readings agree on the powers of two.
\end{remark}

\begin{remark}[Both Bounds Are Attained by the Same Mechanism]
\label{rem:esyncsi-tight}
The matching bounds have a common explanation.  The halving argument of
Proposition~\ref{prop:esyncsi-log} shows that each mistake removes at least
half of the version space, so $\lfloor\log_2M\rfloor$ mistakes suffice; the
cyclic-shift family of Theorem~\ref{thm:esyncsi-theta} is the case in which
each mistake removes exactly half and no more, because each round exposes one
fresh, unconstrained bit of the initial state.  The family is thus extremal for
the halving bound, and the two sides meet exactly rather than up to a
constant.

Solving the associated games exactly confirms both sides.  For
$L=1,\dots,10$ the deterministic minimax mistake count on the cyclic-shift
family is exactly $L=\log_2M$ and the Bayes-optimal expected count under the
uniform prior is exactly $L/2$, with all $M$ states pairwise separated so that
the family is minimal.  In the opposite direction, exhausting all minimal Mealy
machines across nine signatures with $|\Splus|\le4$, $|\mathcal I|\le2$ and
$|\mathcal O|\le4$ --- $46{,}656$ table pairs in the largest, of which
$35{,}640$ are minimal --- the worst-case state-identification cost never
exceeds $\lfloor\log_2M\rfloor$, and attains it at $M=2$ and $M=4$.  Holding
$|\Splus|$ fixed and enlarging $|\mathcal O|$ from $2$ to $4$ leaves the worst
case unchanged, as Remark~\ref{rem:halving-alphabet-free} predicts.

Two consequences follow for the structure of the theory.  First, the collapse
recorded in Remark~\ref{rem:additive-collapses} is now unconditional: it rests
on a two-sided estimate rather than on an upper bound alone, so no choice of
subclass can rescue a mistake-currency additive decomposition.  Second, there
is no dependence on $|\mathcal O|$ at all.  Both bounds equal
$\lfloor\log_2M\rfloor$: the upper bound because an erring learner is left with
a single output class of size at most half the version space
(Remark~\ref{rem:halving-alphabet-free}), the lower bound because the
cyclic-shift witness needs only two output symbols.  Enlarging the alphabet
neither speeds up identification nor slows it down, and the base of the
logarithm is $2$ irrespective of $|\mathcal O|$.
\end{remark}

\begin{remark}[The Additive Form Collapses]
\label{rem:additive-collapses}
Proposition~\ref{prop:esyncsi-log}, sharpened to an equality by
Theorem~\ref{thm:esyncsi-theta}, has a structural consequence for the
additive decomposition.  Since $\EsyncSI(\mathcal C_M)=O(\log M)$ for every
subclass of minimal machines,
\[
M\log M+\EsyncSI(\mathcal C_M)=\Theta(M\log M),
\]
so the two terms never separate and the additive form carries no information
beyond $\Theta(M\log M)$.  Accordingly no additive theorem is asserted with a
state-identification second term.

A genuine two-term decomposition would require measuring the synchronization
component in a currency other than mistakes --- experiment \emph{length}, as in
$\Lsyncu$ of Definition~\ref{def:output-aware-sync}, is the natural candidate,
and Proposition~\ref{prop:active-length-upper} is stated in that currency.  The
mistake-counting decomposition is degenerate.
\end{remark}

\begin{lemma}[Operational Equivalence]
\label{lem:operational-equivalence}
Let the class be a set of deterministic machines and let the protocol be
realizable.  Then a learner guarantees zero future mistakes on all
continuations from a given history if and only if it has attained
objective~(RI) of Definition~\ref{def:residual-knowledge} at that history.

For deterministic realizable classes this is therefore a theorem, not a
hypothesis; it is stated separately because the results below that invoke it
are also meaningful for stochastic classes, where it must be assumed.
\end{lemma}

\begin{proof}
\emph{(RI) implies zero future mistakes.}  This is the prediction-closing
property recorded in Definition~\ref{def:residual-knowledge}: the continuation
function from the current history is determined by the transcript, uniformly
over consistent targets, and realizability places the true target among them,
so predicting by that common function is correct on every subsequent round.

\emph{Zero future mistakes implies (RI).}  Suppose~(RI) has not been attained,
so two machines $A,A'$ in the class are consistent with the transcript and
their continuation functions from the current histories differ.  Let $w$ be a
shortest word on which the two continuation functions disagree, which exists by
Lemma~\ref{lem:moore-separation}, and let $j$ be the first position at which
the emitted symbols differ.  The learner's prediction at that round is a
function of the transcript alone, hence identical in the run against $A$ and in
the run against $A'$, while the two true symbols differ.  Both runs are
realizable, since $A$ and $A'$ both lie in the class and both are consistent
with everything observed.  The learner therefore errs in at least one of them,
contradicting the guarantee of zero future mistakes on all continuations.
\end{proof}

\begin{definition}[Direct-Sum Active Instance]
\label{def:direct-sum-active}
A \textbf{direct-sum active instance} of budget $M$ consists of an
$O(M)$-state machine family together with a decomposition of the learning task
into two components satisfying:
\begin{enumerate}[label=(\roman*)]
\item a synchronization component whose minimax mistake complexity is $S_M$,
measuring \emph{state identification only} and so comparable to
$\EsyncSI$ rather than to $\Esync$;
\item a continuation component whose active realizable mistake complexity is
$C_M$, measuring the residual model search that remains once the state is
known;
\item \textbf{disjoint accounting}: the rounds on which the two components can
force mistakes are disjoint;
\item \textbf{unavoidability}: the objective cannot be met without meeting it
for \emph{both} components, and this is a property of the machine family
rather than an external restriction on the protocol;
\item \textbf{product structure}: the family is parameterized by a pair
$(\theta_1,\theta_2)$ of unknowns, one per component, and \emph{every}
combination occurs, so the family is the full product of its two coordinate
ranges;
\item \textbf{coordinate locality}: the outputs on the rounds of $R_j$ are
determined by $\theta_j$ alone, independently of $\theta_{3-j}$.
\end{enumerate}
Conditions~(iii)--(vi) together constitute the \emph{component-restriction
property}.  It is exactly what the reduction in
Theorem~\ref{thm:active-direct-sum} consumes: (vi) makes a component's
observations well defined once its own parameter is fixed, (v) guarantees that
an arbitrary dummy value of the other parameter yields a legitimate composite
target, (iii) keeps the two mistake counts from overlapping, and (iv) forces
both components to be paid for.
Condition~(iv) is essential in the active protocol, where the learner chooses
its own inputs.  A family in which one component can be bypassed -- for
example by an irreversible transition into the other component -- does not
satisfy~(iv), and for such a family the additive bound fails.
Theorem~\ref{thm:active-two-phase} constructs a family satisfying~(iv) by
making the two components mutually reachable.
\end{definition}

\begin{theorem}[Active Realizable Certified Bounds]
\label{thm:active-certified}
For active realizable learning over $M$-state deterministic machines, the
following three statements hold under \emph{separate} hypotheses, which are
not to be combined implicitly.

\medskip
\noindent\textbf{(I) Upper bound (unconditional).}
Under objective~(RI) of Definition~\ref{def:residual-knowledge},
\[
\MistRI(M)
\ \le\
O\bigl(\Esync(M)\bigr),
\]
and this is sharp: by Definition~\ref{def:sync-mistake-complexity} some
learner attains~(RI) within $\Esync(M)$ mistakes, and~(RI) is
prediction-closing, so no mistake occurs afterwards.  No continuation term is
appended.

\noindent\textbf{(II) Synchronization lower bound (needs
Lemma~\ref{lem:operational-equivalence} alone).}
If any learner guaranteeing zero future mistakes must first acquire
guaranteed residual-class knowledge, then
\[
\MistRI(M)
\ge
\Omega\bigl(\Esync(M)\bigr).
\]
This gives \emph{only} the $\Esync$ term; it does not by itself supply the
$M\log M$ term.

\noindent\textbf{(III) Disjoint-component lower bound (needs direct-sum
saturation, Definition~\ref{def:direct-sum-active}).}
On a subclass $\mathcal C_M$ carrying a direct-sum instance with component
complexities $S_M$ and $C_M$,
\[
\MistRI(M)\ \ge\ S_M+C_M .
\]
This is a lower bound only; no matching upper bound is asserted, and no
identification of $S_M$ with a state-identification complexity is made.  By
Proposition~\ref{prop:esyncsi-log} the state-identification complexity is
$O(\log M)$, so a decomposition with $\EsyncSI$ as its second term would be
degenerate (Remark~\ref{rem:additive-collapses}); the components realized by
the explicit construction of Theorem~\ref{thm:active-two-phase} are both
model-identification components.

\medskip
Parts~(II) and~(III) are stated in terms of the abstract complexity $\Esync$
and are the form that remains informative on \emph{subclasses}
$\mathcal C_M\subseteq\mathcal H_M$, where the two terms of the additive form
can separate.  On the full class $\mathcal H_M$ no hypothesis is needed at all:
Corollary~\ref{cor:active-theta} gives
$\MistRI(M)=\Theta(M\log M)$ outright, and
$M\log M+\Esync(M)=\Theta(\Esync(M))$ there.  An explicit family satisfying the
hypothesis of~(III) is constructed in
Theorem~\ref{thm:active-two-phase}, so~(III) is nonvacuous.
\end{theorem}

\begin{proof}
\emph{Part~(I).}  Fix $\varepsilon>0$.  By
Definition~\ref{def:sync-mistake-complexity} the infimum over learners is
attained to within $\varepsilon$, so there is an active learner
$\mathcal A_\varepsilon$ which, on every $A\in\mathcal H_M$ and every initial
state, attains objective~(RI) after at most $\Esync(M)+\varepsilon$ mistakes.
Run $\mathcal A_\varepsilon$.  At the moment of attainment the continuation
function from the current history is determined by the transcript, uniformly
over all targets still consistent with it; since the protocol is realizable,
the true target is among them, and predicting according to that common
continuation function is correct on every subsequent round.  Hence the total
number of mistakes is at most $\Esync(M)+\varepsilon$.  As $\varepsilon>0$ was
arbitrary and the mistake count is integer valued, the bound
$O(\Esync(M))$ follows.

No continuation term appears, since objective~(RI) already fixes the
continuation function from the current history.

\emph{Part~(II).}  By Lemma~\ref{lem:operational-equivalence}, any
learner that eventually guarantees correct prediction on all continuations must
first attain objective~(RI).  By
Definition~\ref{def:sync-mistake-complexity} the infimum over learners of the
worst-case number of mistakes before that point is $\Esync(M)$, so every
learner incurs at least $\Esync(M)$ mistakes in the worst case:
\[
\MistRI(M)\ \ge\ \Omega(\Esync(M)).
\]
This uses no direct-sum hypothesis and yields \emph{only} the $\Esync$ term.

\emph{Part~(III).}  Assume the direct-sum saturation condition of
Definition~\ref{def:direct-sum-active}.  The Littlestone lower bound is proved
against adversarial inputs and does not transfer automatically to the active
protocol (Remark~\ref{rem:no-automatic-active-lower}), so each component's
cost must be supplied by that hypothesis rather than by counting.

The bound is $S_M+C_M$, not $\max\{S_M,C_M\}$, and the additivity comes from
disjointness rather than from any inequality between maxima and sums.  By
condition~(iv) the objective forces both components; by condition~(iii) the
rounds on which they force mistakes are disjoint; and by
Theorem~\ref{thm:active-direct-sum}, whose restriction argument uses
conditions~(v) and~(vi), each component contributes at least its own minimax
complexity on those rounds.  Summing over the two disjoint round sets gives
\[
\MistRI(M)\ \ge\ S_M+C_M .
\]
No matching upper bound is claimed: Part~(I) bounds $\MistRI(M)$ by
$O(\Esync(\mathcal C_M))$, and identifying $\Esync(\mathcal C_M)$ with
$S_M+C_M$ would require an explicit composite learner, which is not
constructed here.
\end{proof}

\begin{proposition}[Length-Based Active Upper Bound]
\label{prop:active-length-upper}
For fixed finite alphabets --- the hypothesis $\Lsyncu(M)<\infty$ being
automatic by Proposition~\ref{prop:lsyncu-version-space} ---
\[
\MistRI(M)
\le
O\bigl(M\log M+\Lsyncu(M)\bigr),
\]
and in particular
\[
\Esync(M)\ \le\ O\bigl(M\log M+\Lsyncu(M)\bigr).
\]
The universal depth, not the machine-specific $\Lsync(M)$, is the correct
quantity here: a learner facing an unknown machine must fix its experiment in
advance, so it cannot in general realize $L_{\mathrm{sync}}^{\mathrm{adapt}}(A)$
for the particular $A$ it faces.  The two terms are not redundant: $\Lsyncu(M)$
pays for state identification and the $M\log M$ term for the subsequent model
search, these being the two things a universal experiment cannot deliver at
once (Definition~\ref{def:output-aware-sync}).  If the machine is known except
for its initial state, the model-search phase is vacuous and the bound reads
$O(\Lsync(M))$.

In order of magnitude this bound is subsumed by
Theorem~\ref{thm:active-halving}, which attains the same objective in
$O(M\log M)$ mistakes without any synchronization experiment; even for
$\Lsyncu(M)=o(M\log M)$ the total here remains $O(M\log M)$.  Its interest is
structural.  It separates the contribution of the synchronization experiment
from that of the model search, and its first phase is governed by an experiment
fixable in advance, which is the operative quantity when the model-search phase
is externally curtailed.
No lower bound is asserted here; see
Remark~\ref{rem:no-automatic-active-lower}.
\end{proposition}

\begin{proof}
For the upper bound, the learner fixes in advance a universal adaptive
synchronization tree of worst-case depth $\Lsyncu(M)$, as in
Definition~\ref{def:output-aware-sync}; this requires no knowledge of the
target machine.  During this phase it makes at most $\Lsyncu(M)$ prediction
mistakes, one per step in the worst case.  By
Definition~\ref{def:output-aware-sync} the observed transcript then determines
the current state of the target within each machine still consistent with it;
the tables, and hence the residual class, are not yet determined, so
objective~(RI) has not been attained and a second phase is required.

The learner then runs the standard halving algorithm over the set of machines
in $\mathcal H_M$ consistent with the transcript so far, of cardinality at
most $|\mathcal H_M|$.  Each mistake halves that set, so at most
$\log_2|\mathcal H_M|=O(M\log M)$ further mistakes occur before a single
residual class remains, at which point objective~(RI) is attained and, being
prediction-closing, no mistake occurs afterwards.  The two phases together
give
\[
O(M\log M+\Lsyncu(M)).
\]
Here the $M\log M$ term is part of the cost of \emph{attaining}~(RI), not a
continuation term appended after attainment; this is the distinction drawn in
Remark~\ref{rem:objective-bookkeeping}.

For the second display, the composite strategy just described is one
admissible learner attaining~(RI), and it is implementable without knowing the
target.  Since $\Esync(M)$ is an infimum over such learners
(Definition~\ref{def:sync-mistake-complexity}), the bound
$\Esync(M)\le O(M\log M+\Lsyncu(M))$ follows; when
$\Lsyncu(M)=\Omega(M\log M)$ this reads $\Esync(M)=O(\Lsyncu(M))$.

\end{proof}

\begin{proposition}[Machine-Specific Synchronization Depth Is Quadratic]
\label{prop:lsyncu-quadratic}
For fixed finite alphabets and every $M\ge2$,
\[
\Lsync(M)\ \le\ (M-1)^2\ =\ O(M^2).
\]
The bound is on the \emph{machine-specific} depth $\Lsync$, in which the
experiment may depend on the target; see
Remark~\ref{rem:lsync-not-lsyncu} for why it does not transfer to the
universal depth $\Lsyncu$.
\end{proposition}

\begin{proof}
Fix $A\in\mathcal H_M$, minimal without loss of generality by the convention of
Definition~\ref{def:output-aware-sync}, and let $U\subseteq Q_A$ be the set of
states consistent with the transcript so far, initially $U=Q_A$ with
$|U|\le M$.  We exhibit a strategy, fixable in advance, that drives $U$ to a
single observational class.

\emph{Step: one separation episode.}  Suppose $U$ contains two states $s\ne t$
that are not observationally equivalent.  Consider the pair automaton on
unordered pairs of states of $A$: from $\{u,v\}$ and input $x$ there is an edge
to $\{\tau(u,x),\tau(v,x)\}$, and $\{u,v\}$ is \emph{immediately separated} by
$x$ when $\lambda(u,x)\ne\lambda(v,x)$.  Breadth-first search from $\{s,t\}$
reaches an immediately separated pair, since $s\not\sim t$; the pair automaton
has at most $\binom{M}{2}$ nodes, but the standard argument bounds the distance
by $M-1$, because a shortest such path visits pairwise distinct pairs whose
underlying state sets are nested along the search and cannot repeat.  Let $w$
be the corresponding word, $|w|\le M-1$.

Feed $w$.  Transitions are deterministic, so applying an input maps $U$ forward
and $|U|$ never increases.  At the final letter of $w$ the two survivors
descended from $s$ and $t$ emit different outputs, so whichever output is
observed at least one of the two branches is eliminated and $|U|$ strictly
decreases.

\emph{Aggregation.}  Each episode costs at most $M-1$ inputs and strictly
decreases $|U|$.  Since $|U|$ starts at most $M$ and terminates when all
surviving states are observationally equivalent, at most $M-1$ episodes occur,
giving total depth at most $(M-1)^2$.

The strategy runs breadth-first search on the pair automaton \emph{of the
target $A$}, so it presupposes knowledge of $A$ and is admissible for
$\Lsync(A)$.  Taking the supremum over $A\in\mathcal H_M$ gives the stated
bound on $\Lsync(M)$.
\end{proof}

\begin{proposition}[Universal Synchronization Depth, Version-Space Form]
\label{prop:lsyncu-version-space}
Let
\[
N_M=\bigl|\{(A,q):A\in\mathcal H_M,\ q\in Q_A\}\bigr|\ \le\ M\,|\mathcal H_M| .
\]
Then
\[
\Lsync(M)\ \le\ \Lsyncu(M)\ \le\ (M-1)(N_M-1)\ <\ \infty .
\]
In particular $\Lsyncu(M)$ is finite for fixed finite alphabets and fixed $M$,
so the hypothesis $\Lsyncu(M)<\infty$ of
Proposition~\ref{prop:active-length-upper} is automatic.
\end{proposition}

\begin{proof}
Maintain the version space $V\subseteq\{(A,q):A\in\mathcal H_M,\ q\in Q_A\}$ of
machine--state pairs consistent with the transcript so far, initially all of
them.  The whole strategy is a function of the transcript, since $V$ is, and it
is therefore a single decision tree in the sense of
Definition~\ref{def:output-aware-sync}; this is the point on which the
machine-specific bounds of Proposition~\ref{prop:lsyncu-quadratic} fail and
this one does not.

While some machine $A$ occurring in $V$ carries two candidates $(A,q)$ and
$(A,q')$ in $V$ whose states are not observationally equivalent in $A$, select
one such triple.  The learner may compute inside $A$ because $A$ is a
\emph{candidate}, named by the version space, not the unknown target.  Both $q$
and $q'$ are states of the same $M$-state machine, so
Lemma~\ref{lem:moore-separation} applied within $A$ --- rather than to a
disjoint union of two machines --- supplies a word $w$ with
$\lambda_A(q,w)\ne\lambda_A(q',w)$ and $|w|\le M-1$.

Feed $w$ and delete from $V$ every candidate whose predicted output sequence
disagrees with the observation, advancing the survivors along $w$.  At least
one of $(A,q)$, $(A,q')$ is deleted: their predicted sequences differ, so at
most one can agree with whatever is observed.  Hence $|V|$ strictly decreases
in each episode.  Starting from $|V|=N_M$ and terminating when no consistent
machine carries two observationally distinct candidate states, at most
$N_M-1$ episodes occur, each of length at most $M-1$.

At termination the transcript determines the current state of the target within
every machine still consistent with it, which is precisely the objective of
Definition~\ref{def:output-aware-sync}.
\end{proof}

\begin{remark}[Finiteness Is Now Unconditional; the Rate Is Not]
\label{rem:lsyncu-finite-not-rate}
Proposition~\ref{prop:lsyncu-version-space} settles finiteness, which
Remark~\ref{rem:lsync-not-lsyncu} had left open once $|\mathcal I|\ge2$, but it
does not settle the rate.  The bound is polynomial in the number of candidates
$N_M\le M|\mathcal H_M|$ and hence exponential in $M$, since
$|\mathcal H_M|$ is itself exponential; it is therefore far weaker in $M$ than
the machine-specific $\binom{M}{2}$.  Whether $\Lsyncu(M)$ is polynomial in $M$
is Conjecture~\ref{conj:lsyncu-poly}, and whether it is $O(M\log M)$ is the
surviving branch of Open Problem~\ref{open:si-hard-family}.
\end{remark}

\begin{conjecture}[Polynomial Universal Depth]
\label{conj:lsyncu-poly}
$\Lsyncu(M)=O(M^2)$ for fixed finite alphabets.
\end{conjecture}

\begin{remark}[The Machine-Specific Bounds Are on $\Lsync$, Not on $\Lsyncu$]
\label{rem:lsync-not-lsyncu}
Propositions~\ref{prop:lsyncu-quadratic} and~\ref{prop:lsyncu-binomial} bound
the machine-specific depth, in which the separating word is chosen by searching
the pair automaton of the target.  A learner realizing $\Lsyncu$ must fix
\emph{one} decision tree in advance and run it against every member of
$\mathcal H_M$, so it cannot perform that search, and the bounds do not
transfer.

The two quantities genuinely differ, already at $M=2$.  Over
$\mathcal I=\{0,1\}$ and $\mathcal O=\{0,1\}$ take
\[
A_1:\ \lambda(s,0)=s,\ \lambda(s,1)=0,
\qquad
A_2:\ \lambda(s,0)=0,\ \lambda(s,1)=s,
\]
both with $\tau(s,x)=s$.  Each is minimal and each is synchronized by a single
input --- input $0$ for $A_1$, input $1$ for $A_2$ --- so
$\Lsync=1$ on the pair.  A universal tree must commit to its first input before
learning which machine it faces; whichever it plays, the other machine emits
the constant output $0$ and nothing is learned, so a second input is required
and $\Lsyncu=2$ on the same pair.

Consequently the finiteness of $\Lsyncu(M)$, and its rate, are \emph{not}
settled by these propositions.  Finiteness is settled separately, by the
version-space strategy of Proposition~\ref{prop:lsyncu-version-space}, whose
input choices depend on the transcript alone; the rate remains open.
\end{remark}

\begin{lemma}[Separation--Size Tension]
\label{lem:tension}
Let $A$ be a minimal machine with $M$ states and let $U\subseteq Q_A$ with
$|U|\ge2$.  Write
\[
d(U)=\min_{s\ne t\in U}\ \min\{|w| : \lambda(s,w)\ne\lambda(t,w)\}
\]
for the length of the shortest word separating some pair inside $U$.  Then
\[
d(U)\ \le\ M-|U|+1 .
\]
\end{lemma}

\begin{proof}
Write $\sim_k$ for the Moore partition after $k$ rounds, so $\sim_0$ is the
single block $Q_A$ and $\sim_{k+1}$ refines $\sim_k$ by the value of
$\lambda(\cdot,x)$ together with the $\sim_k$-class of $\tau(\cdot,x)$, over
all $x\in\mathcal I$.  By construction $s\sim_k t$ exactly when no word of
length at most $k$ separates $s$ from $t$.

Put $k=d(U)$.  No word of length $k-1$ separates any pair inside $U$, so all
of $U$ lies in a single block of $\sim_{k-1}$.

The refinement is strictly increasing until it stabilizes: if a round fails to
split then the partition is fixed forever, and minimality forces the stable
partition to be discrete.  Since $k-1$ rounds have been taken and the process
has not yet stabilized --- some pair inside $U$ is still unseparated, while
minimality guarantees it will be --- the block count satisfies
\[
\bigl|Q_A/{\sim_{k-1}}\bigr|\ \ge\ k .
\]

On the other hand $U$ occupies one block and the remaining $M-|U|$ states
occupy at most $M-|U|$ further blocks, so
\[
\bigl|Q_A/{\sim_{k-1}}\bigr|\ \le\ 1+\bigl(M-|U|\bigr).
\]
Combining the two displays gives $k\le M-|U|+1$.
\end{proof}

\begin{proposition}[Machine-Specific Depth Is at Most $\binom{M}{2}$]
\label{prop:lsyncu-binomial}
For fixed finite alphabets and every $M\ge2$,
\[
\Lsync(M)\ \le\ \frac{M(M-1)}{2},
\]
The bound is attained at $M=3$ and $M=4$ --- there is a minimal $4$-state
machine over binary alphabets with $\Lsync(A)=6=\binom{4}{2}$ --- but not
beyond: see Remark~\ref{rem:attainment-sporadic}.
\end{proposition}

\begin{proof}
Run the strategy of Proposition~\ref{prop:lsyncu-quadratic}, but account for
each episode with Lemma~\ref{lem:tension} rather than with the crude bound
$M-1$.

Let $U$ be the current uncertainty set.  While $U$ contains two states that
are not observationally equivalent, feed a shortest word realizing $d(U)$.  As
in Proposition~\ref{prop:lsyncu-quadratic}, deterministic transitions never
increase $|U|$, and at the final letter of that word the observed output
eliminates at least one branch, so $|U|$ strictly decreases.  The episode costs
$d(U)\le M-|U|+1$ inputs by Lemma~\ref{lem:tension}.

The successive sizes are therefore at worst $M,M-1,\dots,2$, and the total is
at most
\[
\sum_{m=2}^{M}\bigl(M-m+1\bigr)
=
\sum_{j=1}^{M-1}j
=
\frac{M(M-1)}{2}.
\]
As in Proposition~\ref{prop:lsyncu-quadratic} the separating word is found by
searching the pair automaton of the target, so the bound is on $\Lsync$
(Remark~\ref{rem:lsync-not-lsyncu}).

Attainment at $M=3,4$ is a computational observation rather than a proved
lower bound: exhaustive search over all minimal machines with
$|\mathcal I|=|\mathcal O|=2$ returns maximum adaptive depth exactly $3$ and
$6$ respectively, the latter realized by $3072$ machines.  An explicit witness
at $M=4$ is
\[
\tau(\cdot,0)=(0,0,3,2),\quad
\tau(\cdot,1)=(0,2,3,1),\quad
\lambda(\cdot,0)=(0,1,0,0),\quad
\lambda(\cdot,1)\equiv0 ,
\]
on state set $\{0,1,2,3\}$; the accompanying search program reports its optimal
adaptive depth as $6$.  No hand proof of that lower bound is given here.
\end{proof}

\begin{remark}[Attainment Is Sporadic, and the Extremal Structure Is Linear]
\label{rem:attainment-sporadic}
The bound $\binom{M}{2}$ is attained at $M=3,4$ and, on the evidence, at no
larger $M$.  The distinction matters for
Open Problem~\ref{open:si-hard-family}, because attainment for all $M$ would
give $\Lsyncu(M)=\Theta(M^2)=\omega(M\log M)$ and settle it affirmatively.

Every machine realizing the bound at $M=3,4$ has the same shape.  One input is
a permutation fixing a sink state and cycling the others; the second input is
collapsing and carries the \emph{only} informative output, emitted from a
single \emph{probe} state.  The learner can therefore ask just one question ---
``is the true state at the probe now?'' --- and must rotate candidates past the
probe between questions.  Tracing the optimal play at $M=4$ makes the cost
explicit: the uncertainty set moves
$\{0,1,2,3\}\to\{0,2,3\}\to\{0,1,3\}\to\{0,2\}\to\{0,3\}\to\{0,1\}\to\{0\}$,
alternating probe and rotate.

That mechanism does not scale quadratically.  Generalizing it to arbitrary $M$
--- sink, an $(M-1)$-cycle, and a single probe --- gives adaptive depth exactly
$2M-2$, since each elimination costs one rotation step plus one probe step, and
the ratio to $M\log_2M$ \emph{decreases} in $M$.  Exhaustive search at $M=5$
over the $2{,}839{,}200$ minimal machines whose first input is a permutation
and whose output function has a single probe returns maximum depth $9$, short
of $\binom{5}{2}=10$; hill-climbing over unrestricted minimal machines returns
$9$, $14$ and $14$ at $M=6,7,8$ against $\binom{M}{2}=15$, $21$, $28$, a gap
that widens.

The evidence therefore points away from a quadratic family and towards
$\Lsyncu(M)=O(M)$ for the worst case, which would collapse the length currency
exactly as the mistake currency collapsed in
Remark~\ref{rem:additive-collapses}.  This is evidence, not proof: the search
at $M\ge6$ is heuristic, and the exhaustive $M=5$ search covers a structured
subclass rather than all minimal machines.
\end{remark}

\begin{proposition}[Single-Input Machines Are Linear]
\label{prop:lsyncu-single-input}
Let $|\mathcal I|=1$.  Then for every minimal $A\in\mathcal H_M$ the adaptive
synchronization depth satisfies
\[
L_{\mathrm{sync}}^{\mathrm{adapt}}(A)\ \le\ M-1 ,
\]
and consequently $\Lsync(M)=\Lsyncu(M)\le M-1$ for unary input alphabets.  In
particular $\Lsyncu(M)=O(M)=o(M\log M)$ there.
\end{proposition}

\begin{proof}
With a single input letter the learner has no choice of experiment: the only
strategy is to feed that letter repeatedly, so the adaptive tree degenerates to
a path and $\Lsync$ and $\Lsyncu$ coincide.  This is the one case in which the
distinction of Remark~\ref{rem:lsync-not-lsyncu} is vacuous, because there is
only one universal tree.

Write $a$ for the letter and define the equivalence relations
\[
s\sim_t u
\iff
\lambda(s,a^k)=\lambda(u,a^k)\ \text{ for all } k\le t ,
\]
so that $\sim_1$ is the partition of $Q_A$ by the value of $\lambda(\cdot,a)$
and $\sim_{t+1}$ refines $\sim_t$.  This is exactly Moore's partition
refinement for the unary automaton.

Two observations finish the argument.  First, the refinement is monotone and
\emph{stabilizes at the first round that fails to split}: if
$\sim_{t+1}=\sim_t$ then $\sim_{t'}=\sim_t$ for all $t'\ge t$, since the
refinement step is a function of the current partition alone.  Second,
minimality of $A$ means the stable partition is discrete, with $M$ blocks.

Let $R$ be the number of rounds until stabilization.  The block count is at
least $1$ initially, strictly increases at each of the first $R-1$ rounds by
the stabilization property, and equals $M$ at the end; hence $R\le M-1$.  After
$R$ applications of $a$ the observed output prefix determines the block of
$\sim_R$ containing the initial state, and that block is a singleton, so the
current state is determined.  Therefore
$L_{\mathrm{sync}}^{\mathrm{adapt}}(A)\le R\le M-1$.
\end{proof}

\begin{remark}[The Search for a Witness Needs a Nontrivial Input Alphabet]
\label{rem:witness-needs-two-inputs}
Proposition~\ref{prop:lsyncu-single-input} removes an entire alphabet class
from consideration: any family witnessing
$\Lsyncu(\mathcal C_M)=\omega(M\log M)$ must have $|\mathcal I|\ge2$.  The
reason is structural rather than quantitative.  With one letter there is no
adaptivity to exploit --- the transcript is a fixed function of the initial
state --- so the process is pure partition refinement, and refinement can
happen at most $M-1$ times because each effective round splits at least one
block of a partition of an $M$-element set.  Adaptive depth beyond $M-1$
therefore requires the learner's \emph{choice} of next input to matter, which
needs at least two letters.

Exhaustive computation supports this and locates the extremal behaviour
precisely.  Over all minimal single-input machines with $M\le6$ and binary
output --- $8$, $72$, $960$, $16{,}800$ and $362{,}880$ machines for
$M=2,\dots,6$ --- the depth equals the number of refinement rounds in every
instance, and the maximum is exactly $M-1$; the block count increases strictly
at each round, verified over $9{,}313{,}920$ minimal machines at $M=7$.
\end{remark}

\begin{remark}[Consequences, and What Remains Open]
\label{rem:lsyncu-consequences}
Two consequences follow, and one question is sharpened rather than settled.

First, the finiteness hypothesis in
Proposition~\ref{prop:active-length-upper} is not discharged by \emph{these}
bounds, since they concern $\Lsync$ rather than $\Lsyncu$
(Remark~\ref{rem:lsync-not-lsyncu}); it is discharged in general by
Proposition~\ref{prop:lsyncu-version-space}, and with a bound polynomial in
$M$ rather than in $N_M$ for unary input alphabets by
Proposition~\ref{prop:lsyncu-single-input}.

Second, the quadratic bound does \emph{not} by itself decide whether the
length-form bound is dominated by the halving bound of
Theorem~\ref{thm:active-halving}.  The two rates cross: even the sharpened
$\binom{M}{2}$ of Proposition~\ref{prop:lsyncu-binomial} exceeds $M\log_2M$ for
$M\ge7$, so an upper bound of order $M^2$ leaves room for a family with
$\Lsyncu(\mathcal C_M)=\omega(M\log M)$.  Proposition~\ref{prop:lsyncu-binomial} fixes the constant at $\tfrac12$ for
$\Lsync$, and the bound is achieved at $M=3,4$.  It is not achieved beyond
(Remark~\ref{rem:attainment-sporadic}), and the extremal mechanism at those two
values generalizes to a family of depth $2M-2$ rather than to a quadratic one,
so the evidence favours $\Lsync(M)=\Theta(M)$.  None of this bounds $\Lsyncu$,
by Remark~\ref{rem:lsync-not-lsyncu}.

What the evidence suggests is that no such family exists.  Exhaustive
computation over all minimal machines with $M\le5$ states and
$|\mathcal I|,|\mathcal O|\le3$ gives maximum adaptive depth
$M-1$ for single-input machines and never more than $\binom{M}{2}$; the natural
candidates --- cyclic counters with a single marked state, and the cyclic-shift
family of Theorem~\ref{thm:esyncsi-theta} --- realize $M-1$ and $\log_2M$
respectively, both $o(M\log M)$.  Hill-climbing on the ratio
$\Lsync(A)/(M\log_2M)$ over minimal machines up to $M=20$ never exceeds
$0.78$ and trends downward.  Whether $\Lsyncu(M)=O(M\log M)$ holds in general
remains open (Open Problem~\ref{open:si-hard-family}); what is now settled is
that the answer lies between $M-1$ and $(M-1)^2$, so the question is a
question about the exponent and not about finiteness.
\end{remark}

\begin{lemma}[Separation Across Two Mealy Machines]
\label{lem:moore-separation}
Let $A_1,A_2$ be deterministic Mealy machines over the same alphabets, with
$N$ states in total, and let $q_1,q_2$ be states of $A_1,A_2$ respectively
whose continuation functions differ.  Then some word of length at most $N-1$
produces different output sequences from $q_1$ and $q_2$.
\end{lemma}

\begin{proof}
Work in the disjoint union $A_1\sqcup A_2$, a single Mealy machine with $N$
states.  For $k\ge0$ let $\equiv_k$ identify two states when their output
sequences agree on every input word of length at most $k$.  Each $\equiv_k$ is
an equivalence relation, $\equiv_{k+1}$ refines $\equiv_k$, and the
refinement is the standard Moore construction
\cite[Ch.~II]{sakarovitch2009}.

If $\equiv_{k+1}=\equiv_k$ for some $k$, then the sequence is stationary from
$k$ onwards, since $\equiv_{k+2}$ is determined by $\equiv_{k+1}$ together with
the one-step transition and output maps.  Hence the relations strictly refine
until they stabilize, and $\equiv_0$ has at least one block, so stabilization
occurs at some $k^\ast\le N-1$: each strict refinement increases the block
count by at least one, and the count cannot exceed $N$.

By hypothesis $q_1\not\equiv_\infty q_2$, so $q_1\not\equiv_{k^\ast}q_2$, and a
word of length at most $k^\ast\le N-1$ witnesses the difference.
\end{proof}

\begin{theorem}[Unconditional Active Attainment Bound]
\label{thm:active-halving}
Fix finite alphabets $\mathcal I,\mathcal O$ with $|\mathcal I|,|\mathcal O|\ge2$.
There is an active learner that, on every $A\in\mathcal H_M$ and every initial
state, attains objective~(RI) of Definition~\ref{def:residual-knowledge} after
at most
\[
\log_2\bigl|\mathcal H_M\times Q\bigr|
=
O(M\log M)
\]
prediction mistakes, the implied constant depending only on the alphabets.
Consequently
\[
\Esync(M)\ =\ O(M\log M),
\]
with no synchronizability, operational-equivalence, or direct-sum hypothesis.
\end{theorem}

\begin{proof}
Let $V$ be the version space: the set of pairs $(A,q)$ with
$A\in\mathcal H_M$ and $q$ a state of $A$, consistent with the transcript so
far, where the second coordinate is advanced along the observed inputs.
Initially $|V|=|\mathcal H_M\times Q|$, which is finite.

Call two elements of $V$ \emph{residually distinct} if their continuation
functions differ, i.e.\ if some input word separates their output sequences.
Two residually distinct elements $(A,q)$ and $(B,p)$ of $V$ are separated by a
word of length at most $2M-1$, by Lemma~\ref{lem:moore-separation} applied to
the disjoint union of $A$ and $B$, which has at most $2M$ states.  The familiar
$M-1$ figure is not available here, since it separates two states of a
\emph{single} $M$-state machine whereas the two elements of $V$ may come from
different machines; the disjoint-union bound $2M-1$ is the correct one.  The
exact length is in any case immaterial to the mistake count, which is driven by
halving and not by word length; only finiteness is used.

The learner repeats the following until all elements of $V$ share one
continuation function.  Choose two residually distinct elements of $V$ and a
shortest word $w$ separating them.  Play $w$ one letter at a time; before each
letter, predict the output symbol receiving a plurality of votes among the
elements of $V$; after observing the true symbol, delete from $V$ every element
predicting a different symbol, and advance the states of the survivors.

Two facts complete the argument.  First, each mistake at least halves $|V|$.
Fix the input $x$ played at the current step.  Every element $(A,q)$ of the
version space is a \emph{deterministic} machine together with a state, so it
emits the single symbol $\lambda_A(q,x)$; the version space therefore
partitions into output classes of sizes $c_1\ge c_2\ge\dots$, and the learner
predicts the plurality symbol.  If the prediction is wrong, the observed symbol
identifies one class other than the plurality class, and the survivors are
\emph{that single class}, of size at most $c_2$ --- not the union of all
non-predicted classes.  Since $c_1\ge c_2$ and $c_1+c_2\le|V|$,
\[
2c_2\ \le\ c_1+c_2\ \le\ |V|,
\qquad\text{so}\qquad
c_2\ \le\ \tfrac12|V| .
\]
This is the same counting as in Proposition~\ref{prop:esyncsi-log}, and it
applies verbatim here: what it needs is that each surviving candidate predicts
one deterministic symbol, which holds for machine--state pairs exactly as it
does for states of a known machine.  Consequently the number of mistakes is at
most $\bigl\lfloor\log_2|\mathcal H_M\times Q|\bigr\rfloor$, with no dependence
on $|\mathcal O|$.  Second, the procedure terminates:
playing a separating word $w$ deletes at least one of the two chosen elements,
so $|V|$ strictly decreases on each iteration, and $V$ is finite and never
empty since the realizable target is always retained.

At termination all surviving elements have the same continuation function, and
the target is among them, so the continuation function from the current history
is determined by the transcript: objective~(RI) is attained.  Being
prediction-closing, it incurs no further mistakes.  Finally
$\log_2|\mathcal H_M\times Q|=O\bigl(M|\mathcal I|\log(M|\mathcal O|)\bigr)
=O(M\log M)$ for fixed alphabets, and $\Esync(M)$ is an infimum over learners
by Definition~\ref{def:sync-mistake-complexity}, so
$\Esync(M)\le\log_2|\mathcal H_M\times Q|=O(M\log M)$.
\end{proof}

\begin{remark}[What Becomes Unconditional, and What Does Not]
\label{rem:active-unconditional}
Corollary~\ref{cor:active-theta} settles the \emph{order} of the active
realizable mistake complexity without any of the auxiliary hypotheses used
earlier in this subsection.  In particular
Lemma~\ref{lem:operational-equivalence} and the direct-sum saturation
condition are not needed for the $\Theta(M\log M)$ statement, and the
finiteness of $\Lsyncu(M)$ is not needed either: the halving learner of
Theorem~\ref{thm:active-halving} never performs a synchronization experiment
and is well defined even when no machine in $\mathcal H_M$ admits one.

What those hypotheses still buy is finer information.  Parts~(II) and~(III) of
Theorem~\ref{thm:active-certified} are statements about $\Esync(M)$ as an
abstract complexity, and remain the right form when $\Esync$ is replaced by the
corresponding quantity for a \emph{subclass} of $\mathcal H_M$, where the
counting bound $\log_2|\mathcal H_M|$ need not be $O(M\log M)$ and the two
terms of the additive form can separate.  The additive form is therefore not
vacuous, but on the full class $\mathcal H_M$ it collapses:
$M\log M+\Esync(M)=\Theta(\Esync(M))$.
\end{remark}

\begin{remark}[No Automatic Active Lower Bound]
\label{rem:no-automatic-active-lower}
An active lower bound does not follow from the passive theory, and must be
established by an explicit adversarial construction.  Two standard arguments
fail here.  First, ordinary Littlestone lower bounds are proved against an
\emph{adversarial} input sequence; when the learner chooses the inputs, the
adversary no longer controls the instance selection game, so those
constructions do not transfer.  Second, a counting argument of the form ``each
mistake conveys at most a constant amount of information'' is invalid, because
a \emph{correctly predicted} output also reveals information about the target
and costs the learner nothing.

The obstruction is to those two methods, not to the conclusion.  The gated
family of Definition~\ref{def:gated-active-family} supplies an explicit
construction and yields the unconditional bound $\Omega(M\log M)$ of
Corollary~\ref{cor:active-theta}.  What requires a further hypothesis is the
\emph{separation} of the two terms in the additive form, which uses the
direct-sum saturation condition of Definition~\ref{def:direct-sum-active}.
\end{remark}

\begin{definition}[Gated Active Family $\Gact_M$]
\label{def:gated-active-family}
Let $Q=\{0,1\}^L$, $M=2^L$, and let the state set be $Q\times\{\free,\rd\}$, of
size $2M$.  Write $\rot(v)=v_2\cdots v_Lv_1$ and let $\flipl(v)$ flip the last
bit of $v$.  For an arbitrary map $g:Q\to Q$ the machine $A_g\in\Gact_M$ has
transitions and Mealy outputs
\[
\begin{array}{lll}
(\free,v): &
\mathtt r\mapsto(\free,0^L)/0, \quad
\mathtt e\mapsto(\free,\flipl(v))/0, &
\mathtt d\mapsto(\free,\rot(v))/v_1, \quad
\mathtt c\mapsto(\rd,g(v))/0,\\[3pt]
(\rd,u): &
\mathtt r\mapsto(\free,0^L)/0, \quad
\mathtt e\mapsto(\rd,u)/0, &
\mathtt d\mapsto(\rd,\rot(u))/u_1, \quad
\mathtt c\mapsto(\rd,u)/0 .
\end{array}
\]
The single structural change from the family of
Theorem~\ref{thm:stream-lower-bound} is the last entry: in read mode
$\mathtt c$ is a \emph{no-op}.  The mode is visible in the output alphabet if
desired; this only helps the learner.  Only $\mathtt c$ from free mode depends
on $g$, so $|\Gact_M|=M^M$ and $\Gact_M\subseteq\mathcal H_{2M}$.

The \emph{gating property} is the following.  Since read mode is entered only
by playing $\mathtt c$ from free mode, and $\mathtt c$ does nothing thereafter,
the register in read mode is always a cyclic rotation of $g(v)$ for the single
argument $v$ that was the free-mode state when $\mathtt c$ was played.  That
argument was reached by the learner's own inputs from the fixed state $0^L$
using only the $g$-independent letters $\mathtt r,\mathtt e,\mathtt d$, and is
therefore known to the learner.  Consequently every $g$-dependent output round
reports a bit of $g$ at a learner-known argument, and the set of maps
consistent with any transcript is exactly a product
\[
\prod_{v\in Q}\ S_v,
\qquad
S_v\subseteq Q,
\]
of independent per-argument constraints.
\end{definition}

\begin{theorem}[Gated Family and Unconditional Active Lower Bound]
\label{thm:active-explicit-directsum}
Let $\mathcal I=\{\mathtt r,\mathtt e,\mathtt d,\mathtt c\}$ and
$\mathcal O=\{0,1\}$, and let $\Gact_{M}$ be the \emph{gated} family of
Definition~\ref{def:gated-active-family} on $2M$ states, $M=2^L$.  Consider an
active learner that chooses all inputs and is required to attain
objective~(RI) of Definition~\ref{def:residual-knowledge}.  Then
every \emph{deterministic} such learner incurs at least
\[
M\log_2M
=
\Omega(M\log M)
\]
mistakes in the worst case over $\Gact_M$.  For every \emph{randomized} such
learner there is a fixed target in $\Gact_M$ attaining the bound below, by
averaging over the uniform target distribution; explicitly, every randomized
learner satisfies
\[
\mathbb E[\text{mistakes}]
\ \ge\
\tfrac12M\log_2M
=
\Omega(M\log M),
\]
and, under the uniform target distribution on $Q^Q$, for every
$\eta\in(0,1)$,
\[
\Pr\left[
\text{mistakes}\ \ge\ \bigl(\tfrac12-\eta\bigr)M\log_2M
\right]
\ \ge\
1-\exp\!\left(-\tfrac12\eta^2M\log_2M\right).
\]
This family supplies \emph{one} component; the two-component direct-sum
instance required by Definition~\ref{def:direct-sum-active} is assembled from
two independent copies in Theorem~\ref{thm:active-two-phase}.  The unknown
datum here is the map $g:Q\to Q$, so what is forced is \emph{model}
identification, not state identification within a known skeleton; accordingly
$\Gact_M$ realizes a component of complexity
\[
S_M=\Omega(M\log M)
\]
against objective~(RI), and not the state-identification quantity $\EsyncSI$
of Definition~\ref{def:sync-mistake-complexity},
and $\Esync(M)=\Omega(M\log M)$ on this family.  A two-phase instance,
exhibiting a separate continuation cost $C_M$, is constructed in
Theorem~\ref{thm:active-two-phase}.

Gating doubles the state count, so $\Gact_M$ lies in $\mathcal H_{N}$ with
$N=2M$.  The bound is unaffected in order: writing it against the true budget
$N$ gives
\[
\tfrac{N}{2}\log_2\tfrac{N}{2}
=
\Omega(N\log N),
\]
since $\log_2(N/2)=\log_2N-1\ge\tfrac12\log_2N$ for $N\ge4$.  All statements
below are therefore valid verbatim with $M$ read as the state budget, at the
cost of absolute constants only.
\end{theorem}

\begin{proof}
Throughout, the learner chooses every input; the adversary answers.

\textbf{Step 1: attaining~(RI) requires knowing $g$ everywhere.}
Suppose $g(v)$ is undetermined by the transcript for some $v\in Q$.  From the
current state the learner may play $\mathtt r\,w_v\,\mathtt c\,\mathtt d^{L}$,
where $w_v$ is the transport word of Theorem~\ref{thm:stream-lower-bound};
this reaches $(\free,v)$ using only $g$-independent letters, enters read mode
at $(\rd,g(v))$, and reads out all $L$ bits of $g(v)$.  Two members of
$\Gact_M$ consistent with the transcript disagree on that continuation, so the
continuation function is not determined and~(RI) has not been attained.  Since
$|\Gact_M|=M^M$, attaining~(RI) requires fixing all $M\log_2M$ bits of $g$.

\textbf{Step 2: every $g$-dependent round is a readout at a known argument.}
The letters $\mathtt r,\mathtt e,\mathtt d$ have $g$-independent transitions
and outputs in both modes, and $\mathtt c$ emits the constant $0$ in both
modes, so no round using them at a $g$-independent state reveals anything
about $g$.  The only $g$-dependent outputs are the $\mathtt d$ rounds in read
mode, each emitting one bit of the register.  By the gating property of
Definition~\ref{def:gated-active-family} the register is a rotation of $g(v)$
for a single argument $v$ determined by the learner's own inputs, hence known
to it.  This is the only step at which gating is used, and
Remark~\ref{rem:gating-needed} shows it fails without it.

\textbf{Step 3: the lazy adversary is well defined.}
Maintain a partial assignment recording, for each $v\in Q$ and each
coordinate $j\le L$, either a fixed bit or the symbol $\ast$.  When the learner
plays $\mathtt d$ in read mode with register $\rot^{\,o}(g(v))$ and coordinate
$j=o\bmod L$ still $\ast$, the adversary answers with the complement of the
learner's prediction and records that bit.  The answer is legal precisely
because, by Step~2, the consistent set is at every moment the product
$\prod_{v}S_v$ of per-argument constraints, so fixing one coordinate of $g(v)$
leaves every completion at every other argument available; no earlier round
constrained it.  Each of the $M\log_2M$ bits is therefore revealed on a round
that is a forced mistake, and no bit is ever revealed on a correctly predicted
round.  Rounds where the coordinate is already recorded are predicted
correctly and cost nothing, so no bit is paid for twice.

By Step~1 the learner must reach the all-recorded assignment, so it incurs at
least $M\log_2M$ mistakes.  The learner's control of the input sequence
determines only \emph{when} each bit is paid for, not \emph{whether}.

The argument does not proceed by counting information per mistake, which would
be invalid in the active protocol because correctly predicted rounds may also
be informative (Remark~\ref{rem:no-automatic-active-lower}).  It rests instead
on the structural property of Step~2: in this family every
information-bearing round is a forced mistake.

\textbf{Step 4: randomized learners.}
Here the fixed-stream form of Yao's principle used in
Theorem~\ref{thm:stream-lower-bound}(ii) is \emph{not} available, since an
active learner chooses its own inputs and a fixed stream lower-bounds only
those learners that happen to follow it.  We argue directly on the adaptive
transcript.

We first fix the accounting convention for learners that fail.  By
Definition~\ref{def:sync-mistake-complexity} a learner that never attains~(RI)
is charged $\mistk=+\infty$, so such learners cannot lower the infimum and may
be excluded; equivalently, the bound below is asserted for learners attaining
the objective almost surely, and is vacuously true otherwise.  Without such a
convention no unconditional expectation bound could hold, since a learner could
attain~(RI) only on an event of small probability and idle otherwise.

Draw $g$ uniformly from $Q^Q$; equivalently the $M$ values $g(v)$ are
independent and uniform on $Q$, so the $M\log_2M$ bits of $g$ are
i.i.d.\ fair.  Fix any randomized active learner attaining~(RI) almost surely,
and let $\mathcal F_{t-1}$ be
the $\sigma$-field generated by the transcript through round $t-1$ together
with the learner's internal randomness.  Call round $t$ a \emph{first
emission} if it is a read-mode $\mathtt d$ round whose coordinate $(v,j)$ has
not been emitted at any earlier round, and let $N$ be the set of
first-emission rounds.

By Step~2 the coordinate $(v,j)$ read at round $t$ is
$\mathcal F_{t-1}$-measurable, since the argument $v$ is determined by the
learner's own inputs and $j$ by the number of read-mode shifts since
$\mathtt c$.  By the product structure of Step~2, the transcript through round
$t-1$ constrains only coordinates already emitted, and $(v,j)$ is not among
them; hence $g(v)_j$ is independent of $\mathcal F_{t-1}$ and uniform on
$\{0,1\}$.  The prediction $\hat y_t$ is $\mathcal F_{t-1}$-measurable, so on
the event $\{t\in N\}$,
\[
\Pr\bigl[\hat y_t\ne g(v)_j \,\big|\, \mathcal F_{t-1}\bigr]=\tfrac12 .
\]
By Step~1, attaining~(RI) forces every one of the $m=M\log_2M$ coordinates to
be emitted, so $N$ has exactly $m$ elements.  The rounds of $N$ occur at random
times, so we index them in order of occurrence rather than by $t$: let
$\sigma_1<\sigma_2<\cdots<\sigma_m$ enumerate $N$, each $\sigma_k$ being a
stopping time with respect to $(\mathcal F_t)$, and set
$\mathcal G_k=\mathcal F_{\sigma_k-1}$.  The $\sigma$-fields
$\mathcal G_1\subseteq\cdots\subseteq\mathcal G_m$ form a filtration by the
optional stopping theorem for discrete filtrations.  Put
\[
Y_k
=
\ind{\hat y_{\sigma_k}\ne g(v)_j\text{ at round }\sigma_k}-\tfrac12 ,
\qquad k=1,\dots,m .
\]
By the display above, $\mathbb E[Y_k\mid\mathcal G_k]=0$ and $|Y_k|\le\tfrac12$,
so $(Y_k)_{k\le m}$ is a bounded martingale difference sequence indexed by a
\emph{deterministic} range $1,\dots,m$.  The total number of mistakes is at
least $\sum_{k=1}^mY_k+\tfrac m2$, so taking expectations gives
$\mathbb E[\text{mistakes}]\ge\tfrac m2=\tfrac12M\log_2M$.  Azuma--Hoeffding
applied to $\sum_{k=1}^mY_k$ \cite[Ch.~6]{boucheron2013} gives, for every
$\eta\in(0,1)$,
\[
\Pr\left[\text{mistakes}<\bigl(\tfrac12-\eta\bigr)M\log_2M\right]
\le
\exp\!\left(-\tfrac12\eta^2M\log_2M\right),
\]
which is the stated high-probability form.  No independence among rounds is
assumed, and no input stream is fixed in advance.

\textbf{Step 5: disjoint accounting.}
The readout rounds carrying $g$-information occur in read mode, and the
transport rounds $w_v$ occur in free mode, so the two sets of rounds are
disjoint.  This is condition~(iii) of
Definition~\ref{def:direct-sum-active}.
\end{proof}

\begin{remark}[Why the Passive Family Must Be Gated]
\label{rem:gating-needed}
Gating is what separates the active setting from the passive one.  In the
passive protocol the adversary fixes the input stream, and in the stream of
Theorem~\ref{thm:stream-lower-bound} the letter $\mathtt c$ is always issued
from a state $v$ that the transport word $w_v$ has just established, so every
readout reports $g$ at a known argument.  An active learner is under no such
discipline: playing $\mathtt c$ twice in succession moves the machine to
$g(g(v))$, and a readout there constrains $g$ at an argument that is itself
undetermined.

Under such chaining the set of maps consistent with a transcript is not a
product of independent per-argument constraints.  For $L=2$ a single chained
readout $g(g(0^L))=u$ leaves $48$ of the $256$ maps consistent, whereas the
product of the per-argument projections has $3\cdot4\cdot4\cdot4=192$
elements.  Both the per-bit adversary, which fixes one bit of one $g(v)$ per
readout round and requires every completion to remain available, and the
conditional-uniformity step of the Yao argument, which requires each queried
bit to be fair given the transcript, depend on that product structure.

Making $\mathtt c$ a no-op in read mode, as in
Definition~\ref{def:gated-active-family}, confines every readout argument to
the free-mode state the learner transported to, at the cost of doubling the
state count.  The product structure holds throughout, and with it both
arguments.
\end{remark}

\begin{corollary}[Active Realizable Mistake Complexity, Unconditionally]
\label{cor:active-theta}
For fixed alphabets with $|\mathcal I|\ge4$ and $|\mathcal O|\ge2$, and for
\emph{all} sufficiently large $M$,
\[
\Esync(M)=\Theta(M\log M)
\qquad\text{and}\qquad
\MistRI(M)=\Theta(M\log M),
\]
with no auxiliary hypothesis.  The same holds for every
$|\mathcal I|\ge2$ by the block encoding of
Theorem~\ref{thm:stream-lb-binary}.
\end{corollary}

\begin{proof}
The upper bound $\Esync(M)=O(M\log M)$ is Theorem~\ref{thm:active-halving},
and it holds for every $M$.

For the lower bound, first note that the gated family is defined only for
$M'=2^L$ and occupies $2M'$ states, so it yields a bound along the subsequence
\[
N_L=2^{L+1},
\qquad L=1,2,\dots,
\]
and not immediately at every budget.  Along that subsequence,
$\Gact_{2^L}\subseteq\mathcal H_{N_L}$ forces every active learner
attaining~(RI) to make at least
$2^L\log_22^L=\tfrac12N_L(\log_2N_L-1)$ mistakes by
Theorem~\ref{thm:active-explicit-directsum}.  Since $\Esync$ is an infimum over
learners of a supremum over the whole class, and $\Gact_{2^L}$ is one
subfamily, $\Esync(N_L)\ge\tfrac14N_L\log_2N_L$ for $N_L\ge4$.

To pass from the subsequence to all budgets, observe that
$\mathcal H_M\subseteq\mathcal H_{M+1}$, so $\Esync$ is nondecreasing in $M$,
and that $N_{L+1}/N_L=2$, a bounded ratio.  Lemma~\ref{lem:subsequence-allM},
applied with $\alpha=2$, therefore gives a constant $\beta'>0$ with
\[
\Esync(M)\ \ge\ \beta'M\log M
\]
for all sufficiently large $M$.  Combined with the upper bound,
$\Esync(M)=\Theta(M\log M)$ for all sufficiently large $M$.

For the mistake complexity, Theorem~\ref{thm:active-certified}(I) gives
$\MistRI(M)\le O(\Esync(M))=O(M\log M)$, and
the same family, extended by the identical monotonicity argument, gives the
matching lower bound, since attaining~(RI) is the objective.
\end{proof}


\begin{theorem}[Explicit Two-Phase Direct-Sum Instance]
\label{thm:active-two-phase}
There is a family of $O(M)$-state machines over fixed finite alphabets, and a
pair of switch letters $\mathtt s_1,\mathtt s_2$, realizing a direct-sum active
instance in the sense of Definition~\ref{def:direct-sum-active}: both
\[
S_M=\Omega(M\log M)
\qquad\text{and}\qquad
C_M=\Omega(M\log M)
\]
are forced, on disjoint sets of rounds.  Consequently
\[
\MistRI(M)
\ \ge\
S_M+C_M
=
\Omega(M\log M),
\]
Both components are model-identification components, each of complexity
$\Theta(M\log M)$, so their sum is again $\Theta(M\log M)$ and this family does
\emph{not} exhibit a separation between the two terms of the additive form.
What it establishes is that two independent unknown maps must each be paid for
on disjoint rounds.
\end{theorem}

\begin{proof}
Split the budget as $M=M_1+M_2+O(1)$ with $M_1\asymp M_2\asymp M$.

\emph{Component one.}  On the state block of size $M_1$, install the gated
transport-readout family $\Gact_{M_1}$ of
Definition~\ref{def:gated-active-family}, whose unknown datum is a map
$g:Q_1\to Q_1$.  Determining $g$ in full forces
$S_M=\Omega(M_1\log M_1)=\Omega(M\log M)$ mistakes, by
Theorem~\ref{thm:active-explicit-directsum}.

\emph{Switch.}  Adjoin \emph{two} letters $\mathtt s_1,\mathtt s_2$, where
$\mathtt s_j$ forces entry into block $j$ from anywhere:
\[
\tau(q,\mathtt s_j)=\iota_j
\quad\text{for every state }q,
\qquad
\lambda(q,\mathtt s_j)=0,
\]
$\iota_j$ being the distinguished initial state of block $j$.  Tag the output
alphabet so the learner can see which block is active; this only helps the
learner and does not weaken the bound.

Both letters are idempotent and total, so for each $j$ and each word $u$ the
composite $\mathtt s_j\,u$ reaches block $j$ and executes $u$ from $\iota_j$
irrespective of the current state.  The forcing argument below uses exactly
this: a single word, valid from every configuration, that drives the machine
into a prescribed block.  Totality is also what removes any evasion, since in
the active protocol the learner chooses its own inputs and block one remains
reachable at every time.

\emph{Component two.}  On the state block of size $M_2$, install a second
independent copy of the gated transport-readout family, with its own unknown
map $g':Q_2\to Q_2$.  Restricted to inputs that keep the machine in $Q_2$, the
interaction is exactly a fresh instance of
Theorem~\ref{thm:active-explicit-directsum} on $M_2$ states, forcing
$C_M=\Omega(M_2\log M_2)=\Omega(M\log M)$ mistakes.

\emph{Both components are forced.}  Suppose the learner has attained
objective~(RI) at some time $t$, so that it can predict every continuation with
no further mistakes.  Let $b$ be a bit of $g$, say the $j$-th bit of $g(v)$ for
some $v\in Q_1$, that is not determined by the transcript.  From the current
state -- \emph{whatever} it is, and in particular whichever block it lies in --
the learner may play
\[
\mathtt s_1\,w_v\,\mathtt c\,\mathtt d^{\,L_1},
\]
which reaches $\iota_1$ by the first letter, transports to $(\free,v)$ within
block one by $w_v$, enters read mode, and reads out all $L_1$ bits of $g(v)$,
among them $b$.  Every letter of this word is available at every state, so no
case distinction on the current block is needed.  Since $b$ is unconstrained
by the transcript, two members of the family agree on everything observed so
far and disagree on that continuation, so the continuation function is not
determined and~(RI) has not been attained -- a contradiction.  Hence every one
of the $M_1\log_2M_1$ bits of $g$ is determined by the transcript, and,
by the same argument with $\mathtt s_2$ and $g'$, every one of the
$M_2\log_2M_2$ bits of $g'$.  By
Theorem~\ref{thm:active-explicit-directsum}, Step~3, each such bit was
determined on a round that was a forced mistake.

\emph{Disjointness and additivity.}  A readout of $g$ occurs only while the
machine is in $Q_1$, and a readout of $g'$ only while it is in $Q_2$; the two
blocks are disjoint, so the two sets of forced-mistake rounds are disjoint.
Moreover no observation in one block constrains the unknown map of the other,
the two maps being independent coordinates of the target.  Therefore the total
is at least $S_M+C_M$.

\emph{Conditions~(iv) and~(v).}  The construction satisfies the unavoidability
requirement of Definition~\ref{def:direct-sum-active}: because $\mathtt s_1$
and $\mathtt s_2$ are total, both blocks are reachable from every state at
every time, so neither component can be bypassed, and the argument of the
previous paragraph shows the objective forces both.  Unavoidability is here a
property of the transition structure rather than a restriction imposed on the
protocol, as condition~(iv) requires.

Product structure, condition~(v), also holds: the family is indexed by the pair
$(g,g')\in Q_1^{Q_1}\times Q_2^{Q_2}$, the two maps are chosen independently,
and every combination yields a machine of the family.  The parameter set is
therefore the full Cartesian product of its two coordinate ranges, as
condition~(v) of Definition~\ref{def:direct-sum-active} requires.

The state count is $M_1+M_2+O(1)=O(M)$, so this is a direct-sum active
instance.

The additive conclusion here does \emph{not} route through
Theorem~\ref{thm:active-direct-sum} and so does not need its rectangular
hardness hypothesis.  The adversary above is lazy on each component
independently: every coordinate of $g$ and of $g'$ must be emitted before the
objective is attained, and each coordinate's first emission is a forced mistake
whatever the learner has seen on the other component, because the two maps are
independent and the answering rule commits nothing about one while responding
about the other.  Summing the forced mistakes over the two disjoint coordinate
sets gives $S_M+C_M$ directly.

No matching upper bound of the form $O(S_M+C_M)$ is claimed, and none is
proved here.  The available upper bound is
$\MistRI(M)\le O(\Esync(M))$ from Theorem~\ref{thm:active-certified}(I), stated
in terms of the aggregate synchronization complexity; it is not an upper bound
in the two components separately, and no composite learner is exhibited whose
mistake count is $O(S_M+C_M)$.  For the family constructed here both
quantities happen to be $\Theta(M\log M)$, so the lower and upper bounds do
coincide in order for this family, but that coincidence is a property of the
instance and not a general additive equality.
\end{proof}

\begin{openproblem}[A Non-Degenerate Two-Term Active Decomposition]
\label{open:si-hard-family}
Theorem~\ref{thm:esyncsi-theta} shows that state identification in a known
minimal skeleton costs $\Theta(\log M)$ mistakes, so no
decomposition of the active complexity whose second term is a
state-identification \emph{mistake} complexity can be non-degenerate.  This
settles the mistake currency negatively and leaves the question of currency
itself.

Is there a currency in which a genuine two-term decomposition holds?
The natural candidate is experiment \emph{length}: exhibit a family for which
$\Lsyncu(\mathcal C_M)=\omega(M\log M)$, so that
Proposition~\ref{prop:active-length-upper} is not dominated by the halving
bound of Theorem~\ref{thm:active-halving}, or prove that
$\Lsyncu(M)=O(M\log M)$ always.

Propositions~\ref{prop:lsyncu-quadratic} and~\ref{prop:lsyncu-binomial}
localize the \emph{machine-specific} quantity,
\[
M-1\ \le\ \Lsync(M)\ \le\ \binom{M}{2},
\]
the upper bound attained at $M=3,4$ but apparently not beyond and the lower
bound from the counter family, both in
Remark~\ref{rem:attainment-sporadic}, whose evidence in fact favours
$\Lsync(M)=\Theta(M)$.

For the universal quantity $\Lsyncu$, which is the one appearing in
Proposition~\ref{prop:active-length-upper}, none of this transfers: by
Remark~\ref{rem:lsync-not-lsyncu} a universal tree cannot search the target's
pair automaton, and the two quantities already differ at $M=2$.  Finiteness holds in general by
Proposition~\ref{prop:lsyncu-version-space}, but with a bound exponential in
$M$; for unary alphabets Proposition~\ref{prop:lsyncu-single-input} gives
$M-1$.  What is open is the rate: whether $\Lsyncu(M)$ is polynomial in $M$
(Conjecture~\ref{conj:lsyncu-poly}), and if so whether it exceeds
$M\log M$.

Proposition~\ref{prop:lsyncu-single-input} narrows the search: any witness
must use $|\mathcal I|\ge2$, since unary input alphabets give
$\Lsyncu(M)\le M-1$ --- the one regime where the universal and machine-specific
depths coincide.

The tension between the two effects a witness must combine --- many separation
episodes, and a long word within an episode --- is now precise rather than
heuristic.  Lemma~\ref{lem:tension} states it as
$d(U)\le M-|U|+1$: a set that resists splitting for many rounds must be
\emph{small}, because each round that fails to split it still splits something
else, and there are only $M$ states to divide.  Summing the trade-off over
episodes yields $\binom{M}{2}$.  The remaining question is whether the worst
case over a \emph{family} can be pushed past $M\log M$, which is a statement
about lower bounds and not about the counting above.  The natural extremal
mechanism --- a single probe state with the other candidates rotated past it
--- realizes only $2M-2$, and searches at $M\le8$ find nothing quadratic
(Remark~\ref{rem:attainment-sporadic}), so the likely answer is that no such
family exists and $\Lsyncu(M)=O(M)$.
\end{openproblem}

\begin{remark}[Necessity of the Identification Objective]
\label{rem:active-objective-needed}
Some objective of this kind is unavoidable.  In the bare active protocol, where
the learner chooses every input and is scored only on the rounds it chooses, the
learner may play $\mathtt r$ forever and make no mistakes at all while learning
nothing.  A nontrivial active lower bound therefore requires an identification,
coverage, or guaranteed-future-prediction requirement, and
Theorem~\ref{thm:active-explicit-directsum} supplies the bound relative to the
guaranteed-prediction objective already fixed in
Definition~\ref{def:residual-knowledge}.
\end{remark}

\begin{remark}[Direct-Sum Saturation Alternative]
\label{rem:direct-sum-saturation}
Instead of Lemma~\ref{lem:operational-equivalence}, which requires
determinism, one may assume an
explicit direct-sum saturation condition: there exists a family of
synchronization-hard machines and Littlestone-hard continuation machines whose
sequential composition uses $O(M)$ states and forces any active learner to pay
both costs.  Under that hypothesis, the same additive lower bound follows.
This alternative is made precise in
Definition~\ref{def:direct-sum-active} and
Theorem~\ref{thm:active-direct-sum}, which do not rely on
Lemma~\ref{lem:operational-equivalence}.
\end{remark}

\begin{theorem}[Direct-Sum Additive Lower Bound]
\label{thm:active-direct-sum}
Let a direct-sum active instance of budget $M$ have component complexities
$S_M$ and $C_M$, and assume in addition \emph{rectangular hardness}: there are
distributions $\Pi_1,\Pi_2$ on the two coordinate families such that, for every
composite learner run under the product law $\Pi_1\otimes\Pi_2$, the expected
number of mistakes on component $j$, conditioned on the entire transcript of
the other component, is at least the component bound.  Then
\[
\MistRI(M)
\ge
S_M+C_M.
\]
The bound is a sum, not a maximum, and the additivity comes from the
disjointness of the two round sets together with rectangular hardness.

No two-sided statement follows.  Converting the lower bound into an equality
would require an explicit composite learner attaining $S_M+C_M$, which is not
constructed here.  Moreover, by Proposition~\ref{prop:esyncsi-log} the
state-identification mistake complexity is $O(\log M)$, so a
decomposition in which one component is a state-identification cost is
necessarily degenerate (Remark~\ref{rem:additive-collapses}); the components of
Theorem~\ref{thm:active-two-phase} are both model-identification components,
each of complexity $\Theta(M\log M)$.
\end{theorem}

\begin{proof}
No temporal ordering of the two components is assumed; in the active protocol
the learner may interleave them freely.

Rectangular hardness is what converts two separate lower bounds into a sum, and
it cannot be dispensed with.  Conditions~(iii)--(v) of
Definition~\ref{def:direct-sum-active} give, for each component separately, a
target forcing that component's bound; but those are targets of \emph{different}
runs, and a supremum over each coordinate separately yields
$\max\{S_M,C_M\}$, not $S_M+C_M$.  Under the product law the two conditional
expectations add, by linearity of expectation applied to the conditional
counts, which is the step performed below.

Let $\mathcal A$ be any learner attaining the objective on the composite
instance.  By condition~(iv), unavoidability, $\mathcal A$ must determine the
unknown datum of each component, and this is a property of the machine family
rather than a restriction on the protocol, so it holds however $\mathcal A$
chooses its inputs.

\emph{Restriction.}  For $j=1,2$ let $R_j$ be the set of rounds on which
component $j$ forces a mistake.  We build a learner $\mathcal A_1$ for
component~1 in isolation; the construction for $\mathcal A_2$ is symmetric.

Since $\mathcal A$ chooses its inputs adaptively using outputs from
\emph{both} components, while the outputs of component~2 are not determined by
anything $\mathcal A_1$ observes, the simulation must supply those outputs from
data available to $\mathcal A_1$.  Condition~(v), product structure, is what
makes this possible.  Fix once and for all an arbitrary value $\theta_2^{0}$
of the second unknown.  Define $\mathcal A_1$ on a component-1 instance $\theta_1$ as
follows: run $\mathcal A$ internally; whenever $\mathcal A$ issues an input
whose round belongs to component~2, answer it from $\theta_2^{0}$, which
$\mathcal A_1$ knows; whenever the round belongs to component~1, forward it to
the real interaction and return the observed output.

By condition~(v) the composite target $(\theta_1,\theta_2^{0})$ lies in the
family, so the internal run of $\mathcal A$ is a genuine run against a genuine
target; by condition~(vi) the outputs $\mathcal A_1$ must supply on
component-2 rounds are determined by $\theta_2^{0}$ alone, so it can supply
them, and the outputs on component-1 rounds are determined by $\theta_1$ alone,
so they agree with its own interaction.  The component-1 rounds of the internal
run therefore coincide exactly with the interaction $\mathcal A_1$ has with its
own instance.  Hence $\mathcal A_1$ is
a well-defined component-1 learner making precisely $|R_1|$ mistakes on
$\theta_1$.  Since $\mathcal A$ attains the objective, by condition~(iv) it
determines $\theta_1$, and therefore so does $\mathcal A_1$.  By the definition
of $S_M$ and $C_M$ as the minimax mistake complexities of the components,
\[
|R_1|\ \ge\ S_M,
\qquad
|R_2|\ \ge\ C_M
\]
in the worst case over the composite family.

\emph{Summation.}  By condition~(iii), disjoint accounting,
$R_1\cap R_2=\emptyset$, so the total number of forced-mistake rounds is
$|R_1|+|R_2|$ regardless of the order in which they occur.  Hence
\[
\MistRI(M)
\ge
S_M+C_M.
\]
This is the stated bound.  No upper bound is derived: the unconditional upper bound bounds $\MistRI$ by
$O(\Esync(\mathcal C_M))$, and identifying $\Esync(\mathcal C_M)$ with
$S_M+C_M$ would require exhibiting a learner that attains the sum, which the
restriction argument does not provide.
\end{proof}

\begin{remark}[When Length Gives a Lower Bound]
\label{rem:active-length-lower}
If every step of every optimal synchronization experiment has constant output
ambiguity, then a synchronization-length lower bound may be converted into a
mistake lower bound.  In that case, if
\[
\Esync(M)=\Omega(\Lsync(M)),
\]
then the additive theorem may be written in terms of $\Lsync(M)$:
\[
\MistRI(M)
=
\Theta(M\log M+\Lsync(M)).
\]
Without such an ambiguity hypothesis, synchronization length alone does not
lower-bound prediction mistakes.
\end{remark}

\begin{remark}[Additive Synchronization Cost]
\label{rem:active-additive}
On a subclass for which a direct-sum instance exists, the attainment cost is
additive: it does not replace the $\Theta(M\log M)$ term, and active
synchronization removes the unknown-initial-state penalty only by paying the
cost of the experiment.  On the full class $\mathcal H_M$ the two terms are of
the same order by Corollary~\ref{cor:active-theta}, so additivity carries no
extra information there.
\end{remark}

\begin{remark}[Oracle Integration]
\label{rem:active-oracle}
In the oracle bridge of Section~\ref{sec:oracle}, the passive realizable
estimation rate is
\[
\Est_M^{\mathrm{passive}}(T)
=
\Theta(M\log M),
\]
while on the full class the active realizable estimation rate is, by
Corollary~\ref{cor:active-theta},
\[
\Est_M^{\mathrm{active}}(T)
=
\Theta(M\log M)
\]
unconditionally, so the two protocols enter the oracle bridge with rates of the
same order.  On a subclass the rate is $\Theta(\Esync(\mathcal C_M))$ by
Theorem~\ref{thm:active-certified}(I).  In either case the oracle inequality
itself is unchanged; only the estimation rate input is modified.
\end{remark}

%----------------------------------------------------------------------
\subsection{Agnostic Setting}
\label{subsec:agnostic}
%----------------------------------------------------------------------

In the agnostic setting, the target sequence need not be realized by any
$M$-state Mealy machine.  The learner competes against the best function in
$\mathcal H_M$.

\begin{theorem}[Agnostic Regret]
\label{thm:agnostic}
For the agnostic setting,
\[
\inf_{\mathcal A}
\sup_{(x_t,Y_t)}
R_T^{\mathrm{agn}}(\mathcal A)
=
O\bigl(\sqrt{T\,M\log M}\bigr),
\]
up to constants depending on the alphabets and possible logarithmic
refinements.  The matching lower bound also holds, in the persistent-stream
protocol as well as the reset-word protocol, so
\[
\inf_{\mathcal A}
\sup_{(x_t,Y_t)}
R_T^{\mathrm{agn}}(\mathcal A)
=
\Theta\bigl(\sqrt{T\,M\log M}\bigr).
\]

The two-sided statement is asserted in the regime
\[
T\ \gtrsim\ M\log M ,
\]
that is, for horizons at least of the order of the Littlestone dimension.  For
shorter horizons the lower bound saturates: the adversary cannot force more
than one mistake per round, so the minimax regret is $\Theta(T)$ when
$T=O(M\log M)$, and the correct uniform statement is
$\Theta\bigl(\min\{T,\sqrt{T\,M\log M}\}\bigr)$.
\end{theorem}

\begin{proof}
By Lemma~\ref{lem:littlestone},
\[
\Ldim(\mathcal H_M)=\Theta(M\log M).
\]
Since $\log|\mathcal H_M|=O(M\log M)$, finite-expert aggregation over the
labeled machines gives the unconditional upper bound
\[
R_T^{\mathrm{agn}}
\le
\sqrt{2T\log|\mathcal H_M|}
=
O\bigl(\sqrt{T\,M\log M}\bigr).
\]
For the lower bound, the minimax theorem for online learning
\cite{bendavid2009,daniely2014} gives
$R_T^{\mathrm{agn}}=\Omega(\sqrt{T\,\Ldim(\mathcal H_M)})$ whenever the
Littlestone dimension is witnessed by a mistake tree realizable in the
protocol at hand.  Theorem~\ref{thm:stream-lower-bound} supplies exactly such
a tree: at each of its $M\log_2M$ readout rounds \emph{both} output values
remain consistent with some member of $\mathcal G_M\subseteq\mathcal H_M$,
because the corresponding bit of $g$ is unconstrained at that moment.  The
adversary's decision tree is therefore a complete shattered mistake tree of
depth $M\log_2M$ embedded in a single never-reset stream, whence
\[
\Ldim^{\mathrm{stream}}(\mathcal H_M)\ \ge\ M\log_2M,
\]
and the reverse inequality $\Ldim\le\log_2|\mathcal H_M|=O(M\log M)$ is the
counting bound.  Substituting
$\Ldim^{\mathrm{stream}}(\mathcal H_M)=\Theta(M\log M)$ gives the claim.
\end{proof}

\begin{corollary}[Stream Littlestone Dimension]
\label{cor:stream-ldim}
For fixed alphabets with $|\mathcal I|\ge4$ and $|\mathcal O|\ge2$, the
Littlestone dimension of $\mathcal H_M$ witnessed by mistake trees realizable
inside a single persistent stream satisfies
\[
\Ldim^{\mathrm{stream}}(\mathcal H_M)
=
\Theta(M\log M),
\]
matching the reset-word value.  Thus the two protocols have the same
sequential complexity up to constants.
\end{corollary}

\begin{proof}
The upper bound is the counting bound of Lemma~\ref{lem:littlestone}, since
$\Ldim\le\log_2|\mathcal H_M|=O(M\log M)$.

For the lower bound we exhibit the mistake tree explicitly rather than appeal
to the forcing argument as a black box.  Recall from
Theorem~\ref{thm:stream-lower-bound} that the adversary processes the states
$v^{(1)},\dots,v^{(M)}$ in order, playing the block
$w_{v^{(k)}}\,\mathtt c\,\mathtt d^{L}$ for each, and that the $L$ readout
rounds of block $k$ reveal the $L$ bits of $g(v^{(k)})$.

Build a binary tree of depth $M\log_2M$ whose levels are indexed by the pairs
$(k,j)$ with $1\le k\le M$ and $1\le j\le L$, ordered lexicographically.  The
node at level $(k,j)$ reached by a path $\sigma\in\{0,1\}^{<M\log_2M}$ is
labelled by the input letter that the adversary plays at that round, namely the
$j$-th $\mathtt d$ of block $k$; the transport letters of $w_{v^{(k)}}$ and the
letter $\mathtt c$ are inserted as forced unary segments between levels, since
their outputs are determined by the fixed part of the table and branch nowhere.
The two outgoing edges of the node correspond to the two possible answers for
the bit $g(v^{(k)})_j$.

Each root-to-leaf path $\sigma$ determines a full assignment of all
$M\log_2M$ bits, hence a unique $g_\sigma:Q\to Q$ and a unique member
$f_\sigma$ of the family, which by the consistency step of
Theorem~\ref{thm:stream-lower-bound} realizes every label along $\sigma$.  The
tree is therefore shattered by $\mathcal H_M$, and the whole labelled sequence
occurs inside one persistent stream with no external reset.  Its depth is
$M\log_2M=\Omega(M\log M)$, giving
$\Ldim^{\mathrm{stream}}(\mathcal H_M)=\Omega(M\log M)$.
\end{proof}

\begin{remark}[Logarithmic Refinements]
\label{rem:agnostic-log-factors}
If a refined theorem introduces additional logarithmic factors in $T$, those
factors should be stated explicitly.  The $M$-scaling
\[
\sqrt{M\log M}
\]
is robust.  In the notation of Section~\ref{sec:oracle}, the agnostic
estimation rate is
\[
\Est_M^{\mathrm{agn}}(T)
=
\widetilde\Theta\!\left(\sqrt{T\,M}\right),
\]
with the $\widetilde\Theta$ hiding logarithmic factors such as $\log M$ and
possible $\log T$ refinements.
\end{remark}

\begin{remark}[Log-Loss Bridge for Retention]
\label{rem:logloss-retention}
For retention the compatible online bridge is log-loss.  The approximation term
$\Aapp^{\mathrm{ret}}(M)$ is then a cumulative excess KL relative to the best
stochastic finite-state predictor of budget $M$, and the estimation term is a
mixability-based rate.  It is not a statement about deterministic point-mass
Mealy machines, whose compatible bridge is the realizable mistake bound.
\end{remark}

\begin{remark}[Type-Correct Temporal Rates]
\label{rem:temporal-type-correct}
The bounds in this section are stated for deterministic Mealy hypothesis
classes, which are the type-correct online classes for commitment and for
realizable deterministic prediction.  Retention with KL or log-loss uses
stochastic finite-state predictors and mixability-based estimation rates.
Grounding spectral relaxations use linear finite-rank predictors.  The oracle
bridge imports the appropriate type-correct estimation rate for each regime.
\end{remark}

%======================================================================
\section{Spectral Tail Heuristic}
\label{sec:spectral-heuristic}
%======================================================================

The exact theorems of this manuscript are regime-specific.  Nevertheless, a
common spectral-tail picture organizes many of the analytic converses.  This
section states that picture as a heuristic, not as a theorem dictated solely by
aggregation.

%----------------------------------------------------------------------
\subsection{The Heuristic}
\label{subsec:spectral-heuristic-statement}
%----------------------------------------------------------------------

\begin{heuristic}[Spectral Tail Heuristic]
\label{heur:spectral}
In several approximation regimes, the $M$-state or $M$-rank error is
controlled by the tail of a regime-specific positive operator.  The precise
tail form depends on:
\begin{itemize}
\item operator approximation versus clustering;
\item linear versus affine subspaces;
\item rank count versus state count;
\item compactness, trace-class, domination, and quadratic-regime
assumptions;
\item the aggregation norm and the response representation.
\end{itemize}
Aggregation alone does not determine the tail.
\end{heuristic}

The heuristic serves as a design principle and as a diagnostic.  It becomes a
theorem only after a regime supplies:
\begin{enumerate}[label=(\arabic*)]
\item a response operator $A_\delta$;
\item an effective rank budget $r_{\mathbb T}(M)$;
\item a Schatten or Ky Fan domination inequality;
\item and, for online rates, an achievability bound.
\end{enumerate}

%----------------------------------------------------------------------
\subsection{Why Different Tails Appear}
\label{subsec:spectral-why}
%----------------------------------------------------------------------

The different tail forms reflect different approximation mechanisms.

\paragraph{Operator-norm case: Eckart--Young--Mirsky.}
For the unrestricted linear finite-rank grounding relaxation, the exact gap is
an operator-norm approximation problem:
\[
\Dunres(M)
=
\inf_{\rank B\le M}
\norm{H_\nu-B}_{\mathrm{op}}.
\]
Eckart--Young--Mirsky gives the single singular value
\[
\Dunres(M)
=
\sigma_{M+1}(H_\nu).
\]
The rank budget is $M$, and the error is the first omitted singular value.

The Hankel-structured finite-rank realization gap satisfies the lower bound
\[
\DHankstr(M)
\ge
\sigma_{M+1}(H_\nu),
\]
with equality only under additional structural hypotheses.

\paragraph{Nuclear / squared-error case: covariance residual.}
For quadratic retention, the residual object is a positive semidefinite
covariance residual.  An $M$-cluster approximation induces at most $M-1$
centered centroid directions.  The residual covariance has trace equal to its
nuclear norm, so the lower bound is a trace tail:
\[
\Delta_{\mathrm{ret}}^{\mathrm{quad}}(M)
\ge
\frac12
\sum_{i\ge M}
\lambda_i(\Sigma_\pi)
=
\frac12
\Phi^{(1)}_{M-1}(\Sigma_\pi).
\]
The exponent $1$ appears because the residual covariance is positive
semidefinite, not because the original cost is an $\ell^1$ aggregation.

\paragraph{Commitment case: combinatorial threshold.}
Commitment does not possess a genuine singular-value spectrum.  Its exact gap
is a Boolean rank / Myhill--Nerode threshold:
\[
\Com(M)=0
\iff
M\ge\kappaMN.
\]
One may view this as a combinatorial step spectrum: the zero-error gap drops
from positive to zero when $M$ crosses the Boolean realization rank.  This is
not an ordinary Schatten $p\to0$ limit.

\paragraph{Index offset.}
The grounding tail begins at $\sigma_{M+1}$ because a rank-$M$ linear
approximation uses an $M$-dimensional linear subspace.  The retention tail
begins at $i\ge M$ because $M$ centroids span at most $M-1$ centered
directions after the global mean is removed.  Thus the same state budget $M$
produces different effective ranks in operator approximation and clustering.

\paragraph{Frobenius / Schatten-$2$ case.}
If a regime satisfies a Schatten-$2$ domination hypothesis, the conditional
template gives a Frobenius tail:
\[
\Delta_{\mathbb T}(M)
\ge
\frac1{c_2}
\left(
\sum_{i>r_{\mathbb T}(M)}
\sigma_i(A_\delta)^2
\right)^{1/2}.
\]
This is a conditional consequence of the Schatten template, not a verified
manuscript regime unless such a domination bridge is explicitly supplied.

%----------------------------------------------------------------------
\subsection{Conditional Cross-Term Conjectures}
\label{subsec:spectral-conjectures}
%----------------------------------------------------------------------

Some cross-term tail forms require additional hypotheses.  They are not
consequences of aggregation alone.

\begin{conjecture}[Conditional Cross Terms]
\label{conj:cross-terms}
The following are plausible under additional hypotheses:
\begin{enumerate}[label=(\arabic*),leftmargin=2.0em]
\item An $\mathcal L^1$-type grounding bound of nuclear-tail form
\[
\Delta_{\mathrm{grd}}^{L^1}(M)
\gtrsim
\sum_{i>M}
\sigma_i(H_\nu)
\]
requires $H_\nu$ to be trace class and requires a nuclear-norm
domination hypothesis relating the task cost to the trace norm of
$H_\delta-H_{\widehat\delta}$.

\item An $\mathcal L^\infty$-type retention bound of single-eigenvalue form
\[
\Delta_{\mathrm{ret}}^{L^\infty}(M)
\gtrsim
\frac12
\lambda_M(\Sigma_\pi)
\]
requires a worst-case quadratic surrogate or additional geometric
assumptions aligning the worst-case state with the top Fisher
eigendirection.

\item A Frobenius / Schatten-$2$ converse
\[
\Delta_{\mathbb T}(M)
\gtrsim
\left(
\sum_{i>r(M)}
\sigma_i(A_\delta)^2
\right)^{1/2}
\]
requires an explicit Schatten-$2$ domination bridge.

\item Intermediate R\'enyi orders $\alpha\notin\{0,1,\infty\}$ may organize
interpolation inequalities, but they do not define operational
finite-state regimes unless a matching task theory, domination bridge,
and rank-control structure are supplied.
\end{enumerate}
\end{conjecture}

These conjectures are compatible with the Schatten template and the exponent
correspondence.  They are not asserted as theorems.

%----------------------------------------------------------------------
\subsection{Scope and Conditions}
\label{subsec:spectral-scope}
%----------------------------------------------------------------------

\begin{itemize}[leftmargin=1.6em]
\item \textbf{Aggregation alone is insufficient.}
The tail shape is influenced by aggregation, but also by affine versus linear
approximation, clustering versus operator approximation, state count versus
rank count, and regime-specific domination or quadratic-regime assumptions.

\item \textbf{No Boolean singular values.}
Commitment does not have a genuine Schatten singular-value spectrum.  Its
threshold is a Boolean rank / Myhill--Nerode realizability threshold.

\item \textbf{No unconditional cross-regime ordering.}
The spectral-tail heuristic does not imply inequalities among grounding,
retention, and commitment gaps.  The Independence Theorem remains intact.

\item \textbf{Converses are not achievability rates.}
Spectral tails are lower bounds on approximation gaps.  They become online
regret upper rates only under explicit spectral-rate assumptions, as in
Section~\ref{sec:oracle}.

\item \textbf{Numerical checks are illustrations.}
Numerical evaluations can illustrate singular-value tails, covariance tails,
and threshold behavior on finite instances.  They do not prove domination
hypotheses, Boolean-rank identities, relaxation gaps, or independence.
\end{itemize}

The spectral-tail picture is an organizing heuristic once a response operator,
an effective rank budget, and a domination bridge are supplied.  It is not a
single theorem derived from aggregation alone, and it does not unify the
numerical gaps across regimes.

%======================================================================
\section{Conditional Representation Theorem}
\label{sec:conditional-rep}
%======================================================================

The schema does not by itself force KL divergence, Boolean equality, or
operator norm.  Those cost functionals arise under additional analytical
axioms.  The following theorem records the conditional representation results
that justify the canonical regime costs.

%----------------------------------------------------------------------
\subsection{Theorem Statement}
\label{subsec:conditional-rep-statement}
%----------------------------------------------------------------------

\begin{theorem}[Conditional Representation]
\label{thm:master-rigorous}
Let $\mathbb T$ be a task theory on a finite input alphabet $\mathcal I$.

\medskip
\noindent\textbf{Clause I, Retention.}
Suppose $\mathbf R=\mathbf{FinStoch}$, so that residual objects are finite
probability distributions.  Suppose $\mathcal E$ satisfies:
\begin{enumerate}[label=(\alph*)]
\item \textbf{Faithfulness.}
\[
\mathcal E(P\|Q)\ge0,
\]
and
\[
\mathcal E(P\|Q)=0
\iff
P=Q.
\]

\item \textbf{Regularity.}
$\mathcal E$ is continuous on full-support pairs and lower semicontinuous on
the boundary.

\item \textbf{Data processing.}
For every stochastic kernel $K$,
\[
\mathcal E(P\|Q)
\ge
\mathcal E(PK\|QK).
\]

\item \textbf{Exact chain rule.}
For all finite joint distributions,
\[
\mathcal E(P_{XY}\|Q_{XY})
=
\mathcal E(P_X\|Q_X)
+
\mathbb E_{P_X}
\bigl[
\mathcal E(P_{Y|X}\|Q_{Y|X})
\bigr].
\]

\item \textbf{$f$-divergence form (independent structural axiom).}
Assume outright that $\mathcal E$ lies in the class of $f$-divergences, i.e.
\[
\mathcal E(P\|Q)
=
\sum_i q_i f\!\left(\frac{p_i}{q_i}\right)
\]
for a convex generator $f$ with $f(1)=0$, extended lower semicontinuously at
the boundary.  This is an \emph{assumed} hypothesis, not a consequence of
clauses~(a)--(d): faithfulness, regularity, and data processing are all
satisfied by the R\'enyi divergences, which are not $f$-divergences
(Remark~\ref{rem:csiszar-sumform-needed}).  Deriving the $f$-divergence form
from information monotonicity requires further structural axioms such as
expansibility, symmetry, and sum-form decomposability
\cite{csiszar1978,amari2009}, which are not assumed here.  Given clause~(e), the chain rule pins down the generator by
Lemma~\ref{lem:csiszar-representation}.

Given this clause, data processing is not used in the derivation: every
$f$-divergence with convex $f$ satisfies it automatically.  It appears in
the hypothesis list because the clauses are stated as independent axioms on
$\mathcal E$, but the proof below invokes only faithfulness, the
$f$-divergence form, and the chain rule.
\end{enumerate}
Then there exists a constant $c>0$ such that
\[
\mathcal E(P\|Q)
=
c\,D_{\mathrm{KL}}(P\|Q).
\]

\medskip
\noindent\textbf{Clause II, Commitment.}
Suppose $\mathbf R=\mathbf{Set}$ and $\mathcal E$ is faithful and
$\{0,\infty\}$-valued:
\[
\mathcal E(x,y)=0
\iff
x=y.
\]
Then
\[
\mathcal E(x,y)=
\begin{cases}
0,&x=y,\\
\infty,&x\neq y.
\end{cases}
\]

\medskip
\noindent\textbf{Clause III, Grounding.}
Suppose $\operatorname{ob}(\mathbf R)$ is a concrete $C^*$-algebra of bounded
operators and
\[
\mathcal E(T,S)=\|T-S\|,
\]
where $\|\cdot\|$ satisfies the $C^*$-identity
\[
\|T^*T\|=\|T\|^2.
\]
Then
\[
\|T\|=\|T\|_{\mathrm{op}}.
\]

\medskip
\noindent\textbf{Consequence.}
Under these representations, the canonical gaps are respectively:
\begin{itemize}
\item the full-KL information-bottleneck objective;
\item the Myhill--Nerode index;
\item the Eckart--Young--Mirsky spectral lower bound for the unrestricted
linear Hankel relaxation.
\end{itemize}
\end{theorem}

%----------------------------------------------------------------------
\subsection{Csisz\'ar Representation Lemma}
\label{subsec:csiszar-lemma}
%----------------------------------------------------------------------

\begin{lemma}[Csisz\'ar Representation]
\label{lem:csiszar-representation}
The $f$-divergence class and its characterizations are due to Csisz\'ar
\cite{csiszar1978}; the failure of information monotonicity alone to pin down
the class, and the role of the additional sum-form axioms, are discussed by
Amari \cite{amari2009}.  The chain-rule selection of the logarithmic generator
below is stated in the form used here.

Let $\mathcal E(P\|Q)$ be a functional on pairs of probability distributions
over a finite alphabet satisfying:
\begin{enumerate}[label=(\alph*)]
\item \textbf{Faithfulness.}
$\mathcal E(P\|Q)\ge0$, with equality iff $P=Q$.
\item \textbf{Regularity.}
$\mathcal E$ is continuous on full-support pairs and lower semicontinuous on
the boundary.
\item \textbf{Data processing.}
For every stochastic kernel $K$,
\[
\mathcal E(P\|Q)
\ge
\mathcal E(PK\|QK).
\]

\item \textbf{$f$-divergence form.}
$\mathcal E$ is an $f$-divergence: there is a convex
$f:[0,\infty)\to\mathbb R\cup\{+\infty\}$ with $f(1)=0$, extended lower
semicontinuously at the boundary, such that for all finite alphabets
\[
\mathcal E(P\|Q)
=
\sum_i q_i\,f\!\left(\frac{p_i}{q_i}\right).
\]
\end{enumerate}
Recall that the generator is determined only modulo an affine term: replacing
$f(t)$ by $f(t)+a(t-1)$ leaves $\mathcal E$ unchanged.

If, in addition, $\mathcal E$ satisfies the exact $P$-weighted chain rule for
all finite joint distributions, then, modulo such an affine term,
$f(t)=c\,t\log t$ for some $c>0$; equivalently the normalized convex
representative is $f(t)=c(t\log t-t+1)$, and hence
\[
\mathcal E(P\|Q)
=
c\,D_{\mathrm{KL}}(P\|Q).
\]
\end{lemma}

\begin{proof}
Hypothesis~(d) places $\mathcal E$ in the $f$-divergence class; it is assumed,
not derived, and cannot be replaced by faithfulness, regularity, and data
processing alone (Remark~\ref{rem:csiszar-sumform-needed}).  Faithfulness
excludes the identically zero generator and forces $f$ to be strictly convex
at $t=1$.

For the second assertion, apply the chain rule to a joint distribution
$P_{XY}$ and its marginal-conditional factorization.  The chain rule requires
\[
\sum_{x,y} q_x q_{y|x}\,
f\!\left(\frac{p_x p_{y|x}}{q_x q_{y|x}}\right)
=
\sum_x q_x\,f\!\left(\frac{p_x}{q_x}\right)
+
\sum_x p_x
\sum_y q_{y|x}\,
f\!\left(\frac{p_{y|x}}{q_{y|x}}\right).
\]
Since the generator matters only modulo an affine term $a(t-1)$, write
$f(t)=g(t)+a(t-1)$ and normalize by $g(1)=0$.

\emph{Product additivity is not enough.}  Specializing to product
distributions, $p_{y|x}=p_y$ and $q_{y|x}=q_y$ independent of $x$, yields the
additivity
$\mathcal E(p\otimes p'\|q\otimes q')=\mathcal E(p\|q)+\mathcal E(p'\|q')$,
which is too weak to select the generator.  The reverse
Kullback--Leibler divergence, with generator $g(t)=-\log t$, is an
$f$-divergence and is additive over products, yet it is not a multiple of
$D_{\mathrm{KL}}$.  Any argument resting on product additivity alone therefore
cannot reach the conclusion, and we must use the chain rule in its genuinely
conditional form, where $q_{y|x}$ varies with $x$.

\emph{An exact two-variable identity from a non-product conditional.}  Let $X$
be binary with $q_X=(\alpha,1-\alpha)$ and $p_X=(u\alpha,1-u\alpha)$ for
$u>0$ and $\alpha$ small enough that both are probability vectors.  Fix a
conditional reference $q_{Y|X}$ that is the \emph{same} law $q'$ on both
branches, and a conditional $p_{Y|X}$ that equals $p'$ on the branch $x=1$ and
equals $q'$ on the branch $x=2$.  Write $t_j=p'_j/q'_j$.

On the branch $x=2$ the conditional divergence vanishes, since
$p_{Y|x=2}=q_{Y|x=2}$ and $g(1)=0$.  The joint ratios are $u\,t_j$ on the
branch $x=1$, with weights $\alpha q'_j$, and $\rho\coloneqq(1-u\alpha)/(1-\alpha)$
on the branch $x=2$, with weights $(1-\alpha)q'_j$.  The $P$-weighted chain
rule
\[
\mathcal E(P_{XY}\|Q_{XY})
=
\mathcal E(P_X\|Q_X)
+
\sum_xp_x\,\mathcal E(P_{Y|x}\|Q_{Y|x})
\]
therefore reads
\[
\alpha\sum_jq'_j\,g(u\,t_j)
+
(1-\alpha)\,g(\rho)
=
\Bigl[\alpha\,g(u)+(1-\alpha)g(\rho)\Bigr]
+
u\alpha\sum_jq'_j\,g(t_j),
\]
where the bracket is $\mathcal E(P_X\|Q_X)$.  The terms $(1-\alpha)g(\rho)$
cancel \emph{exactly} --- no expansion in $\alpha$ is needed --- and dividing
by $\alpha>0$ leaves the identity
\[
\sum_jq'_j\,g(u\,t_j)
=
g(u)+u\sum_jq'_j\,g(t_j),
\tag{$\dagger$}
\]
valid for every $u>0$, every reference $q'$, and every $p'$.

\emph{One-sided specialization; no differentiability is assumed.}
Take $q'=(\beta,1-\beta)$ and $p'=(v\beta,1-v\beta)$ with $v>0$ and
$\beta\in(0,1)$ small enough that $p'$ is a probability vector, so that
$t_1=v$ and $\sigma(\beta)\coloneqq t_2=(1-v\beta)/(1-\beta)$.  Then
$(\dagger)$ reads
\[
\beta\,g(uv)+(1-\beta)\,g(u\sigma)
=
g(u)+u\beta\,g(v)+u(1-\beta)\,g(\sigma).
\tag{$\dagger\dagger$}
\]
Both sides tend to $g(u)$ as $\beta\downarrow0$.  Subtract $g(u)$, divide by
$\beta$, and let $\beta\downarrow0$.  This uses no differentiability
hypothesis: a convex function on an open interval has finite one-sided
derivatives at every point, and $\beta$ approaches $0$ from above only, so
every limit occurring is one-sided.  Since
$\sigma=1+\beta(1-v)+O(\beta^2)$, the argument $u\sigma$ approaches $u$ from
above when $v<1$ and from below when $v>1$, and $\sigma$ approaches $1$ from
the same side.  Writing $g'_+$ and $g'_-$ for the right and left derivatives,
\[
g(uv)-g(u)+u(1-v)\,g'_{\pm}(u)
=
u\,g(v)+u(1-v)\,g'_{\pm}(1),
\qquad
\begin{cases}+, & v<1,\\ -, & v>1,\end{cases}
\tag{$\ddagger$}
\]
for every $u>0$; the case $v=1$ is trivial.

\emph{Symmetry on each side of $1$.}  The left side of $(\ddagger)$ is
symmetric in $u$ and $v$.  For $u,v\in(0,1)$ both $(\ddagger)$ with the $+$
sign and its transpose hold, and equating the two expressions for $g(uv)$
gives
\[
\frac{u\,g'_+(u)-g(u)-g'_+(1)}{u-1}
=
\frac{v\,g'_+(v)-g(v)-g'_+(1)}{v-1}.
\]
Each side depends on one variable only, so there is a constant $c_1$ with
\[
u\,g'_+(u)=g(u)+c_1(u-1)+g'_+(1),
\qquad u\in(0,1).
\tag{$\ddagger_+$}
\]
The same argument with the $-$ sign for $u,v>1$ gives a constant $c_2$ with
$u\,g'_-(u)=g(u)+c_2(u-1)+g'_-(1)$ on $(1,\infty)$.

\emph{Regularity bootstrap.}  Equation $(\ddagger_+)$ exhibits $g'_+(u)$ as
$[g(u)+c_1(u-1)+g'_+(1)]/u$, which is continuous on $(0,1)$ because convex
functions are continuous on open intervals.  A convex function whose right
derivative is continuous is $C^1$, since then
$g'_-(u)=\lim_{t\uparrow u}g'_+(t)=g'_+(u)$ at every point.  Hence $g$ is
$C^1$ on $(0,1)$ and satisfies $u\,g'(u)-g(u)=c_1(u-1)+\alpha$ with
$\alpha\coloneqq g'_+(1)$.  Dividing by $u^2$ makes the left side
$(g(u)/u)'$, and integrating with $g(1)=0$ gives
\[
g(u)=c_1\bigl(u\log u-u+1\bigr)+\alpha(u-1),
\qquad u\in(0,1].
\]
Then $g'_-(1)=\lim_{u\uparrow1}g'(u)=\alpha$, and the same bootstrap applied
on $(1,\infty)$ yields
$g(u)=c_2(u\log u-u+1)+\alpha(u-1)$ for $u\ge1$.

\emph{Gluing across $1$.}  Apply $(\ddagger)$ with the $-$ sign at
$u\in(0,1)$ and $v=1/u>1$, so $uv=1$ and $g(uv)=0$.  Substituting the two
branch formulas and $g'_-(u)=c_1\log u+\alpha$, the affine terms cancel and
\[
0=(c_1-c_2)\bigl(\log u-u+1\bigr)
\qquad\text{for all }u\in(0,1).
\]
Since $\log u<u-1$ for $u\ne1$, this forces $c_1=c_2\eqqcolon c$, so
\[
g(t)=c\,(t\log t-t+1)+\alpha(t-1),
\qquad t>0.
\]
The affine term does not change the divergence, so the normalized convex
representative is $f(t)=c(t\log t-t+1)$.  Convexity forces $c\ge0$ and
faithfulness excludes $c=0$, since $c=0$ gives $\mathcal E\equiv0$; hence
$c>0$ and $\mathcal E=c\,D_{\mathrm{KL}}$.  See
\cite[Sec.~2.1]{csiszar1978} and \cite[Ch.~3]{cover} for the classical
characterization within the finite-alphabet $f$-divergence class.
\end{proof}

\begin{remark}[Smoothness of the Generator Is Derived, Not Assumed]
\label{rem:csiszar-automatic-smoothness}
The proof assumes only convexity and derives regularity from the chain rule:
the displayed formula shows that every convex generator satisfying the
$P$-weighted chain rule is $C^\infty$ on $(0,\infty)$, with matching one-sided
derivatives at $t=1$.  A $C^1$ hypothesis, sometimes imposed in
characterizations of this kind, is therefore a conclusion here rather than a
prerequisite.

Two steps carry the weight.  The passage to $(\ddagger)$ is a one-sided limit
in the probability weight $\beta$, which is legitimate for any convex $g$
because $\beta\downarrow0$ only; differentiating $(\dagger\dagger)$ at
$\beta=0$ would not be, since convexity gives differentiability only off a
countable set.  The bootstrap then converts $(\ddagger_+)$ from a constraint
\emph{on} $g'_+$ into a formula \emph{for} it, and a convex function with
continuous right derivative is $C^1$.

The gluing step is not decorative.  A generator equal to
$c_1(t\log t-t+1)+\alpha(t-1)$ on $(0,1]$ and to
$c_2(t\log t-t+1)+\alpha(t-1)$ on $[1,\infty)$ with $c_1\ne c_2$ is convex,
is $C^1$ at $t=1$, and satisfies both branch equations; only the identity
$0=(c_1-c_2)(\log u-u+1)$ excludes it.
\end{remark}

\begin{remark}[Why the Conditional Structure Is Indispensable]
\label{rem:csiszar-conditional-needed}
The proof above uses a conditional $P_{Y|X}$ that differs across the two
branches of $X$.  This is not a matter of convenience.  Restricting the chain
rule to product distributions yields only additivity over products, and that
property does not characterize the logarithmic generator: reverse
Kullback--Leibler divergence, $g(t)=-\log t$, satisfies
\[
\mathcal E(p\otimes p'\|q\otimes q')=\mathcal E(p\|q)+\mathcal E(p'\|q')
\]
for all factors, and is an $f$-divergence, but is not proportional to
$D_{\mathrm{KL}}$.  Substituting $g(t)=-\log t$ into the exact identity
$(\dagger)$, by contrast, fails immediately: with $u=2.3$,
$q'=(0.3,0.7)$ and $p'=(0.55,0.45)$ the two sides of $(\dagger)$ differ by
$-0.1657$, whereas for $g(t)=t\log t-t+1$ they agree to machine precision.

The separation is systematic rather than incidental to that triple.  Over
random $(u,q',p')$ with $|\mathcal O|\le5$, the maximum discrepancy in
$(\dagger)$ is $5.5\times10^{-40}$ in $40$-digit arithmetic for the forward
generator, while for reverse Kullback--Leibler, $\chi^2$, Hellinger, the
Jensen--Shannon generator and total variation it is $12.6$, $2.0\times10^4$,
$3.1$, $3.2$ and $4.5$ respectively.  Sweeping the $\alpha$-divergence family
$g_\alpha(t)=(t^\alpha-1-\alpha(t-1))/(\alpha(\alpha-1))$ through
$\alpha\in\{-1,0,\tfrac12,0.9,0.99,1,1.01,\tfrac32,2,3\}$, the defect in
$(\dagger)$ vanishes only at $\alpha=1$ and is bounded away from zero on either
side, so no neighbouring member of the family survives the conditional chain
rule.  The subsequent symmetry step is equally decisive: the quotient
$(u\,g'(u)-g(u))/(u-1)$ evaluated at $u=1.3,2.1,3.7,5.9,9.4$ is constant to
$0$ for $t\log t-t+1$ and varies by $0.61$, $8.1$, $0.22$ and $0.27$ for the
other differentiable generators listed.
The $P$-weighting of the conditional term --- the factor $p_x$ rather than
$q_x$ in front of $\mathcal E(P_{Y|x}\|Q_{Y|x})$ --- is what breaks the
symmetry between forward and reverse divergence, and it is visible only when
$P_{Y|x}$ genuinely depends on $x$.
\end{remark}

\begin{remark}[The Sum-Form Hypothesis Cannot Be Dropped]
\label{rem:csiszar-sumform-needed}
Faithfulness, regularity, and data processing alone do \emph{not}
characterize $f$-divergences.  The R\'enyi divergences
\[
D_\alpha(P\|Q)
=
\frac1{\alpha-1}
\log\sum_i p_i^\alpha q_i^{1-\alpha},
\qquad
\alpha\in(0,\infty)\setminus\{1\},
\]
are faithful, continuous on full-support pairs, and satisfy the data
processing inequality, yet they are not $f$-divergences: the outer logarithm
destroys the sum form.  Concretely, every $f$-divergence is jointly convex in
the pair $(P,Q)$, whereas $D_\alpha$ fails joint convexity for $\alpha>1$.
More generally, any strictly increasing reparametrization of an
$f$-divergence inherits faithfulness and data processing while leaving the
$f$-divergence class.  The $f$-divergence hypothesis is therefore not removable, and clause~(e) of
Theorem~\ref{thm:master-rigorous} is an additional assumption rather
than a consequence of the others.  Sum-form decomposability alone is likewise
insufficient without further symmetry and expansibility axioms; we therefore
assume membership in the $f$-divergence class outright, which is the weakest
formulation that makes the chain-rule argument valid.
\end{remark}

\begin{remark}[Role of the Lemma]
\label{rem:csiszar-role}
The lemma separates two logical steps, which have different status.  The
first---membership in the $f$-divergence class---is an \emph{assumed}
structural axiom, hypothesis~(d).  It does \emph{not} follow from faithfulness,
regularity, and data processing: the R\'enyi divergences satisfy all three and
are not $f$-divergences (Remark~\ref{rem:csiszar-sumform-needed}).  The
second---the chain rule forces the logarithmic generator---is a genuine
characterization \emph{within} the $f$-divergence class, and is what the lemma
proves.  Neither step is derivable from the enriched-category schema alone: the
schema supplies the finite-state syntax, while the analytical axioms supply the
cost functional.
\end{remark}

%----------------------------------------------------------------------
\subsection{Proof of Clause I}
\label{subsec:proof-clause-I}
%----------------------------------------------------------------------

\begin{proof}[Proof of Clause I]
By clause~(e), $\mathcal E$ is an $f$-divergence:
\[
\mathcal E(P\|Q)
=
\sum_i q_i f\!\left(\frac{p_i}{q_i}\right),
\]
with $f$ convex and $f(1)=0$.

The exact chain rule forces $f$ to be logarithmic.  Applying the chain rule
to finite alphabets and to binary refinements yields the standard functional
equation characterizing those $f$-divergences that are additive under
conditioning.  Under continuity, the solutions are exactly
\[
f(t)=c(t\log t-t+1),
\qquad c>0.
\]
Faithfulness excludes $c=0$.  Substitution gives
\[
\mathcal E(P\|Q)
=
c
\sum_i p_i\log\frac{p_i}{q_i}
=
cD_{\mathrm{KL}}(P\|Q).
\]
\end{proof}

\begin{remark}[Why Faithfulness Is Necessary]
\label{rem:clauseI-faithfulness}
Without faithfulness, the identically zero functional
\[
\mathcal E(P\|Q)=0
\]
satisfies data processing, joint convexity, lower semicontinuity, and the
chain rule, but corresponds to $c=0$.  Clause~I therefore requires
faithfulness, or an equivalent nontriviality assumption.
\end{remark}

\begin{remark}[The Chain Rule Is External]
\label{rem:chain-rule-external}
The chain rule in Clause I is an analytical axiom on the cost bimodule.  It is
not derivable from the $\mathbf V$-category axioms alone.  The enriched schema
supplies the finite-state syntax; the chain rule supplies the information
geometry that selects KL divergence.
\end{remark}

\begin{remark}[Retention Vertex]
\label{rem:clauseI-retention-vertex}
Clause I explains why retention is the exact $\alpha=1$ vertex in the exponent
correspondence.  The cost is KL / relative entropy.  The spectral converse for
the quadratic restricted retention gap is a separate Schatten-$1$ /
trace-norm statement and is not derived from Clause I alone.
\end{remark}

%----------------------------------------------------------------------
\subsection{Proof of Clause II}
\label{subsec:proof-clause-II}
%----------------------------------------------------------------------

\begin{proof}[Proof of Clause II]
Since $\mathcal E$ is $\{0,\infty\}$-valued and faithful, the only possible
zero value occurs at equality:
\[
\mathcal E(x,y)=0
\iff
x=y.
\]
If $x\neq y$, faithfulness forces $\mathcal E(x,y)\neq0$.  Because the only
other available value is $\infty$, one must have
\[
\mathcal E(x,y)=\infty.
\]
Thus
\[
\mathcal E(x,y)=
\begin{cases}
0,&x=y,\\
\infty,&x\neq y.
\end{cases}
\]
\end{proof}

\begin{remark}[Normalized $0/1$ Version]
\label{rem:clauseII-01}
The structural Boolean cost uses $\{0,\infty\}$.  The normalized operational
cost replaces $\infty$ by $1$:
\[
\mathcal E_{0/1}(x,y)=
\begin{cases}
0,&x=y,\\
1,&x\neq y.
\end{cases}
\]
The Myhill--Nerode threshold is unchanged.
\end{remark}

\begin{remark}[Commitment Vertex]
\label{rem:clauseII-commitment-vertex}
Clause II is the cost-face justification for the Booleanized symmetrized
$\alpha=0$ vertex.  It does not identify commitment with ordinary Schatten
rank or with raw one-sided $D_0$.  The relevant rank is deterministic Boolean
realization rank, i.e.\ the Myhill--Nerode index.
\end{remark}

%----------------------------------------------------------------------
\subsection{Proof of Clause III}
\label{subsec:proof-clause-III}
%----------------------------------------------------------------------

\begin{proof}[Proof of Clause III]
In a $C^*$-algebra, the norm is uniquely determined by the $C^*$-identity:
\[
\|T\|^2
=
\|T^*T\|.
\]
For a positive element $T^*T$, its norm equals its spectral radius:
\[
\|T^*T\|
=
r(T^*T).
\]
In a concrete operator algebra, the spectral radius of the positive operator
$T^*T$ is
\[
r(T^*T)
=
\sup_{\|v\|=1}
\langle v,T^*Tv\rangle
=
\sup_{\|v\|=1}
\|Tv\|^2
=
\|T\|_{\mathrm{op}}^2.
\]
Therefore
\[
\|T\|^2
=
\|T\|_{\mathrm{op}}^2,
\]
and hence
\[
\|T\|
=
\|T\|_{\mathrm{op}}.
\]
\end{proof}

\begin{remark}[Grounding Cost]
\label{rem:clauseIII-grounding}
Clause III justifies using operator-norm distance as the canonical grounding
cost when residual objects form a concrete $C^*$-algebra.  The exact spectral
converse for the unrestricted linear Hankel relaxation then uses
Eckart--Young--Mirsky:
\[
\Dunres(M)
=
\sigma_{M+1}(H_\nu).
\]
The Hankel-structured finite-rank realization gap satisfies
\[
\DHankstr(M)
\ge
\sigma_{M+1}(H_\nu),
\]
with equality only under additional structural hypotheses.  The $C^*$-norm
identification does not by itself prove Hankel domination, Hardy-space
embedding, or exact symbolic deterministic grounding.
\end{remark}

\begin{remark}[What Clause III Supports at the Grounding Face]
\label{rem:clauseIII-grounding-vertex}
Clause III supports the operator-norm / Schatten-$\infty$ face of the
grounding regime.  The R\'enyi $\alpha=\infty$ vertex, max-divergence, is
related to this face under bounded-overlap hypotheses but is not identical to
the $C^*$-norm distance.
\end{remark}

%----------------------------------------------------------------------
\subsection{Scope of the Conditional Representation Theorem}
\label{subsec:conditional-rep-scope}
%----------------------------------------------------------------------

\begin{remark}[Not an Unconditional Derivation]
\label{rem:conditional-not-unconditional}
The Conditional Representation Theorem is conditional by design.  It says that
if a regime's cost satisfies the relevant analytical axioms, then the canonical
cost form follows.  It does not derive KL divergence, Boolean equality, or
operator norm from the enriched-category schema alone.
\end{remark}

\begin{remark}[Regime-Specific Consequences]
\label{rem:conditional-consequences}
The theorem supports the following regime-specific identifications:
\begin{enumerate}[label=(\arabic*)]
\item Retention costs satisfying data processing, convexity, and the chain
rule are scaled KL divergences.
\item Faithful $\{0,\infty\}$ commitment costs are Boolean equality costs.
\item Grounding costs arising from concrete $C^*$-norms are operator-norm
costs.
\end{enumerate}
The approximation gaps built from these costs remain governed by the separate
meta-theorems: information bottleneck for retention, Myhill--Nerode for
commitment, and Eckart--Young--Mirsky for the unrestricted linear Hankel
grounding relaxation.
\end{remark}

%======================================================================
\section{Type Discipline for Approximation Statements}
\label{sec:type-discipline}
%======================================================================

The regimes of this manuscript share one combinatorial resource, the index of
a finite right congruence on histories, but they do \emph{not} share a
semantic category.  Most of the errors that a comparative theory of this kind
invites are not false lemmas but \emph{type errors}: a quantity proved in one
regime, protocol, or aggregation is silently reused in another.  This section
fixes a type signature that every approximation statement in this article
carries, and records the signatures of the principal results, so that a
mismatch between regimes is visible in the statement itself.

\begin{definition}[Type signature of an approximation claim]
\label{def:type-signature}
A claim about a finite-state approximation gap is \textbf{well typed} when the
following seven coordinates are fixed and stated:
\[
\Bigl\langle
\underbrace{\mathsf R}_{\text{regime}},\
\underbrace{\mathsf M}_{\text{machine type}},\
\underbrace{\mathsf F}_{\text{feasible set}},\
\underbrace{\mathsf A}_{\text{aggregation}},\
\underbrace{\mathsf S}_{\text{support}},\
\underbrace{\mathsf P}_{\text{protocol}},\
\underbrace{\mathsf B}_{\text{analytical bridge}}
\Bigr\rangle,
\]
with admissible values
\begin{enumerate}[label=(\arabic*),leftmargin=2.2em]
\item $\mathsf R\in\{\text{commitment},\text{retention},\text{grounding},
\text{joint}\}$;
\item $\mathsf M\in\{\text{deterministic Mealy},\text{stochastic lumped
predictor},\text{unifilar controlled }\epsilon\text{-machine},
\text{linear finite-rank realization}\}$;
\item $\mathsf F\in\{\text{right congruences},
\text{support-relative right congruences},\text{lumpable quotients},
\text{unifilar-lumpable quotients},\text{history factors},
\text{unrestricted rank-}M,\text{Hankel-restricted}\}$;
\item $\mathsf A\in\{\text{worst case},\text{stationary Ces\`aro},
\text{discounted prefix},\text{finite-horizon cumulative}\}$;
\item $\mathsf S\in\{\text{all histories},\text{stationary support }\Splus,
\text{positive-mass reachable}\}$;
\item $\mathsf P\in\{\text{none},\text{reset word},\text{persistent stream},
\text{active query},\text{active identification}\}$;
\item $\mathsf B$ the analytical engine actually invoked, e.g.\
Myhill--Nerode, information bottleneck, Ky Fan, Eckart--Young--Mirsky,
Adamjan--Arov--Krein, Schur test, Le Cam.
\end{enumerate}
\end{definition}

\begin{remark}[Role of Each Coordinate]
\label{rem:type-discipline-coordinates}
Each coordinate marks a hypothesis whose omission changes the truth value of a
statement in this setting.
\begin{itemize}[leftmargin=1.6em]
\item $\mathsf A$ and $\mathsf S$: the Boolean zero-error dichotomy is a
biconditional under worst-case aggregation, but under an average-case law it
holds only with the Nerode relation restricted to $\operatorname{supp}\mu$
(Meta-Theorem~\ref{meta:boolean} and
Remark~\ref{rem:boolean-aggregation-qualifier}).
\item $\mathsf R$ and the meaning of ``predictive state'': equality of
one-step laws is strictly coarser than causal-state equivalence and need not
be a congruence (Definition~\ref{def:z-predictive-equivalence},
Example~\ref{ex:onestep-not-congruence}).
\item $\mathsf F$: ``rank $\le M$'' is three different feasible sets ---
algebraic, bounded-operator, and Hankel-restricted --- with
different values (Remark~\ref{rem:algebraic-vs-operator-rank},
Theorems~\ref{thm:spectral-grounding} and~\ref{thm:aak-equality}).
\item $\mathsf M$ for retention: the input-driven model of
Definition~\ref{def:controlled-markov} is a proper subclass of the unifilar
model of Definition~\ref{def:unifilar-machine}, and lumpability is output-aware
in general (Definition~\ref{def:unifilar-lumpable}).  Conflating the two
changes the feasible set, and with it the value of the gap and the
zero-retention threshold
(Remarks~\ref{rem:unifilar-support-not-automatic}
and~\ref{rem:controlled-zero-not-kernels}).
\item $\mathsf P$: reset-word and persistent-stream mistake bounds are
different theorems; the transfer required an explicit construction
(Theorem~\ref{thm:stream-lower-bound}), and active bounds require an
identification objective (Remark~\ref{rem:active-objective-needed}).
\item $\mathsf B$: Eckart--Young--Mirsky and Adamjan--Arov--Krein solve
different problems, and the latter needs a shift-intertwining embedding, not
merely a unitary (Theorem~\ref{thm:aak-equality}).  That theorem is scalar:
for $|\Sigma|>1$ the equality $\DHankstr(M)=\sigma_{M+1}(H_\nu)$ is
conditional on a multi-shift Nehari/AAK theorem that is not proved here
(Theorem~\ref{thm:aak-multiletter}, Open Problem~\ref{open:hankel-multiletter}).
\end{itemize}
\end{remark}

\begin{table}[htbp]
\centering
\small
\caption{Type signatures of the principal approximation results.}
\label{tab:type-signatures}
\begin{tabular}{@{}p{3.5cm}p{2.0cm}p{2.4cm}p{2.2cm}p{3.0cm}@{}}
\toprule
\textbf{Result} & \textbf{Regime} & \textbf{Feasible set} &
\textbf{Aggregation / support} & \textbf{Bridge} \\
\midrule
Myhill--Nerode threshold,
Theorem~\ref{thm:commitment-spec}
& commitment & right congruences & worst case, all histories
& Myhill--Nerode \\[2pt]
Boolean dichotomy~(i),
Meta-Theorem~\ref{meta:boolean}
& any separated & right congruences & worst case, all histories
& separatedness \\[2pt]
Boolean dichotomy~(ii),
Meta-Theorem~\ref{meta:boolean}
& any separated & right congruences & discounted-prefix $\mu$-average,
$\operatorname{supp}\mu$ & separatedness \\[2pt]
IB identity,
Meta-Theorem~\ref{meta:ib}
& retention & lumpable quotients & stationary, $\Splus$
& information bottleneck \\[2pt]
Zero retention,
Theorem~\ref{thm:retention-zero}
& retention & lumpable quotients & stationary, $\Splus$
& $Z$-predictive distinctness \\[2pt]
Controlled IB identity,
Theorem~\ref{thm:controlled-ib}
& retention & unifilar-lumpable quotients
& stationary, $\Splus$, i.i.d.\ input
& fiberwise information bottleneck \\[2pt]
Zero controlled retention,
Theorem~\ref{thm:controlled-zero}
& retention & unifilar-lumpable quotients
& stationary, $\Splus$, i.i.d.\ input
& stable kernel refinement \\[2pt]
Quadratic converse,
Theorem~\ref{thm:spectral-retention}
& retention & lumpable quotients & stationary, $\Splus$
& Ky Fan / ANOVA \\[2pt]
Unrestricted grounding,
Theorem~\ref{thm:spectral-grounding}
& grounding & unrestricted rank-$M$ & operator norm on $H_\nu$
& Eckart--Young--Mirsky \\[2pt]
Structured grounding,
Theorem~\ref{thm:aak-equality}
& grounding & Hankel-restricted, $|\Sigma|=1$ & worst case, all histories
& Adamjan--Arov--Krein, given a shift-intertwining embedding \\[2pt]
Determinism gap,
Theorem~\ref{thm:exp-gap}
& grounding & \emph{algebraic} rank & worst case, all histories
& PFA succinctness, rate $\mathcal S(k)$ \\[2pt]
Safe retention,
Corollary~\ref{cor:price-safety}
& joint & history factors & stationary, $\Splus$
& information bottleneck \\
\bottomrule
\end{tabular}
\end{table}

\begin{table}[htbp]
\centering
\small
\caption{Type signatures of the temporal results.  Protocol is the
distinguishing coordinate.}
\label{tab:type-signatures-temporal}
\begin{tabular}{@{}p{4.2cm}p{3.0cm}p{2.4cm}p{3.4cm}@{}}
\toprule
\textbf{Result} & \textbf{Protocol} & \textbf{Alphabets} &
\textbf{Status} \\
\midrule
Reset-word complexity
& reset word & $|\mathcal I|,|\mathcal O|\ge2$
& $\Theta(M\log M)$, exact \\[2pt]
Persistent-stream upper bound,
Theorem~\ref{thm:passive-realizable}
& persistent stream & $|\mathcal I|,|\mathcal O|\ge2$
& $O(M\log M)$, unconditional \\[2pt]
Persistent-stream lower bound,
Theorem~\ref{thm:stream-lower-bound}
& persistent stream & $|\mathcal I|\ge4$
& $\Omega(M\log M)$, deterministic; expectation and
high-probability for randomized \\[2pt]
Binary-input lower bound,
Theorem~\ref{thm:stream-lb-binary}
& persistent stream & $|\mathcal I|=2$
& $\Omega(M\log M)$, constant $c<\tfrac13$ \\[2pt]
Agnostic regret,
Theorem~\ref{thm:agnostic}
& both & $|\mathcal I|,|\mathcal O|\ge2$
& $\Theta(\sqrt{TM\log M})$ \\[2pt]
Active upper bound,
Theorem~\ref{thm:active-certified}(I)
& active query & fixed finite
& $O(\Esync)$ under~(RI), unconditional \\[2pt]
Active $\Esync$ lower bound,
Theorem~\ref{thm:active-certified}(II)
& active identification & fixed finite
& conditional on operational equivalence \\[2pt]
Active additive lower bound,
Theorem~\ref{thm:active-certified}(III)
& active identification & $|\mathcal I|\ge4$
& conditional on direct-sum saturation; witnessed by the two-component
construction of Theorem~\ref{thm:active-two-phase} \\[2pt]
Active complexity, Corollary~\ref{cor:active-theta}
& active identification & $|\mathcal I|\ge2$
& $\Theta(M\log M)$, unconditional \\
\bottomrule
\end{tabular}
\end{table}

\begin{remark}[How to read a gap symbol]
\label{rem:gap-symbol-convention}
Where ambiguity is possible, a gap carries its distinguishing coordinates
explicitly.  Thus
\[
\Delta_{\mathbb T}^{\sup}(M),
\qquad
\Delta_{\mathbb T}^{\mu}(M),
\qquad
\Dunres(M),
\qquad
\DHankstr(M),
\qquad
\operatorname{Mistakes}^{\text{stream}}_{\mathrm{realizable}}(M)
\]
are distinct objects, and no inequality between symbols of different type is
asserted without an explicit bridge.  In particular, the absence of a
universal cross-regime ordering
(Theorem~\ref{thm:schatten-nogo}) is a statement about incomparable types, not
a gap in the analysis.
\end{remark}

%======================================================================
\section{Scope of Results and Open Problems}
\label{sec:epistemic}
%======================================================================

The schema provides a common structural syntax.  It does not derive the
Myhill--Nerode theorem, the Eckart--Young--Mirsky theorem, or the
information-bottleneck objective from categorical axioms alone.  The
analytical engines remain domain-specific.

The model distinguishes:
\begin{itemize}[leftmargin=1.6em]
\item shared combinatorial syntax;
\item regime-specific semantics;
\item exact theorems;
\item conditional theorems;
\item local or surrogate theorems;
\item heuristics;
\item open problems.
\end{itemize}

The classification below states, for each principal result, the hypotheses
under which it holds, and lists the problems that remain open.

%----------------------------------------------------------------------
\subsection{Status Categories}
\label{subsec:status-categories}
%----------------------------------------------------------------------

The lists below use the following status labels.

\begin{itemize}[leftmargin=1.6em]
\item \textbf{Exact.}
A theorem proved under the stated hypotheses, with no additional domination
or achievability assumption required.

\item \textbf{Conditional.}
A theorem proved under an explicit additional hypothesis, such as Schatten
domination, stable decay, bounded overlap, or spectral-rate achievability.

\item \textbf{Local.}
A theorem valid in a small-radius or second-order regime under explicit
regularity assumptions.

\item \textbf{Heuristic.}
An organizing principle or diagnostic, not a theorem without further
hypotheses.

\item \textbf{Open.}
A problem left unresolved here, stated explicitly as an open problem.
\end{itemize}

%----------------------------------------------------------------------
\subsection{Exact Results}
\label{subsec:exact-results}
%----------------------------------------------------------------------

The following results are exact under the hypotheses stated in the body of the
manuscript.

\begin{itemize}[leftmargin=1.6em]
\item Retention zero iff $M\ge|\Splus|$, assuming distinct predictive
distributions on the stationary support.
\item Full-KL information-bottleneck identity
\[
\RetKL(\phi)=I(S;Z\mid K_\phi)
=
I(S;Z)-I(K_\phi;Z).
\]
\item Quadratic restricted spectral converse
\[
\RetQuad(M)
\ge
\frac12
\sum_{i\ge M}
\lambda_i(\Sigma_\pi).
\]
\item Unifilar retention: the controlled information-bottleneck identity
(Theorem~\ref{thm:controlled-ib}), the zero-retention threshold at the index of
the stable kernel refinement (Theorem~\ref{thm:controlled-zero}), and the
fiberwise quadratic, probability-coordinate, and interior-Fisher spectral
converses
(Corollaries~\ref{cor:controlled-quad-spectral}--\ref{cor:controlled-fisher}).
\item Gaussian quadratic restricted retention NP-completeness for reset
machines of synchronization depth one.
\item Tagged depth-two NP-hardness in the same regime.
\item Commitment exact threshold
\[
\Com(M)=0
\iff
M\ge\kappa_{\det}(F).
\]
\item Safety controller complexity $\kappa_{\mathrm{ctrl}}(\Lambda)$.
\item Strategic-spread lower bound with explicit myopic adversary.
\item Stateless-game exact corollary $\ComGame(M)=\alpha/(1-\gamma)$.
\item Booleanized symmetrized $\alpha=0$ commitment vertex.
\item Unrestricted linear Hankel relaxation converse
\[
\Dunres(M)
=
\sigma_{M+1}(H_\nu).
\]
\item Hankel-structured lower bound
\[
\DHankstr(M)\ge\sigma_{M+1}(H_\nu).
\]

\item Stochasticity floor for deterministic Mealy grounding, and its exact
observable refinement: $\sigma_\gamma(\nu)\le F_\gamma(\nu)$ with
$F_\gamma(\nu)$ the exact input-only deterministic floor and
$\lim_{M\to\infty}\Delta_{\mathrm{grd}}(M;\gamma)=F_\gamma(\nu)$
(Theorem~\ref{thm:observable-floor}).
\item Finite-section operator certificates for linear grounding: computable
two-sided bounds on $\sigma_{M+1}(H_\nu)=\Dunres(M)$ from finite Hankel
truncations $H_N$, with an explicit lower certificate
$\DHankstr(M)\ge\sigma_{M+1}(H_N)$ for the structured problem
(Proposition~\ref{prop:grounding-finite-section}).
\item Exact zero threshold for the Hankel-structured linear problem:
$\DHankstr(M)=0\iff\rank(H)\le M$
(Proposition~\ref{prop:grounding-structured-zero}).
\item Finite-horizon symbolic certificates: a two-sided sandwich
$\Delta^{(T)}_{\mathrm{grd}}(M;\gamma)\le\Delta_{\mathrm{grd}}(M;\gamma)\le
\Delta^{(T)}_{\mathrm{grd}}(M;\gamma)+\gamma^T/(1-\gamma)$, decidable and
rational under rational finite-state and rational-$\gamma$ hypotheses
(Proposition~\ref{prop:symbolic-finite-horizon}).
\item Refinement monotonicity of the tracking deficit: $\phi$ refining $\psi$
implies $D(\phi)\le D(\psi)$, so $\min_{|\mathcal K|\le M}D(\phi)$ is
nonincreasing in $M$ (Correction to item (iii) of
Proposition~\ref{prop:grounding-tracking}).
\item Joint zero-error memory theorem via meet of right congruences.
\item Comparative state-complexity sandwich.
\item Independence of regime obstructions.
\item Interaction theorem for exact commitment and approximate retention.
\item Price of Safety in mutual-information form.
\item Reset-word realizable mistake complexity $\Theta(M\log M)$.
\item \textbf{Persistent-stream realizable mistake complexity
$\Theta(M\log M)$}, with the lower bound proved by an explicit transport-plus-
readout adversary on a single never-reset stream
(Theorem~\ref{thm:stream-lower-bound}).
\item Active upper bound $O(\Esync(M))$ under objective~(RI)
(Theorem~\ref{thm:active-certified}(I)), and
$O(M\log M+\Lsyncu(M))$ in universal length form
(Proposition~\ref{prop:active-length-upper}).
\item \textbf{Active realizable mistake complexity
$\Theta(M\log M)$, unconditionally}: an active halving learner attains
objective~(RI) within $\log_2|\mathcal H_M\times Q|=O(M\log M)$ mistakes
(Theorem~\ref{thm:active-halving}), matching the gated-family lower bound, so
$\Esync(M)=\Theta(M\log M)$ with no synchronizability,
operational-equivalence, or direct-sum hypothesis
(Corollary~\ref{cor:active-theta}).
\item \textbf{Explicit gated family} $\Gact_M$ with
$\Esync(M)=\Omega(M\log M)$, giving an active lower bound
$\Omega(M\log M)$ relative to the residual-identification objective
(Theorem~\ref{thm:active-explicit-directsum}).  The unknown datum is a
transition map, so the bound is for model identification; it does not witness a
state-identification component.
\item Conditional additive active lower bound under direct-sum saturation
(Theorem~\ref{thm:active-direct-sum}), which separates the two terms; the
unconditional active lower bound itself is
Corollary~\ref{cor:active-theta}.
\item \textbf{Agnostic minimax regret $\Theta(\sqrt{T M\log M})$} in both the
reset-word and persistent-stream protocols
(Theorem~\ref{thm:agnostic}, Corollary~\ref{cor:stream-ldim}).
\item Fixed-budget bias--variance decomposition.
\item Adaptive agnostic oracle inequality under bounded losses.
\item Realizable $0/1$ and mixable log-loss oracle inequalities under stated
conditions.
\item Linear pinching Price-of-Safety surrogate nonnegativity and Ky Fan
majorization.
\item Retention as exact $\alpha=1$ vertex.
\item Commitment as Booleanized symmetrized $\alpha=0$ vertex.
\item Grounding as $\alpha=\infty$ supremum/operator-norm face, related but
not identical to max-divergence.
\end{itemize}

%----------------------------------------------------------------------
\subsection{Conditional, Local, and Surrogate Results}
\label{subsec:conditional-results}
%----------------------------------------------------------------------

The following results are conditional, local, or surrogate statements.

\begin{itemize}[leftmargin=1.6em]
\item Local Fisher asymptotic lower bound for full KL under regularity
assumptions.
\item Conditional unified Schatten converse under response representation and
domination.
\item Frobenius / Schatten-$2$ converse under explicit Schatten-$2$
domination.
\item Stable-decay residual-norm spectral bound.
\item Max-divergence / operator-norm local equivalence under bounded-overlap
hypotheses.
\item Superpolynomial determinism gap for formal Hankel rank: the reduction is
exact, but the statement is conditional on the external bounded-error
succinctness family (Theorem~\ref{thm:exp-gap}).
\item Budget laws under spectral-rate assumptions.
\item Linear block-local Price-of-Safety surrogate.  It is unconditionally
nonnegative and bounded by the Ky Fan coupling term, but its identification
with the discrete Price of Safety requires a two-sided bound on the signed
discrepancy $\rho_{\mathrm{safe}}-\rho_{\mathrm{free}}$, neither term of which
has a determined sign (Section~\ref{subsec:pos-discrete}).
\item Spectral-tail behaviour, which is a heuristic organizing principle
requiring a regime-specific response operator and domination bridge rather
than following from aggregation
(Heuristic~\ref{heur:spectral}).
\item The vertex correspondence $\alpha=p$, which is a labelling consistency
between R\'enyi order and Schatten exponent and carries no derivation of the
Mirsky converses (Section~\ref{sec:exponent}).
\item The commitment vertex, which is a Boolean-rank and Myhill--Nerode
threshold rather than an ordinary Schatten $p\to0$ limit
(Theorem~\ref{thm:commitment-alpha0}).
\item The Hankel-structured converse, which attains the Eckart--Young--Mirsky
value only under a Hardy-space embedding
(Theorems~\ref{thm:aak-equality} and~\ref{thm:aak-multiletter}).
\item The block-mode tracking construction, which realizes a deterministic
Mealy machine of the claimed stationary error only under the explicit
input-driven synchronization and right-congruence implementability
hypothesis of Proposition~\ref{prop:block-mode-implementability}; absent that
hypothesis the same error formula is only the loss of a genie that observes
the latent block label.
\end{itemize}

%----------------------------------------------------------------------
\subsection{Open Problems}
\label{subsec:open-problems}
%----------------------------------------------------------------------

The following open problems remain.

\begin{enumerate}[leftmargin=2.0em]
\item \textbf{Exact symbolic deterministic grounding gap (partially resolved).}
\emph{Hypotheses.} $0<\gamma<1$, finite input alphabet, target channel $\nu$.
\emph{Established baseline.} Theorem~\ref{thm:observable-floor} identifies
the exact input-only deterministic floor $F_\gamma(\nu)$ and proves
$\lim_{M\to\infty}\Delta_{\mathrm{grd}}(M;\gamma)=F_\gamma(\nu)$, with
$\sigma_\gamma(\nu)\le F_\gamma(\nu)$ possibly strict.
Proposition~\ref{prop:symbolic-finite-horizon} supplies, for every finite
horizon $T$, a two-sided certified interval for
$\Delta_{\mathrm{grd}}^{\mathrm{symb}}(M)=\Delta_{\mathrm{grd}}(M;\gamma)$ of
width at most $\gamma^T/(1-\gamma)$, decidable and rational for rational
finite-state channels and rational $\gamma$.  What remains open is a
\emph{closed-form value at finite $M$}, not merely a limit or a certified
interval, and its relation to the unrestricted linear finite-rank Hankel gap
$\Dunres(M)=\sigma_{M+1}(H_\nu)$, to the stochasticity floor $\sigma_\gamma$,
and to the superpolynomial determinism gap of
Theorem~\ref{thm:exp-gap} (see Open Problem~\ref{open:symbolic-grounding}).
\emph{Success criterion.} A formula or algorithm computing
$\Delta_{\mathrm{grd}}^{\mathrm{symb}}(M)$ exactly for every finite $M$ (not
just $M\to\infty$ or a fixed horizon $T$), together with either a proof or a
counterexample to $\Delta_{\mathrm{grd}}^{\mathrm{symb}}(M)\ge\sigma_{M+1}(H_\nu)$
for all $M$.

\item \textbf{Hankel-structured equality conditions.}
Characterize necessary and sufficient conditions under which
\[
\DHankstr(M)
=
\sigma_{M+1}(H_\nu).
\]
Theorem~\ref{thm:aak-equality} settles the sufficient direction: a Hardy-space
symbol with $\varphi\in H^\infty+C(\mathbb T)$ forces equality for
\emph{every} $M$, with an attained optimum of rational symbol degree at most
$M$.  Corollary~\ref{cor:hankel-strict} converts this into a necessary
condition on the operator: a strict gap at any $M$ rules out Hardy-space
representability.  What remains open is the intrinsic characterization -- an
$\ell^2((\mathcal I\times\mathcal O)^*)$-side criterion, stated directly in
terms of $\nu$, deciding when such a symbol exists.

\item \textbf{Relaxation gaps for the Price of Safety.}
Prove explicit bounds relating the linear block-local Price-of-Safety
surrogate
\[
\PoSlin(M)
\]
to the discrete right-congruence Price of Safety
\[
\PoSquad(M)
\]
or the full-KL Price of Safety.

\item \textbf{Exact and fixed-alphabet complexity of full-KL retention.}
The unrestricted full-KL decision problem is: given $\{P_s\}$, $\{\pi_s\}$, a
budget $M\in\mathbb N$, lumpability constraints, and $\tau$, decide whether
\[
\min_{\substack{\phi\ \text{lumpable}\\|\mathcal K|\le M}}
\sum_k\sum_{s:\phi(s)=k}\pi_s\KL\bigl(P_s\|\bar P_k\bigr)
\le\tau .
\]
The budget $M$ is present because it is what makes the problem a genuine
finite-state compression decision rather than a bare clustering feasibility
question at the natural number of distinct rows; the case $M\ge|\Splus|$ is
trivial by Theorem~\ref{thm:retention-zero}, so the decision problem is of
interest exactly for $M<|\Splus|$.
Because the optimal block representative is the mixture centroid
$\bar P_k$ (Lemma~\ref{lem:mixture-centroid}), the objective is exactly the
weighted \emph{Jensen--Shannon} clustering cost, so the problem is
\emph{constrained Jensen--Shannon clustering with lumpability constraints}.
The Gaussian quadratic regime of Theorem~\ref{thm:retention-reset-np} is its
second-order Fisher special case, where the centroid degenerates to a mean and
the cost to squared Euclidean distance, which is how the $k$-means reduction
applies there.  That reduction is carried out in Theorem~\ref{thm:full-kl-promise-np}: the
obstruction --- that KL centroids are not rational in the inputs, so the
promise gap cannot be produced by denominator clearing --- is met by embedding
$k$-means into a shrinking interior neighbourhood of the uniform law, where the
Jensen--Shannon cost is a controlled perturbation of the Euclidean one and the
gap is preserved analytically.  Four questions remain.

\begin{enumerate}[label=(\alph*),leftmargin=2.0em]
\item \textbf{Membership in NP.}  The objective compares a rational threshold
against a sum of terms $p\log(p/q)$ with $p,q$ rational.  Deciding such
inequalities in polynomial time is not known, so exact NP-completeness is open
even though hardness is proved.

Proposition~\ref{prop:full-kl-algebraic-form} locates the difficulty exactly.
For a fixed partition the whole sum collapses,
$\RetKL(\phi)=D^{-1}\log R$ with $D$ a positive integer and $R$ a positive
rational presented as a product of powers, so the comparison with $\tau$ is
equivalent to deciding $R\ge e^{D\tau}$.  The number of logarithmic terms is
therefore not the obstruction; the single rational-versus-exponential
comparison is.  That comparison is decidable, by
Lindemann--Weierstrass as recorded in
Remark~\ref{rem:full-kl-exp-membership}, and what is missing is a polynomial
bound on the precision needed to resolve it, equivalently a polynomial lower
bound on $|R-e^{D\tau}|$.  Remark~\ref{rem:full-kl-slp} shows that $R$ must be
kept in factored form for this to be a meaningful question, its expanded
numerator and denominator being of exponential bit size.  Consequently the
problem lies in NP as soon as this $\textsc{PosSLP}$-type comparison admits
polynomial-time verifiable certificates.
\item \textbf{Fixed output alphabet.}  The embedding of
Theorem~\ref{thm:full-kl-promise-np} uses $|\mathcal O|=2d$, growing with the
dimension of the $k$-means instance.  Whether the problem is hard for a fixed
alphabet, binary in particular, is open.

Remark~\ref{rem:output-alphabet-2d} determines how far the alphabet can be
reduced within this scheme, and the answer is arithmetic rather than a
universal constant: $|\mathcal O|$ may be taken as small as the least
$n\ge d+1$ whose zero-sum subspace in $\mathbb Q^{n}$ contains $d$ pairwise
orthogonal rational vectors of equal norm.  The value $n=2d$ always qualifies,
by the doubling map; $n=d+1$ qualifies for $d=3,7,15,\dots$, by the
non-constant rows of a Hadamard matrix of order $d+1$; and $n=d+1$ fails for
$d=2$, by a discriminant obstruction.  What no choice of scheme can avoid is
$|\mathcal O|\ge d+1$, since the zero-sum subspace of
$\mathbb R^{|\mathcal O|}$ has dimension $|\mathcal O|-1$ and must accommodate
$d$ independent directions.  A fixed alphabet therefore caps the embeddable
dimension, and a fixed-alphabet hardness proof would have to start from a
source problem hard in bounded dimension rather than from Euclidean
$k$-means.
\item \textbf{Hardness without a promise.}  The reduction separates instances
whose values differ by at least the amplified gap; exact comparison against a
rational threshold is a different question.
\item \textbf{Nonvacuous lumpability.}  The reduction makes every partition
lumpable.  The complexity under genuinely restrictive right-congruence
constraints is not addressed.
\end{enumerate}

\item \textbf{Domination bridges.}
Identify task-theoretic conditions under which contextwise aggregation
implies Schatten domination, or under which frame-type lower bounds hold.

\item \textbf{Intermediate exponent regimes.}
Determine whether intermediate R\'enyi orders
\[
\alpha\in(0,\infty)\setminus\{1\}
\]
can be made operational as finite-state approximation regimes with matching
rank-control and domination structures.

\item \textbf{Active learning without synchronization.}
Develop active realizable and agnostic bounds when no output-aware
synchronizing word exists, or when synchronization is only partial.

\item \textbf{Hardy-space embeddings for grounding.}
Identify when a Hankel operator admits a Hardy-space embedding making
Adamjan--Arov--Krein applicable, and compare the resulting converses with
the default Eckart--Young--Mirsky converse.

\item \textbf{Deterministic commitment rate-distortion.}
\emph{Hypotheses.} A fixed history law $\mu$ on $\mathcal I^{\mathbb N}$ and a
causal length-preserving specification $F$ (Definition~\ref{def:causal-lp-spec}).
\emph{Established baseline.} Theorem~\ref{thm:com-rd-formula} gives the exact
fixed-budget value of $\ComRD(M)$ over right-congruence quotients of a
one-step observable, and its zero threshold at
$M=\kappa_{\mathrm{pair}}(F,\mu)$, the one-step determination index; Remark~\ref{rem:com-rd-scope} shows
this is generically distinct from, and no more than, the worst-case Boolean
gap $\Comex(M)$.  What is open is the \emph{full multi-step} rate-distortion
law: the exact value of the $M$-state quotient loss for approximating the
entire residual process $(F_u)_{u\in\mathcal I^*}$ under $\mu$-average loss,
not just its one-step marginal, together with a converse comparable in
strength to the Blahut--Arimoto-type characterization available for classical
rate-distortion.
\emph{Success criterion.} A formula (or a provably tight two-sided bound with
matching rate) for the $M$-state $\mu$-average multi-step commitment
distortion, together with a decidable finite-horizon certificate analogous to
Proposition~\ref{prop:symbolic-finite-horizon}.

\item \textbf{Typed joint quotients across the linear/symbolic boundary.}
\emph{Hypotheses.} Two tasks, one governed by a right-congruence quotient
budget (commitment, retention, or symbolic grounding under
Proposition~\ref{prop:block-mode-implementability}) and one governed by the
unrestricted linear finite-rank Hankel relaxation of
Theorem~\ref{thm:spectral-grounding}.
\emph{Established baseline.} Theorem~\ref{thm:comparative} and
Remark~\ref{rem:comparative-grounding-excluded} prove the product sandwich
\emph{only} when every component's budget is a right-congruence index, and
explicitly exclude the unrestricted linear coordinate: no representative
map is known that combines an $M$-state quotient with a rank-$M'$ Hankel
truncation into one object of controlled joint budget.
What is open is whether \emph{any} nontrivial joint budget notion
$M_{\mathrm{joint}}$, mixing a quotient index and an operator rank, admits
even a one-sided analogue of the sandwich of Theorem~\ref{thm:comparative}.
\emph{Success criterion.} Either (a) a bridge construction -- a map from
rank-$M'$ Hankel approximants to $O(f(M'))$-state quotients (or conversely)
that preserves admissibility and control loss by a stated factor, from which
a mixed sandwich would follow by composing with Theorem~\ref{thm:comparative};
or (b) a proof, e.g.\ via an explicit family of channels, that no such bridge
with subexponential $f$ can exist.

\item \textbf{Strategic memory versus information.}
\emph{Hypotheses.} A commitment task admitting both the distributional gap
$\ComRD(M)$ and the simultaneous-move strategic gap $\ComGame(M)$ of
Section~\ref{subsec:quant-simul}.
\emph{Established baseline.} Remark~\ref{rem:com-rd-scope} proves these two
quantities are posed on different protocols (averaging under a fixed law
versus play against an online adversary) and asserts no general inequality
between them; Section~\ref{subsec:quant-simul} computes $\ComGame(M)$ exactly
in the stateless-game case.  What is open is whether, under an explicit
correspondence between the adversary's input distribution and a history law
$\mu$ (e.g.\ the adversary's minimax-optimal mixed strategy), $\ComGame(M)$
and $\ComRD(M)$ become comparable -- whether strategic memory (states needed
to play well against an adaptive adversary) is provably bounded above or
below by informational memory (states needed to match a fixed one-step
predictive law).
\emph{Success criterion.} Either a proved inequality
$\ComGame(M)\le g(\ComRD(M'))$ or $\ComRD(M)\le g(\ComGame(M'))$ for an
explicit function $g$ and budget map $M\mapsto M'$, under a stated
correspondence between the adversary's strategy and $\mu$; or an explicit
task family showing the two quantities are order-independent (each can be
made arbitrarily larger than the other at matching budgets), closing the
question negatively.

\item \textbf{Necessity of the aggregation penalty for nested classes.}
Proposition~\ref{prop:aggregation-necessary} settles necessity for arbitrary
experts; what remains is the nested case.
Theorem~\ref{thm:oracle-minimax-lower} proves a two-sided minimax
characterization of the adaptive oracle envelope under the explicit
approximation and estimation floors of
Assumption~\ref{ass:oracle-floors}, matching the adaptive upper bound up to
the aggregation penalty $\PsiAgg_M(T)$ and universal constants.  The floors are
discharged for the realizable persistent-stream regime by
Proposition~\ref{prop:floors-instance} and assumed in the stochastic regimes.
What remains open is whether $\PsiAgg_M(T)$ is itself necessary, i.e.\ whether
adapting to an unknown budget is strictly harder than the oracle problem, and
whether horizon-efficient packings exist for the retention, grounding, and
commitment regimes.  See Remark~\ref{rem:oracle-minimax-residual}.
\end{enumerate}

%----------------------------------------------------------------------
\subsection{Classification of Results}
\label{subsec:final-epistemic}
%----------------------------------------------------------------------

The schema supplies a shared finite-state syntax, while the analytical engines
remain regime-specific.  Results are classified as exact, conditional, local,
heuristic, or open according to the hypotheses under which they are proved.

%======================================================================
\section{Conclusion}
\label{sec:conclusion}
%======================================================================

Finite-index right congruences on histories provide a common combinatorial
resource for the three regimes, while the analytical engines governing them
remain regime-specific.  The principal conclusions are the following.

\begin{itemize}[leftmargin=1.8em]
\item \textbf{Shared syntax, regime-specific semantics.}
Regime-typed histories, finite-index right-action quotients, and variational
gaps form a common finite-state syntax: input histories for commitment,
feasible joint histories for unifilar retention, the joint alphabet for
grounding.  Commitment, retention, and grounding attach different semantic
objects to states: residual functions, predictive distributions or kernels,
and operators or kernels.

\item \textbf{Stationary aggregation.}
The average-case theory uses stationary Ces\`aro aggregation or discounted
prefix measures.  Finite prefixes do not define a probability measure over
all lengths.

\item \textbf{Retention has two distinct gaps.}
The full-KL gap satisfies an exact finite-state information-bottleneck
decomposition over stationary lumpable quotients:
\[
\RetKL(M)
=
I(S;Z)
-
\max_{\substack{
\phi\ \text{lumpable}\\
|\mathcal K|\le M
}}
I(K_\phi;Z).
\]
Zero retention occurs exactly when $M$ is at least the stationary support size
of distinct predictive states.  The Gaussian quadratic restricted gap satisfies
a spectral lower bound with effective rank $M-1$.  Full KL has a precise local
Fisher asymptotic lower bound under regularity assumptions.

The same analysis holds for general unifilar controlled machines: the
controlled information-bottleneck decomposition conditional on the input, and
fiberwise quadratic, probability-coordinate, and interior-Fisher spectral
converses, with the input-driven theory as the output-independent
specialization.  The zero-retention threshold, by contrast, does not transfer
verbatim: it is the index of the coarsest unifilar-lumpable refinement of the
predictive-kernel partition, which can strictly exceed the number of distinct
kernels (Theorem~\ref{thm:controlled-zero}).

\item \textbf{Grounding spectral converses are linear finite-rank.}
The unrestricted linear finite-rank relaxation satisfies
\[
\Dunres(M)
=
\sigma_{M+1}(H_\nu)
\]
by Eckart--Young--Mirsky.  The Hankel-structured finite-rank realization gap
satisfies
\[
\DHankstr(M)
\ge
\sigma_{M+1}(H_\nu),
\]
with equality only under additional structural hypotheses, and finite
Hankel truncations supply computable two-sided certificates for both
quantities (Proposition~\ref{prop:grounding-finite-section}).  The structured
problem has an exact zero threshold at $\rank(H)\le M$
(Proposition~\ref{prop:grounding-structured-zero}).  For the symbolic
deterministic (Mealy-machine) grounding gap, the exact input-only
deterministic floor $F_\gamma(\nu)$ and its finite-state limit are now
theorems (Theorem~\ref{thm:observable-floor}), with decidable finite-horizon
certified intervals (Proposition~\ref{prop:symbolic-finite-horizon}); a
closed-form value at finite $M$, and its relation to the linear relaxation,
remain open (Open Problem~\ref{open:symbolic-grounding}).  AAK is used only
under separate Hardy-space embeddings.

\item \textbf{Commitment is Boolean rank.}
Exact commitment is governed by Myhill--Nerode index.  The exponent
correspondence identifies commitment with the Booleanized symmetrized
$\alpha=0$ vertex, not with raw $D_0$ and not with ordinary Schatten rank.

\item \textbf{The spectral-tail picture is a heuristic.}
Spectral tails arise from regime-specific response operators, effective rank
budgets, and domination bridges.  Aggregation alone does not determine the
tail.  No unconditional cross-regime ordering follows.

\item \textbf{The Schatten template is conditional.}
Mirsky's theorem unifies the analytic converse form once a Schatten
domination bridge is supplied.  Grounding is the $p=\infty$ instance for the
unrestricted linear finite-rank relaxation; retention is the $p=1$ instance
with effective rank $M-1$; commitment is the Boolean-rank vertex.

\item \textbf{The oracle bridge is adaptive and type-correct.}
Offline approximation and online estimation compose into a bias--variance
decomposition.  Adaptive aggregation attains a penalized minimum without
knowledge of the optimal budget.  Spectral tails govern online rates only
under explicit achievability assumptions.

\item \textbf{The oracle floors are discharged in one regime only.}
Assumption~\ref{ass:oracle-floors} is a minimax separation hypothesis.  It is
established unconditionally for the realizable persistent-stream regime by
Proposition~\ref{prop:floors-instance}, where the packing is the gated family
$\Gact_M$ and $\Delta_M=M\log M$; there
Theorem~\ref{thm:oracle-minimax-lower} is a theorem rather than a reduction.
For the retention, grounding, and commitment regimes the floors are not
established, and there the theorem remains a conditional reduction.  The
discharged case is realizable, so it certifies the estimation axis only; by
Remark~\ref{rem:floors-two-axes} the sum-form envelope follows from the minimum
form through the discrete sandwich, and by Lemma~\ref{lem:kappa-bounded} the
resulting constant $(1+\kappa)^{-1}$ is absolute rather than
instance-dependent.  See Section~\ref{subsec:open-problems}.

\item \textbf{The matching minimax lower bound is conditional on explicit
floors.}
An adaptive oracle upper bound and a conditional continuous bias--variance
sandwich hold in general, and under the floors of
Assumption~\ref{ass:oracle-floors} a matching minimax lower bound holds as well
(Theorem~\ref{thm:oracle-minimax-lower}); the envelope is therefore optimal up
to the aggregation penalty $\PsiAgg_M(T)$.  The constant in the sum-form
version is absolute, since Lemma~\ref{lem:kappa-bounded} bounds the jump ratio
by $2^\alpha(1+\log2)^\beta$ for a regularly varying estimation envelope.  The
floors themselves are proved for the realizable persistent-stream regime
(Proposition~\ref{prop:floors-instance}) and assumed elsewhere.  Whether the
penalty $\PsiAgg_M(T)$ is necessary, and whether the floors hold in the
stochastic regimes at horizons $T$ of the order at which the envelopes cross,
remain open.

\item \textbf{The linear Price-of-Safety surrogate is exact but distinct.}
Pinching gives an exact majorization law for a linear second-order,
state-indexed, block-local surrogate.  It does not compute the discrete
right-congruence Price of Safety without relaxation-error control.

\item \textbf{The exponent correspondence is a vertex correspondence.}
R\'enyi order, Schatten exponent, and aggregation label share vertex labels
$\{0,1,\infty\}$, but they are distinct layers.  The equality $\alpha=p$ is
a formal consistency, not a derivation.

\item \textbf{Temporal claims are protocol-specific.}
For the reset-word protocol, realizable finite-automaton learning has
$\Theta(M\log M)$ mistake complexity.  The same $\Theta(M\log M)$ holds for a
single \emph{persistent} Mealy stream that is never reset: the upper bound is
finite-class counting, and the matching lower bound is proved in
Theorem~\ref{thm:stream-lower-bound} by an adversary that makes state
navigation free and extracts $\Omega(M\log M)$ bits through forced-mistake
readouts.  Finite-class aggregation gives agnostic regret
\[
O\bigl(\sqrt{T M\log M}\bigr).
\]
For active learning, all bounds are stated against the residual-identification
objective, an identification requirement being unavoidable in a protocol where
the learner selects its own inputs.  The sharp upper bound is
\[
\MistRI(M)
\le
O(\Esync(M)),
\]
with the universal length form $O(M\log M+\Lsyncu(M))$ when the learner must
fix its synchronization experiment in advance.  On the full class
$\mathcal H_M$ the complexity is settled outright,
\[
\MistRI(M)
=
\Theta(M\log M),
\]
the upper bound by an active halving learner over the version space of
machine-state pairs, the lower bound by the gated adversarial family.  The
$\Omega(M\log M)$ lower bound comes from that explicit construction,
not from ordinary Littlestone dimension, which does not transfer to a protocol
in which the learner chooses its own inputs.

A direct-sum instance forces its two components on disjoint rounds, giving
$\MistRI(M)\ge S_M+C_M$, but no non-degenerate two-term decomposition of the
active complexity is available in the mistake currency: state identification in
a known minimal skeleton costs only $O(\log M)$ mistakes
(Proposition~\ref{prop:esyncsi-log}).  Whether a genuine decomposition exists
in the experiment-length currency is
Open Problem~\ref{open:si-hard-family}.

\item \textbf{NP-hardness is regime-qualified.}
The Gaussian quadratic restricted retention decision problem is NP-complete
already for minimal reset machines of synchronization depth one, and a tagged
depth-two construction also gives NP-hardness.  For unrestricted full-KL
retention over a finite output alphabet,
Theorem~\ref{thm:full-kl-promise-np} gives NP-hardness under a promise and
Corollary~\ref{cor:full-kl-apx} upgrades this to APX-hardness, so no
polynomial-time approximation scheme exists unless $\mathsf P=\mathsf{NP}$.
Two qualifications attach to those two results, and both have been located
precisely.  The output alphabet produced by the embedding has size $2d$, so
hardness for a \emph{fixed} alphabet is not established; by
Remark~\ref{rem:output-alphabet-2d} the size can be reduced to $d+1$ exactly
when the zero-sum subspace of $\mathbb Q^{d+1}$ admits an orthogonal rational
basis of equal norms, which holds for $d=3,7,15,\dots$ and fails for $d=2$, but
it can never be taken below $d+1$, so the alphabet grows with the dimension
under any variant of the construction.  Membership in NP is open; by
Proposition~\ref{prop:full-kl-algebraic-form} the objective is $D^{-1}\log R$
for a factored rational $R$, so the question is the single comparison
$R\ge e^{D\tau}$, decidable by Lindemann--Weierstrass but lacking a polynomial
precision bound.
\end{itemize}

Each result above is stated under the hypotheses given in the corresponding
section, and its type signature in the sense of
Definition~\ref{def:type-signature} fixes the regime, feasible set,
aggregation, support, protocol, and analytical bridge to which it applies.

%======================================================================
\begin{thebibliography}{99}
%======================================================================

\bibitem{cover}
T.~M.~Cover and J.~A.~Thomas,
\emph{Elements of Information Theory}, 2nd ed.,
Wiley, 2006.

\bibitem{sakarovitch2009}
J.~Sakarovitch,
\emph{Elements of Automata Theory},
Cambridge University Press, 2009.

\bibitem{peller2003}
V.~V.~Peller,
\emph{Hankel Operators and Their Applications},
Springer, 2003.

\bibitem{rabin1963}
M.~O.~Rabin,
``Probabilistic automata,''
\emph{Information and Control}, vol.~6, no.~3, pp.~230--245, 1963.

\bibitem{ambainis1996}
A.~Ambainis,
``The complexity of probabilistic versus deterministic finite automata,''
in \emph{Algorithms and Computation (ISAAC~1996)},
Lecture Notes in Computer Science, vol.~1178,
Springer, 1996, pp.~233--238.

\bibitem{freivalds1981}
R.~Freivalds,
``Probabilistic two-way machines,''
in \emph{Mathematical Foundations of Computer Science 1981},
Lecture Notes in Computer Science, vol.~118,
Springer, 1981, pp.~33--45.

\bibitem{ishigami1993}
Y.~Ishigami and S.~Tani,
``The VC-dimensions of finite automata with $n$ states,''
in \emph{Algorithmic Learning Theory},
1993, pp.~328--341.

\bibitem{littlestone1988}
N.~Littlestone,
``Learning quickly when irrelevant attributes abound: A new linear-threshold
algorithm,''
\emph{Machine Learning}, vol.~2, no.~4, pp.~285--318, 1988.

\bibitem{bendavid2009}
S.~Ben-David, D.~P\'al, and S.~Shalev-Shwartz,
``Agnostic online learning,''
in \emph{Proceedings of COLT}, 2009.

\bibitem{daniely2014}
A.~Daniely, S.~Sabato, S.~Ben-David, and S.~Shalev-Shwartz,
``Multiclass learnability and the ERM principle,''
\emph{Journal of Machine Learning Research},
vol.~15, pp.~41--76, 2014.

\bibitem{aak1971}
V.~M.~Adamjan, D.~Z.~Arov, and M.~G.~Kre\u{\i}n,
``Analytic properties of Schmidt pairs for a Hankel operator and the
generalized Schur--Takagi problem,''
\emph{Math. USSR Sb.}, vol.~15, no.~1, pp.~31--73, 1971.

\bibitem{eckart1936}
C.~Eckart and G.~Young,
``The approximation of one matrix by another of lower rank,''
\emph{Psychometrika}, vol.~1, no.~3, pp.~211--218, 1936.

\bibitem{mirsky1960}
L.~Mirsky,
``Symmetric gauge functions and unitarily invariant norms,''
\emph{Quarterly Journal of Mathematics}, vol.~11, no.~1, pp.~50--58, 1960.

\bibitem{kyfan1951}
K.~Fan,
``Maximum properties and inequalities for the eigenvalues of completely
continuous operators,''
\emph{Proceedings of the National Academy of Sciences},
vol.~37, no.~11, pp.~760--762, 1951.

\bibitem{csiszar1978}
I.~Csisz\'ar,
``Information measures: A critical survey,''
in \emph{Transactions of the 7th Prague Conference on Information Theory},
Academia, Prague, 1978, pp.~73--86.

\bibitem{amari2009}
S.-I. Amari,
``$\alpha$-divergence is unique, belonging to both $f$-divergence and
Bregman divergence classes,''
\emph{IEEE Trans. Inform. Theory}, vol.~55, no.~11, pp.~4925--4931, 2009.

\bibitem{tsybakov2009}
A.~B.~Tsybakov,
\emph{Introduction to Nonparametric Estimation},
Springer, 2009.

\bibitem{vovk1990}
V.~Vovk,
``Aggregating strategies,''
in \emph{Proceedings of the Third Annual Workshop on Computational Learning
Theory}, 1990, pp.~371--386.

\bibitem{tishby1999}
N.~Tishby, F.~C.~Pereira, and W.~Bialek,
``The information bottleneck method,''
in \emph{Proceedings of the 37th Annual Allerton Conference on Communication,
Control, and Computing}, 1999, pp.~368--377.

\bibitem{bialek2001}
W.~Bialek, I.~Nemenman, and N.~Tishby,
``Predictability, complexity, and learning,''
\emph{Neural Computation}, vol.~13, no.~11, pp.~2409--2463, 2001.

\bibitem{boucheron2013}
S.~Boucheron, G.~Lugosi, and P.~Massart,
\emph{Concentration Inequalities: A Nonasymptotic Theory of Independence},
Oxford University Press, 2013.

\bibitem{cesabianchi2006}
N.~Cesa-Bianchi and G.~Lugosi,
\emph{Prediction, Learning, and Games},
Cambridge University Press, 2006.

\bibitem{derooij2014}
S.~de~Rooij, T.~van~Erven, P.~D.~Gr\"unwald, and W.~M.~Koolen,
``Follow the leader if you can, hedge if you must,''
\emph{Journal of Machine Learning Research},
vol.~15, pp.~1281--1316, 2014.

\bibitem{awasthi2015}
P.~Awasthi, M.~Charikar, R.~Krishnaswamy, and A.~K.~Sinop,
``The hardness of approximation of Euclidean $k$-means,''
in \emph{31st International Symposium on Computational Geometry (SoCG 2015)},
LIPIcs vol.~34, Schloss Dagstuhl, 2015, pp.~754--767.

\bibitem{lee2017}
E.~Lee, M.~Schmidt, and J.~Wright,
``Improved and simplified inapproximability for $k$-means,''
\emph{Information Processing Letters}, vol.~120, pp.~40--43, 2017.

\bibitem{aloise2009}
D.~Aloise, A.~Deshpande, P.~Hansen, and P.~Popat,
``NP-hardness of Euclidean sum-of-squares clustering,''
\emph{Machine Learning}, vol.~75, no.~2, pp.~245--248, 2009.

\bibitem{shalizi2001}
C.~R.~Shalizi and J.~P.~Crutchfield,
``Computational mechanics: Pattern and prediction, structure and simplicity,''
\emph{Journal of Statistical Physics}, vol.~104, no.~3--4, pp.~817--879, 2001.

\bibitem{lawvere1973}
F.~W.~Lawvere,
``Metric spaces, generalized logic, and closed categories,''
\emph{Rendiconti del Seminario Matematico e Fisico di Milano},
vol.~43, pp.~135--166, 1973.

\end{thebibliography}

%======================================================================
\section*{Data and Code Availability}
%======================================================================

No empirical dataset was used.  The results are theoretical, and the
computational material reported here is of two kinds, held to different
standards.

Machine-checked statements are formalized in Lean~4 against Mathlib and are
verified by the compiler; these are stated as such in
Remark~\ref{rem:lean-formalization}.

Exhaustive enumerations, numerical evaluations and search results are
\emph{computational observations}.  They are reported with the enumeration
convention, the arithmetic used, and the resulting counts, and they illustrate
or falsify statements rather than proving them; where a theorem depends on one,
the dependence is stated in the theorem.  The corresponding programs, the
extremal machine tables, and the exact outputs accompany the manuscript as
supplementary material.

%======================================================================
\section*{Funding}
%======================================================================

The author received no funding for this work.

%======================================================================
\section*{Competing Interests}
%======================================================================

The author declares no competing interests.

%======================================================================
\section*{Acknowledgments and Declaration of Generative AI}
%======================================================================

The author used generative AI systems, specifically Qwen, Claude, and
ChatGPT, to assist with drafting and manuscript preparation.  The author
reviewed, edited, and verified the content and takes full responsibility for
the final manuscript.

\end{document}